% In this file you should put the actual content of the blueprint.
% It will be used both by the web and the print version.
% It should *not* include the \begin{document}
%
% If you want to split the blueprint content into several files then
% the current file can be a simple sequence of \input. Otherwise It
% can start with a \section or \chapter for instance.

\part{Tight sample complexity bounds for entropic best policy identification}
\label{part:main}

\textbf{Overview}

This blueprint guides the Lean formalization of the paper
\emph{Tight Sample Complexity Bounds for Entropic Best Policy Identification} by Amer Essakine
and Claire Vernade (COLT 2026, \cite{essakine2026tight}). The paper studies best-policy
identification in finite-horizon tabular MDPs under the entropic risk measure
$V^\pi_h(s) = \beta^{-1}\log \E[\exp(\beta \sum_{i \ge h} r_i(S_i, A_i)) \mid S_h = s]$, in the
online (forward) interaction model: the learner plays one deterministic policy per episode from
a fixed initial state $s_1$, observes the states of the episode (the rewards are known and
deterministic), and stops after a random number $\tau$ of episodes to output a policy
$\hat\pi$ which must be $\eps$-optimal with probability at least $1 - \delta$. Its results are:
\begin{enumerate}
  \item[(i)] a lower bound (Theorem~\ref{thm:lower_bound}, the paper's Theorem~3): on a hard
        instance $\mathcal M_0$, every $(\eps,\delta)$-PAC algorithm uses in expectation
        $\Omega\big(\frac{(e^{|\beta| \Gmax(\mathcal M_0)} - 1)^2}{e^{|\beta| \Gmax(\mathcal M_0)}}
        \frac{S A H}{(e^{|\beta|\eps} - 1)^2} \log\frac1\delta\big)$ episodes, where
        $\Gmax(\mathcal M)$ is the maximal achievable return;
  \item[(ii)] the algorithm Entropic-BPI (KL-based exploration bonuses in the exponential
        space and an entropic stopping rule) and its matching upper bound
        (Theorems~\ref{thm:upper_bound_pac} and~\ref{thm:upper_bound_complexity}, the paper's
        Theorem~4): it is $(\eps,\delta)$-PAC and stops after
        $\tilde O\big(\frac{e^{2|\beta|\eps}}{(e^{|\beta|\eps} - 1)^2}
        \frac{(e^{|\beta| \Gmax(\mathcal M)} - 1)^2}{e^{|\beta| \Gmax(\mathcal M)}} S A H\big)$
        episodes with probability at least $1 - \delta$, closing the $e^{|\beta| H}$ gap of the
        previous bounds.
\end{enumerate}

\textbf{Goals of the formalization.} Two goals drive the organization of this document.
First, formalize the two main results of the paper. Second, build a reusable library, to be
contributed to \emph{Lean Machine Learning} (LML) and Mathlib: finite-horizon MDPs and their
values (on top of LML's MDPs and its algorithm--environment framework), entropic values,
empirical models, identification algorithms for episodic environments, the change of measure at
a stopping time, and the concentration inequalities (Sanov, a Bernoulli maximal inequality, the
self-normalized Bernstein inequality, the KL-Bernstein inequality) are introduced in the
generality in which they are useful, even when the paper only needs a special case.

\textbf{Deviations from the paper.} The blueprint departs from the paper whenever the paper's
statement is imprecise or a different route is simpler to formalize while giving the same (or
a stronger) statement; every such point is recorded in a remark next to the statement it
concerns. The main ones:
\begin{itemize}
  \item Never-visited state-action pairs get the trivial bounds (the paper's $\alpha(0)/0 = \infty$),
        the greedy policy for $\beta < 0$ is $\argmin_a \Ul$ (as in the appendix), and the
        clips of the optimistic backups and of the certificate at step $h$ are $e^{\beta(H+1-h)}$
        (the ranges of $U^*_h$; the paper's exponents are off by one).
  \item The third term of the exploration bonus carries the KL rate $\alpha$ (as in the
        statement of the paper's Lemma~7), not $\alpha^*$: with $\alpha^*$ the certificate
        lemmas fail (Chapter~\ref{chap:algorithm}, Definition~\ref{def:bonus}).
  \item The Bernstein concentration event uses step-dependent ranges and a non-strict
        inequality (with the paper's strict one, the event is empty for $\beta < 0$), the
        pseudo-counts are the random predictable sums of conditional visit probabilities, and
        the rate terms are $\min\{\alpha(n)/n, 1\}$ with value $1$ for unvisited pairs; the
        additive term of the certificate recursion uses $e^{\beta(H+1-h)}$ (resp. $1$ for
        $\beta < 0$) so that it dominates the clip of unvisited pairs. The KL event is obtained
        from a time-uniform mixture-martingale bound rather than from Sanov's bound and a union
        bound over $n$ (which diverges with the paper's rate).
  \item The paper's Lemma~27 needs $\alpha \ge \alpha^{\mathrm{cnt}} \ge 1$, the counting
        argument has constant $16$ (not $3$), and the paper's Lemma~29 is false as stated: the
        factor $(C + D)$ in its logarithm must be squared (Chapter~\ref{chap:pre_analysis}).
  \item Theorem~4 is stated in its detailed appendix form, with its two conclusions as two
        separate statements, and with the two-term bound that the proof actually gives (the
        paper's ``in particular'' simplification assumes that the first term dominates); its
        absolute constants are generous so that the formal proof has slack, and its
        hypothesis $\eps \le 2/(|\beta| H S)$ is kept although the proof does not use it.
  \item The PAC property is a property of the algorithm (LML's \texttt{IdentAlg.IsPAC}) for the
        class of MDPs with the known reward function of the algorithm; Theorem~3 assumes PAC
        for the class with the reward function of the hard instances (a stronger theorem).
  \item The hard instance of Theorem~3 uses a possibly non-full tree (the first $S-3$ nodes of
        the $A$-ary tree in breadth-first order), which the paper leaves to
        \cite{domingues2021episodic}; the leaves are the childless nodes. Condition~A is
        stated for the effective horizon $H' = H - \lfloor H/3 \rfloor - d$ of the construction
        (the paper's condition with $H$ does not imply admissibility) and with a margin
        $p_+ \le (1 + p_-)/2$: the paper's bound on the binary divergence $\klbin(p_-, p_+)$
        is false near the boundary of its condition. The gap $\Delta = p_+ - p_-$ is chosen so
        that non-realizing policies are strictly more than $\eps$-suboptimal. The constant
        $1/1650$ becomes $1/5000$ (the proof gives $1/3456$).
  \item The change of measure (Lemma~17 of the paper) is proved in LML generality for two
        stationary environments and any stopping time of the history filtration, and applied
        to the episode environment, whose actions are policies.
  \item Lemma~7 (and~12) is stated without a sign assumption on $\Zt - Z^*$, the optimism lemma
        being proved jointly by backward induction; the $\eps$-optimality at stopping for
        $\beta < 0$ is proved through $Z^\pi_1 \ge Z^*_1 \ge \Zl_1$ (the paper's argument uses
        an unproved lower bound on $Z^\pi_1$ and has a sign typo).
\end{itemize}

\textbf{Reading guide.}
Part~\ref{part:main} follows the paper: Chapter~\ref{chap:mdp} sets up the model (finite-horizon
MDPs, entropic values, the episode environment and PAC identification),
Chapter~\ref{chap:empirical} the empirical transitions and the concentration events,
Chapter~\ref{chap:algorithm} the Entropic-BPI algorithm, Chapters~\ref{chap:analysis_pos}
and~\ref{chap:analysis_neg} its analysis for $\beta > 0$ and $\beta < 0$,
Chapter~\ref{chap:sample_complexity} the sample complexity theorem and
Chapter~\ref{chap:lower_bound} the lower bound.
Part~\ref{part:prereq} collects the prerequisites that are not in Mathlib or LML, decomposed
into small lemmas: finite-horizon MDP theory (occupancy measures, unrolling of recursions,
entropic variances), concentration inequalities, the change of measure, and elementary
inequalities.

\textbf{Conventions.} The paper's steps are $h = 1, \dots, H$ with the terminal step $H + 1$;
the blueprint keeps this indexing, and the Lean remarks give the $0$-indexed translation
(a step-$h$ quantity of the paper is at Lean index $h - 1$, and the paper's
$e^{\beta(H - h)}$ is $e^{\beta k}$ with $k = H - 1 - h$ the number of steps after the Lean
step $h$). Episodes are indexed by $t = 0, 1, \dots$ (LML rounds): $t$ is the number of
episodes already played and $\pi^{t+1}$ the policy computed from them. All random objects live
on a probability space $(\Omega, \mathcal F, \Prob)$; $\Prob^\pi_{(s,h)}$, $\E^\pi_{(s,h)}$ denote the
law of the trajectory of the policy $\pi$ started at the state $s$ at step $h$, and
$\E^\pi := \E^\pi_{(s_1, 1)}$.

\textbf{Status.} The formalization proceeds in two phases. Phase~1 states the headline results
of the paper in Lean, over definitions written in library generality: each such statement
carries a \verb+\lean+ tag naming its Lean declaration and is marked as formalized, with its
proof left to phase~2. Each headline statement is also a standalone challenge file for
\texttt{leanprover/comparator} (directory \texttt{comparator/}). Phase~2 will formalize the
proofs, following this document.

\textbf{Blueprint author.} Rémy Degenne. Anyone is welcome to contribute!

\chapter{Model: finite-horizon MDPs, entropic values and best-policy identification}
\label{chap:mdp}

This chapter fixes the objects of the paper (Section~2 of \cite{essakine2026tight}, together
with the exponential Bellman equation of its Section~4) in the form in which they are
formalized on top of the LML framework of sequential algorithms and environments
(\texttt{Learning.Algorithm}, \texttt{Learning.Environment}, \texttt{Learning.IsAlgEnvSeq},
\texttt{Learning.stationaryEnv}) and of LML's MDPs (\texttt{Learning.MDP}). Its definitions are
library material (namespace \texttt{Learning.MDP}, files \texttt{Episodic.lean},
\texttt{Entropic.lean} of \texttt{Essakine2026Tight/LeanMachineLearning/ReinforcementLearning/MDP/}).

\textbf{Indexing convention.} The paper's steps are $h = 1, \dots, H$ with the terminal step
$H + 1$: a state-action pair is visited at the steps $1, \dots, H$ and the state $S_{H+1}$ is
the last state of an episode. The blueprint keeps this $1$-indexing throughout. Lean is
$0$-indexed: the steps of the trajectory laws and of the values are $h \in \N$ (the horizon
$H$ is an argument of the values, and the values are trivial from the step $H$ on), while the
steps of the $H$-step policies played by the learner and of the quantities computed by the
algorithm are $h \in \texttt{Fin H}$ or $h \in \texttt{Fin (H + 1)}$; the paper's step $h$ is
Lean's $h - 1$. The quantity
$e^{\beta(H - h)}$ of the paper (the range of $Z^*_{h+1}$, see Lemma~\ref{lem:opt_exp_value_range})
is $e^{\beta k}$ with $k = H - 1 - h$ the number of steps \emph{after} the Lean step $h$, and the
clip $e^{\beta(H + 1 - h)}$ of Chapter~\ref{chap:algorithm} is $e^{\beta(k + 1)}$. Backward
recursions are indexed in Lean by the number $j$ of steps to the end: $j = 0$ is the terminal
step $H + 1$ and the paper's step $h$ has $j = H + 1 - h$. Episodes are indexed by
$t = 0, 1, \dots$: $t$ is the number of episodes already played and $\pi^{t+1}$ is the policy
computed from them (LML rounds are $0$-indexed: the policy $\pi^{t+1}$ is the action of the
round $t$).

\section{MDPs and policies}
\label{sec:mdp_mdps}

\begin{definition}[Finite MDP]\label{def:mdp}
\lean{MDP, Learning.MDP.env, Learning.MDP.policyAlg, Learning.MDP.policyMeasure}\leanok
A (stationary) \emph{Markov decision process} on a finite state space $\Ss$ and a finite action
space $\As$ is given by a transition kernel $P(\cdot \mid s, a)$, a probability measure on
$\Ss$ for every $(s, a) \in \Ss \times \As$, and a reward kernel $R(\cdot \mid s, a)$, a
probability measure on $\R$ for every $(s, a)$ (the law of the reward received when playing $a$
in $s$). Its \emph{environment} from an initial law $\mu_0$ on $\Ss$ is the LML environment
whose observation at each round is the current state (drawn from $\mu_0$ at the first round
and from $P(\cdot \mid s, a)$ after the round $(s, a)$) and whose feedback to the action $a$ in
the state $s$ is a reward drawn from $R(\cdot \mid s, a)$.
\end{definition}

\textbf{Lean remark.} This is LML's \texttt{Learning.MDP 𝓢 𝓐 𝓡} (fields \texttt{M.P},
\texttt{M.R}, Markov kernels on \texttt{𝓢 × 𝓐}), copied from the \texttt{mdp} branch of LML to
\texttt{Essakine2026Tight/LeanMachineLearning/ReinforcementLearning/MDP/Basic.lean}, and its
environment \texttt{MDP.env M μ₀ : Environment 𝓢 𝓐 𝓡}. The state and action spaces of this
blueprint are finite types with the discrete $\sigma$-algebra (\texttt{Fintype},
\texttt{MeasurableSingletonClass}), so every function on them is measurable and the kernels are
determined by their transition probabilities.

\begin{definition}[Finite-horizon MDP]\label{def:episodic_mdp}
\lean{Learning.MDP.Episodic.EpisodicMDP, Learning.MDP.Episodic.EpisodicMDP.ofDet, Learning.MDP.Episodic.EpisodicMDP.transVec, Learning.MDP.Episodic.EpisodicMDP.meanReward, Learning.MDP.Episodic.EpisodicMDP.RewardsIn}\leanok
\uses{def:mdp}
A \emph{finite-horizon MDP} with horizon $H \ge 1$ is
$\mathcal M = (\Ss, \As, H, (p_h)_{h \le H}, (R_h)_{h \le H})$ where $\Ss$, $\As$ are finite,
$p_h(\cdot \mid s, a)$ is a probability vector on $\Ss$ (the transition kernel at step $h$) and
$R_h(\cdot \mid s, a)$ a probability measure on $\R$ (the reward kernel at step $h$), for
$h \in \{1, \dots, H\}$ and $(s, a) \in \Ss \times \As$. We write $S := \abs{\Ss}$, $A := \abs{\As}$,
$(p_h f)(s, a) := \sum_{s'} p_h(s' \mid s, a) f(s') = \E_{S' \sim p_h(\cdot \mid s, a)}[f(S')]$
for $f : \Ss \to \R$, and $\Var_{p_h}(f)(s, a) := p_h\big((f - (p_h f)(s, a))^2\big)(s, a)$;
more generally $q f := \sum_{s'} q(s') f(s')$ and $\Var_q(f) := q(f - qf)^2$ for a probability
vector $q$ on $\Ss$. The \emph{mean reward} is $r_h(s, a) := \int x\, R_h(dx \mid s, a)$.
\end{definition}

\textbf{Lean remark.} \texttt{EpisodicMDP S A} with fields
\texttt{trans : ℕ → Kernel (S × A) S} and \texttt{reward : ℕ → Kernel (S × A) ℝ}
(Markov kernels indexed by all the steps $h \in \N$, the horizon $H$ being an argument of the
values and of the episode environment; the state and action spaces are arbitrary measurable
spaces, the finiteness of the paper entering through the hypotheses of the statements that
need it; the reward and the next state are independent given $(s, a)$, which is the
paper's model and which the deterministic rewards of Definition~\ref{def:reward_fn} make
harmless). \texttt{M.transVec h s a : S → ℝ} is the vector of transition probabilities
$p_h(\cdot \mid s, a)$ and \texttt{M.meanReward h s a} the mean reward; the operators
$q \mapsto qf$ and $\Var_q$ on finite-state vectors are \texttt{vecExp} and \texttt{vecVar}
(\texttt{Vec.lean}), with \texttt{vecExp (M.transVec h s a) f} for $(p_h f)(s, a)$. A
finite-horizon MDP with the transition kernels \texttt{trans} and a deterministic reward function
\texttt{r} is \texttt{EpisodicMDP.ofDet trans r}.

\begin{definition}[Reward function; the class $\mathcal M_r$]\label{def:reward_fn}
\lean{Learning.MDP.Episodic.EpisodicMDP.HasRewardFn}\leanok
\uses{def:episodic_mdp}
A \emph{reward function} is a map $r : \{1, \dots, H\} \times \Ss \times \As \to [0, 1]$,
written $r_h(s, a)$. The MDP $\mathcal M$ \emph{has reward function $r$} if all its rewards are
deterministic with values given by $r$: $R_h(\cdot \mid s, a) = \delta_{r_h(s, a)}$ for all $h$,
$s$, $a$. For a fixed reward function $r$, $\mathcal M_r$ denotes the class of all finite-horizon
MDPs on $(\Ss, \As, H)$ with reward function $r$ (their transition kernels are arbitrary). All
the MDPs of the paper are in some $\mathcal M_r$: the rewards are known, deterministic and in
$[0, 1]$; only the transitions are unknown to the learner.
\end{definition}

\textbf{Lean remark.} \texttt{M.HasRewardFn r} is the predicate
$\forall h\, s\, a,\ \texttt{M.reward h (s, a) = Measure.dirac (r h s a)}$; the class
$\mathcal M_r$ is the subtype \texttt{\{M : EpisodicMDP S A // M.HasRewardFn r\}}, with
\texttt{r : ℕ → S → A → ℝ}. The
condition $r \in [0, 1]$ is a separate hypothesis \texttt{∀ h s a, r h s a ∈ Set.Icc 0 1} of
the statements that need it (all the results of Chapters~\ref{chap:algorithm}
to~\ref{chap:lower_bound}).

\begin{definition}[Deterministic policies]\label{def:policy}
\lean{Learning.MDP.Episodic.Policy}\leanok
\uses{def:episodic_mdp}
A \emph{(deterministic, non-stationary) policy} is $\pi = (\pi_h)_{h \le H}$ with
$\pi_h : \Ss \to \As$; $\pi_h(s)$ is the action played in the state $s$ at step $h$. There are
finitely many policies ($A^{SH}$). For $g : \Ss \times \As \to \R$ we write
$(\pi_h g)(s) := g(s, \pi_h(s))$ for the composition with the policy at step $h$.
\end{definition}

\textbf{Lean remark.} The trajectory laws and the values are defined for a policy
\texttt{π : ℕ → S → A} on all the steps. As for a stochastic process in Mathlib, the policy is not
bundled with its measurability: the definitions are total (the trajectory law of a policy which
is not measurable at every step is the zero measure, as \texttt{Measure.map} of a non-measurable
function), and the lemmas take the hypothesis \texttt{hπ : ∀ h, Measurable (π h)}. The learner
plays $H$-step policies \texttt{Policy S A H := Fin H → S → A}, a finite type (\texttt{Fintype}
instance by \texttt{Pi.instFintype}) with the discrete $\sigma$-algebra, extended to all the steps
by \texttt{Policy.extend} (the policy on the first $H$ steps, a fixed action beyond;
\texttt{Policy.extend\_fin}, \texttt{Policy.extend\_injective}, and
\texttt{Policy.measurable\_extend} on a countable discrete state space); the values of
\texttt{π.extend} only depend on \texttt{π} (\texttt{expValue\_congr}). The paper also allows
stochastic policies in the notation $(\pi_h g)(s) = \E_{A \sim \pi_h(\cdot \mid s)} g(s, A)$;
only deterministic policies are needed (the learner plays deterministic policies and the
optimal policies are deterministic, Lemma~\ref{lem:opt_bellman}).

\section{Trajectory laws through the time-augmented MDP}
\label{sec:mdp_trajectories}

The trajectory of a policy in a finite-horizon MDP is defined as the trajectory of a
stationary policy in a stationary MDP, the \emph{time-augmented} MDP, whose state records the
current step; this reuses LML's trajectory measures and makes the Markov property a consequence
of the Ionescu--Tulcea construction.

\begin{definition}[Time-augmented MDP]\label{def:layered_mdp}
\lean{Learning.MDP.Episodic.EpisodicMDP.augmented, Learning.MDP.Episodic.EpisodicMDP.augTrans, Learning.MDP.Episodic.EpisodicMDP.augReward, Learning.MDP.Episodic.augPolicy, Learning.MDP.Episodic.measurable_augPolicy, Learning.MDP.Episodic.Policy.extend, Learning.MDP.Episodic.Policy.measurable_extend}\leanok
\uses{def:episodic_mdp, def:policy, def:mdp}
The \emph{time-augmented MDP} of $\mathcal M$ is the stationary MDP (Definition~\ref{def:mdp})
on the state space $\Ss \times \N$ with actions $\As$: from a state $(s, h)$ and an action $a$,
the reward has law $R_h(\cdot \mid s, a)$ and the next state is $(s', h + 1)$ with
$s' \sim p_h(\cdot \mid s, a)$. A policy $\pi = (\pi_h)_{h \in \N}$ of $\mathcal M$ becomes the
stationary policy $\pi^{\mathrm{aug}}(s, h) := \pi_h(s)$; an $H$-step policy is extended by a
fixed action beyond the horizon.
\end{definition}

\textbf{Lean remark.} \texttt{M.augmented : MDP (S × ℕ) A ℝ} (kernels \texttt{M.augTrans},
\texttt{M.augReward}) and \texttt{augPolicy π : S × ℕ → A}, measurable when every
\texttt{π h} is (\texttt{measurable\_augPolicy}, by
\texttt{measurable\_from\_prod\_countable\_left}). The steps are all of $\N$: there is no
terminal layer, the trajectory runs forever and the finite horizon only enters through the
return of the first $H - h$ rounds (Definition~\ref{def:episode_return}). An $H$-step policy
\texttt{π : Policy S A H} is \texttt{π.extend}, which plays a default action of the nonempty
type \texttt{A} beyond the horizon.

\begin{definition}[Trajectory law from a state at a step]\label{def:step_law}
\lean{Learning.MDP.Episodic.stepLaw}\leanok
\uses{def:layered_mdp}
For a policy $\pi$, a state $s$ and a step $h \in \N$, the \emph{trajectory law}
$\Prob^\pi_{(s, h)}$ is the law of the trajectory of the stationary policy $\pi^{\mathrm{aug}}$ in
the time-augmented MDP started at the state $(s, h)$, i.e.\ the sequence of rounds
$((S_{h+j}, h + j), A_{h+j}, R_{h+j})_{j \ge 0}$ (observation, action, feedback). We write
$(S_i, A_i, R_i)_{i \ge h}$ for this trajectory (so $S_h = s$, $A_i = \pi_i(S_i)$ and
$R_i \sim R_i(\cdot \mid S_i, A_i)$ for $i \ge h$) and $\E^\pi_{(s, h)}$ for the expectation under
$\Prob^\pi_{(s, h)}$; $\Prob^\pi := \Prob^\pi_{(s_1, 1)}$ and $\E^\pi := \E^\pi_{(s_1, 1)}$ when an
initial state $s_1$ is fixed.
\end{definition}

\textbf{Lean remark.} \texttt{stepLaw M π (s, h) := M.augmented.policyMeasure (augPolicy π) \_ (s, h)}
for \texttt{π : ℕ → S → A} measurable at every step (\texttt{stepLaw\_of\_measurable}; it is
\texttt{0} otherwise) and \texttt{h : ℕ}, a probability measure
(\texttt{isProbabilityMeasure\_stepLaw}) on
\texttt{ℕ → Round (S × ℕ) A ℝ} (LML \texttt{MDP.policyMeasure}, the trajectory measure
\texttt{trajMeasure} of the stationary policy algorithm \texttt{policyAlg} in
\texttt{MDP.env M.augmented (Measure.dirac (s, h))}); \texttt{stepLawKernel M π h} is the same
law as a kernel from the state. The round $j$ of the trajectory corresponds to the paper's step
$h + j$.

\begin{definition}[Episode return and state sequence]\label{def:episode_return}
\lean{Learning.MDP.Episodic.episodeReturn, Learning.MDP.Episodic.episodeStates, Learning.MDP.Episodic.Traj}\leanok
\uses{def:step_law}
The \emph{return from step $h$} of a trajectory is $R^\pi_h := \sum_{i = h}^H R_i$, the sum of
the first $H + 1 - h$ rewards of the trajectory from $(s, h)$ ($0$ if $h > H$). The
\emph{state sequence} of an episode started at $(s_1, 1)$ is $(S_1, \dots, S_{H+1})$: the
observation states of the first $H + 1$ rounds. We call the pair (state sequence, rewards) the
\emph{episode}.
\end{definition}

\textbf{Lean remark.} \texttt{episodeReturn n : (ℕ → Round (S × ℕ) A ℝ) → ℝ},
\texttt{traj ↦ ∑ k ∈ range n, (traj k).feedback}, the sum of the first \texttt{n} rewards: the
return from the ($0$-indexed) step \texttt{h} is \texttt{episodeReturn (H - h)}, and
\texttt{episodeStates H : \_ → Traj S H}, \texttt{traj ↦ fun i : Fin (H + 1) ↦ (traj i).obs.1},
where \texttt{Traj S H := Fin (H + 1) → S}. Both are measurable (finite sums and coordinate
projections).

\begin{lemma}[Markov property of the trajectory of a stationary policy]\label{lem:policy_markov}
\lean{Learning.MDP.hasCondDistrib_shiftRound_trajMeasure, Learning.MDP.hasLaw_step_zero_policyMeasure, Learning.MDP.hasCondDistrib_obs_one_trajMeasure}\leanok
\uses{def:mdp}
Let $M$ be an MDP whose state space has measurable singletons, $\pi$ a measurable
stationary deterministic policy and $\mu_0$ an initial law, and let $(O_j, A_j, Y_j)_{j \ge 0}$ be
the canonical trajectory of $\pi$ in the environment of $M$ from $\mu_0$. Then the first round
$(O_0, A_0, Y_0)$ has law $\mu_0(ds)\,\delta_{\pi(s)}(da)\,R(dy \mid s, a)$, the state $O_1$ has
conditional law $P(\cdot \mid O_0, A_0)$ given the first round, and the shifted trajectory
$(O_{j+1}, A_{j+1}, Y_{j+1})_{j \ge 0}$ has conditional law $\Prob^\pi_{O_1}$ (the trajectory law of
$\pi$ from the state $O_1$) given $((O_0, A_0, Y_0), O_1)$.
\end{lemma}

\begin{proof}
\leanok
Let $Q$ be the law of $((O_0, A_0, Y_0), O_1)$ followed by an independent trajectory drawn from
$\Prob^\pi_{O_1}$, a measure on (round, state, trajectory), and let $\Phi$ map such a triple to the
trajectory obtained by prepending the round to the trajectory. Under $Q$, the process
$\Phi$ has the conditional laws of the canonical trajectory: its round $0$ has the law of the
first round; its round $1$, given round $0$, has observation $P(\cdot \mid O_0, A_0)$ and then the
action and reward of $\pi$ from that state (the initial round of $\Prob^\pi_{O_1}$); its round
$j + 2$, given the previous rounds, has the step kernel of the pair at round $j + 1$, which is the
step kernel of $\Prob^\pi_{O_1}$ at round $j + 1$ because the step kernels at positive rounds do
not depend on the initial law. By uniqueness of the Ionescu--Tulcea measure with given
conditional laws (LML \texttt{Kernel.hasLaw\_trajMeasureFin}), $Q \circ \Phi^{-1}$ is the law of
the canonical trajectory. Since the observation of round $0$ of $\Prob^\pi_{s'}$ is $s'$ almost
surely, $\Phi$ is $Q$-almost surely inverted by (first round, state of round $1$, shift), so the
law of this triple under the trajectory law is $Q$, which is the claim.
\end{proof}

\textbf{Lean remark.} File \texttt{MDP/Markov.lean}, in LML generality (\texttt{Learning.MDP},
arbitrary state space with \texttt{MeasurableSingletonClass}). The trajectory law of $\pi$ as a
kernel from the initial state is \texttt{policyKernel M hπ}, built with Mathlib's
\texttt{Kernel.traj} from the step kernels (\texttt{iicStepKernel}) so that its measurability
is Mathlib's; the first round and the state of round $1$ are
\texttt{firstRoundState}, the shift is \texttt{shiftRound}, the auxiliary measure is
\texttt{shiftMeasure} and the prepending map is \texttt{consRound}. The general
\texttt{HasCondDistrib} tools used are in \texttt{Mathlib/Probability/HasCondDistrib.lean}
(\texttt{hasCondDistrib\_snd\_compProd}, \texttt{HasCondDistrib.comp\_hasLaw},
\texttt{HasCondDistrib.compProd\_left}).

\begin{lemma}[Markov property of the trajectory law]\label{lem:step_law_markov}
\lean{Learning.MDP.Episodic.hasLaw_stepLaw_succ, Learning.MDP.Episodic.ae_obs_zero_stepLaw, Learning.MDP.Episodic.ae_action_zero_stepLaw, Learning.MDP.Episodic.hasLaw_feedback_zero_stepLaw, Learning.MDP.Episodic.ae_feedback_eq_policy_stepLaw, Learning.MDP.Episodic.episodeReturn_succ_eq_add_shiftRound, Learning.MDP.Episodic.ae_stepLaw_succ_of_ae, Learning.MDP.Episodic.lintegral_stepLaw_succ, Learning.MDP.Episodic.integral_stepLaw_succ}\leanok
\uses{def:step_law, def:layered_mdp}
Let $h \in \N$, $s \in \Ss$ and $\pi$ a policy. Under $\Prob^\pi_{(s, h)}$:
\begin{enumerate}
  \item[(i)] the first round is $S_h = s$, $A_h = \pi_h(s)$, and $(R_h, S_{h+1})$ has law
        $R_h(\cdot \mid s, \pi_h(s)) \otimes p_h(\cdot \mid s, \pi_h(s))$;
  \item[(ii)] the conditional law of the shifted trajectory $(S_i, A_i, R_i)_{i \ge h+1}$ given
        the first round is $\Prob^\pi_{(S_{h+1}, h + 1)}$: for every measurable bounded (or
        nonnegative) $F$ on trajectories,
        \[
          \E^\pi_{(s, h)}\big[F\big((S_i, A_i, R_i)_{i \ge h + 1}\big)\big]
          = \E_{(R, s') \sim R_h(\cdot \mid s, \pi_h(s)) \otimes p_h(\cdot \mid s, \pi_h(s))}
            \Big[ \E^\pi_{(s', h + 1)}[F] \Big],
        \]
        and more generally
        $\E^\pi_{(s, h)}[G(R_h, S_{h+1})\, F((S_i, A_i, R_i)_{i \ge h+1})]
        = \E_{(R, s')}[G(R, s')\, \E^\pi_{(s', h+1)}[F]]$ for bounded measurable $G$.
\end{enumerate}
\end{lemma}

\begin{proof}
\uses{def:step_law, def:layered_mdp, lem:policy_markov}
\leanok
The trajectory measure of an algorithm in an environment is the Ionescu--Tulcea measure built
from the initial law of the observation and the step kernels (LML \texttt{trajMeasure}). For the
stationary policy $\pi^{\mathrm{aug}}$ in the time-augmented MDP started at $(s, h)$: the
observation of the round $0$ is $(s, h)$ almost surely (initial law $\delta_{(s, h)}$), the action
is the deterministic $\pi^{\mathrm{aug}}(s, h) = \pi_h(s)$ (the policy algorithm is deterministic),
and the feedback and next observation are drawn from the augmented step kernel at
$((s, h), \pi_h(s))$,
which is $R_h(\cdot \mid s, \pi_h(s))$ for the reward and $s' \mapsto (s', h + 1)$ with
$s' \sim p_h(\cdot \mid s, \pi_h(s))$ for the next state, independently: this is (i). For (ii), the
canonical trajectory process is an algorithm-environment sequence for the policy algorithm and
the augmented environment (LML \texttt{isAlgEnvSeq\_trajMeasure}); the conditional law of the
future rounds $j \ge 1$ given the round $0$ is, by the Ionescu--Tulcea construction, the
trajectory measure of the same algorithm-environment pair started at the observation of round
$1$, that is $(S_{h+1}, h + 1)$ (the policy algorithm and the MDP environment are stationary:
their kernels at round $j$ only depend on the last observation and action, not on $j$ nor on
the earlier rounds). Since the observation of round $1$ is $(S_{h+1}, h + 1)$ and the law of
$(R_h, S_{h+1})$ is the product kernel of (i), the displayed identities follow by integrating
the conditional law, first for $F$ an indicator of a measurable rectangle of the trajectory
space and then by the monotone class theorem.
\end{proof}

\textbf{Lean remark.} File \texttt{MDP/StepLaw.lean}, derived from
Lemma~\ref{lem:policy_markov} for the time-augmented MDP. For a Lean step \texttt{h : ℕ} (the
paper's $h + 1$), \texttt{hasLaw\_stepLaw\_succ} states that under \texttt{stepLaw M π (s, h)}
the triple (reward of round $0$, state of round $1$, shifted trajectory) has law
\texttt{(M.reward h (s, π h s)).prod (M.trans h (s, π h s)) ⊗ₘ Kernel.prodMkLeft ℝ (stepLawKernel M π (h + 1))},
which contains (i) and (ii); the almost sure statements and integrals along one step are
\texttt{ae\_stepLaw\_succ\_of\_ae}, \texttt{lintegral\_stepLaw\_succ} and
\texttt{integral\_stepLaw\_succ}; with a reward function the first reward is
$r_h(s, \pi_h(s))$ almost surely (\texttt{ae\_feedback\_eq\_policy\_stepLaw}). The return of
$n + 1$ rounds is the first reward plus the return of $n$ rounds of the shifted trajectory
(\texttt{episodeReturn\_succ\_eq\_add\_shiftRound}, an identity of functions): this is
$R^\pi_h = R_h + R^\pi_{h+1}$, used by every backward recursion, and the recursions stop at the
step $H$, where the return is the empty sum (no terminal layer is needed).

\begin{definition}[Additive values]\label{def:eval_value}
\uses{def:step_law, def:episode_return}
For $f : \{1, \dots, H\} \times \Ss \times \As \to \R$, a policy $\pi$, a step $h$ and a state $s$,
the \emph{additive value} of $\pi$ for $f$ is
$\E^\pi_{(s, h)}\big[\sum_{i = h}^H f_i(S_i, A_i)\big]$; with $f = r$ (a reward function) it is the
usual (risk-neutral) value $\E^\pi_{(s, h)}[R^\pi_h]$ of the MDP with reward function $r$.
\end{definition}

\textbf{Lean remark.} Not formalized (it would be
\texttt{∫ traj, ∑ k ∈ range (H - h), f (h + k) (traj k).obs.1 (traj k).action ∂(stepLaw M π (s, h))}).
It is not needed for the results of the paper (the entropic values are not additive) and
serves as the risk-neutral reference point.

\section{Exponential and entropic values}
\label{sec:mdp_values}

Throughout, $\beta \ne 0$ is the risk parameter. The paper works in the \emph{exponential space}:
all the recursions are for $Z = e^{\beta V}$, which satisfies a linear Bellman equation, and the
entropic values are recovered by a logarithm.

\begin{definition}[Exponential values]\label{def:exp_value}
\lean{Learning.MDP.Episodic.expValue, Learning.MDP.Episodic.expQ}\leanok
\uses{def:step_law, def:episode_return, def:episodic_mdp}
For $\beta \ne 0$, a policy $\pi$, a step $h \in \{1, \dots, H + 1\}$ and a state $s$, the
\emph{exponential value} is
\[
  Z^\pi_h(s) := \E^\pi_{(s, h)}\big[e^{\beta R^\pi_h}\big] ,
  \qquad\text{so that } Z^\pi_{H+1}(s) = 1 ,
\]
and for $h \le H$ and $a \in \As$ the \emph{exponential state-action value} is
\[
  U^\pi_h(s, a) := \E_{(R, S') \sim R_h(\cdot \mid s, a) \otimes p_h(\cdot \mid s, a)}
    \big[e^{\beta R}\, Z^\pi_{h+1}(S')\big] .
\]
When $\mathcal M$ has the reward function $r$ (Definition~\ref{def:reward_fn}),
$U^\pi_h(s, a) = e^{\beta r_h(s, a)}\,(p_h Z^\pi_{h+1})(s, a)$.
\end{definition}

\textbf{Lean remark.} \texttt{expValue M H β π h s := ∫ traj, Real.exp (β * episodeReturn (H - h) traj) ∂(stepLaw M π (s, h))}
for \texttt{h : ℕ} and \texttt{π : ℕ → S → A}, and
\texttt{expQ M H β π h s a := (∫ x, Real.exp (β * x) ∂M.reward h (s, a)) * ∫ s', expValue M H β π (h + 1) s' ∂M.trans h (s, a)};
with a reward function, \texttt{expQ = Real.exp (β * r h s a) * ∫ s', expValue M H β π (h + 1) s' ∂M.trans h (s, a)}
(\texttt{expQ\_of\_hasRewardFn}), and on a finite state space
\texttt{= Real.exp (β * r h s a) * vecExp (M.transVec h s a) (expValue M H β π (h + 1))}
(\texttt{expQ\_of\_hasRewardFn\_eq\_vecExp}). The terminal value $Z^\pi_{H+1} = 1$ is a lemma
(\texttt{expValue\_of\_le}: the return of $H - h = 0$ rounds is $0$ for $h \ge H$), not a
separate definition.

\begin{definition}[Entropic values]\label{def:entropic_value}
\lean{Learning.MDP.Episodic.entropicValue, Learning.MDP.Episodic.entropicQ}\leanok
\uses{def:exp_value}
The \emph{entropic value} and \emph{entropic state-action value} of $\pi$ are
\[
  V^\pi_h(s) := \frac1\beta \log Z^\pi_h(s)
  = \frac1\beta \log \E^\pi_{(s, h)}\Big[\exp\Big(\beta \sum_{i = h}^H R_i\Big)\Big],
  \qquad
  Q^\pi_h(s, a) := \frac1\beta \log U^\pi_h(s, a) ,
\]
so that $V^\pi_{H+1} = 0$ and $Z^\pi_h = e^{\beta V^\pi_h}$, $U^\pi_h = e^{\beta Q^\pi_h}$ (the
logarithms are well defined by Lemma~\ref{lem:exp_value_range}).
\end{definition}

\textbf{Lean remark.} \texttt{entropicValue M H β π h s := β⁻¹ * Real.log (expValue M H β π h s)},
\texttt{entropicQ}. The entropic risk measure with $\beta > 0$ is risk-seeking and with
$\beta < 0$ risk-averse \cite{borkar2001risk}; the formalization treats both signs with the
same definitions.

\begin{lemma}[Exponential Bellman equation]\label{lem:exp_bellman}
\lean{Learning.MDP.Episodic.expValue_succ, Learning.MDP.Episodic.expValue_of_le, Learning.MDP.Episodic.lexpValue_succ, Learning.MDP.Episodic.lexpValue_lt_top, Learning.MDP.Episodic.expQ_of_hasRewardFn, Learning.MDP.Episodic.expQ_of_hasRewardFn_eq_vecExp, Learning.MDP.Episodic.integrable_exp_mul_of_rewardsIn, Learning.MDP.Episodic.expValue_congr}\leanok
\uses{def:exp_value, lem:step_law_markov}
Assume that the rewards are bounded: $R_h(\cdot \mid s, a)$ is supported in a fixed interval
$[a, b]$ for all $h, s, a$ (true for a reward function with values in $[0, 1]$). Then
$Z^\pi_{h} = 1$ for $h \ge H + 1$ and for every policy $\pi$, step $h \le H$ and state $s$,
\[
  Z^\pi_h(s) = U^\pi_h(s, \pi_h(s)),
\]
i.e.\ $Z^\pi_h(s) = e^{\beta r_h(s, \pi_h(s))}\,(p_h Z^\pi_{h+1})(s, \pi_h(s))$ when $\mathcal M$
has the reward function $r$.
\end{lemma}

\begin{proof}
\uses{lem:step_law_markov, def:exp_value, def:episode_return}
\leanok
Under $\Prob^\pi_{(s, h)}$, $R^\pi_h = R_h + R^\pi_{h+1}$ where $R^\pi_{h+1} = \sum_{i \ge h+1} R_i$
is a function of the shifted trajectory, so
$e^{\beta R^\pi_h} = e^{\beta R_h}\, e^{\beta R^\pi_{h+1}}$. Lemma~\ref{lem:step_law_markov}(ii) with
$G(R, s') = e^{\beta R}$ and $F = e^{\beta R^\pi_{h+1}}$ (bounded, the rewards being bounded)
gives
$Z^\pi_h(s) = \E_{(R, S') \sim R_h \otimes p_h(\cdot \mid s, \pi_h(s))}[e^{\beta R} Z^\pi_{h+1}(S')]
= U^\pi_h(s, \pi_h(s))$, since $\E^\pi_{(s', h+1)}[e^{\beta R^\pi_{h+1}}] = Z^\pi_{h+1}(s')$ by
definition. The deterministic-reward form is Definition~\ref{def:exp_value}.
\end{proof}

\textbf{Lean remark.} The identity is first proved for the Lebesgue integrals
$L^\pi_h(s) := \int^- e^{\beta R^\pi_h}\,d\Prob^\pi_{(s,h)} \in [0, \infty]$ (\texttt{lexpValue}),
without any hypothesis: $L^\pi_h(s) = \big(\int^- e^{\beta x}R_h(dx \mid s, a)\big)\sum_{s'} p_h(s' \mid s, a)L^\pi_{h+1}(s')$
with $a = \pi_h(s)$ (\texttt{lexpValue\_succ}). The boundedness hypothesis
(\texttt{M.RewardsIn (Set.Icc a b)}) makes all $L^\pi_h(s)$ finite
(\texttt{lexpValue\_lt\_top}), and then $Z^\pi_h = (L^\pi_h)^{\mathrm{toReal}}$ satisfies the
Bellman equation (\texttt{expValue\_succ}). Without such a hypothesis the Bochner-integral
identity can fail (a non-integrable exponential return has Bochner integral $0$), so the
hypothesis is not an artifact. The backward inductions on the step \texttt{h : ℕ} are
\texttt{horizon\_induction} (the case $h \ge H$, then the step from $h + 1$ to $h < H$).

\begin{lemma}[Ranges of the exponential values]\label{lem:exp_value_range}
\lean{Learning.MDP.Episodic.expValue_mem_Icc, Learning.MDP.Episodic.expQ_mem_Icc, Learning.MDP.Episodic.ae_episodeReturn_mem_Icc, Learning.MDP.Episodic.expValue_pos}\leanok
\uses{def:exp_value, lem:step_law_markov, lem:exp_bellman}
Assume that the rewards of $\mathcal M$ are in $[0, 1]$ ($R_h(\cdot \mid s, a)$ is supported in
$[0, 1]$ for all $h, s, a$; in particular when $\mathcal M$ has a reward function). Then for every
policy $\pi$, step $h \le H + 1$ and state $s$, $R^\pi_h \in [0, H + 1 - h]$
$\Prob^\pi_{(s, h)}$-a.s., and
\[
  Z^\pi_h(s) \in \big[\min(1, e^{\beta(H + 1 - h)}),\ \max(1, e^{\beta(H + 1 - h)})\big],
  \qquad
  U^\pi_h(s, a) \in \big[\min(1, e^{\beta(H + 1 - h)}),\ \max(1, e^{\beta(H + 1 - h)})\big]
\]
(the second for $h \le H$); in particular $Z^\pi_h(s) > 0$ and $U^\pi_h(s, a) > 0$. For
$\beta > 0$ the intervals are $[1, e^{\beta(H+1-h)}]$, for $\beta < 0$ they are
$[e^{\beta(H+1-h)}, 1]$.
\end{lemma}

\begin{proof}
\uses{lem:step_law_markov, def:exp_value, def:episode_return}
\leanok
By Lemma~\ref{lem:step_law_markov}(i) applied at each round (induction on the rounds, using (ii)
to move to the next step), the reward $R_i$ of the round at step $i \le H$ has law
$R_i(\cdot \mid S_i, A_i)$, hence lies in $[0, 1]$ a.s., and the rewards of the rounds at
steps $i > H$ are $0$. So $R^\pi_h = \sum_{i = h}^H R_i \in [0, H + 1 - h]$ a.s. The function
$x \mapsto e^{\beta x}$ is monotone, so $e^{\beta R^\pi_h}$ lies a.s.\ between $1 = e^{0}$ and
$e^{\beta(H + 1 - h)}$, and so does its expectation $Z^\pi_h(s)$ (monotonicity of the integral of a
probability measure). For $U^\pi_h(s, a)$: $e^{\beta R}$ lies between $1$ and $e^\beta$ for
$R \in [0, 1]$, $Z^\pi_{h+1}(S')$ lies between $1$ and $e^{\beta(H - h)}$ by the first part at the
step $h + 1$, and the product of two numbers lying respectively between $1$ and $e^{\beta a}$
and between $1$ and $e^{\beta b}$ ($a, b \ge 0$) lies between $1$ and $e^{\beta(a + b)}$; the
expectation preserves the bounds. Positivity: the lower bounds are positive.
\end{proof}

\textbf{Lean remark.} The ranges are stated with \texttt{min 1 (Real.exp (β * (H - h)))} and
\texttt{max 1 (Real.exp (β * (H - h)))} for the Lean step \texttt{h : ℕ} (\texttt{H - h} steps to
go, truncated subtraction; for \texttt{expQ} at \texttt{h < H} the same exponent). The
Lean proof of the ranges of $Z^\pi$ and $U^\pi$ is analytic, by backward induction on the
exponential Bellman equation (Lemma~\ref{lem:exp_bellman}), with the product rule
\texttt{mul\_mem\_Icc\_min\_max\_exp}; the almost sure bound on the return is proved by
induction on the number of rounds with \texttt{episodeReturn\_succ\_eq\_add\_shiftRound} and
\texttt{ae\_stepLaw\_succ\_of\_ae}. Positivity \texttt{expValue\_pos} only needs bounded rewards.

\begin{lemma}[Monotonicity of the exponential Bellman operator]\label{lem:exp_value_monotone}
\lean{Learning.MDP.vecExp_mono, Learning.MDP.vecExp_nonneg, Learning.MDP.vecExp_const, Learning.MDP.vecExp_add, Learning.MDP.vecExp_const_mul, Learning.MDP.vecExp_mem_Icc, Learning.MDP.abs_vecExp_le, Learning.MDP.transVec_nonneg, Learning.MDP.sum_transVec, Learning.MDP.Episodic.EpisodicMDP.sum_transVec, Learning.MDP.Episodic.integral_eq_vecExp_transVec}\leanok
\uses{def:episodic_mdp}
For every step $h \le H$ and $(s, a)$: $f \mapsto (p_h f)(s, a)$ is linear on $\R^{\Ss}$,
positive ($f \ge 0 \Rightarrow p_h f \ge 0$), monotone ($f \le g \Rightarrow p_h f \le p_h g$),
$p_h \mathbf 1 = 1$, and $\abs{(p_h f)(s, a)} \le \max_{s'} \abs{f(s')}$. Consequently, for every
$c \in \R$ (in particular $c = e^{\beta r_h(s, a)} > 0$), $f \le g$ implies
$e^{c}(p_h f)(s, a) \le e^{c}(p_h g)(s, a)$; the same holds for every probability vector $q$ on
$\Ss$ in place of $p_h(\cdot \mid s, a)$.
\end{lemma}

\begin{proof}
\uses{def:episodic_mdp}
\leanok
$(p_h f)(s, a) = \sum_{s'} p_h(s' \mid s, a) f(s')$ with nonnegative weights summing to $1$:
linearity, positivity and $p_h \mathbf 1 = 1$ are immediate, monotonicity is positivity applied
to $g - f$, and the bound is the triangle inequality. Multiplying by $e^c > 0$ preserves the
order.
\end{proof}

\textbf{Lean remark.} These are the lemmas of \texttt{Vec.lean}, stated for a vector
\texttt{p : S → ℝ} with the hypotheses \texttt{∀ s, 0 ≤ p s} and \texttt{∑ s, p s = 1} where
needed; \texttt{M.transVec h s a} satisfies both (\texttt{transVec\_nonneg},
\texttt{EpisodicMDP.sum\_transVec}), and \texttt{integral\_eq\_vecExp\_transVec} identifies
$\int f\,dp_h(\cdot \mid s, a)$ with \texttt{vecExp (M.transVec h s a) f}.

\begin{definition}[Optimal values]\label{def:opt_value}
\lean{Learning.MDP.Episodic.optEntropicValue, Learning.MDP.Episodic.optExpValue, Learning.MDP.Episodic.optExpQ}\leanok
\uses{def:entropic_value, def:exp_value}
The \emph{optimal entropic value} is $V^*_h(s) := \sup_\pi V^\pi_h(s)$ (a supremum over all
measurable policies, finite for bounded rewards), the \emph{optimal exponential value} is $Z^*_h(s) := e^{\beta V^*_h(s)}$, and the
\emph{optimal exponential state-action value} is, for $h \le H$,
\[
  U^*_h(s, a) := \E_{(R, S') \sim R_h(\cdot \mid s, a) \otimes p_h(\cdot \mid s, a)}
    \big[e^{\beta R} Z^*_{h+1}(S')\big]
  = e^{\beta r_h(s, a)}\,(p_h Z^*_{h+1})(s, a) \text{ for a reward function } r .
\]
The paper writes $Q^*_h := \beta^{-1} \log U^*_h$; then $Z^*_{H+1} = 1$ and $V^*_{H+1} = 0$.
\end{definition}

\textbf{Lean remark.} \texttt{optEntropicValue M H β h s := ⨆ π : \{π : ℕ → S → A // ∀ k, Measurable (π k)\}, entropicValue M H β π h s}
(a \texttt{ciSup} in $\R$ over the policies measurable at every step, so that the zero measure
of a non-measurable policy never enters the supremum; bounded above for bounded rewards:
\texttt{bddAbove\_entropicValue}; it is attained, Lemma~\ref{lem:opt_bellman}),
\texttt{optExpValue M H β h s := Real.exp (β * optEntropicValue M H β h s)},
\texttt{optExpQ M H β h s a} defined as \texttt{expQ} with \texttt{optExpValue M H β (h + 1)} in
place of the value of a policy.

\begin{lemma}[Optimal exponential value as an extremum]\label{lem:opt_exp_value_eq}
\lean{Learning.MDP.Episodic.expValue_le_optExpValue, Learning.MDP.Episodic.optExpValue_le_expValue, Learning.MDP.Episodic.expValue_eq_exp_mul_entropicValue, Learning.MDP.Episodic.entropicValue_le_optEntropicValue, Learning.MDP.Episodic.bddAbove_entropicValue, Learning.MDP.Episodic.entropicValue_le_of_pos, Learning.MDP.Episodic.entropicValue_le_of_neg}\leanok
\uses{def:opt_value, lem:exp_value_range}
Assume $\beta \ne 0$ and that the rewards are in $[0, 1]$. Then for every $h$ and $s$ the values
$V^\pi_h(s)$ are bounded above uniformly in $\pi$ (by $H - h$ for $\beta > 0$, by $0$ for
$\beta < 0$), so $V^*_h(s)$ is finite and $V^\pi_h(s) \le V^*_h(s)$ for every policy $\pi$; hence
$Z^\pi_h(s) \le Z^*_h(s)$ ($\beta > 0$), resp.\ $Z^*_h(s) \le Z^\pi_h(s)$ ($\beta < 0$), for every
policy $\pi$. (The supremum is attained, by the greedy policy of Lemma~\ref{lem:opt_bellman}.)
\end{lemma}

\begin{proof}
\uses{def:opt_value, def:entropic_value, lem:exp_value_range}
\leanok
By Lemma~\ref{lem:exp_value_range}, $Z^\pi_h(s) > 0$ and $Z^\pi_h(s)$ lies between $1$ and
$e^{\beta(H - h)}$, so $V^\pi_h(s) = \beta^{-1} \log Z^\pi_h(s)$ is bounded above uniformly in
$\pi$ and the supremum $V^*_h(s)$ is a real number with $V^\pi_h(s) \le V^*_h(s)$. The map
$x \mapsto e^{\beta x}$ is increasing for $\beta > 0$ and decreasing for $\beta < 0$, and
$Z^\pi_h(s) = e^{\beta V^\pi_h(s)}$ (Definition~\ref{def:entropic_value}, using $Z^\pi_h(s) > 0$),
which gives the inequalities.
\end{proof}

\begin{lemma}[Optimal exponential Bellman equation]\label{lem:opt_bellman}
\lean{Learning.MDP.Episodic.optExpValueRec, Learning.MDP.Episodic.optExpQRec, Learning.MDP.Episodic.optPolicy, Learning.MDP.Episodic.expValue_optPolicy, Learning.MDP.Episodic.expValue_le_optExpValueRec, Learning.MDP.Episodic.optExpValueRec_le_expValue, Learning.MDP.Episodic.optEntropicValue_eq, Learning.MDP.Episodic.optExpValue_eq_optExpValueRec, Learning.MDP.Episodic.optExpValue_eq_expValue_optPolicy, Learning.MDP.Episodic.exists_expValue_eq_optExpValue, Learning.MDP.Episodic.exists_entropicValue_eq_optEntropicValue, Learning.MDP.Episodic.optExpValue_succ_eq, Learning.MDP.Episodic.optExpQ_le_optExpValue, Learning.MDP.Episodic.optExpValue_le_optExpQ, Learning.MDP.Episodic.optExpValue_of_le, Learning.MDP.Episodic.optExpQ_of_hasRewardFn, Learning.MDP.Episodic.optExpQ_of_hasRewardFn_eq_vecExp, Learning.MDP.Episodic.expQ_le_optExpQ, Learning.MDP.Episodic.optExpQ_le_expQ}\leanok
\uses{def:opt_value, lem:exp_bellman, lem:exp_value_monotone, lem:opt_exp_value_eq}
Let $\beta \ne 0$ and assume that the rewards of $\mathcal M$ are in $[0, 1]$ (for instance,
$\mathcal M$ has a reward function $r$ with values in $[0, 1]$), with finitely many actions and
countably many states. Then the supremum $V^*_h(s)$ is attained,
$Z^*_{h} = 1$ for $h \ge H + 1$, $U^*_h(s, a) = e^{\beta r_h(s, a)}(p_h Z^*_{h+1})(s, a)$, and for
every $h \le H$ and $s$,
\[
  Z^*_h(s) = \max_{a \in \As} U^*_h(s, a) \quad (\beta > 0),
  \qquad
  Z^*_h(s) = \min_{a \in \As} U^*_h(s, a) \quad (\beta < 0).
\]
Moreover the \emph{greedy policy} $\pi^*_h(s) \in \argmax_a U^*_h(s, a)$ ($\beta > 0$), resp.\
$\argmin_a U^*_h(s, a)$ ($\beta < 0$), satisfies $Z^{\pi^*}_h(s) = Z^*_h(s)$ for all $h$ and $s$:
it is optimal from every state at every step, and in particular $V^*_h(s) = V^{\pi^*}_h(s)$.
\end{lemma}

\begin{proof}
\uses{lem:exp_bellman, lem:exp_value_monotone, lem:opt_exp_value_eq, def:opt_value, lem:exp_value_range}
\leanok
Let $\beta > 0$ (the case $\beta < 0$ is identical with $\min$ in place of $\max$, $\argmin$ in
place of $\argmax$ and all inequalities reversed). Define the backward recursion
$\tilde Z_h(s) := 1$ for $h \ge H + 1$ and
$\tilde Z_h(s) := \max_a \tilde U_h(s, a)$ with
$\tilde U_h(s, a) := e^{\beta r_h(s, a)} (p_h \tilde Z_{h+1})(s, a)$ for $h \le H$ (the maximum
over the finite set of actions exists), and let $\pi^*_h(s) \in \argmax_a \tilde U_h(s, a)$ be a
greedy policy of this recursion.
\begin{enumerate}
  \item[(a)] $Z^{\pi^*} = \tilde Z$: by backward induction on $h$, using the exponential
        Bellman equation (Lemma~\ref{lem:exp_bellman}) for $\pi^*$ and the definition of $\pi^*_h$.
  \item[(b)] For every policy $\pi$, $Z^\pi_h(s) \le \tilde Z_h(s)$: by backward induction on $h$,
        $Z^\pi_h(s) = e^{\beta r_h(s, \pi_h(s))}(p_h Z^\pi_{h+1})(s, \pi_h(s))
        \le e^{\beta r_h(s, \pi_h(s))}(p_h \tilde Z_{h+1})(s, \pi_h(s)) = \tilde U_h(s, \pi_h(s))
        \le \tilde Z_h(s)$, by the Bellman equation, the induction hypothesis and monotonicity
        (Lemma~\ref{lem:exp_value_monotone}).
\end{enumerate}
By (b), $V^\pi_h(s) \le \beta^{-1}\log \tilde Z_h(s) = V^{\pi^*}_h(s)$ for every $\pi$, so the supremum
$V^*_h(s)$ is at most $V^{\pi^*}_h(s)$, and $V^{\pi^*}_h(s) \le V^*_h(s)$ by
Lemma~\ref{lem:opt_exp_value_eq}: the supremum is attained at $\pi^*$ and
$Z^* = Z^{\pi^*} = \tilde Z$. The Bellman equation and the formula for $U^*_h$ then follow from
the recursion defining $\tilde Z$ and $\tilde U$.
\end{proof}

\textbf{Lean remark.} The recursion is \texttt{optExpValueRec} / \texttt{optExpQRec} (a
well-founded recursion on \texttt{H - h}) and the greedy policy is
\texttt{optPolicy : ℕ → S → A}, built with \texttt{Finite.exists\_max} /
\texttt{Finite.exists\_min} on the finite action type and measurable on a countable discrete
state space (\texttt{measurable\_optPolicy}); (a) is \texttt{expValue\_optPolicy}, (b) is
\texttt{expValue\_le\_optExpValueRec} (\texttt{optExpValueRec\_le\_expValue} for $\beta < 0$),
and the conclusions are \texttt{optEntropicValue\_eq}, \texttt{optExpValue\_eq\_optExpValueRec}
and \texttt{optExpValue\_succ\_eq}. This is where the countability of the state space and the
finiteness of the action space enter (\texttt{[Countable S] [Finite A] [Nonempty A]}); the
supremum over all measurable policies needs no finiteness of the set of policies.

\begin{lemma}[Ranges of the optimal exponential values]\label{lem:opt_exp_value_range}
\lean{Learning.MDP.Episodic.optExpValue_mem_Icc, Learning.MDP.Episodic.optExpQ_mem_Icc, Learning.MDP.Episodic.optExpValue_pos}\leanok
\uses{def:opt_value, lem:exp_value_range, lem:opt_exp_value_eq}
For $\beta \ne 0$ and rewards in $[0, 1]$, for every $h \le H + 1$, $s$ and $a$,
\[
  Z^*_h(s) \in \big[\min(1, e^{\beta(H + 1 - h)}),\ \max(1, e^{\beta(H + 1 - h)})\big],
  \qquad
  U^*_h(s, a) \in \big[\min(1, e^{\beta(H + 1 - h)}),\ \max(1, e^{\beta(H + 1 - h)})\big]
\]
($h \le H$ for the second); in particular $Z^*_{h+1}$ takes values in $[1, e^{\beta(H - h)}]$
for $\beta > 0$ and in $[e^{\beta(H - h)}, 1]$ for $\beta < 0$, an interval of length
$\abs{e^{\beta(H - h)} - 1}$, and $Z^*_h > 0$, $U^*_h > 0$.
\end{lemma}

\begin{proof}
\uses{lem:exp_value_range, lem:opt_exp_value_eq, def:opt_value}
\leanok
By Lemma~\ref{lem:opt_exp_value_eq}, $Z^*_h(s) = Z^{\pi^\circ}_h(s)$ for some policy $\pi^\circ$, and
Lemma~\ref{lem:exp_value_range} gives the range. For $U^*_h(s, a)$, argue exactly as in the proof
of Lemma~\ref{lem:exp_value_range}: $e^{\beta R}$ lies between $1$ and $e^{\beta}$, $Z^*_{h+1}(S')$
between $1$ and $e^{\beta(H - h)}$, and the expectation of the product lies between $1$ and
$e^{\beta(H + 1 - h)}$.
\end{proof}

\begin{lemma}[Value gap in the exponential space]\label{lem:value_gap_log}
\lean{Learning.MDP.Episodic.optEntropicValue_sub_entropicValue, Learning.MDP.Episodic.optEntropicValue_sub_entropicValue_le_of_pos, Learning.MDP.Episodic.optEntropicValue_sub_entropicValue_le_of_neg, Learning.MDP.Episodic.optEntropicValue_eq_inv_mul_log}\leanok
\uses{def:opt_value, def:entropic_value, lem:exp_value_range, lem:opt_exp_value_eq}
Assume rewards in $[0, 1]$ and $\beta \ne 0$. For every policy $\pi$ and state $s$,
\[
  V^*_1(s) - V^\pi_1(s) = \frac1\beta \log \frac{Z^*_1(s)}{Z^\pi_1(s)}
  = \frac1\beta \log\Big(1 + \frac{Z^*_1(s) - Z^\pi_1(s)}{Z^\pi_1(s)}\Big) ,
\]
and consequently, for $\eps \ge 0$:
\begin{enumerate}
  \item[(i)] if $\beta > 0$ and $Z^*_1(s) - Z^\pi_1(s) \le (e^{\beta\eps} - 1)\, Z^\pi_1(s)$, then
        $V^*_1(s) - V^\pi_1(s) \le \eps$;
  \item[(ii)] if $\beta < 0$ and $Z^\pi_1(s) - Z^*_1(s) \le (1 - e^{\beta\eps})\, Z^\pi_1(s)$, then
        $V^*_1(s) - V^\pi_1(s) \le \eps$.
\end{enumerate}
The same holds at every step $h$ in place of $1$.
\end{lemma}

\begin{proof}
\uses{def:entropic_value, def:opt_value, lem:exp_value_range, lem:opt_exp_value_eq}
\leanok
$Z^\pi_1(s) > 0$ and $Z^*_1(s) > 0$ (Lemma~\ref{lem:exp_value_range},
Lemma~\ref{lem:opt_exp_value_range} through Lemma~\ref{lem:opt_exp_value_eq}), so
$V^*_1 - V^\pi_1 = \beta^{-1}(\log Z^*_1 - \log Z^\pi_1) = \beta^{-1}\log(Z^*_1/Z^\pi_1)$ and
$Z^*_1/Z^\pi_1 = 1 + (Z^*_1 - Z^\pi_1)/Z^\pi_1$. (i) The hypothesis is $Z^*_1(s) \le e^{\beta\eps} Z^\pi_1(s)$,
so $\log(Z^*_1/Z^\pi_1) \le \beta\eps$ and dividing by $\beta > 0$ gives the claim. (ii) The
hypothesis is $e^{\beta\eps} Z^\pi_1(s) \le Z^*_1(s)$, i.e.\ $Z^\pi_1/Z^*_1 \le e^{-\beta\eps} = e^{\abs\beta\eps}$,
so $V^*_1 - V^\pi_1 = \abs\beta^{-1}\log(Z^\pi_1/Z^*_1) \le \eps$ (for $\beta < 0$,
$\beta^{-1}\log(Z^*/Z^\pi) = \abs\beta^{-1}\log(Z^\pi/Z^*)$).
\end{proof}

\textbf{Lean remark.} This is the identity behind the stopping rule of Section~4 of the
paper (``Stopping rule'') and of its Appendix~B.2; both directions of the sign of $\beta$ are
needed. The formal statements hold at every step \texttt{h : ℕ} and only assume
$Z^\pi_h(s) > 0$ (not the range of the rewards); $Z^*_h(s) = e^{\beta V^*_h(s)} > 0$ holds by
definition.

\section{The maximal return}
\label{sec:mdp_gmax}

\begin{definition}[Maximal achievable return]\label{def:gmax}
\lean{Learning.MDP.Episodic.maxReturn}\leanok
\uses{def:step_law, def:episode_return}
For an initial state $s_1$, the \emph{maximal achievable return} of $\mathcal M$ (Definition~1 of
the paper) is
\[
  \Gmax(\mathcal M) := \sup_{\pi}\ \operatorname{ess\,sup}_{\Prob^\pi_{(s_1, 1)}} R^\pi_1 ,
\]
the supremum over the policies of the essential supremum of the return of an episode under the
trajectory law of the policy from $(s_1, 1)$.
\end{definition}

\textbf{Lean remark.} \texttt{maxReturn M H s₁ := ⨆ π : \{π : ℕ → S → A // ∀ k, Measurable (π k)\}, essSup (episodeReturn H) (stepLaw M π (s₁, 0))}
(a supremum over the measurable policies, bounded by $H$ for rewards in $[0, 1]$:
\texttt{bddAbove\_essSup\_episodeReturn}; \texttt{essSup} in $\R$, a conditionally complete
lattice, is the Mathlib \texttt{MeasureTheory.essSup}). The paper's ``supremum over all sources of
randomness'' is the essential supremum under the trajectory law: the two agree for the
finitely supported laws of the paper (finite state space, deterministic rewards) and the
essential supremum is the version that makes sense for reward kernels.

\begin{lemma}[The maximal return is between $0$ and $H$]\label{lem:gmax_le_horizon}
\lean{Learning.MDP.Episodic.maxReturn_nonneg, Learning.MDP.Episodic.maxReturn_le, Learning.MDP.Episodic.ae_episodeReturn_le_maxReturn, Learning.MDP.Episodic.essSup_episodeReturn_mem_Icc}\leanok
\uses{def:gmax, lem:exp_value_range}
For rewards in $[0, 1]$: $0 \le \Gmax(\mathcal M) \le H$, and for every policy $\pi$,
$R^\pi_1 \le \Gmax(\mathcal M)$ $\Prob^\pi_{(s_1, 1)}$-almost surely.
\end{lemma}

\begin{proof}
\uses{def:gmax, lem:exp_value_range}
\leanok
By Lemma~\ref{lem:exp_value_range}, $R^\pi_1 \in [0, H]$ a.s.\ under every $\Prob^\pi_{(s_1, 1)}$, so
each essential supremum lies in $[0, H]$ (the essential supremum of a function bounded a.s.\ by
$H$ is at most $H$, and it is at least $0$ since $R^\pi_1 \ge 0$ a.s.\ and the measure is a
probability measure); the supremum over the finite set of policies is then in $[0, H]$. The
last claim is the defining property of the essential supremum
($R^\pi_1 \le \operatorname{ess\,sup} R^\pi_1$ a.s., Mathlib \texttt{ae\_le\_essSup}) together with
$\operatorname{ess\,sup} R^\pi_1 \le \Gmax(\mathcal M)$ (\texttt{le\_ciSup}).
\end{proof}

\section{Episodes as rounds, BPI algorithms and the PAC property}
\label{sec:mdp_bpi}

The learner interacts with $\mathcal M$ in the online model: at each episode it chooses a
policy, plays it from the fixed initial state $s_1$ and observes the states of the episode.
Since the rewards are known and deterministic (Definition~\ref{def:reward_fn}), the state
sequence determines the episode; the whole interaction is an LML algorithm-environment
sequence in which one round is one episode.

\begin{definition}[Episode environment]\label{def:episode_env}
\lean{Learning.MDP.Episodic.statesLaw, Learning.MDP.Episodic.statesKernel, Learning.MDP.Episodic.statesEnv}\leanok
\uses{def:step_law, def:episode_return, def:policy}
Fix an initial state $s_1$. The \emph{episode environment} of $\mathcal M$ from $s_1$ is the
stationary LML environment (\texttt{stationaryEnv}) without observations, whose actions are the
policies $\pi \in \As^{\{1, \dots, H\} \times \Ss}$ and whose feedback to the action $\pi$ is the
state sequence $(S_1, \dots, S_{H+1})$ of an episode of $\pi$ from $s_1$, drawn from the
\emph{states kernel} $\nu_{\mathcal M}(\pi) :=$ the law of $(S_1, \dots, S_{H+1})$ under
$\Prob^\pi_{(s_1, 1)}$. A run of an algorithm in this environment is a sequence
$(X_t, Y_t)_{t \ge 0}$ where $X_t = \pi^{t+1}$ is the policy of the episode $t + 1$ and
$Y_t = (s^{t+1}_1, \dots, s^{t+1}_{H+1})$ its state sequence; we write $a^{t+1}_h := \pi^{t+1}_h(s^{t+1}_h)$
for the action played at step $h$ of the episode $t + 1$.
\end{definition}

\textbf{Lean remark.} \texttt{statesLaw M H s₁ π := (stepLaw M π (s₁, 0)).map (episodeStates H)}
for a policy \texttt{π : ℕ → S → A}, \texttt{statesKernel M H s₁ : Kernel (Policy S A H) (Traj S H)},
\texttt{π ↦ statesLaw M H s₁ π.extend} (measurable on the countable discrete type of
$H$-step policies), and \texttt{statesEnv M H s₁ := stationaryEnv (statesKernel M H s₁) : Environment Unit (Policy S A H) (Traj S H)}.
The observation type is \texttt{Unit}. The rewards are not part of the feedback: an
algorithm for this environment is a learner who knows $r$; this is the reason why the PAC
property below is stated for the class $\mathcal M_r$.

\begin{lemma}[Law of an episode given the past]\label{lem:episode_env_law}
\lean{Learning.MDP.Episodic.hasCondDistrib_feedback_history_action_statesEnv, Learning.MDP.Episodic.hasCondDistrib_feedback_statesEnv}\leanok
\uses{def:episode_env, lem:step_law_markov}
Let $(X, Y)$ be an algorithm-environment sequence for any algorithm and the episode
environment of $\mathcal M$ from $s_1$, on $(\Omega, \Prob)$. For every $t$, the conditional law of
$Y_t$ given the history $(X_i, Y_i)_{i < t}$ and $X_t$ is $\nu_{\mathcal M}(X_t)$: for every
bounded measurable $F$ on $\Ss^{H+1}$,
$\E[F(Y_t) \mid (X_i, Y_i)_{i < t}, X_t] = \E^{X_t}_{(s_1, 1)}[F(S_1, \dots, S_{H+1})]$ a.s. The
within-episode form (the law of $s^{t+1}_{h+1}$ given the past and $s^{t+1}_{\le h}$) is
Lemma~\ref{lem:pmdp_episode_env_step}.
\end{lemma}

\begin{proof}
\uses{def:episode_env}
\leanok
This is the defining property of an algorithm-environment sequence for a stationary environment
(LML \texttt{IsObliviousEnv.hasCondDistrib\_feedback\_history\_action} together with
\texttt{feedbackCondAction\_stationaryEnv}: the feedback kernel of \texttt{stationaryEnv ν} at every
round is $\nu$ applied to the action).
\end{proof}

\textbf{Lean remark.} \texttt{hasCondDistrib\_feedback\_history\_action\_statesEnv}:
\texttt{HasCondDistrib (Y t) (fun ω ↦ ((history O X Y t ω, O t ω), X t ω)) ((statesKernel M s₁).prodMkLeft \_) P},
and \texttt{hasCondDistrib\_feedback\_statesEnv} given the policy only.

\begin{definition}[BPI algorithms, runs and stopping times]\label{def:bpi_alg}
\lean{Learning.MDP.Episodic.BPIAlg}\leanok
\uses{def:episode_env}
A \emph{best-policy identification (BPI) algorithm} is an LML identification algorithm
(\texttt{IdentAlg}) in the episode environment: a triple $\mathcal A = (\mathrm{alg}, \mathcal T, \rho)$
of a \emph{sampling rule} $\mathrm{alg}$, an LML algorithm with actions in the policies and
feedback in $\Ss^{H+1}$ (at the round $t$ it chooses the policy $\pi^{t+1}$ of the episode
$t + 1$ as a function, possibly randomized, of the history of the first $t$ episodes), a
\emph{stopping rule} $\mathcal T$, a measurable set of finite histories of variable length
(``stop after $t$ episodes when their history belongs to $\mathcal T$''), and an \emph{output
rule} $\rho$, a Markov kernel from histories of variable length to policies. A \emph{run} of
$\mathcal A$ in the episode environment of $\mathcal M$ from $s_1$ on $(\Omega, \Prob)$ is a triple
$(X, Y, \hat\pi)$ where $(X, Y)$ is an algorithm-environment sequence for $\mathrm{alg}$ and the
episode environment and the output $\hat\pi$ has conditional law $\rho(\hist_\tau)$ given the
history $\hist_\tau$ at the stopping time. The \emph{number of episodes} (sample complexity) is
the stopping time
\[
  \tau := \min\{t \in \N : (X_i, Y_i)_{i < t} \in \mathcal T\} \in \N \cup \{\infty\} ,
\]
a stopping time of the history filtration (LML \texttt{isStoppingTime\_stoppingTime}); the law of
$(\hist_\tau, \hat\pi)$ is the same in all runs, and we write $\Prob_{\mathcal M}$ for it when a
statement only depends on it (LML \texttt{IdentAlg.outputMeasure} for the law of $\hat\pi$).
\end{definition}

\textbf{Lean remark.} \texttt{BPIAlg S A H := IdentAlg Unit (Policy S A H) (Traj S H) (Policy S A H)};
the stopping rule is the field \texttt{stopSet : Set (Σ n, Hist Unit (Policy S A H) (Traj S H) n)},
the output rule \texttt{output : Kernel (Σ n, Hist \_ \_ \_ n) (Policy S A H)}, the stopping time
\texttt{IdentAlg.stoppingTime A O X Y : Ω → ℕ∞} (\texttt{hittingAfter} of the process of
histories \texttt{sigmaHistory}), runs are \texttt{IdentAlg.IsRun A env O X Y out P} (with
\texttt{O : ℕ → Ω → Unit} trivial), and \texttt{IdentAlg.stoppedHist} is $\hist_\tau$
(\texttt{stoppedHist\_mem\_stopSet\_of\_ne\_top}: it belongs to the stopping rule when
$\tau < \infty$). When $\tau = \infty$ the stopped history is a history of an arbitrary
number of rounds; all statements about the output are on the event $\{\tau < \infty\}$ or
about the law of the output.

\begin{definition}[$(\eps, \delta)$-PAC algorithms for entropic BPI]\label{def:entropic_pac}
\lean{Learning.MDP.Episodic.IsEntropicPAC}\leanok
\uses{def:bpi_alg, def:opt_value, def:entropic_value, def:reward_fn}
Let $r$ be a reward function, $\beta \ne 0$, $\eps > 0$, $\delta \in (0, 1)$ and $s_1 \in \Ss$. A
BPI algorithm $\mathcal A$ is \emph{$(\eps, \delta)$-PAC for entropic BPI with reward function
$r$, parameter $\beta$ and initial state $s_1$} (Definition~2 of the paper) if for every MDP
$\mathcal M \in \mathcal M_r$, the output $\hat\pi$ of $\mathcal A$ in the episode environment of
$\mathcal M$ from $s_1$ satisfies
\[
  \Prob_{\mathcal M}\big( V^*_1(s_1) - V^{\hat\pi}_1(s_1) > \eps \big) \le \delta ,
\]
where the values are those of $\mathcal M$. Equivalently: in every run of $\mathcal A$ in the episode
environment of every $\mathcal M \in \mathcal M_r$, $\Prob(V^*_1(s_1) - V^{\hat\pi}_1(s_1) \le \eps) \ge 1 - \delta$.
\end{definition}

\textbf{Lean remark.} \texttt{IsEntropicPAC 𝒜 r β ε δ s₁ := 𝒜.IsPAC (fun M : \{M // M.HasRewardFn r\} ↦ statesEnv M.1 s₁) (fun M π ↦ ε < optEntropicValue M.1 β 0 s₁ - entropicValue M.1 β π 0 s₁) δ}.
LML's \texttt{IdentAlg.IsPAC env bad δ} says that, for every parameter $\theta$ (here every
$\mathcal M \in \mathcal M_r$), the law \texttt{outputMeasure} of the output in \texttt{env θ} gives
mass at most $\delta$ to the bad outputs; for a run $(X, Y, \hat\pi)$ on $(\Omega, \Prob)$,
\texttt{IsPAC.measureReal\_bad\_of\_isRun} gives $\Prob(\hat\pi \text{ bad}) \le \delta$
(\texttt{IsRun.hasLaw\_output}: the output of a run has law \texttt{outputMeasure}). Conversely,
to prove that an algorithm is PAC one proves the bound for one run, the canonical run on the
Ionescu--Tulcea space (LML \texttt{IdentAlg.runMeasure}, \texttt{isRun\_runMeasure}), which
is what Theorem~\ref{thm:upper_bound_pac} does. Two points of the encoding are decisions
recorded in the overview (deviation~7): the PAC property is a property of the \emph{algorithm}
(quantified over the environments of the class), not of a run, and the class is
$\mathcal M_r$ for the \emph{known} reward function $r$ of the algorithm and arbitrary
transition kernels; the lower bound (Theorem~\ref{thm:lower_bound}) only assumes PAC for the
class with the reward function of its hard instances, a weaker hypothesis than PAC for all
reward functions, hence a stronger theorem. Note also that the paper's Definition~2 does not
require $\tau < \infty$: an algorithm that never stops has an output law given by the output
kernel applied to an infinite history in LML, and finiteness of $\tau$ is a separate
conclusion (Theorem~\ref{thm:upper_bound_complexity}).

\chapter{Empirical transitions and concentration events}
\label{chap:empirical}

This chapter defines the statistics that any learner in the episode environment computes from
the observed episodes (Section~2 of \cite{essakine2026tight}, ``Empirical MDP''), the
exploration rates and the three concentration events of Appendix~A of the paper, and proves
Lemma~5 of the paper: the good event has probability at least $1 - \delta$.

\textbf{Standing setting.} $\mathcal M$ is a finite-horizon MDP with a reward function $r$
(Definition~\ref{def:reward_fn}), $s_1$ an initial state, and $(X, Y)$ is an
algorithm-environment sequence for an \emph{arbitrary} algorithm and the episode environment of
$\mathcal M$ from $s_1$ (Definition~\ref{def:episode_env}), on a probability space
$(\Omega, \Prob)$: $X_t = \pi^{t+1}$ is the policy of the episode $t + 1$ and
$Y_t = (s^{t+1}_1, \dots, s^{t+1}_{H+1})$ its state sequence, $a^{t+1}_h := \pi^{t+1}_h(s^{t+1}_h)$.
The \emph{history of the first $t$ episodes} is $\hist_t := (X_i, Y_i)_{i < t}$ and we write
$\mathcal F_t := \sigma(\hist_t, X_t)$ for the $\sigma$-algebra generated by the history of the
first $t$ episodes together with the policy of the episode $t + 1$ ($t \ge 0$; the filtration
$(\mathcal F_t)_{t \ge 0}$ is nondecreasing since $X_t$ is part of $\hist_{t+1}$). Nothing in this
chapter depends on the algorithm; Chapter~\ref{chap:algorithm} specializes to Entropic-BPI,
for which $X_t$ is a function of $\hist_t$ and $\mathcal F_t = \sigma(\hist_t)$.

\textbf{Lean remark.} The history of the first $t$ episodes of the run is
\texttt{histAt X Y t ω : Hist Unit (Policy S A H) (Traj S H) t} (\texttt{Run.lean}), LML's
\texttt{history O X Y t} with the trivial observations; all the quantities below are first
defined as functions of a history \texttt{hist : Hist Unit (Policy S A H) (Traj S H) t}
(\texttt{Empirical.lean}, namespace \texttt{Learning.MDP}) and evaluated on the run through
\texttt{histAt}. The $\sigma$-algebra $\mathcal F_t$ is generated by the pair
\texttt{(histAt X Y t, X t)}; LML's history filtration \texttt{IsAlgEnvSeq.filtration n} is
generated by the rounds $0, \dots, n$, i.e.\ by $\hist_{n+1}$, so that $\mathcal F_t$ sits between
\texttt{filtration (t - 1)} and \texttt{filtration t}. The concentration statements are
formulated with conditional distributions (\texttt{HasCondDistrib}) given
\texttt{(histAt X Y t, X t)}, which avoids naming the filtration.

\section{Counts and empirical transitions}
\label{sec:empirical_counts}

\begin{definition}[Visit counts]\label{def:counts}
\lean{Learning.MDP.EmpiricalModel, Learning.MDP.EmpiricalModel.ofEpisodeAt, Essakine2026Tight.stepModel, Essakine2026Tight.visitCount, Essakine2026Tight.histAt, Essakine2026Tight.visitCountAt}\leanok
\uses{def:episode_env}
For $t \ge 0$, $h \in \{1, \dots, H\}$, $(s, a) \in \Ss \times \As$ and $s' \in \Ss$, the
\emph{visit counts} after $t$ episodes are
\[
  n_h^t(s, a) := \sum_{i = 1}^{t} \indic\{(s^i_h, a^i_h) = (s, a)\},
  \qquad
  n_h^t(s, a, s') := \sum_{i = 1}^{t} \indic\{(s^i_h, a^i_h, s^i_{h+1}) = (s, a, s')\},
\]
where $a^i_h = \pi^i_h(s^i_h)$. They are functions of the history $\hist_t$ of the first $t$
episodes (hence $\mathcal F_t$-measurable, and even $\sigma(\hist_t)$-measurable),
$n_h^t(s, a) = \sum_{s'} n_h^t(s, a, s')$, $n_h^0 = 0$, and $n_h^{t+1}(s, a) - n_h^t(s, a) \in \{0, 1\}$.
\end{definition}

\textbf{Lean remark.} \texttt{EmpiricalModel S A := (S × A → ℕ) × (S × A → ℝ) × (S × A → S → ℕ)}
(visit counts, reward sums, transition counts; an additive commutative monoid, so that the
model of a history is a finite sum), \texttt{EmpiricalModel.ofEpisodeAt π τ r h} the model of
the single transition at step \texttt{h} of the episode with policy \texttt{π}, states
\texttt{τ} and rewards \texttt{r}, \texttt{stepModel hist h := ∑ i, ofEpisodeAt (hist i).action (hist i).feedback 0 h}
(the rewards are not observed, hence set to $0$), \texttt{visitCount hist h s a := (stepModel hist h).count (s, a)};
on the run, \texttt{visitCountAt X Y h s a t ω := visitCount (histAt X Y t ω) h s a}.
Lemma~\ref{lem:count_le_episodes} ($n_h^t(s, a) \le t$, $\sum_{s, a} n_h^t(s, a) = t$) is in
Chapter~\ref{chap:pre_mdp}.

\begin{definition}[Empirical transitions]\label{def:emp_trans}
\lean{Essakine2026Tight.empTrans, Essakine2026Tight.empTransAt}\leanok
\uses{def:counts}
The \emph{empirical transition probabilities} after $t$ episodes are
\[
  \phat_h^t(s' \mid s, a) :=
  \begin{cases}
    n_h^t(s, a, s') / n_h^t(s, a) & \text{if } n_h^t(s, a) > 0, \\
    1/S & \text{otherwise},
  \end{cases}
\]
and $\phat_h^t(s, a) := \phat_h^t(\cdot \mid s, a) \in \Sigma_{\Ss}$ denotes the corresponding
probability vector; as in Definition~\ref{def:episodic_mdp}, $(\phat_h^t f)(s, a) := \sum_{s'} \phat_h^t(s' \mid s, a) f(s')$
and $\Var_{\phat_h^t}(f)(s, a) := \phat_h^t\big((f - \phat_h^t f(s, a))^2\big)(s, a)$.
\end{definition}

\textbf{Lean remark.} \texttt{empTrans hist h s a : S → ℝ} (\texttt{if visitCount hist h s a = 0 then fun \_ ↦ (Fintype.card S : ℝ)⁻¹ else (stepModel hist h).empTrans (s, a)}),
\texttt{empTransAt X Y h s a t ω := empTrans (histAt X Y t ω) h s a}. The value $1/S$ for a
never-visited pair is the paper's convention; it never matters, since every use of $\phat$
for a pair with $n = 0$ is replaced by a trivial bound (deviation~1 of the overview).

\begin{lemma}[The empirical transitions form a probability vector]\label{lem:emp_trans_simplex}
\uses{def:emp_trans}
For all $t, h, (s, a)$: $\phat_h^t(\cdot \mid s, a)$ is a probability vector on $\Ss$
(nonnegative entries summing to $1$). Consequently $f \mapsto (\phat_h^t f)(s, a)$ is linear,
positive and monotone, $\phat_h^t \mathbf 1 = 1$, $\abs{(\phat_h^t f)(s, a)} \le \max_{s'} \abs{f(s')}$,
$\Var_{\phat_h^t}(f)(s, a) \ge 0$, and if $f$ takes values in $[m, M]$ then so does $(\phat_h^t f)(s, a)$.
\end{lemma}

\begin{proof}
\uses{def:emp_trans, def:counts}
If $n_h^t(s, a) > 0$, the entries $n_h^t(s, a, s')/n_h^t(s, a)$ are nonnegative and sum to
$\sum_{s'} n_h^t(s, a, s')/n_h^t(s, a) = 1$ (Definition~\ref{def:counts}); otherwise the entries are
$1/S$, which sum to $1$. The consequences hold for every probability vector $q$: linearity,
positivity and $q \mathbf 1 = 1$ are immediate from $q f = \sum_{s'} q(s') f(s')$ with
$q(s') \ge 0$, $\sum q(s') = 1$; the bound and the preservation of $[m, M]$ follow from the
convexity of the combination; $\Var_q(f) = q(f - qf)^2 \ge 0$ by positivity.
\end{proof}

\textbf{Lean remark.} \texttt{isProbVec\_empTrans}, then the \texttt{vecExp}/\texttt{vecVar}
lemmas of \texttt{Vec.lean} (Lemma~\ref{lem:exp_value_monotone}) apply to
\texttt{empTrans hist h s a}. The KL divergence of Definition~\ref{def:event_kl} views this
vector as a measure through \texttt{weightedMeasure : (S → ℝ) → Measure S}
(\texttt{Essakine2026Tight/Mathlib/}), $\sum_{s'} q(s') \delta_{s'}$.

\section{Pseudo-counts}
\label{sec:empirical_pseudo}

\begin{definition}[Pseudo-counts]\label{def:pseudo_counts}
\lean{Essakine2026Tight.pseudoCount}\leanok
\uses{def:counts, def:occupancy, lem:episode_env_law}
The \emph{pseudo-count} of $(s, a)$ at step $h$ after $t$ episodes is the sum of the conditional
probabilities of the visits given the past:
\[
  \nbar_h^t(s, a) := \sum_{i = 1}^{t} p^{\pi^i}_h(s, a),
  \qquad
  p^\pi_h(s, a) := \Prob^\pi_{(s_1, 1)}(S_h = s, A_h = a)
\]
(the occupancy measure of $\pi$, Definition~\ref{def:occupancy}). It is a random variable, a
function of the policies $\pi^1, \dots, \pi^t$ played, hence $\mathcal F_{t-1}$-measurable for
$t \ge 1$ (predictable), and $\nbar_h^0 = 0$.
\end{definition}

\begin{remark}[Deviation 10: random pseudo-counts]\label{rem:empirical_pseudo_counts}
The paper defines the pseudo-counts as the expectations $\E[n_h^t(s, a)]$ (following
\cite{menard2021fast}), and then uses two properties: Lemma~20 of the paper (the Bernoulli
maximal inequality, Theorem~\ref{thm:bernoulli_concentration}), which controls
$n_h^t(s, a)$ by $\sum_{i \le t} \Prob(\text{visit at episode } i \mid \text{past})$, a
\emph{random} sum; and the counting argument of its Appendix~B.1 (Lemma~\ref{lem:sum_counts_bound}),
which telescopes the increments $\nbar_h^{t+1} - \nbar_h^t = p^{\pi^{t+1}}_h(s, a)$, again the random
conditional visit probabilities. Both are properties of the random predictable sum, not of its
expectation (the two differ as soon as the algorithm is adaptive). The blueprint therefore
\emph{defines} the pseudo-counts as the random sums; Lemma~\ref{lem:pseudo_counts_increments}(iv)
shows that the paper's quantity is their expectation. Since the policies are deterministic,
the paper's ``$\nu_{\hat\pi}(u) > 1/2$'' in the lower bound is ``$\hat\pi$ realizes $u$''.
\end{remark}

\textbf{Lean remark.} \texttt{pseudoCount M s₁ X h s a t ω := ∑ i ∈ Finset.range t, occupancy M (X i ω) s₁ h s a}
(the outline's \texttt{pseudoCount X Y P h s a t}, now depending on the MDP through the
occupancy measures and not on the measure \texttt{P}; the expectation version
$\E[n_h^t(s, a)]$ is not used).

\begin{lemma}[Increments of the pseudo-counts]\label{lem:pseudo_counts_increments}
\uses{def:pseudo_counts, def:counts, lem:episode_env_law, def:occupancy, lem:occupancy_sum_one}
For all $t \ge 0$, $h \le H$ and $(s, a)$:
\begin{enumerate}
  \item[(i)] $\nbar_h^{t+1}(s, a) - \nbar_h^t(s, a) = p^{\pi^{t+1}}_h(s, a) \in [0, 1]$, with
        $p^{\pi}_h(s, a) = \indic\{a = \pi_h(s)\}\, \Prob^\pi_{(s_1, 1)}(S_h = s)$;
  \item[(ii)] $0 \le \nbar_h^t(s, a) \le t$ and $\sum_{s, a} \nbar_h^t(s, a) = t$;
  \item[(iii)] $\Prob\big((s^{t+1}_h, a^{t+1}_h) = (s, a) \mid \mathcal F_t\big) = p^{\pi^{t+1}}_h(s, a)$
        almost surely;
  \item[(iv)] $\E[n_h^t(s, a)] = \E[\nbar_h^t(s, a)]$: the paper's pseudo-count is the expectation
        of $\nbar_h^t(s, a)$.
\end{enumerate}
\end{lemma}

\begin{proof}
\uses{def:pseudo_counts, def:counts, lem:episode_env_law, lem:occupancy_sum_one}
(i) is the definition; the formula for $p^\pi_h(s, a)$ and $\sum_{s, a} p^\pi_h(s, a) = 1$ are
Lemma~\ref{lem:occupancy_sum_one}, which gives $p^\pi_h(s, a) \in [0, 1]$. (ii) follows from
(i) by induction on $t$. (iii) By Lemma~\ref{lem:episode_env_law}, the conditional law of
$Y_t = (s^{t+1}_1, \dots, s^{t+1}_{H+1})$ given $\mathcal F_t = \sigma(\hist_t, X_t)$ is the law of
the state sequence of the policy $X_t = \pi^{t+1}$ from $s_1$; the event
$\{(s^{t+1}_h, a^{t+1}_h) = (s, a)\} = \{s^{t+1}_h = s\} \cap \{\pi^{t+1}_h(s) = a\}$ is the
event $\{S_h = s, A_h = a\}$ of this state sequence (with $A_h = \pi^{t+1}_h(S_h)$), whose
probability under $\Prob^{\pi^{t+1}}_{(s_1, 1)}$ is $p^{\pi^{t+1}}_h(s, a)$. (iv) Taking
expectations in (iii) (tower rule), $\E[\indic\{(s^{i}_h, a^{i}_h) = (s, a)\}] = \E[p^{\pi^{i}}_h(s, a)]$
for every $i \ge 1$; sum over $i = 1, \dots, t$.
\end{proof}

\section{Exploration rates}
\label{sec:empirical_rates}

\begin{definition}[Exploration rates]\label{def:rates}
\lean{Essakine2026Tight.alphaKL, Essakine2026Tight.alphaStar, Essakine2026Tight.alphaCnt}\leanok
\uses{def:episodic_mdp}
For $\delta \in (0, 1)$ and $n \in \N$, with $S = \abs{\Ss}$, $A = \abs{\As}$ (Appendix~A of the
paper, eq.~(5)):
\[
  \alpha(n, \delta) := \log\frac{3 S A H}{\delta} + S \log\big(8 e (n + 1)\big),
  \qquad
  \alpha^*(n, \delta) := \log\frac{3 S A H}{\delta} + \log\big(8 e (n + 1)\big),
  \qquad
  \alpha^{\mathrm{cnt}}(\delta) := \log\frac{3 S A H}{\delta} .
\]
We write $\alpha(n)$, $\alpha^*(n)$ when $\delta$ is fixed.
\end{definition}

\textbf{Lean remark.} \texttt{alphaKL S A H δ n}, \texttt{alphaStar S A H δ n},
\texttt{alphaCnt S A H δ} (\texttt{Essakine2026Tight/EV2026/Rates.lean}), functions of the
cardinalities \texttt{S A : ℕ}, of \texttt{H : ℕ} and of \texttt{δ : ℝ}; they are applied
with \texttt{Fintype.card S}, \texttt{Fintype.card A}.

\begin{lemma}[Properties of the rates]\label{lem:rates_props}
\uses{def:rates}
For $\delta \in (0, 1)$ (more generally $\delta \in (0, 1]$):
\begin{enumerate}
  \item[(i)] $0 \le \alpha^{\mathrm{cnt}}(\delta) \le \alpha^*(n, \delta) \le \alpha(n, \delta)$ for all $n$,
        and $\alpha^{\mathrm{cnt}}(\delta) \ge \log 3 > 1$;
  \item[(ii)] $n \mapsto \alpha(n, \delta)$ and $n \mapsto \alpha^*(n, \delta)$ are nondecreasing;
  \item[(iii)] $n \mapsto \alpha(n, \delta)/n$ and $n \mapsto \alpha^*(n, \delta)/n$ are nonincreasing on
        $n \ge 1$: $\alpha(n, \delta)/n \le \alpha(n', \delta)/n'$ for $n \ge n' \ge 1$, and the same for $\alpha^*$;
  \item[(iv)] $\alpha(n, \delta) \le \log(3SAH/\delta) + S \log(8e(n + 1))$ and
        $\alpha^*(n, \delta) \le \log(3SAH/\delta) + \log(8e(n + 1))$ (equalities, recorded as the
        only form in which the rates enter the final bound).
\end{enumerate}
\end{lemma}

\begin{proof}
\uses{def:rates}
(i) $S, A, H \ge 1$ and $\delta \le 1$ give $3SAH/\delta \ge 3$, so $\alpha^{\mathrm{cnt}} \ge \log 3 > 1 > 0$;
$\log(8e(n + 1)) \ge \log(8e) > 0$ and $S \ge 1$ give the two other inequalities.
(ii) $n \mapsto \log(8e(n + 1))$ is nondecreasing. (iii) Write $\alpha(n)/n = c/n + S \log(n + 1)/n$
with $c := \log(3SAH/\delta) + S \log(8e) \ge 0$; $c/n$ is nonincreasing, and $x \mapsto \log(1 + x)/x$
is nonincreasing on $(0, \infty)$ because $x \mapsto \log(1 + x)$ is concave with value $0$ at
$0$ (for $0 < x' \le x$, concavity gives $\log(1 + x') \ge (x'/x)\log(1 + x) + (1 - x'/x)\log 1$).
The same argument applies to $\alpha^*$ with $S$ replaced by $1$. (iv) is the definition.
\end{proof}

\section{Concentration events}
\label{sec:empirical_events}

The three events below are defined for the run $(X, Y)$ and a confidence level
$\delta \in (0, 1)$; they are events of $\Omega$ (countable intersections of measurable sets).

\begin{definition}[KL concentration event]\label{def:event_kl}
\lean{Essakine2026Tight.eventKL, Essakine2026Tight.klThreshold}\leanok
\uses{def:emp_trans, def:counts, def:rates}
\[
  \evKL := \Big\{ \forall t \in \N,\ \forall h \le H,\ \forall (s, a) \in \Ss \times \As :\
    \KL\big(\phat_h^t(s, a) \,\big\|\, p_h(s, a)\big) \le \frac{\alpha(n_h^t(s, a), \delta)}{n_h^t(s, a)} \Big\},
\]
where $\KL(q \| p) := \sum_{s'} q(s') \log(q(s')/p(s'))$ for probability vectors ($+\infty$ if
$q$ is not absolutely continuous with respect to $p$), and the threshold is $+\infty$ when
$n_h^t(s, a) = 0$ (the paper's $\alpha(0)/0 = \infty$): the constraint only concerns visited pairs.
\end{definition}

\textbf{Lean remark.} \texttt{eventKL X Y M δ : Set Ω} is
\texttt{\{ω | ∀ t h s a, klDiv (weightedMeasure (empTransAt X Y h s a t ω)) (M.trans h (s, a)) ≤ klThreshold S A H δ (visitCountAt X Y h s a t ω)\}}
with \texttt{klThreshold S A H δ n : ℝ≥0∞ := if n = 0 then ⊤ else ENNReal.ofReal (alphaKL S A H δ n / n)}
and Mathlib's \texttt{InformationTheory.klDiv} (with values in \texttt{ℝ≥0∞}, infinite when the
first measure is not absolutely continuous with respect to the second; on $\evKL$ every
visited pair has $\phat_h^t(s, a) \ll p_h(s, a)$). The paper's Appendix~A writes the threshold
$\alpha(n_h^t(s, a), \delta)$ without the division by $n_h^t(s, a)$; the confidence sets of its
Appendix~B and Lemma~19 have the division, which is the intended form.

\begin{definition}[Bernstein concentration event]\label{def:event_bern}
\lean{Essakine2026Tight.eventBern}\leanok
\uses{def:emp_trans, def:counts, def:rates, def:episodic_mdp}
Let $f = (f_h)_{h \le H + 1}$ with $f_h : \Ss \to \R$ and let $b = (b_h)_{h \le H}$ with $b_h \ge 0$
(in the applications, $b_h$ bounds the oscillation $\max f_{h+1} - \min f_{h+1}$ of $f_{h+1}$).
\[
  \evBern(f, b) := \Big\{ \forall t, h, (s, a) \text{ with } n := n_h^t(s, a) > 0 :\
    \big| (\phat_h^t - p_h) f_{h+1}(s, a) \big|
    \le \sqrt{2 \Var_{p_h}(f_{h+1})(s, a)\, \frac{\alpha^*(n, \delta)}{n}} + 3 b_h \frac{\alpha^*(n, \delta)}{n} \Big\} .
\]
\end{definition}

\begin{remark}[Deviation 12: step-dependent ranges, non-strict inequality]\label{rem:empirical_bern_event}
The paper's event $\mathcal E_f$ uses a single range $b$ for all steps and a strict inequality.
Two changes. (a) The ranges are step-dependent: Lemmas~7 and~12 of the paper
(Lemmas~\ref{lem:concentration_pos} and~\ref{lem:concentration_neg}) apply the event to
$f = Z^*$ with the range $e^{\beta(H - h)}$ (resp.\ $1 - e^{\beta(H - h)}$) of $Z^*_{h+1}$ \emph{at
the step $h$}, which is what makes the bonus of step $h$ scale with $e^{\beta(H - h)}$; a
uniform $b = e^{\beta H}$ would lose this. (b) The inequality is non-strict: with a strict
inequality the event would be \emph{empty} in the case $\beta < 0$, because at the step
$h = H$ the function $Z^*_{H+1} \equiv 1$ is constant and $b_H = 1 - e^{0} = 0$, so both sides
vanish for every visited pair. The concentration inequality (Corollary~\ref{cor:bernstein_explicit})
gives the non-strict form, and the analysis only uses the event as an upper bound.
\end{remark}

\textbf{Lean remark.} \texttt{eventBern X Y M δ f b : Set Ω} with \texttt{f : Fin (H + 1) → S → ℝ}
and \texttt{b : Fin H → ℝ}: for all \texttt{t h s a} with \texttt{0 < visitCountAt X Y h s a t ω},
\texttt{|vecExp (empTransAt …) (f h.succ) - vecExp (M.transVec h s a) (f h.succ)| ≤ √(2 * vecVar (M.transVec h s a) (f h.succ) * (alphaStar … n / n)) + 3 * b h * (alphaStar … n / n)}.

\begin{definition}[Counts concentration event]\label{def:event_cnt}
\lean{Essakine2026Tight.eventCnt}\leanok
\uses{def:counts, def:pseudo_counts, def:rates}
\[
  \evCnt := \Big\{ \forall t \in \N,\ \forall h \le H,\ \forall (s, a) :\
    n_h^t(s, a) \ge \tfrac12\, \nbar_h^t(s, a) - \alpha^{\mathrm{cnt}}(\delta) \Big\} .
\]
\end{definition}

\textbf{Lean remark.} \texttt{eventCnt M s₁ X Y δ : Set Ω}, with the random pseudo-counts of
Definition~\ref{def:pseudo_counts} (the outline's \texttt{eventCnt X Y δ P} loses the
dependence on \texttt{P}).

\begin{definition}[The good event]\label{def:good_event}
\lean{Essakine2026Tight.goodEvent}\leanok
\uses{def:event_kl, def:event_bern, def:event_cnt, def:opt_value, lem:opt_exp_value_range}
Let $Z^* = (Z^*_h)_{h \le H + 1}$ be the optimal exponential values of $\mathcal M$
(Definition~\ref{def:opt_value}). The \emph{good event} is
\[
  \evGood^+ := \evKL \cap \evBern(Z^*, b^+) \cap \evCnt \quad (\beta > 0),
  \qquad
  \evGood^- := \evKL \cap \evBern(Z^*, b^-) \cap \evCnt \quad (\beta < 0),
\]
with $b^+_h := e^{\beta(H - h)}$ and $b^-_h := 1 - e^{\beta(H - h)}$ for $h \le H$. By
Lemma~\ref{lem:opt_exp_value_range}, $Z^*_{h+1}$ takes values in $[1, e^{\beta(H - h)}]$ for
$\beta > 0$ and in $[e^{\beta(H - h)}, 1]$ for $\beta < 0$: an interval of length at most $b^\pm_h$,
so that $b^\pm_h \ge \max Z^*_{h+1} - \min Z^*_{h+1}$ as required by Lemma~\ref{lem:event_bern_prob}.
\end{definition}

\textbf{Lean remark.} \texttt{goodEvent M s₁ X Y β δ := eventKL X Y M δ ∩ eventBern X Y M δ (optExpValue M β) (fun h ↦ if 0 < β then Real.exp (β * k) else 1 - Real.exp (β * k)) ∩ eventCnt M s₁ X Y δ}
with \texttt{k = H - 1 - h} for the Lean step \texttt{h : Fin H} (the paper's $H - h$ for its
$1$-indexed step); both signs in one definition, the analysis chapters instantiating
\texttt{0 < β} or \texttt{β < 0}.

\section{Probability of the good event}
\label{sec:empirical_good_event}

The three probability bounds share one reduction: for a fixed triple $(h, s, a)$, the next
states observed at the successive visits of $(s, a)$ at step $h$ behave like an i.i.d.\
sample from $p_h(\cdot \mid s, a)$ (Lemma~\ref{lem:visits_iid}), and the empirical transition
after $t$ episodes is the empirical distribution of the first $n_h^t(s, a)$ of them
(Lemma~\ref{lem:emp_trans_reindex}). The time-uniform inequalities are then applied to this
sample, indexed by the number $n$ of visits, and transferred to all episodes $t$ at once.
Throughout this section, for a fixed $(h, s, a)$ we write $p := p_h(\cdot \mid s, a)$,
\[
  T_k := \min\{t \ge 1 : n_h^t(s, a) \ge k\} \in \N \cup \{\infty\}
  \quad\text{(the episode of the $k$-th visit)},
  \qquad
  W_k := s^{T_k}_{h+1} \text{ on } \{T_k < \infty\}
\]
(the next state observed at the $k$-th visit), and $\hat q_n(x_1, \dots, x_n) := \frac1n \sum_{k \le n} \delta_{x_k}$
for the empirical distribution of a finite sequence in $\Ss$. Each $T_k$ is a stopping time of
$(\mathcal F_t)$ (since $n_h^t(s, a)$ is $\mathcal F_t$-measurable), $T_k \le T_{k+1}$, and
$T_{n_h^t(s, a)} \le t$.

\begin{lemma}[Domination of the observed transitions by an i.i.d.\ sample, time-uniform form]\label{lem:empirical_visits_domination}
\uses{lem:visits_iid, def:counts, def:episode_env}
Fix $(h, s, a)$ and let $\xi_1, \xi_2, \dots$ be i.i.d.\ with law $p = p_h(\cdot \mid s, a)$ on some
auxiliary probability space $(\Omega', \Prob')$. Then:
\begin{enumerate}
  \item[(i)] for every $n \ge 1$ and every set $B \subseteq \Ss^n$,
        $\Prob\big(T_n < \infty,\ (W_1, \dots, W_n) \in B\big) \le \Prob'\big((\xi_1, \dots, \xi_n) \in B\big) = p^{\otimes n}(B)$;
  \item[(ii)] for every sequence of sets $B_n \subseteq \Ss^n$, $n \ge 1$,
        \[
          \Prob\big(\exists n \ge 1 :\ T_n < \infty,\ (W_1, \dots, W_n) \in B_n\big)
          \le \Prob'\big(\exists n \ge 1 :\ (\xi_1, \dots, \xi_n) \in B_n\big) .
        \]
\end{enumerate}
\end{lemma}

\begin{proof}
\uses{lem:visits_iid, def:counts}
(i) Lemma~\ref{lem:visits_iid}(i) gives the inequality for product sets
$B = B^1 \times \dots \times B^n$: $\Prob(W_1 \in B^1, \dots, W_n \in B^n, T_n < \infty) \le \prod_i p(B^i)$.
Since $\Ss^n$ is finite, an arbitrary $B$ is the disjoint union of its singletons, which are
product sets, and both sides are additive in $B$.

(ii) First-entrance decomposition. For $n \ge 1$ let
$B'_n := \{x \in \Ss^n : x \in B_n \text{ and } (x_1, \dots, x_k) \notin B_k \text{ for all } k < n\}$
and $E_n := \{T_n < \infty, (W_1, \dots, W_n) \in B'_n\}$. On $\{T_n < \infty\}$ all $T_k$, $k \le n$,
are finite, so $E_n$ is the event that $n$ is the \emph{first} index at which
``$T_n < \infty$ and $(W_1, \dots, W_n) \in B_n$'' holds (if it holds at $k < n$ then $T_k < \infty$
and $(W_1, \dots, W_k) \in B_k$, excluded by $B'_n$; conversely the first index at which it holds
is in some $E_n$). Hence the $E_n$ are pairwise disjoint and their union is the event of the
left-hand side, so its probability is $\sum_n \Prob(E_n) \le \sum_n p^{\otimes n}(B'_n)$ by (i).
Under $\Prob'$, the events $\{(\xi_1, \dots, \xi_n) \in B'_n\}$ are, by the same argument, the
pairwise disjoint events ``$n$ is the first index with $(\xi_1, \dots, \xi_n) \in B_n$'', whose
union is $\{\exists n, (\xi_1, \dots, \xi_n) \in B_n\}$; so
$\sum_n p^{\otimes n}(B'_n) = \Prob'(\exists n, (\xi_1, \dots, \xi_n) \in B_n)$.
\end{proof}

\textbf{Lean remark.} The auxiliary space can be taken to be $\Ss^{\N}$ with the product measure
\texttt{Measure.infinitePi (fun \_ ↦ weightedMeasure p)} (Mathlib \texttt{Measure.infinitePi}),
and the events of (ii) are measurable as countable unions. The lemma is the bridge between
the i.i.d.\ statements of Chapter~\ref{chap:pre_concentration} and the run; it is the
episodic analogue of LML's \texttt{Online/Bandit/RewardByCountMeasure.lean} (the ``arm'' is
$(h, s, a)$ and the ``reward'' the next state), to be contributed to LML with
Lemma~\ref{lem:visits_iid}.

\begin{lemma}[Time-uniform KL concentration of the empirical distribution]\label{lem:empirical_kl_time_uniform}
\uses{lem:ville, lem:num_types}
Let $p$ be a probability vector on a finite set $\mathcal X$ with $m := \abs{\mathcal X} \ge 1$
elements, $\xi_1, \xi_2, \dots$ i.i.d.\ with law $p$, and $\hat q_n := \hat q_n(\xi_1, \dots, \xi_n)$.
For every $\delta' \in (0, 1)$,
\[
  \Prob\Big( \exists n \ge 1 :\ n\, \KL(\hat q_n \| p) \ge \log\frac1{\delta'} + (m - 1)\log(n + 1) + 1 + \tfrac12 \log n \Big) \le \delta' ,
\]
and in particular
$\Prob\big( \exists n \ge 1 : \KL(\hat q_n \| p) \ge (\log(1/\delta') + m \log(e(n + 1)))/n \big) \le \delta'$.
\end{lemma}

\begin{proof}
\uses{lem:ville, lem:num_types}
\emph{Support.} Almost surely all $\xi_k$ lie in the support $\mathcal X_0 := \{x : p(x) > 0\}$;
on this event $\hat q_n$ is supported in $\mathcal X_0$ and $\KL(\hat q_n \| p) = \sum_{x \in \mathcal X_0} \hat q_n(x)\log(\hat q_n(x)/p(x)) < \infty$.
Replacing $\mathcal X$ by $\mathcal X_0$ ($m_0 := \abs{\mathcal X_0} \le m$ elements) only decreases
the threshold, so we assume $p(x) > 0$ for all $x$. Write $N_n(x) := n \hat q_n(x) = \#\{k \le n : \xi_k = x\}$,
so that, with $0 \log 0 := 0$,
\[
  n\, \KL(\hat q_n \| p) = \sum_x N_n(x) \log\frac{N_n(x)}{n} - \sum_{k \le n} \log p(\xi_k)
\]
(since $\sum_{k \le n} \log p(\xi_k) = \sum_x N_n(x) \log p(x)$).

\emph{A mixture martingale.} Let $\nu$ be the uniform probability measure on the simplex
$\Delta := \{q \in [0, 1]^{\mathcal X} : \sum_x q(x) = 1\}$ (normalized Lebesgue measure on the
$(m - 1)$-dimensional simplex, whose Lebesgue measure is $1/(m - 1)!$; the point mass if $m = 1$)
and
\[
  M_n := \int_\Delta \prod_{k \le n} \frac{q(\xi_k)}{p(\xi_k)}\, \nu(dq)
  = \Big(\prod_{k \le n} p(\xi_k)\Big)^{-1} \int_\Delta \prod_x q(x)^{N_n(x)}\, \nu(dq),
  \qquad M_0 = 1 .
\]
For a fixed $q$, $L_n(q) := \prod_{k \le n} q(\xi_k)/p(\xi_k)$ is a nonnegative martingale for the
natural filtration $\mathcal G_n := \sigma(\xi_1, \dots, \xi_n)$ with $L_0(q) = 1$:
$\E[L_{n+1}(q) \mid \mathcal G_n] = L_n(q)\, \E[q(\xi)/p(\xi)] = L_n(q) \sum_x q(x) = L_n(q)$. By
Fubini for nonnegative integrands, $M_n = \int L_n(q)\, \nu(dq)$ is a nonnegative martingale with
$M_0 = 1$, and Ville's inequality (Lemma~\ref{lem:ville}) gives $\Prob(\exists n, M_n \ge 1/\delta') \le \delta'$.

\emph{Lower bound on $M_n$.} The Dirichlet integral gives, for nonnegative integers $N(x)$ with
$\sum_x N(x) = n$, $\int_\Delta \prod_x q(x)^{N(x)}\, \nu(dq) = (m - 1)!\, \prod_x N(x)! / (n + m - 1)!$
(the integral of $\prod_x q(x)^{a_x - 1}$ over the simplex with respect to Lebesgue measure is
$\prod_x \Gamma(a_x)/\Gamma(\sum_x a_x)$, by induction on $m$ with the Beta integral). Hence
\[
  \log M_n - n\, \KL(\hat q_n \| p)
  = \log\frac{(m - 1)!}{(n + m - 1)!} + \sum_x \big[\log N_n(x)! - N_n(x)\log N_n(x)\big] + n \log n .
\]
Now $N! \ge (N/e)^N$ gives $\sum_x [\log N_n(x)! - N_n(x)\log N_n(x)] \ge -\sum_x N_n(x) = -n$;
and $(n + m - 1)!/(m - 1)! = \binom{n + m - 1}{m - 1}\, n!$ with $\binom{n + m - 1}{m - 1} \le (n + 1)^{m - 1}$
(Lemma~\ref{lem:num_types}) and $n! \le e \sqrt n\,(n/e)^n$ (Stirling's upper bound). Therefore
\[
  \log M_n - n\, \KL(\hat q_n \| p)
  \ge -(m - 1)\log(n + 1) - 1 - \tfrac12 \log n - n\log n + n - n + n\log n
  = -(m - 1)\log(n + 1) - 1 - \tfrac12 \log n .
\]

\emph{Conclusion.} On the event $\{\forall n, M_n < 1/\delta'\}$, of probability at least
$1 - \delta'$, for every $n \ge 1$:
$n\, \KL(\hat q_n \| p) \le \log M_n + (m - 1)\log(n + 1) + 1 + \frac12 \log n < \log(1/\delta') + (m - 1)\log(n + 1) + 1 + \frac12\log n$.
The simplified form follows from $1 + \frac12 \log n \le \log(e(n + 1))$ and $\log(n + 1) \le \log(e(n + 1))$.
\end{proof}

\begin{remark}[Why a union bound over $n$ of Sanov's bound does not suffice]\label{rem:empirical_sanov_union}
The paper proves the KL event by applying its Lemma~19 (Sanov's bound for a fixed sample size,
Theorem~\ref{thm:sanov_hp}) with the confidence $\delta/(3SAH)$ and a union bound over
$(h, s, a)$; the union over the sample sizes $n$ (equivalently over the episodes $t$) is left
implicit. It cannot be done with the rate $\alpha(n, \delta) = \log(3SAH/\delta) + S \log(8e(n + 1))$:
Sanov's bound gives $\Prob(\KL(\hat q_n \| p) > \alpha(n)/n) \le (n + 1)^{S - 1} e^{-\alpha(n)} = \frac{\delta}{3SAH}\,(8e)^{-S}\,\frac{1}{n + 1}$,
whose sum over $n \ge 1$ diverges. A union bound over $n$ with the summable weights
$1/(n(n + 1))$ would require the larger rate $\log(3SAH/\delta) + (S + 1)\log(8e(n + 1))$, i.e.\
$S + 1$ in place of $S$ in every downstream constant. Lemma~\ref{lem:empirical_kl_time_uniform}
(the mixture-martingale argument behind the time-uniform KL bound used by
\cite{menard2021fast}, whose rate is exactly $\alpha$) is the correct route and keeps the paper's
rates; it uses Ville's inequality and the type count of Chapter~\ref{chap:pre_concentration}
in place of Theorem~\ref{thm:sanov_hp}.
\end{remark}

\textbf{Lean remark.} The statement is library material for \texttt{Essakine2026Tight/Mathlib/}
(then Chapter~\ref{chap:pre_concentration}): a sample \texttt{ξ : ℕ → Ω → 𝒳} with
\texttt{iIndepFun} and \texttt{HasLaw (ξ k) (weightedMeasure p)}, the empirical measure as a
\texttt{weightedMeasure} and \texttt{klDiv}. Ingredients: Mathlib's
\texttt{Real.pow\_div\_factorial\_le\_exp} for $N^N/N! \le e^N$, the Stirling bound
$n! \le e\sqrt n\,(n/e)^n$ from \texttt{Stirling.stirlingSeq'\_antitone} and
\texttt{Stirling.stirlingSeq\_one} ($n!/(\sqrt{2n}(n/e)^n)$ is antitone with value $e/\sqrt 2$ at
$1$), and the Dirichlet integral, which is not in Mathlib (prove it by induction on $m$ with
\texttt{Real.betaIntegral}, or define $\nu$ as the law of a normalized vector of i.i.d.\
exponential variables and compute the moments). The martingale is with respect to the
natural filtration of \texttt{ξ}, and Lemma~\ref{lem:ville} is used in the form
$\Prob(\sup_n M_n \ge 1/\delta') \le \delta'$.

\begin{lemma}[Probability of the KL event]\label{lem:event_kl_prob}
\uses{def:event_kl, lem:empirical_kl_time_uniform, lem:empirical_visits_domination, lem:emp_trans_reindex, lem:rates_props}
For $\delta \in (0, 1)$, $\Prob(\evKL) \ge 1 - \delta/3$.
\end{lemma}

\begin{proof}
\uses{def:event_kl, lem:empirical_kl_time_uniform, lem:empirical_visits_domination, lem:emp_trans_reindex, def:rates, def:counts, def:emp_trans}
Fix $(h, s, a)$, put $p := p_h(\cdot \mid s, a)$, $\delta' := \delta/(3SAH)$, and let
$B_n := \{x \in \Ss^n : \KL(\hat q_n(x) \| p) > \alpha(n, \delta)/n\}$ for $n \ge 1$.

\emph{Reduction to the visits.} Let $\omega \notin \evKL$: there are $t$, $h$, $(s, a)$ with
$\KL(\phat_h^t(s, a) \| p_h(s, a)) > \alpha(n, \delta)/n$ where $n := n_h^t(s, a)$; since the threshold is
$+\infty$ for $n = 0$, $n \ge 1$, hence $T_n \le t < \infty$ and, by Lemma~\ref{lem:emp_trans_reindex},
$\phat_h^t(\cdot \mid s, a) = \hat q_n(W_1, \dots, W_n)$: so $(W_1, \dots, W_n) \in B_n$. Therefore
\[
  \Omega \setminus \evKL \subseteq \bigcup_{(h, s, a)} \big\{\exists n \ge 1 : T_n < \infty,\ (W_1, \dots, W_n) \in B_n\big\} .
\]

\emph{The i.i.d.\ bound.} With $\xi_1, \xi_2, \dots$ i.i.d.\ of law $p$ on a finite set of $S$
elements, Lemma~\ref{lem:empirical_kl_time_uniform} with $m = S$ and the confidence $\delta'$
gives $\Prob'(\exists n \ge 1 : \KL(\hat q_n(\xi) \| p) \ge (\log(1/\delta') + S\log(e(n + 1)))/n) \le \delta'$,
and $\log(1/\delta') + S\log(e(n + 1)) = \log(3SAH/\delta) + S \log(e(n + 1)) \le \alpha(n, \delta)$
(Definition~\ref{def:rates}: $e(n + 1) \le 8e(n + 1)$). Hence
$\Prob'(\exists n \ge 1 : (\xi_1, \dots, \xi_n) \in B_n) \le \delta'$.

\emph{Transfer and union bound.} Lemma~\ref{lem:empirical_visits_domination}(ii) gives
$\Prob(\exists n \ge 1 : T_n < \infty, (W_1, \dots, W_n) \in B_n) \le \delta'$ for each $(h, s, a)$, and
the union bound over the $SAH$ triples gives
$\Prob(\Omega \setminus \evKL) \le SAH \cdot \delta/(3SAH) = \delta/3$.
\end{proof}

\begin{lemma}[Probability of the Bernstein event]\label{lem:event_bern_prob}
\uses{def:event_bern, cor:bernstein_explicit, lem:empirical_visits_domination, lem:emp_trans_reindex, lem:rates_props}
Let $\delta \in (0, 1)$, $f = (f_h)_{h \le H + 1}$ and $b = (b_h)_{h \le H}$ with
$b_h \ge \max_{s'} f_{h+1}(s') - \min_{s'} f_{h+1}(s')$ for every $h \le H$. Then
$\Prob(\evBern(f, b)) \ge 1 - \delta/3$.
\end{lemma}

\begin{proof}
\uses{def:event_bern, cor:bernstein_explicit, thm:bernstein_self_normalized, lem:empirical_visits_domination, lem:emp_trans_reindex, def:rates, def:emp_trans, def:counts}
Fix $(h, s, a)$ and put $p := p_h(\cdot \mid s, a)$, $g := f_{h+1}$, $\bar g := pg$,
$\sigma^2 := \Var_p(g) = p(g - \bar g)^2$, $\delta' := \delta/(3SAH)$, and
\[
  B_n := \Big\{x \in \Ss^n : \big|(\hat q_n(x) - p) g\big| > \sqrt{2 \sigma^2 \alpha^*(n, \delta)/n} + 3 b_h\, \alpha^*(n, \delta)/n\Big\} .
\]

\emph{Reduction to the visits.} Exactly as in Lemma~\ref{lem:event_kl_prob}: if
$\omega \notin \evBern(f, b)$, some $t, h, (s, a)$ with $n := n_h^t(s, a) \ge 1$ violate the inequality,
$T_n \le t < \infty$, and $\phat_h^t(\cdot \mid s, a) = \hat q_n(W_1, \dots, W_n)$
(Lemma~\ref{lem:emp_trans_reindex}), so $(W_1, \dots, W_n) \in B_n$; thus
$\Omega \setminus \evBern(f, b) \subseteq \bigcup_{(h, s, a)} \{\exists n \ge 1 : T_n < \infty, (W_1, \dots, W_n) \in B_n\}$.

\emph{The i.i.d.\ bound.} Let $\xi_1, \xi_2, \dots$ be i.i.d.\ with law $p$ and
$Y_k := g(\xi_k) - \bar g$. The $Y_k$ are i.i.d., centered, with $\E[Y_k^2] = \sigma^2$, and
$\abs{Y_k} \le \max g - \min g \le b_h$ since $\bar g \in [\min g, \max g]$. If $b_h = 0$ then $g$ is
constant, $(\hat q_n(x) - p) g = 0$ for every $x$, and $B_n = \emptyset$ (the right-hand side is
nonnegative). If $b_h > 0$, apply Corollary~\ref{cor:bernstein_explicit} (the explicit form of the
self-normalized Bernstein inequality, Theorem~\ref{thm:bernstein_self_normalized}) to the
sequence $(Y_k)$ with its natural filtration, the constant weights $w_k := 1 \in [0, 1]$ (trivially
predictable), the bound $b_h$ and the confidence $\delta'$: here
$S_n = \sum_{k \le n} Y_k = n\,(\hat q_n(\xi) - p) g$ and $V_n = \sum_{k \le n} \E[Y_k^2 \mid \mathcal G_{k-1}] = n\sigma^2$,
so with probability at least $1 - \delta'$, for all $n \ge 1$,
\[
  n\,\big|(\hat q_n(\xi) - p) g\big| \le \sqrt{2 n \sigma^2 L_n} + 3 b_h L_n,
  \qquad L_n := \log\frac{4e(2n + 1)}{\delta'} .
\]
Since $4e(2n + 1) \le 8e(n + 1)$, $L_n \le \log(8e(n + 1)) + \log(3SAH/\delta) = \alpha^*(n, \delta)$
(Definition~\ref{def:rates}); dividing by $n$ and using the monotonicity of the right-hand side
in $L_n$, on the same event $(\xi_1, \dots, \xi_n) \notin B_n$ for all $n \ge 1$. Hence
$\Prob'(\exists n \ge 1 : (\xi_1, \dots, \xi_n) \in B_n) \le \delta'$.

\emph{Transfer and union bound.} Lemma~\ref{lem:empirical_visits_domination}(ii) and the union
bound over the $SAH$ triples give $\Prob(\Omega \setminus \evBern(f, b)) \le SAH \cdot \delta/(3SAH) = \delta/3$.
\end{proof}

\begin{remark}[The index of the Bernstein inequality]\label{rem:empirical_bernstein_index}
The paper's proof of Lemma~5 applies the self-normalized Bernstein inequality to the process
indexed by the \emph{episodes} $t$, with the predictable weights $w_t = \indic\{(s^t_h, a^t_h) = (s, a)\}$;
this gives the deviation bound with the logarithmic term $\log(4e(2t + 1)/\delta')$ in the
episode index $t$, and the paper then ``uses $n_h^t(s, a) \le t$ and the monotonicity of the
logarithm'' to replace it by $\alpha^*(n_h^t(s, a), \delta)$, which goes in the wrong direction
($\log(4e(2t+1)) \ge \log(4e(2n+1))$). Indexing the process by the number of visits, as above
and as in \cite{menard2021fast}, gives the term in $n$ directly; the domination lemma
(Lemma~\ref{lem:empirical_visits_domination}) is what allows to use the i.i.d.\ version of the
inequality. The paper's proof also chooses $\delta_{h, s, a} = \delta/(SAH)$ in the union bound,
which gives $\delta$ rather than the claimed $\delta/3$; the rates of Definition~\ref{def:rates}
contain $\log(3SAH/\delta)$, i.e.\ the choice $\delta_{h, s, a} = \delta/(3SAH)$ used here.
\end{remark}

\begin{lemma}[Probability of the counts event]\label{lem:event_cnt_prob}
\uses{def:event_cnt, thm:bernoulli_concentration, lem:pseudo_counts_increments, lem:rates_props}
For $\delta \in (0, 1)$, $\Prob(\evCnt) \ge 1 - \delta/3$.
\end{lemma}

\begin{proof}
\uses{def:event_cnt, thm:bernoulli_concentration, lem:pseudo_counts_increments, def:pseudo_counts, def:counts, def:rates}
Fix $(h, s, a)$ and $\delta' := \delta/(3SAH)$. For $i \ge 1$ let $X'_i := \indic\{(s^i_h, a^i_h) = (s, a)\}$,
the visit indicator of the episode $i$, and consider the filtration $(\mathcal F_i)_{i \ge 0}$.
$X'_i$ is a function of $(X_{i-1}, Y_{i-1})$, which is part of $\hist_i$, hence $\mathcal F_i$-measurable;
and by Lemma~\ref{lem:pseudo_counts_increments}(iii) (applied with $t = i - 1$),
$P_i := \Prob(X'_i = 1 \mid \mathcal F_{i-1}) = p^{\pi^i}_h(s, a)$ almost surely, which is
$\mathcal F_{i-1}$-measurable (a function of $X_{i-1} = \pi^i$). Theorem~\ref{thm:bernoulli_concentration}
(the Bernoulli maximal inequality, Lemma~20 of the paper, from \cite{dann2017unifying}) applied
to $(X'_i)$, $(P_i)$ and $(\mathcal F_i)$ with the confidence $\delta'$ gives
\[
  \Prob\Big(\exists t \ge 1 : \sum_{i \le t} X'_i < \frac12 \sum_{i \le t} P_i - \log\frac1{\delta'}\Big) \le \delta' .
\]
Now $\sum_{i \le t} X'_i = n_h^t(s, a)$ (Definition~\ref{def:counts}), $\sum_{i \le t} P_i = \nbar_h^t(s, a)$
(Definition~\ref{def:pseudo_counts}) and $\log(1/\delta') = \log(3SAH/\delta) = \alpha^{\mathrm{cnt}}(\delta)$;
for $t = 0$ the inequality $n_h^0 = 0 \ge 0 - \alpha^{\mathrm{cnt}}(\delta)$ holds since
$\alpha^{\mathrm{cnt}}(\delta) \ge 0$ (Lemma~\ref{lem:rates_props}). So the constraint of $\evCnt$ for the
triple $(h, s, a)$ fails with probability at most $\delta'$, and the union bound over the $SAH$
triples gives $\Prob(\Omega \setminus \evCnt) \le \delta/3$.
\end{proof}

\textbf{Lean remark.} The filtration is only needed inside Theorem~\ref{thm:bernoulli_concentration};
its hypotheses are supplied as: \texttt{X' i} is measurable with respect to the $\sigma$-algebra
generated by \texttt{histAt X Y i}, and the conditional law of \texttt{X' i} given
\texttt{(histAt X Y (i - 1), X (i - 1))} is the Bernoulli law of parameter
\texttt{occupancy M (X (i - 1)) s₁ h s a} (Lemma~\ref{lem:pseudo_counts_increments}(iii) as a
\texttt{HasCondDistrib} statement, from Lemma~\ref{lem:episode_env_law}). The martingale
$\exp(\sum_{i \le t}(P_i/2 - X'_i))$ of the proof of Theorem~\ref{thm:bernoulli_concentration} is
adapted to $(\mathcal F_t)$.

\begin{lemma}[Lemma 5: the good event has probability at least $1 - \delta$]\label{lem:good_event}
\uses{def:good_event, lem:event_kl_prob, lem:event_bern_prob, lem:event_cnt_prob, lem:opt_exp_value_range}
For $\beta > 0$, $\Prob(\evGood^+) \ge 1 - \delta$, and for $\beta < 0$, $\Prob(\evGood^-) \ge 1 - \delta$.
\end{lemma}

\begin{proof}
\uses{def:good_event, lem:event_kl_prob, lem:event_bern_prob, lem:event_cnt_prob, lem:opt_exp_value_range}
By Lemma~\ref{lem:opt_exp_value_range}, for every $h \le H$ the function $Z^*_{h+1}$ takes values
in $[1, e^{\beta(H - h)}]$ if $\beta > 0$ and in $[e^{\beta(H - h)}, 1]$ if $\beta < 0$, so
$\max Z^*_{h+1} - \min Z^*_{h+1} \le \abs{e^{\beta(H - h)} - 1} \le b^\pm_h$ (for $\beta > 0$,
$e^{\beta(H-h)} - 1 \le e^{\beta(H - h)} = b^+_h$; for $\beta < 0$, $1 - e^{\beta(H-h)} = b^-_h$). Hence
Lemma~\ref{lem:event_bern_prob} applies to $(Z^*, b^\pm)$ and gives
$\Prob(\evBern(Z^*, b^\pm)) \ge 1 - \delta/3$; with Lemmas~\ref{lem:event_kl_prob}
and~\ref{lem:event_cnt_prob}, the union bound gives
$\Prob(\Omega \setminus \evGood^\pm) \le \delta/3 + \delta/3 + \delta/3 = \delta$.
\end{proof}

\textbf{Lean remark.} \texttt{measureReal\_goodEvent\_ge : 1 - δ ≤ P.real (goodEvent M s₁ X Y β δ)}
for any algorithm-environment sequence \texttt{(X, Y)} in \texttt{statesEnv M s₁}, with the
hypotheses $\delta \in (0, 1)$, \texttt{M.HasRewardFn r} and $r \in [0, 1]$ (needed for the ranges of
$Z^*$); it is applied in Chapters~\ref{chap:analysis_pos} and~\ref{chap:analysis_neg} to runs
of Entropic-BPI, but holds for every algorithm.

\chapter{The Entropic-BPI algorithm}
\label{chap:algorithm}

This chapter defines the algorithm of the paper (Section~4, Algorithm~1, and the equations of
Appendix~B: (bonus positive), (optimism\_equation\_positive), (width\_certificate\_positive)
and their negative versions) as an identification algorithm in the episode environment, and
the auxiliary ``ring'' quantities of its analysis. The algorithm knows the reward function
$r$ and the parameters $\beta \ne 0$, $\delta \in (0, 1)$, $\eps > 0$, $s_1$; the transition
kernels are unknown to it.

\textbf{Standing setting and abbreviations.} $r$ is a reward function with values in $[0, 1]$
(Definition~\ref{def:reward_fn}), $\beta \ne 0$, $\delta \in (0, 1)$, $\eps > 0$ and $s_1 \in \Ss$.
All the quantities of Sections~\ref{sec:algorithm_backups}--\ref{sec:algorithm_alg} are
computed after $t$ episodes from the history $\hist_t$ of the first $t$ episodes: they only
depend on the counts $n_h^t(s, a)$ (Definition~\ref{def:counts}), the empirical transitions
$\phat_h^t(\cdot \mid s, a)$ (Definition~\ref{def:emp_trans}), the rates $\alpha, \alpha^*$ of
Definition~\ref{def:rates} and $r$; when $t$, $h$ and $(s, a)$ are clear we abbreviate
$n := n_h^t(s, a)$, $\phat := \phat_h^t(\cdot \mid s, a)$ and write $\phat g$ for
$(\phat_h^t g)(s, a)$. The ring quantities of Section~\ref{sec:algorithm_ring} also use the true
kernels $p_h$ and exist only for the analysis.

\textbf{Lean remark.} Every quantity of the chapter is defined as a function of a history
\texttt{hist : Hist Unit (Policy S A H) (Traj S H) t} (file \texttt{Essakine2026Tight/EV2026/Algorithm.lean},
namespace \texttt{Essakine2026Tight}) and evaluated on a run $(X, Y)$ through
\texttt{histAt X Y t ω} (\texttt{Run.lean}: \texttt{UtildeAt}, \texttt{UlowerAt},
\texttt{ZtildeAt}, \texttt{ZlowerAt}, \texttt{bonusAt}, \texttt{greedyAt}, \texttt{certAt},
\texttt{ringZAt}, \texttt{ringUAt}). The backward recursions are indexed by the number
\texttt{j : ℕ} of steps to the end: \texttt{j = 0} is the terminal step $H + 1$ and the paper's
step $h$ has \texttt{j = H + 1 - h}; the Lean step of a pair is \texttt{stepOf H k : Fin H}, the
step with \texttt{k = j - 1 = H - h} steps after it, and the values at the step \texttt{j} are
computed from the values at \texttt{j - 1} by one application of the one-step functions
\texttt{bonus}, \texttt{stepU}, \texttt{ringUAux} (which take the step data
\texttt{k n r phat Zt Zl} as arguments and are independent of histories).

\section{Bonus and backups}
\label{sec:algorithm_backups}

\begin{definition}[Exploration bonus]\label{def:bonus}
\lean{Essakine2026Tight.bonus, Essakine2026Tight.bonusAt}\leanok
\uses{def:counts, def:emp_trans, def:rates, lem:emp_trans_simplex}
Let $t \ge 0$, $h \le H$ and $(s, a)$ with $n = n_h^t(s, a) \ge 1$, and let $\Zt^t_{h+1}, \Zl^t_{h+1} : \Ss \to \R$
be the optimistic and pessimistic values of the next step (Definition~\ref{def:backups}). The
\emph{bonus} is, for $\beta > 0$ (eq.~(bonus positive) of the paper),
\[
  b_h^t(s, a) := 2\sqrt2\, \sqrt{\Var_{\phat}\big(\Zt^t_{h+1}\big)\, \frac{\alpha^*(n)}{n}}
  + 5\, e^{\beta(H - h)}\, \frac{\alpha(n)}{n}
  + 4 H\, e^{\beta(H - h)}\, \frac{\alpha(n)}{n} ,
\]
and for $\beta < 0$ (eq.~(bonus\_negative)),
\[
  b_h^t(s, a) := 2\sqrt2\, \sqrt{\Var_{\phat}\big(\Zl^t_{h+1}\big)\, \frac{\alpha^*(n)}{n}}
  + 5\, \big(1 - e^{\beta(H - h)}\big)\, \frac{\alpha(n)}{n}
  + 4 H\, \big(1 - e^{\beta(H - h)}\big)\, \frac{\alpha(n)}{n} ,
\]
where $\Var_{\phat}(g) = \Var_{\phat_h^t}(g)(s, a)$ is the variance under the empirical
transitions (Definition~\ref{def:emp_trans}). The bonus is not defined (and never used) for a
pair with $n_h^t(s, a) = 0$. The third term carries the KL rate $\alpha$ (as in the statement of
the paper's Lemma~7), not the Bernstein rate $\alpha^*$ of the paper's equations
(bonus positive) and (bonus\_negative): with $\alpha^*$, which can be as small as $\alpha/S$,
the certificate lemmas (Lemmas~\ref{lem:certificate_pos} and~\ref{lem:certificate_neg}) fail
(Remark~\ref{rem:pos_constants_bonus}).
\end{definition}

\textbf{Lean remark.} \texttt{bonus nA H β δ k n phat Zt Zl : ℝ} with \texttt{k = H - h} the
number of steps after the step (so $e^{\beta(H - h)}$ is \texttt{Real.exp (β * k)}),
\texttt{nA = Fintype.card A}, the rates applied with \texttt{Fintype.card S}, and
\texttt{if 0 < β then … else …} selecting the sign; \texttt{vecVar phat Zt} is $\Var_{\phat}(\Zt)$.
The function is total in Lean and takes the value $0$ at \texttt{n = 0} (Lean's $\alpha(0)/0 = 0$);
this value never appears in a statement, since every statement about the bonus assumes
\texttt{0 < n} (deviation~1). The constants $2\sqrt2$, $5$, $4H$ are those of Lemma~7 of the
paper (Lemma~\ref{lem:concentration_pos}), of which the bonus is by definition the right-hand
side minus the last term; they are verified there.

\begin{definition}[Optimistic and pessimistic backups]\label{def:backups}
\lean{Essakine2026Tight.stepU, Essakine2026Tight.optZStep, Essakine2026Tight.optZ, Essakine2026Tight.stepUAt, Essakine2026Tight.ZtildeAt, Essakine2026Tight.ZlowerAt, Essakine2026Tight.UtildeAt, Essakine2026Tight.UlowerAt}\leanok
\uses{def:bonus, def:counts, def:emp_trans, def:reward_fn}
Let $t \ge 0$. The \emph{optimistic and pessimistic exponential values} $\Zt^t_h, \Zl^t_h : \Ss \to \R$
($h \le H + 1$) and \emph{backups} $\Ut^t_h, \Ul^t_h : \Ss \times \As \to \R$ ($h \le H$) are defined by the
backward recursion $\Zt^t_{H+1} := 1$, $\Zl^t_{H+1} := 1$ and, for $h = H, \dots, 1$, with the
\emph{clip} $c_h := e^{\beta(H + 1 - h)}$:

\emph{Case $\beta > 0$} (eq.~(optimism\_equation\_positive)). For a pair with $n = n_h^t(s, a) \ge 1$
and $b = b_h^t(s, a)$,
\begin{align*}
  \Ut^t_h(s, a) &:= \min\Big\{ c_h,\ e^{\beta r_h(s, a)} \Big[ \phat \Zt^t_{h+1} + b + \frac1H\, \phat\big(\Zt^t_{h+1} - \Zl^t_{h+1}\big) \Big] \Big\}, \\
  \Ul^t_h(s, a) &:= \max\Big\{ 1,\ e^{\beta r_h(s, a)} \Big[ \phat \Zl^t_{h+1} - b - \frac1H\, \phat\big(\Zt^t_{h+1} - \Zl^t_{h+1}\big) \Big] \Big\},
\end{align*}
and for a never-visited pair ($n_h^t(s, a) = 0$) $\Ut^t_h(s, a) := c_h$, $\Ul^t_h(s, a) := 1$; then
$\Zt^t_h(s) := \max_{a} \Ut^t_h(s, a)$ and $\Zl^t_h(s) := \max_a \Ul^t_h(s, a)$.

\emph{Case $\beta < 0$} (eq.~(optimism\_equation\_negative)). The clips are exchanged: for $n \ge 1$,
\begin{align*}
  \Ut^t_h(s, a) &:= \min\Big\{ 1,\ e^{\beta r_h(s, a)} \Big[ \phat \Zt^t_{h+1} + b + \frac1H\, \phat\big(\Zt^t_{h+1} - \Zl^t_{h+1}\big) \Big] \Big\}, \\
  \Ul^t_h(s, a) &:= \max\Big\{ c_h,\ e^{\beta r_h(s, a)} \Big[ \phat \Zl^t_{h+1} - b - \frac1H\, \phat\big(\Zt^t_{h+1} - \Zl^t_{h+1}\big) \Big] \Big\},
\end{align*}
for a never-visited pair $\Ut^t_h(s, a) := 1$, $\Ul^t_h(s, a) := c_h$; and
$\Zt^t_h(s) := \min_a \Ut^t_h(s, a)$, $\Zl^t_h(s) := \min_a \Ul^t_h(s, a)$.
\end{definition}

\begin{remark}[Deviations 1 and 3: unvisited pairs and the clips]\label{rem:algorithm_clips}
(a) The paper handles never-visited pairs through the convention $\alpha(0)/0 = \infty$, which
makes the bonus infinite and the minimum/maximum equal to the clips; the blueprint defines the
backups of such pairs directly as the clips, so that Lean's $\alpha(0)/0 = 0$ never enters.
(b) The paper writes the optimistic clip as $e^{\beta(H - h - 1)}$ (both in Section~4 and in
eq.~(optimism\_equation\_positive)) and, for $\beta < 0$, the pessimistic clip as
$e^{\beta(H - h - 1)}$. With the paper's $1$-indexed steps this is off by two: the target of the
optimistic backup is $U^*_h$, which ranges in $[1, e^{\beta(H + 1 - h)}]$ for $\beta > 0$
(Lemma~\ref{lem:opt_exp_value_range}), and optimism ($U^*_h \le \Ut^t_h$, Lemma~\ref{lem:optimism_pos})
forces the clip to dominate $e^{\beta(H + 1 - h)}$: at the last step $h = H$ the paper's clip
$e^{-\beta} < 1 \le U^*_H(s, a) = e^{\beta r_H(s, a)}$ would break optimism for every pair. The
recursion is consistent with the clip $c_h = e^{\beta(H + 1 - h)}$: the bracket of $\Ut^t_h$ is
at most $e^{\beta}\cdot e^{\beta(H - h)}$ plus the bonus terms, and $e^{\beta(H + 1 - h)}$ is the
exact range of the values it approximates (Lemma~\ref{lem:backups_range}). In Lean the clip is
\texttt{Real.exp (β * (k + 1))} for \texttt{k = H - h} steps after the step (the paper's
exponent $H - h - 1$ is $k - 1$).
\end{remark}

\textbf{Lean remark.} \texttt{stepU nA H β δ k n r phat Zt Zl : ℝ × ℝ} computes the pair
$(\Ut, \Ul)(s, a)$ of one pair from its step data, with the clips
\texttt{up := if 0 < β then Real.exp (β * (k + 1)) else 1},
\texttt{low := if 0 < β then 1 else Real.exp (β * (k + 1))} and \texttt{if n = 0 then (up, low) else (min up …, max low …)};
\texttt{optZ r β δ hist : ℕ → (S → ℝ) × (S → ℝ)} is the backward recursion
($\Zt, \Zl$ at the step with \texttt{j} steps to the end; \texttt{optZ 0 = (1, 1)}, and
\texttt{optZ (k + 1)} takes the \texttt{max} (\texttt{0 < β}) or \texttt{min} over
\texttt{A} of the \texttt{stepU} of the step \texttt{stepOf H k}); \texttt{stepUAt r β δ hist h s a}
is the pair of backups at the step \texttt{h : Fin H}. The finite maxima and minima are
\texttt{Finset.univ.sup'}/\texttt{inf'} or LML's \texttt{Function.max}/\texttt{min} on a
finite nonempty type.

\begin{lemma}[Ranges and order of the backups]\label{lem:backups_range}
\uses{def:backups, def:bonus, lem:rates_props, lem:emp_trans_simplex}
For every $t$, $h \le H$, $(s, a)$ with $n_h^t(s, a) \ge 1$: $b_h^t(s, a) \ge 0$. For every $t$,
$h \le H$ and $(s, a)$:
\begin{enumerate}
  \item[(i)] if $\beta > 0$: $1 \le \Ul^t_h(s, a) \le \Ut^t_h(s, a) \le e^{\beta(H + 1 - h)}$;
  \item[(ii)] if $\beta < 0$: $e^{\beta(H + 1 - h)} \le \Ul^t_h(s, a) \le \Ut^t_h(s, a) \le 1$;
\end{enumerate}
and for every $h \le H + 1$ and $s$, the same bounds hold for $\Zl^t_h(s) \le \Zt^t_h(s)$:
$1 \le \Zl^t_h(s) \le \Zt^t_h(s) \le e^{\beta(H + 1 - h)}$ if $\beta > 0$, and
$e^{\beta(H + 1 - h)} \le \Zl^t_h(s) \le \Zt^t_h(s) \le 1$ if $\beta < 0$ (at $h = H + 1$ all
values are $1$). In particular $\Zt^t_{h+1}$ and $\Zl^t_{h+1}$ take values in $[1, e^{\beta(H - h)}]$
($\beta > 0$), resp.\ $[e^{\beta(H - h)}, 1]$ ($\beta < 0$), and $\Zt^t_h, \Zl^t_h > 0$.
\end{lemma}

\begin{proof}
\uses{def:backups, def:bonus, lem:rates_props, lem:emp_trans_simplex}
Fix $t$. We prove by backward induction on $h = H + 1, H, \dots, 1$ the statement $P(h)$:
\emph{for all $s$, $\Zl^t_h(s) \le \Zt^t_h(s)$ and both lie in $I_h := [\min(1, e^{\beta(H + 1 - h)}), \max(1, e^{\beta(H + 1 - h)})]$},
proving along the way the claims about $b$, $\Ul$ and $\Ut$ at the step $h$. $P(H + 1)$: both
values are $1 \in I_{H+1} = \{1\}$. Assume $P(h + 1)$ for some $h \le H$ and fix $(s, a)$.

\emph{The bonus is nonnegative.} For $n \ge 1$: $\Var_{\phat}(\cdot) \ge 0$ (Lemma~\ref{lem:emp_trans_simplex}),
$\alpha(n), \alpha^*(n) \ge 0$ (Lemma~\ref{lem:rates_props}), $e^{\beta(H - h)} > 0$, and for $\beta < 0$
also $1 - e^{\beta(H - h)} \ge 0$; so $b_h^t(s, a)$ is a sum of nonnegative terms.

\emph{The correction term is nonnegative.} $w := \frac1H \phat(\Zt^t_{h+1} - \Zl^t_{h+1}) \ge 0$ by
$P(h + 1)$ and the positivity of $\phat$ (Lemma~\ref{lem:emp_trans_simplex}).

\emph{Case $\beta > 0$.} Put $c := e^{\beta(H + 1 - h)} \ge 1$ (as $H + 1 - h \ge 1$). If $n = 0$,
$\Ul = 1 \le c = \Ut$. If $n \ge 1$, let $X := e^{\beta r_h(s, a)}[\phat\Zt^t_{h+1} + b + w]$ and
$Y := e^{\beta r_h(s, a)}[\phat\Zl^t_{h+1} - b - w]$, so that $\Ut = \min\{c, X\}$, $\Ul = \max\{1, Y\}$.
Then: $X \ge 1$, since $e^{\beta r} \ge 1$, $\phat \Zt^t_{h+1} \ge 1$ (by $P(h + 1)$, $\Zt^t_{h+1} \ge 1$,
and $\phat$ preserves lower bounds) and $b, w \ge 0$; $Y \le c$, since
$Y \le e^{\beta r_h(s, a)} \phat \Zl^t_{h+1} \le e^{\beta} e^{\beta(H - h)} = c$ (as $r \le 1$ and
$\Zl^t_{h+1} \le \Zt^t_{h+1} \le e^{\beta(H - h)}$); and $Y \le X$ (as $\Zl^t_{h+1} \le \Zt^t_{h+1}$ and
$b, w \ge 0$). Hence $\max\{1, Y\} \le \min\{c, X\}$: indeed $1 \le c$, $1 \le X$, $Y \le c$ and
$Y \le X$. So $1 \le \Ul \le \Ut \le c$. Taking the maximum over $a$ of both sides preserves
the order and the bounds: $1 \le \Zl^t_h(s) = \max_a \Ul \le \max_a \Ut = \Zt^t_h(s) \le c$, which is $P(h)$.

\emph{Case $\beta < 0$.} Put $c := e^{\beta(H + 1 - h)} \le 1$. If $n = 0$, $\Ul = c \le 1 = \Ut$. If
$n \ge 1$, with $X, Y$ as above, $\Ut = \min\{1, X\}$ and $\Ul = \max\{c, Y\}$. Then $X \ge c$: since
$\beta < 0$ and $r \le 1$, $e^{\beta r} \ge e^{\beta}$, and $\phat\Zt^t_{h+1} \ge e^{\beta(H - h)}$ by
$P(h + 1)$, so $X \ge e^{\beta} e^{\beta(H - h)} + 0 = c$; $Y \le 1$: $e^{\beta r} \le 1$ and
$\phat\Zl^t_{h+1} \le 1$; $Y \le X$ as before; and $c \le 1$. Hence $c \le \Ul \le \Ut \le 1$, and the
minimum over $a$ gives $P(h)$.

The final claims restate $P(h + 1)$ and positivity of the lower bounds.
\end{proof}

\textbf{Lean remark.} The induction is on \texttt{j} for \texttt{optZ r β δ hist j}
(\texttt{optZ\_le}, \texttt{optZ\_mem\_Icc}), and the one-step statement is a lemma about
\texttt{stepU} with hypotheses on \texttt{Zt}, \texttt{Zl} (\texttt{stepU\_fst\_le},
\texttt{stepU\_mem\_Icc}); the bonus nonnegativity is \texttt{bonus\_nonneg} (for
\texttt{0 < n}, or unconditionally, the value at \texttt{n = 0} being $0$).

\section{Greedy policy and certificate}
\label{sec:algorithm_certificate}

\begin{definition}[Greedy policy]\label{def:greedy}
\lean{Essakine2026Tight.greedy, Essakine2026Tight.greedyAt}\leanok
\uses{def:backups}
The policy played at the episode $t + 1$ is the \emph{greedy policy} of the history $\hist_t$:
\[
  \pi^{t+1}_h(s) := \argmax_{a \in \As} \Ut^t_h(s, a) \quad (\beta > 0),
  \qquad
  \pi^{t+1}_h(s) := \argmin_{a \in \As} \Ul^t_h(s, a) \quad (\beta < 0),
\]
for $h \le H$ and $s \in \Ss$, with a fixed tie-breaking rule. Consequently
$\Zt^t_h(s) = \Ut^t_h(s, \pi^{t+1}_h(s))$ for $\beta > 0$ and $\Zl^t_h(s) = \Ul^t_h(s, \pi^{t+1}_h(s))$ for
$\beta < 0$ (Definition~\ref{def:backups}).
\end{definition}

\begin{remark}[Deviation 2: the greedy policy for $\beta < 0$]\label{rem:algorithm_greedy_neg}
The algorithm box of the paper (Section~4) writes $\argmin_a \Ut^t_h(s, a)$ for $\beta < 0$, while
its Appendix~B.2 (after eq.~(optimism\_equation\_negative)) defines $\pi^{t+1}_h(s) = \argmin_a \Ul^t_h(s, a)$
and its analysis uses the latter: the stopping rule compares the certificate to $\Zl^t_1(s_1)$
and Lemma~16 (Lemma~\ref{lem:certificate_neg}) bounds $\Ur - \Ul$ at the greedy action, both
of which need $\Zl^t_h(s) = \Ul^t_h(s, \pi^{t+1}_h(s))$. For $\beta < 0$ the pessimistic $\Ul$ is
the \emph{optimistic} bound on the entropic value $Q^*$ (the exponential is decreasing), so
$\argmin_a \Ul$ is the optimistic greedy choice; this is the definition adopted.
\end{remark}

\textbf{Lean remark.} \texttt{greedy r β δ hist : Policy S A H := fun h s ↦ if 0 < β then argmax (fun a ↦ (stepUAt r β δ hist h s a).1) else argmin (fun a ↦ (stepUAt r β δ hist h s a).2)}
with LML's \texttt{argmax}/\texttt{argmin} on a finite nonempty type (a definite choice of a
maximizer, \texttt{le\_argmax}); \texttt{greedyAt r β δ X Y t ω := greedy r β δ (histAt X Y t ω)}.
The map \texttt{hist ↦ greedy r β δ hist} is measurable because the history type is countable
(\texttt{measurable\_of\_countable}).

\begin{definition}[Certificate]\label{def:certificate}
\lean{Essakine2026Tight.certStep, Essakine2026Tight.cert, Essakine2026Tight.certAt}\leanok
\uses{def:backups, def:bonus, def:greedy}
Let $t \ge 0$ and $\pi := \pi^{t+1}$. The \emph{certificate} of $\pi$ is the family of functions
$(\pi_h \cert^t_h) : \Ss \to \R$, $h \le H + 1$, defined by the backward recursion
$\cert^t_{H+1} := 0$ and, for $h = H, \dots, 1$ and $s \in \Ss$, with $a := \pi_h(s)$,
$n := n_h^t(s, a)$, $\phat := \phat_h^t(\cdot \mid s, a)$ and the clip
\[
  c_h := e^{\beta(H + 1 - h)} \ (\beta > 0), \qquad c_h := 1 \ (\beta < 0):
\]
\[
  (\pi_h \cert^t_h)(s) := \cert^t_h(s, a),
  \qquad
  \cert^t_h(s, a) :=
  \begin{cases}
    \min\Big\{ c_h,\ e^{\beta r_h(s, a)} \Big[ 3\, b_h^t(s, a) + \Big(1 + \frac3H\Big)\, \phat\,\big(\pi_{h+1}\cert^t_{h+1}\big) \Big] \Big\} & \text{if } n \ge 1, \\[4pt]
    c_h & \text{if } n = 0
  \end{cases}
\]
(eq.~(width\_certificate\_positive) and (width\_certificate\_negative) of the paper, with
$\pi_{H+1}\cert^t_{H+1} := 0$). Only the composed functions $\pi_h \cert^t_h$ enter the
recursion and the algorithm; $\cert^t_h(s, a)$ for $a \ne \pi_h(s)$ is never used.
\end{definition}

\begin{remark}[The clip of the certificate]\label{rem:algorithm_cert_clip}
For $\beta > 0$ the paper clips the certificate at $e^{\beta(H - h)}$ in Appendix~B (and at
$e^{\beta(H - h - 1)}$ in Section~4). The certificate must dominate $\Zt^t_h - \Zr^t_h$
(Lemma~\ref{lem:cert_dominates_gap_pos}), which is at most $e^{\beta(H + 1 - h)} - 1$ for a visited
pair and equal to $c_h - \min\{e^{\beta r}(p_h \Zr^t_{h+1}), 1\}$-type quantities for an unvisited
one; with the paper's clip the domination fails at $h = H$ as soon as $e^{\beta} - 1 > 1$. The
clip $c_h = e^{\beta(H + 1 - h)}$ (the same as for $\Ut$, deviation~3) dominates every gap of
Lemma~\ref{lem:backups_range}. For $\beta < 0$ the clip $1$ dominates $1 - e^{\beta(H + 1 - h)}$.
\end{remark}

\textbf{Lean remark.} \texttt{cert r β δ hist : ℕ → S → ℝ}, the backward recursion of the
composed certificate ($\pi_h \cert^t_h$ at the step with \texttt{j} steps to the end;
\texttt{cert 0 = 0}), whose step \texttt{k + 1} computes the greedy action
\texttt{a := greedy r β δ hist h s} at \texttt{h := stepOf H k}, the bonus from
\texttt{optZ r β δ hist k}, the clip \texttt{if 0 < β then Real.exp (β * (k + 1)) else 1}, and
\texttt{if n = 0 then clip else min clip (…)}; \texttt{certAt r β δ X Y t ω h}.

\begin{lemma}[Range of the certificate]\label{lem:certificate_range}
\uses{def:certificate, lem:backups_range, lem:emp_trans_simplex}
For every $t$, $h \le H + 1$ and $s$: $0 \le (\pi^{t+1}_h \cert^t_h)(s) \le c_h$, with
$c_{H+1} := 0$; in particular $\pi^{t+1}_{h+1}\cert^t_{h+1}$ takes values in $[0, e^{\beta(H - h)}]$
($\beta > 0$) or $[0, 1]$ ($\beta < 0$) for $h \le H$.
\end{lemma}

\begin{proof}
\uses{def:certificate, lem:backups_range, lem:emp_trans_simplex}
Backward induction on $h$. At $h = H + 1$ the certificate is $0$. Let $h \le H$ and assume
$0 \le \pi_{h+1}\cert^t_{h+1} \le c_{h+1}$; fix $s$, $a = \pi_h(s)$. If $n = 0$ the value is $c_h \in [0, c_h]$
($c_h > 0$). If $n \ge 1$: $b_h^t(s, a) \ge 0$ (Lemma~\ref{lem:backups_range}),
$\phat(\pi_{h+1}\cert^t_{h+1}) \ge 0$ by the induction hypothesis and positivity of $\phat$
(Lemma~\ref{lem:emp_trans_simplex}), and $e^{\beta r} > 0$, so the bracket is nonnegative and
$\min\{c_h, \cdot\} \in [0, c_h]$. For $\beta > 0$, $c_{h+1} = e^{\beta(H - h)}$; for $\beta < 0$, $c_{h+1} = 1$
(and $0$ at $h = H$, contained in $[0, 1]$).
\end{proof}

\section{Stopping rule and the algorithm}
\label{sec:algorithm_alg}

\begin{definition}[Stopping rule]\label{def:stopping_rule}
\lean{Essakine2026Tight.stopCond}\leanok
\uses{def:certificate, def:backups}
The \emph{stopping condition} after $t$ episodes is
\[
  (\pi^{t+1}_1 \cert^t_1)(s_1) \le \frac{e^{\beta\eps} - 1}{e^{\beta\eps}}\, \Zt^t_1(s_1)
  \quad (\beta > 0, \text{ eq.~(computable\_stopping\_condition)}),
\]
\[
  (\pi^{t+1}_1 \cert^t_1)(s_1) \le \big(1 - e^{\beta\eps}\big)\, \Zl^t_1(s_1)
  \quad (\beta < 0, \text{ Algorithm~1 and Appendix~B.2}).
\]
Entropic-BPI stops after $t$ episodes if the condition holds for the history $\hist_t$ (and not
before). Both sides are computable from $\hist_t$ and $r$ by the backward recursions of
Definitions~\ref{def:backups} and~\ref{def:certificate}.
\end{definition}

\textbf{Lean remark.} \texttt{stopCond r β δ ε s₁ hist : Prop := if 0 < β then cert r β δ hist H s₁ ≤ (Real.exp (β * ε) - 1) / Real.exp (β * ε) * (optZ r β δ hist H).1 s₁ else cert r β δ hist H s₁ ≤ (1 - Real.exp (β * ε)) * (optZ r β δ hist H).2 s₁}
(the step $1$ has \texttt{j = H} steps to the end). For $\beta > 0$ the condition is
equivalent to $\pi\cert \le (e^{\beta\eps} - 1)(\Zt^t_1(s_1) - \pi\cert)$, the form in which it is
used (Lemma~\ref{lem:pac_at_stopping_pos}). Appendix~B.2 of the paper writes the $\beta < 0$
condition as $\pi \cert \le (e^{\beta\eps} - 1)\Zl^t_1(s_1)$, whose right-hand side is negative for
$\beta < 0$ (the algorithm would never stop); the algorithm box and the proof of the sample
complexity use $(1 - e^{\beta\eps})\Zl^t_1(s_1)$, which is the intended condition. Note that the
condition fails after $0$ episodes: all pairs are unvisited, so $\pi^1\cert^0_1(s_1) = c_1$ while
$\Zt^0_1(s_1) = e^{\beta H}$ (resp.\ $\Zl^0_1(s_1) = e^{\beta H}$) and the factors
$(e^{\beta\eps} - 1)/e^{\beta\eps}$, $1 - e^{\beta\eps}$ are less than $1$; hence $\tau \ge 1$.

\begin{definition}[Algorithm 1: Entropic-BPI]\label{def:entropic_bpi}
\lean{Essakine2026Tight.entropicBPI}\leanok
\uses{def:bpi_alg, def:greedy, def:stopping_rule, def:backups, def:certificate}
\emph{Entropic-BPI} with parameters $(r, \beta, \delta, \eps, s_1)$ is the BPI algorithm
(Definition~\ref{def:bpi_alg}) whose
\begin{enumerate}
  \item[(i)] \emph{sampling rule} is deterministic: after $t$ episodes, with history $\hist_t$, it
        plays the greedy policy $\pi^{t+1}$ of $\hist_t$ (Definition~\ref{def:greedy}) at the
        episode $t + 1$;
  \item[(ii)] \emph{stopping rule} is the set of histories $\hist_t$ (of any length $t$) satisfying
        the stopping condition of Definition~\ref{def:stopping_rule};
  \item[(iii)] \emph{output rule} is deterministic: it outputs the greedy policy of the history at
        which it stops, i.e.\ $\hat\pi := \pi^{\tau + 1}$ when it stops after $\tau$ episodes.
\end{enumerate}

\medskip\noindent\textbf{Algorithm 1 (Entropic-BPI).}
\begin{enumerate}
  \item[] \textbf{Input:} the reward function $r$, $\beta \ne 0$, $\delta \in (0, 1)$, $\eps > 0$, the
        initial state $s_1$.
  \item[] \textbf{Initialization:} $t := 0$, empty history ($n_h^0 \equiv 0$, $\phat_h^0 \equiv 1/S$).
  \item[] \textbf{Repeat} (after $t$ episodes, from the history $\hist_t$):
  \item[(1)] \emph{Terminal initialization:} $\Zt^t_{H+1} := 1$, $\Zl^t_{H+1} := 1$, $\cert^t_{H+1} := 0$.
  \item[(2)] \emph{Backward pass, for $h = H, H - 1, \dots, 1$:} for all $(s, a)$ with $n_h^t(s, a) \ge 1$
        compute the bonus $b_h^t(s, a)$ (Definition~\ref{def:bonus}); for all $(s, a)$ compute the
        backups $\Ut^t_h(s, a)$, $\Ul^t_h(s, a)$ and then $\Zt^t_h$, $\Zl^t_h$
        (Definition~\ref{def:backups}); for all $s$ set the greedy action $\pi^{t+1}_h(s)$
        (Definition~\ref{def:greedy}); for all $s$ compute the certificate $(\pi^{t+1}_h\cert^t_h)(s)$
        (Definition~\ref{def:certificate}).
  \item[(3)] \emph{Stopping test:} if the stopping condition of Definition~\ref{def:stopping_rule}
        holds, \textbf{stop}: $\tau := t$, \textbf{output} $\hat\pi := \pi^{t+1}$.
  \item[(4)] \emph{Otherwise} play the episode $t + 1$ with the policy $\pi^{t+1}$ from $s_1$,
        observe its state sequence $(s^{t+1}_1, \dots, s^{t+1}_{H+1})$, append it to the history,
        set $t := t + 1$ and go to (1).
\end{enumerate}
\end{definition}

\textbf{Lean remark.} \texttt{entropicBPI r β δ ε s₁ : BPIAlg S A H} with
\texttt{alg := detAlgorithm (fun \_ p ↦ greedy r β δ p.1) (fun \_ ↦ measurable\_of\_countable \_)}
(LML \texttt{detAlgorithm}: the next action is a measurable function of the pair (history,
observation); the observation is \texttt{()}), \texttt{stopSet := \{h : Σ n, Hist Unit (Policy S A H) (Traj S H) n | stopCond r β δ ε s₁ h.2\}}
(measurable since the history type is countable: \texttt{(Set.to\_countable \_).measurableSet}),
and \texttt{output := Kernel.deterministic (fun h ↦ greedy r β δ h.2) (measurable\_of\_countable \_)}.
This is the LML main branch's \texttt{IdentAlg} with a \texttt{stopSet} of histories of
variable length and one \texttt{output} kernel on them (the previous statement-only
formalization used a stopping predicate \texttt{stop n hist} and a family \texttt{output n}).

\begin{lemma}[Runs of Entropic-BPI]\label{lem:entropic_bpi_run}
\uses{def:entropic_bpi, def:bpi_alg, def:greedy, def:stopping_rule, def:episode_env}
Let $(X, Y, \hat\pi)$ be a run of Entropic-BPI in the episode environment of an MDP
$\mathcal M \in \mathcal M_r$ from $s_1$, on $(\Omega, \Prob)$, and let $\tau$ be its stopping time.
Almost surely:
\begin{enumerate}
  \item[(i)] for every $t \ge 0$, $X_t = \pi^{t+1}$, the greedy policy of the history $\hist_t = (X_i, Y_i)_{i < t}$
        of the first $t$ episodes (Definition~\ref{def:greedy});
  \item[(ii)] $\tau = \min\{t \ge 0 : \hist_t \text{ satisfies the stopping condition}\}$: on
        $\{\tau < \infty\}$ the stopping condition holds for $\hist_\tau$ and fails for $\hist_t$
        for every $t < \tau$; on $\{\tau = \infty\}$ it fails for every $t$;
  \item[(iii)] on $\{\tau < \infty\}$, $\hat\pi = \pi^{\tau + 1}$, the greedy policy of $\hist_\tau$.
\end{enumerate}
\end{lemma}

\begin{proof}
\uses{def:entropic_bpi, def:bpi_alg, def:greedy, def:stopping_rule}
(i) $(X, Y)$ is an algorithm-environment sequence for the sampling rule of Entropic-BPI, a
deterministic algorithm whose next action after the history $\hist_t$ is $\pi^{t+1} = \mathrm{greedy}(\hist_t)$.
The defining property of an algorithm-environment sequence is that the conditional law of
$X_t$ given $\hist_t$ (and the trivial observation) is the policy kernel, here the Dirac mass
at $\mathrm{greedy}(\hist_t)$; a random variable whose conditional law given $Z$ is the Dirac mass
at $g(Z)$ equals $g(Z)$ almost surely. A countable intersection over $t$ gives (i) almost
surely for all $t$ simultaneously.

(ii) The stopping time of an identification algorithm is the hitting time of its stopping rule
by the process $t \mapsto \hist_t$ (Definition~\ref{def:bpi_alg}): $\tau$ is the least $t$ with
$\hist_t$ in the stopping rule, i.e.\ satisfying the stopping condition, and $+\infty$ if there is
none. Minimality gives the failure of the condition for all $t < \tau$, and $\tau < \infty$ gives
$\hist_\tau$ in the stopping rule.

(iii) The output has conditional law $\rho(\hist_\tau)$ given the stopped history $\hist_\tau$,
where $\rho$ is the deterministic kernel $\mathrm{hist} \mapsto \delta_{\mathrm{greedy}(\mathrm{hist})}$; as in (i),
$\hat\pi = \mathrm{greedy}(\hist_\tau)$ almost surely. On $\{\tau < \infty\}$ the stopped history is the
history of the first $\tau$ episodes, so $\hat\pi = \mathrm{greedy}(\hist_\tau) = \pi^{\tau + 1}$
(which is also $X_\tau$ by (i)).
\end{proof}

\textbf{Lean remark.} (i) is LML's \texttt{IsDeterministicAlg.action\_ae\_all\_eq} (or
\texttt{action\_detAlgorithm\_ae\_all\_eq}) for \texttt{detAlgorithm}: \texttt{∀ᵐ ω ∂P, ∀ t, X t ω = greedy r β δ (histAt X Y t ω)};
its proof is \texttt{hasCondDistrib\_action} with the deterministic policy kernel and
\texttt{ae\_eq\_of\_hasCondDistrib\_deterministic}. (ii) is the characterization of
\texttt{hittingAfter} (Mathlib: \texttt{hittingAfter\_mem\_set\_of\_ne\_top},
\texttt{hittingAfter\_le\_iff}, and LML's \texttt{IdentAlg.stoppedHist\_mem\_stopSet\_of\_ne\_top});
the value \texttt{τ ω : ℕ∞} is compared with \texttt{t} through \texttt{ENat} coercions. (iii)
is \texttt{IsRun.hasCondDistrib\_output} with \texttt{Kernel.deterministic}, giving
\texttt{out =ᵐ[P] fun ω ↦ greedy r β δ (A.stoppedHist O X Y ω).2}, and
\texttt{stoppedHist ω = ⟨τ ω, histAt X Y (τ ω) ω⟩} on \texttt{τ ω ≠ ⊤} (LML
\texttt{stoppedValue} of \texttt{sigmaHistory}). Statements about the output are only made on
$\{\tau < \infty\}$, which has full probability under the good event
(Lemma~\ref{lem:stopping_time_bound_pos}).

\section{The ring quantities of the analysis}
\label{sec:algorithm_ring}

The analysis compares the computable values to the values of the played policy through an
auxiliary backward recursion, which follows the true exponential Bellman recursion of the
greedy policy but is clipped by the empirical pessimistic (resp.\ optimistic) backup. It uses
the unknown kernels $p_h$ and is never computed by the algorithm.

\begin{definition}[Ring values and backups]\label{def:ring}
\lean{Essakine2026Tight.ringUAux, Essakine2026Tight.ringZStep, Essakine2026Tight.ringZ, Essakine2026Tight.ringU, Essakine2026Tight.ringZAt, Essakine2026Tight.ringUAt}\leanok
\uses{def:backups, def:bonus, def:greedy, def:episodic_mdp, def:reward_fn}
Let $\mathcal M \in \mathcal M_r$ with transition kernels $(p_h)$, $t \ge 0$ and $\pi := \pi^{t+1}$. The
\emph{ring values} $\Zr^t_h : \Ss \to \R$, $h \le H + 1$, and \emph{ring backups}
$\Ur^t_h : \Ss \times \As \to \R$, $h \le H$, are defined by $\Zr^t_{H+1} := 1$ and, for $h = H, \dots, 1$
and every $(s, a)$, with $n := n_h^t(s, a)$, $\phat := \phat_h^t(\cdot \mid s, a)$, $b := b_h^t(s, a)$
and $e^{\beta r} := e^{\beta r_h(s, a)}$:

\emph{Case $\beta > 0$} (Appendix~B.1, ``Stopping rule''). For $n \ge 1$,
\begin{align*}
  \Ur^t_{h, \mathrm{pes}}(s, a) &:= \max\Big\{ 1,\ e^{\beta r} \Big[ \phat \Zr^t_{h+1} - b - \frac1H\, \phat\big(\Zt^t_{h+1} - \Zr^t_{h+1}\big) \Big] \Big\}, \\
  \Ur^t_h(s, a) &:= \min\Big\{ e^{\beta r}\, (p_h \Zr^t_{h+1})(s, a),\ \Ur^t_{h, \mathrm{pes}}(s, a) \Big\},
\end{align*}
and for a never-visited pair $\Ur^t_h(s, a) := \min\{ e^{\beta r}\, (p_h \Zr^t_{h+1})(s, a),\ 1\}$
(the pessimistic backup being $\Ul^t_h(s, a) = 1$ there).

\emph{Case $\beta < 0$} (Appendix~B.2). For $n \ge 1$,
\begin{align*}
  \Ur^t_{h, \mathrm{opt}}(s, a) &:= \min\Big\{ 1,\ e^{\beta r} \Big[ \phat \Zr^t_{h+1} + b + \frac1H\, \phat\big(\Zr^t_{h+1} - \Zl^t_{h+1}\big) \Big] \Big\}, \\
  \Ur^t_h(s, a) &:= \max\Big\{ e^{\beta r}\, (p_h \Zr^t_{h+1})(s, a),\ \Ur^t_{h, \mathrm{opt}}(s, a) \Big\},
\end{align*}
and for a never-visited pair $\Ur^t_h(s, a) := \max\{ e^{\beta r}\, (p_h \Zr^t_{h+1})(s, a),\ 1\}$
(the optimistic backup being $\Ut^t_h(s, a) = 1$ there).

In both cases $\Zr^t_h(s) := \Ur^t_h(s, \pi_h(s))$. Here $\Zt^t_{h+1}$, $\Zl^t_{h+1}$ and $b$ are those
of Definitions~\ref{def:backups} and~\ref{def:bonus} for the same history, and
$(p_h \Zr^t_{h+1})(s, a) = \sum_{s'} p_h(s' \mid s, a)\Zr^t_{h+1}(s')$ uses the true kernel.
\end{definition}

\textbf{Lean remark.} \texttt{ringUAux nA H β δ k n r p phat Zr Zt Zl : ℝ} is the one-step
function (arguments: the true transition vector \texttt{p = M.transVec h s a}, the empirical
one, the next-step ring value and the next-step optimistic and pessimistic values), with the
sign selected by \texttt{if 0 < β} and the unvisited case by \texttt{if n = 0};
\texttt{ringZ M β δ hist : ℕ → S → ℝ} is the backward recursion ($\Zr$ at the step with
\texttt{j} steps to the end, \texttt{ringZ 0 = 1}, evaluated at the greedy action
\texttt{greedy r β δ hist h s} of the same history, where \texttt{r} is the reward function of
\texttt{M}), \texttt{ringU M β δ hist h s a} the ring backup at the step \texttt{h}, and
\texttt{ringZAt}, \texttt{ringUAt} their values on a run. The lemmas about these quantities
(Lemmas~\ref{lem:ring_pos}, \ref{lem:certificate_pos}, \ref{lem:cert_dominates_gap_pos} and
their negative counterparts) are in Chapters~\ref{chap:analysis_pos} and~\ref{chap:analysis_neg}.

\chapter{Analysis of Entropic-BPI for \texorpdfstring{$\beta > 0$}{beta > 0}}
\label{chap:analysis_pos}

This chapter reproduces Appendix~B.1 of \cite{essakine2026tight} (Lemmas~7--11 and the proof
of the sample complexity for $\beta > 0$), decomposed into lemmas that are each a plausible
single Lean declaration, with every constant recomputed. The structure is the one of
\cite{menard2021fast} transported to the exponential space of \cite{fei2021exponential}: a
concentration inequality for the optimal exponential value (Lemma~\ref{lem:concentration_pos})
gives optimism (Lemma~\ref{lem:optimism_pos}); an auxiliary ``ring'' value
$\Zr^t$ bridges the pessimistic value $\Zl^t$ and the value $Z^{\pi^{t+1}}$ of the greedy
policy (Lemmas~\ref{lem:ring_pos}, \ref{lem:certificate_pos}), so that the computable
certificate $\pi^{t+1}\cert^t$ dominates the gap $Z^* - Z^{\pi^{t+1}}$
(Lemmas~\ref{lem:cert_dominates_gap_pos}, \ref{lem:stopping_pos}) and the stopping rule is
sound (Lemma~\ref{lem:pac_at_stopping_pos}); finally the certificate is bounded under the true
kernel (Lemma~\ref{lem:cert_true_recursion_pos}), unrolled along the trajectory of $\pi^{t+1}$
(Lemmas~\ref{lem:cert_unrolled_pos}, \ref{lem:ratio_bound_pos}) and summed over the episodes
(Lemmas~\ref{lem:progress_pos}, \ref{lem:stopping_time_bound_pos}, \ref{lem:pac_pos}).

\section{Standing setting and notation}
\label{sec:pos_setting}

Throughout this chapter: $\beta > 0$; $\mathcal M$ is a finite-horizon MDP
(Definition~\ref{def:episodic_mdp}) with reward function $r$ with values in $[0, 1]$
(Definition~\ref{def:reward_fn}); $\delta \in (0, 1)$ and $\eps > 0$; $(X, Y, \mathrm{out})$ is a run
of Entropic-BPI (Definition~\ref{def:entropic_bpi}) with parameters $(r, \beta, \delta, \eps, s_1)$
in the episode environment of $\mathcal M$ from $s_1$ (Definition~\ref{def:episode_env}), on a
probability space $(\Omega, \mathcal F, \Prob)$, with stopping time $\tau$. For every $t \ge 0$ the
quantities $n_h^t$, $\phat_h^t$, $b_h^t$, $\Ut_h^t$, $\Ul_h^t$, $\Zt_h^t$, $\Zl_h^t$, $\pi^{t+1}$,
$\cert_h^t$, $\Ur_h^t$, $\Zr_h^t$ are those of Definitions~\ref{def:counts}, \ref{def:emp_trans},
\ref{def:bonus}, \ref{def:backups}, \ref{def:greedy}, \ref{def:certificate}, \ref{def:ring}
computed from the history of the first $t$ episodes of the run, and $\rho_h^t$, $\rho^{*t}_h$
are the rate terms of Definition~\ref{def:rate_min}. Inside a statement about a fixed $t$ we
write $\pi := \pi^{t+1}$ and, for a fixed $(h, s, a)$,
\[
  n := n_h^t(s,a), \quad \phat := \phat_h^t(\cdot \mid s, a), \quad b := b_h^t(s,a), \quad
  \alpha := \alpha(n, \delta), \quad \alpha^* := \alpha^*(n, \delta), \quad
  \rho := \rho_h^t(s,a), \quad \rho^* := \rho^{*t}_h(s,a),
\]
and for a function $f : \Ss \to \R$, $\phat f := \sum_{s'} \phat(s' \mid s,a) f(s')$,
$(p_h f)(s,a) := \sum_{s'} p_h(s' \mid s,a) f(s')$, $\Var_{\phat}(f)$ and $\Var_{p_h}(f)(s,a)$
the variances of $f$ under these probability vectors (Lemma~\ref{lem:emp_trans_simplex}).
We also write
\[
  B_h := e^{\beta (H - h)} \quad \text{(the range parameter of step $h$)}, \qquad
  c_h := e^{\beta (H + 1 - h)} \quad \text{(the clip of step $h$)},
\]
so that $B_h = c_{h+1}$, $1 \le B_h \le c_h$ and $B_H = 1$. The ranges used all the time are:
$Z^*_{h+1}, U^*_h \in [1, c_h]$ and $Z^*_{h+1} \in [1, B_h]$ (Lemma~\ref{lem:opt_exp_value_range});
$Z^\pi_{h+1} \in [1, B_h]$ for every policy $\pi$ (Lemma~\ref{lem:exp_value_range});
$1 \le \Zl_{h+1}^t \le \Zt_{h+1}^t \le B_h$ and $1 \le \Ul_h^t \le \Ut_h^t \le c_h$
(Lemma~\ref{lem:backups_range}); $0 \le \pi_{h+1}\cert_{h+1}^t \le c_{h+1} = B_h$
(Lemma~\ref{lem:certificate_range}). The good event is $\evGood^+ = \evKL \cap \evBern(Z^*, B) \cap \evCnt$
with $B = (B_h)_{h \le H}$ (Definition~\ref{def:good_event}); ``on $\evGood^+$'' means ``for every
$\omega \in \evGood^+$'', and the statements are pointwise in $\omega$ (no measure theory is
involved before Lemma~\ref{lem:pac_pos}).

\textbf{Lean remark.} All the objects of this chapter are functions of $\omega$ through the
history $\mathrm{histAt}\ X\ Y\ t\ \omega$ (\texttt{Essakine2026Tight/EV2026/Run.lean}:
\texttt{UtildeAt}, \texttt{ZtildeAt}, \texttt{certAt}, \texttt{ringZAt}, \ldots), and the
events $\evGood^+$ etc.\ are subsets of $\Omega$; a lemma ``on $\evGood^+$, for all $t, h, s, a$,
$\ldots$'' is a Lean statement with the hypothesis \texttt{hω : ω ∈ goodEvent X Y M s₁ β δ} and
the arguments \texttt{t h s a} (\texttt{EV2026/AnalysisPos.lean}, section \texttt{Run}). No hypothesis
that $(X, Y)$ is a run of Entropic-BPI is needed before Lemma~\ref{lem:stopping_time_bound_pos}: all the
quantities are functions of the history, and the lemmas are about these functions. Every lemma is
first proved for a fixed history \texttt{hist} under the pointwise consequences of the good event at
that history (Definition~\ref{def:pos_hist_concentration}), and the run version is a one-line
corollary (Lemma~\ref{lem:pos_hist_concentration}); both names are listed in the \texttt{\textbackslash lean}
tags. The standing hypotheses are \texttt{hM : M.HasRewardFn r}, \texttt{hr : ∀ h s a, r h s a ∈ Set.Icc 0 1},
\texttt{hβ : 0 < β}, \texttt{hδ : 0 < δ} and \texttt{hδ1 : δ ≤ 1} ($\delta = 1$ is allowed), and
$\eps \ge 0$ suffices in Lemma~\ref{lem:pac_at_stopping_pos}. The
step $h$ of the paper is the Lean index $h - 1$; the backward recursions are indexed by the
number $j = H + 1 - h$ of steps to the end, so ``backward induction on $h$ from $H + 1$ down to
$1$'' is an induction on $j$ from $0$ to $H$, and the ranges $B_h = e^{\beta(H-h)}$,
$c_h = e^{\beta(H+1-h)}$ are \texttt{Real.exp (β * (j - 1))} and \texttt{Real.exp (β * j)}.
The pair $(s, a)$ with $n_h^t(s,a) = 0$ (never visited) is a separate case of every
backward induction below (deviation~1 of Chapter~\ref{chap:mdp}): its backups are the clips
$\Ut = c_h$, $\Ul = 1$, its certificate is $c_h$ and its rate terms are $\rho = \rho^* = 1$.

\section{Concentration at a history and sign-independent cores}
\label{sec:pos_cores}

This section records the auxiliary objects of the Lean proofs of this chapter and of
Chapter~\ref{chap:analysis_neg} (\texttt{EV2026/AnalysisCommon.lean}): the pointwise content of the
good event at a history, the concentration lemmas of Chapter~\ref{chap:pre_concentration} for
probability vectors, and the parts of the proofs below that do not depend on the sign of $\beta$.

\begin{definition}[Concentration at a history]\label{def:pos_hist_concentration}
\lean{Essakine2026Tight.HistConcentration, Essakine2026Tight.bernRange}\leanok
\uses{def:event_kl, def:event_bern, def:good_event, def:counts, def:emp_trans, def:rates, def:opt_value}
A history of $t$ episodes satisfies the concentration inequalities (for $\mathcal M$, $\beta$, $\delta$)
if for all $h \le H$ and $(s, a)$, $\KL(\phat_h^t(\cdot \mid s, a) \,\|\, p_h(\cdot \mid s, a)) \le \alpha(n, \delta)/n$
with $n = n_h^t(s, a)$ (the threshold being $\infty$ for $n = 0$), and for every visited pair
$\abs{(\phat_h^t - p_h)Z^*_{h+1}(s, a)} \le \sqrt{2\Var_{p_h}(Z^*_{h+1})(s, a)\alpha^*(n, \delta)/n} + 3B_h\alpha^*(n, \delta)/n$,
where $B_h = e^{\beta(H - h)}$ if $\beta > 0$ and $B_h = 1 - e^{\beta(H - h)}$ otherwise.
\end{definition}

\begin{lemma}[The good event gives concentration at every history]\label{lem:pos_hist_concentration}
\lean{Essakine2026Tight.histConcentration_histAt, Essakine2026Tight.HistConcentration.klDiv_le_ofReal}\leanok
\uses{def:pos_hist_concentration, def:good_event}
For $\omega \in \evGood^\pm$ and every $t$, the history of the first $t$ episodes satisfies the
concentration inequalities; at a visited pair the KL threshold is $\alpha(n, \delta)/n$.
\end{lemma}

\begin{proof}
\uses{def:good_event, def:event_kl, def:event_bern}
\leanok
These are the first two components of the good event, read at the episode $t$.
\end{proof}

\begin{lemma}[KL inequalities for probability vectors]\label{lem:pos_vec_kl}
\lean{Essakine2026Tight.integral_eq_vecExp_of_measureReal, Essakine2026Tight.variance_eq_vecVar_of_measureReal, Essakine2026Tight.abs_vecExp_sub_le_of_klDiv_le, Essakine2026Tight.vecVar_le_two_mul_vecVar_add_of_klDiv_le, Essakine2026Tight.vecVar_le_two_mul_vecVar_add_of_klDiv_le', Essakine2026Tight.vecVar_le_two_mul_vecVar_add, Essakine2026Tight.vecVar_le_mul_vecExp}\leanok
\uses{lem:kl_bernstein, lem:kl_transport_var, lem:transport_var, lem:var_le_range_mean, lem:pconc_weighted_measure}
Let $p$ be a probability vector on a finite set, $\nu$ a probability measure on it with vector
$q = (\nu\{x\})_x$, $a \ge 0$, $b \ge 0$, $c \in \R$, and $f, g$ with values in $[c, c + b]$. If
$\KL(\sum_x p(x)\delta_x \,\|\, \nu) \le a$, then
$\abs{pf - qf} \le \sqrt{2\Var_q(f)a} + \frac23 ba$, $\Var_q(f) \le 2\Var_p(f) + 4b^2a$ and
$\Var_p(f) \le 2\Var_q(f) + 4b^2a$. Without the KL assumption,
$\Var_p(f) \le 2\Var_p(g) + 2b\,p\abs{f - g}$, and $\Var_p(f) \le b\,pf$ if $c = 0$. (The integral and
the variance under $\nu$ are $qf$ and $\Var_q(f)$.)
\end{lemma}

\begin{proof}
\uses{lem:kl_bernstein, lem:kl_transport_var, lem:transport_var, lem:var_le_range_mean, lem:pconc_weighted_measure}
\leanok
If $b = 0$, $f$ is constant and all the claims are trivial. Otherwise apply
Lemmas~\ref{lem:kl_bernstein}, \ref{lem:kl_transport_var}, \ref{lem:transport_var}
and~\ref{lem:var_le_range_mean} to $f - c$ and $g - c$ (values in $[0, b]$), under the weighted
measure of $p$ (Lemma~\ref{lem:pconc_weighted_measure}) and $\nu$; the means shift by $c$ and the
variances are invariant. On a finite set, $\int f\,d\nu = \sum_x\nu\{x\}f(x)$
(\texttt{integral\_fintype}) and $\Var_\nu(f) = \sum_x \nu\{x\}(f(x) - \nu f)^2$.
\end{proof}

\begin{lemma}[Core of Lemmas~7 and~12]\label{lem:pos_core_concentration}
\lean{Essakine2026Tight.abs_vecExp_sub_le_bonus_add}\leanok
\uses{lem:pos_vec_kl}
Let $p, \nu, q$ be as in Lemma~\ref{lem:pos_vec_kl}, $0 \le \alpha^\star \le \alpha$ (the rates divided by
$n$), $B \ge 0$, $H \ge 1$, and $f, g$ with values in $[c, c + B]$. If $\KL(p \,\|\, \nu) \le \alpha$ and
$\abs{pf - qf} \le \sqrt{2\Var_q(f)\alpha^\star} + 3B\alpha^\star$, then
\[
  \abs{pf - qf} \le 2\sqrt2\sqrt{\Var_p(g)\alpha^\star} + 5B\alpha + 4HB\alpha + \frac1H p\abs{g - f} .
\]
\end{lemma}

\begin{proof}
\uses{lem:pos_vec_kl, lem:sqrt_mul_le}
\leanok
The computation of the proof of Lemma~\ref{lem:concentration_pos}: with $D := p\abs{g - f}$,
$2\Var_q(f)\alpha^\star \le 8\Var_p(g)\alpha^\star + 8BD\alpha^\star + 8B^2\alpha\alpha^\star$, and this is at most
$a_1^2 + a_2^2 + a_3^2 \le (a_1 + a_2 + a_3)^2$ for $a_1 := 2\sqrt2\sqrt{\Var_p(g)\alpha^\star}$,
$a_2 := D/H + 2HB\alpha^\star$ (as $a_2^2 \ge 4(D/H)(2HB\alpha^\star)$) and $a_3 := 2\sqrt2 B\alpha$ (as
$\alpha^\star \le \alpha$). Hence $\abs{pf - qf} \le a_1 + D/H + (2H + 2\sqrt2 + 3)B\alpha$, and
$2\sqrt2 + 3 \le 5 + 2H$.
\end{proof}

\begin{lemma}[Core of the third term of Lemmas~11 and~16]\label{lem:pos_core_third_term}
\lean{Essakine2026Tight.vecExp_sub_vecExp_le_of_klDiv_le}\leanok
\uses{lem:pos_vec_kl}
Let $p, \nu, q$ be as in Lemma~\ref{lem:pos_vec_kl}, $\alpha \ge 0$, $B \ge 0$, $H \ge 1$, and $f$ with
values in $[0, B]$. If $\KL(p \,\|\, \nu) \le \alpha$, then $qf - pf \le \frac1H pf + (5 + 4H)B\alpha$.
\end{lemma}

\begin{proof}
\uses{lem:pos_vec_kl, lem:sqrt_mul_le}
\leanok
As in the proof of Lemma~\ref{lem:certificate_pos}: $2\Var_q(f)\alpha \le 4B\alpha\,pf + 8B^2\alpha^2$ by
Lemma~\ref{lem:pos_vec_kl}, which is at most $(pf/H + HB\alpha)^2 + (2\sqrt2B\alpha)^2$, so
$qf - pf \le pf/H + (H + 2\sqrt2 + \frac23)B\alpha \le pf/H + (5 + 4H)B\alpha$.
\end{proof}

\begin{lemma}[Cores of the certificate recursion]\label{lem:pos_core_recursion}
\lean{Essakine2026Tight.vecExp_le_one_add_mul_vecExp_add, Essakine2026Tight.bonusTerm_le_of_klDiv_le, Essakine2026Tight.three_mul_add_le}\leanok
\uses{lem:pos_vec_kl}
Let $p, \nu, q$ be as in Lemma~\ref{lem:pos_vec_kl}, $0 \le \alpha^\star \le \alpha$, $B \ge 0$, $H \ge 1$ and
$\KL(p \,\|\, \nu) \le \alpha$.
\begin{enumerate}
  \item[(i)] If $g$ has values in $[0, B]$, then $pg \le (1 + \frac1H)qg + 2HB\alpha$.
  \item[(ii)] If $G, Z$ have values in $[c, c + B]$ and $q\abs{G - Z} \le \Delta$, then
        $2\sqrt2\sqrt{\Var_p(G)\alpha^\star} + 5B\alpha + 4HB\alpha \le 6\sqrt{\Var_q(Z)\alpha^\star} + \frac1H\Delta + 23HB\alpha$.
  \item[(iii)] (Arithmetic.) If $\sigma, \Delta, \rho, B' \ge 0$, $B, B' \le C$, $b \le 6\sigma + \frac1H\Delta + 23HB\rho$
        and $\hat g \le (1 + \frac1H)\Delta + 2HB'\rho$, then
        $3b + (1 + \frac3H)\hat g \le 36\sigma + (1 + \frac{13}H)\Delta + 84HC\rho$.
\end{enumerate}
\end{lemma}

\begin{proof}
\uses{lem:pos_vec_kl, lem:sqrt_mul_le}
\leanok
(i) and (ii) are steps (i) and (ii) of the proof of Lemma~\ref{lem:cert_true_recursion_pos}, with every
square root bounded by comparing squares as in Lemma~\ref{lem:pos_core_concentration}. (iii): $(1 + \frac3H)(1 + \frac1H) \le 1 + \frac7H$,
$(1 + \frac3H)2HB'\rho \le 8HB'\rho$ and $69HB\rho + 8HB'\rho \le 77HC\rho \le 84HC\rho$.
\end{proof}

\section{Concentration and optimism}
\label{sec:pos_optimism}

\begin{lemma}[Concentration of the optimal exponential value; Lemma~7 of the paper]\label{lem:concentration_pos}
\lean{Essakine2026Tight.abs_vecExp_sub_le_bonus_add_of_pos, Essakine2026Tight.abs_vecExp_sub_le_bonusAt_add_of_pos}\leanok
\uses{def:good_event, def:backups, def:emp_trans, def:counts, def:rates, def:opt_value, def:entropic_bpi}
On $\evGood^+$, for all $t \ge 0$, $h \in \{1, \dots, H\}$ and $(s,a)$ with $n = n_h^t(s,a) \ge 1$,
\[
  \abs{(p_h - \phat) Z^*_{h+1}(s,a)}
  \le 2\sqrt2 \sqrt{\Var_{\phat}(\Zt^t_{h+1}) \frac{\alpha^*}{n}}
  + 5 B_h \frac{\alpha}{n} + 4 H B_h \frac{\alpha}{n}
  + \frac1H \phat \abs{\Zt^t_{h+1} - Z^*_{h+1}} .
\]
The sum of the first three terms of the right-hand side is the bonus $b = b_h^t(s,a)$
(Definition~\ref{def:bonus}, with the third term $4 H B_h \alpha(n,\delta)/n$; see
Remark~\ref{rem:pos_constants_bonus}), so the conclusion reads
$\abs{(p_h - \phat) Z^*_{h+1}(s,a)} \le b + \frac1H \phat\abs{\Zt^t_{h+1} - Z^*_{h+1}}$.
\end{lemma}

\begin{proof}
\uses{def:good_event, def:event_bern, def:event_kl, lem:opt_exp_value_range, lem:backups_range, lem:kl_transport_var, lem:transport_var, lem:sqrt_add_le, lem:sqrt_mul_le, lem:rates_props, lem:emp_trans_simplex, lem:pos_core_concentration, lem:pos_hist_concentration}
\leanok
Fix $\omega \in \evGood^+$, $t$, $h$, $(s,a)$ with $n \ge 1$. Write $f := Z^*_{h+1} - 1$ and
$g := \Zt^t_{h+1} - 1$; both take values in $[0, B_h - 1] \subseteq [0, B_h]$
(Lemmas~\ref{lem:opt_exp_value_range} and~\ref{lem:backups_range}) and variances are invariant
under the shift by $1$.

\emph{Bernstein event.} Since $\omega \in \evBern(Z^*, B)$ with $B_h \ge \max Z^*_{h+1} - \min Z^*_{h+1}$
and $n \ge 1$ (Definitions~\ref{def:event_bern}, \ref{def:good_event}),
\[
  \abs{(p_h - \phat) Z^*_{h+1}(s,a)}
  \le \sqrt{2 \Var_{p_h}(Z^*_{h+1})(s,a) \frac{\alpha^*}{n}} + 3 B_h \frac{\alpha^*}{n} .
\]

\emph{Transport of the variance.} Since $\omega \in \evKL$, $\KL(\phat \,\|\, p_h(\cdot \mid s,a)) \le \alpha/n$
(Definition~\ref{def:event_kl}), so Lemma~\ref{lem:kl_transport_var} (first inequality, with
$p = \phat$, $q = p_h(\cdot \mid s,a)$, $f$, $b = B_h$) gives
$\Var_{p_h}(Z^*_{h+1}) \le 2 \Var_{\phat}(Z^*_{h+1}) + 4 B_h^2 \alpha/n$, and
Lemma~\ref{lem:transport_var} (first part, with $f$, $g$, $b = B_h$) gives
$\Var_{\phat}(Z^*_{h+1}) \le 2\Var_{\phat}(\Zt^t_{h+1}) + 2 B_h \phat\abs{Z^*_{h+1} - \Zt^t_{h+1}}$.
Hence
\[
  2 \Var_{p_h}(Z^*_{h+1}) \frac{\alpha^*}{n}
  \le 8 \Var_{\phat}(\Zt^t_{h+1}) \frac{\alpha^*}{n}
  + 8 B_h \, \phat\abs{\Zt^t_{h+1} - Z^*_{h+1}} \, \frac{\alpha^*}{n}
  + 8 B_h^2 \frac{\alpha}{n}\frac{\alpha^*}{n} .
\]

\emph{Square roots.} By Lemma~\ref{lem:sqrt_add_le} (twice), the square root of the right-hand
side is at most the sum of the three square roots. The first is
$2\sqrt2\sqrt{\Var_{\phat}(\Zt^t_{h+1})\alpha^*/n}$. For the second, Lemma~\ref{lem:sqrt_mul_le}
with $x = \phat\abs{\Zt^t_{h+1} - Z^*_{h+1}}$, $y = 8 B_h \alpha^*/n$ and $\eta = 2/H$ gives
\[
  \sqrt{x y} \le \frac{\eta x}{2} + \frac{y}{2\eta}
  = \frac1H \phat\abs{\Zt^t_{h+1} - Z^*_{h+1}} + 2 H B_h \frac{\alpha^*}{n} .
\]
For the third, Lemma~\ref{lem:sqrt_mul_le} with $\eta = 1$ gives
$\sqrt{\alpha \alpha^*} \le (\alpha + \alpha^*)/2$, so
$\sqrt{8 B_h^2 \alpha\alpha^*/n^2} \le \sqrt2 B_h (\alpha + \alpha^*)/n$.

\emph{Collecting.} Altogether, using $\alpha^* \le \alpha$ (Lemma~\ref{lem:rates_props}) and
$H \ge 1$,
\begin{align*}
  \abs{(p_h - \phat) Z^*_{h+1}(s,a)}
  &\le 2\sqrt2\sqrt{\Var_{\phat}(\Zt^t_{h+1})\frac{\alpha^*}{n}}
     + \frac1H \phat\abs{\Zt^t_{h+1} - Z^*_{h+1}}
     + \big(\sqrt2 \alpha + 3 \alpha^*\big)\frac{B_h}{n}
     + \big(2H + \sqrt2\big) B_h \frac{\alpha^*}{n} \\
  &\le 2\sqrt2\sqrt{\Var_{\phat}(\Zt^t_{h+1})\frac{\alpha^*}{n}}
     + \frac1H \phat\abs{\Zt^t_{h+1} - Z^*_{h+1}}
     + 5 B_h \frac{\alpha}{n} + 4 H B_h \frac{\alpha}{n},
\end{align*}
since $\sqrt2 + 3 \le 5$ and $2H + \sqrt2 \le 4H$.
\end{proof}

\begin{remark}[The constants of the bonus]\label{rem:pos_constants_bonus}
\uses{def:bonus, lem:concentration_pos}
The proof of Lemma~7 in \cite{essakine2026tight} takes the square root of
$\Var_{p_h}(Z^*_{h+1})\alpha^*/n$ instead of $2\Var_{p_h}(Z^*_{h+1})\alpha^*/n$ and thereby
loses a factor $\sqrt2$ on the last two terms; the constants $2\sqrt2$, $5$ and $4H$ of its
statement are nevertheless correct, by the splitting $\sqrt{\alpha\alpha^*} \le (\alpha + \alpha^*)/2$
and the choice $\eta = 2/H$ above (the paper's $\eta = 1/H$ would give $\sqrt2 \cdot 4H$).
Two points about the bonus:
\begin{enumerate}
  \item[(i)] The third term of the bonus is stated here with the KL rate $\alpha(n,\delta)$,
        as in the statement of Lemma~7 of the paper (its equation (bonus positive) writes
        $\alpha^*$ there instead). Lemma~\ref{lem:concentration_pos} with $\alpha^*$ in the
        third term is true as well (the proof above gives the bound with $(2H + \sqrt2)\alpha^*$ in
        that place), but the certificate bound of Lemma~\ref{lem:certificate_pos} needs
        $b \ge (5 + 4H) B_h \alpha/n$, and this fails for the version with $\alpha^*$ (the
        KL-Bernstein inequality that controls $(p_h - \phat)(Z^*_{h+1} - \Zr^t_{h+1})$ only has the
        KL radius $\alpha/n$ at its disposal, and $\alpha^*(n,\delta)$ can be as small as
        $\alpha(n,\delta)/S$). \textbf{Consequently Definition~\ref{def:bonus} must use
        $4 H e^{\beta(H-h)}\alpha(n,\delta)/n$ as its third term} (for both signs of $\beta$);
        the sample complexity is unaffected, since the analysis of
        Lemma~\ref{lem:cert_true_recursion_pos} bounds this term by the additive rate term
        $84 H c_h \rho$, which is already in $\alpha$.
  \item[(ii)] The bonus is \emph{defined} as the right-hand side of
        Lemma~\ref{lem:concentration_pos} minus its last term, so the lemma is the statement
        ``$\abs{(p_h - \phat)Z^*_{h+1}} \le b + \frac1H\phat\abs{\Zt^t_{h+1} - Z^*_{h+1}}$'', which is
        the only form used afterwards (Lemmas~\ref{lem:optimism_pos} and~\ref{lem:certificate_pos}).
        The statement carries no sign assumption on $\Zt^t_{h+1} - Z^*_{h+1}$ (deviation~13);
        the absolute value is removed in Lemma~\ref{lem:optimism_pos} by its induction hypothesis.
\end{enumerate}
\end{remark}

\begin{lemma}[Optimism and pessimism; Lemma~8 of the paper]\label{lem:optimism_pos}
\lean{Essakine2026Tight.optExpValue_mem_Icc_optZ_of_pos, Essakine2026Tight.optExpQ_mem_Icc_stepUAt_of_pos, Essakine2026Tight.optExpQ_le_stepUAt_fst_of_pos, Essakine2026Tight.stepUAt_snd_le_optExpQ_of_pos, Essakine2026Tight.optExpValue_mem_Icc_ZlowerAt_ZtildeAt_of_pos, Essakine2026Tight.optExpQ_mem_Icc_UlowerAt_UtildeAt_of_pos, Essakine2026Tight.optExpValue_mem_Icc_of_pos, Essakine2026Tight.optZ_mem_of_pos}\leanok
\uses{def:good_event, def:backups, def:opt_value, def:entropic_bpi}
On $\evGood^+$, for all $t \ge 0$, $h \in \{1, \dots, H + 1\}$, $s \in \Ss$ and $a \in \As$,
\[
  \Ul_h^t(s,a) \le U^*_h(s,a) \le \Ut_h^t(s,a) \quad (h \le H)
  \qquad\text{and}\qquad
  \Zl_h^t(s) \le Z^*_h(s) \le \Zt_h^t(s) .
\]
\end{lemma}

\begin{proof}
\uses{lem:concentration_pos, lem:opt_exp_value_range, lem:opt_bellman, lem:emp_trans_simplex, def:backups, lem:backups_range}
\leanok
Fix $\omega \in \evGood^+$ and $t$. We prove by backward induction on $h$ (from $H + 1$ down to
$1$) the statement $P(h)$: ``for all $s$, $\Zl_h^t(s) \le Z^*_h(s) \le \Zt_h^t(s)$'', proving the
$U$-inequalities at step $h \le H$ along the way from $P(h+1)$.

\emph{Base case} $h = H + 1$: $\Zl_{H+1}^t = Z^*_{H+1} = \Zt_{H+1}^t = 1$
(Definitions~\ref{def:backups}, \ref{def:opt_value}).

\emph{Induction step.} Let $h \le H$, assume $P(h+1)$ and fix $(s,a)$. By
Lemma~\ref{lem:opt_bellman}, $U^*_h(s,a) = e^{\beta r_h(s,a)} (p_h Z^*_{h+1})(s,a)$ and
$U^*_h(s,a) \in [1, c_h]$ (Lemma~\ref{lem:opt_exp_value_range}).

\emph{Never-visited pair} ($n = 0$): $\Ut_h^t(s,a) = c_h \ge U^*_h(s,a) \ge 1 = \Ul_h^t(s,a)$.

\emph{Visited pair} ($n \ge 1$). Upper bound: $\Ut_h^t(s,a) = \min\{c_h, X\}$ with
$X := e^{\beta r_h(s,a)}[\phat\Zt^t_{h+1} + b + \frac1H\phat(\Zt^t_{h+1} - \Zl^t_{h+1})]$, and
$U^*_h(s,a) \le c_h$, so it suffices to show $U^*_h(s,a) \le X$. Now
\[
  X - U^*_h(s,a) = e^{\beta r_h(s,a)}\Big[ \phat(\Zt^t_{h+1} - Z^*_{h+1}) + (\phat - p_h)Z^*_{h+1}(s,a)
  + b + \tfrac1H \phat(\Zt^t_{h+1} - \Zl^t_{h+1}) \Big].
\]
By $P(h+1)$, $\Zt^t_{h+1} - Z^*_{h+1} \ge 0$ pointwise, so
$\phat\abs{\Zt^t_{h+1} - Z^*_{h+1}} = \phat(\Zt^t_{h+1} - Z^*_{h+1})$ and
Lemma~\ref{lem:concentration_pos} gives
$(\phat - p_h)Z^*_{h+1}(s,a) \ge -b - \frac1H\phat(\Zt^t_{h+1} - Z^*_{h+1})$. Hence
\[
  X - U^*_h(s,a) \ge e^{\beta r_h(s,a)}\Big[\big(1 - \tfrac1H\big)\phat(\Zt^t_{h+1} - Z^*_{h+1})
  + \tfrac1H \phat(Z^*_{h+1} - \Zl^t_{h+1})\Big] \ge 0 ,
\]
because $H \ge 1$, both functions inside are nonnegative by $P(h+1)$, and $\phat$ is a
positive linear functional (Lemma~\ref{lem:emp_trans_simplex}). Lower bound:
$\Ul_h^t(s,a) = \max\{1, Y\}$ with
$Y := e^{\beta r_h(s,a)}[\phat\Zl^t_{h+1} - b - \frac1H\phat(\Zt^t_{h+1} - \Zl^t_{h+1})]$; since
$U^*_h(s,a) \ge 1$ it suffices to show $Y \le U^*_h(s,a)$:
\[
  U^*_h(s,a) - Y = e^{\beta r_h(s,a)}\Big[\phat(Z^*_{h+1} - \Zl^t_{h+1}) + (p_h - \phat)Z^*_{h+1}(s,a)
  + b + \tfrac1H\phat(\Zt^t_{h+1} - \Zl^t_{h+1})\Big]
  \ge e^{\beta r_h(s,a)}\big(1 + \tfrac1H\big)\phat(Z^*_{h+1} - \Zl^t_{h+1}) \ge 0,
\]
using again Lemma~\ref{lem:concentration_pos} in the form
$(p_h - \phat)Z^*_{h+1}(s,a) \ge -b - \frac1H\phat(\Zt^t_{h+1} - Z^*_{h+1})$ and
$-\frac1H\phat(\Zt^t_{h+1} - Z^*_{h+1}) + \frac1H\phat(\Zt^t_{h+1} - \Zl^t_{h+1}) = \frac1H\phat(Z^*_{h+1} - \Zl^t_{h+1})$.

\emph{From $U$ to $Z$.} By Lemma~\ref{lem:opt_bellman} ($\beta > 0$), $Z^*_h(s) = \max_a U^*_h(s,a)$;
by Definition~\ref{def:backups}, $\Zt_h^t(s) = \max_a\Ut_h^t(s,a)$ and $\Zl_h^t(s) = \max_a \Ul_h^t(s,a)$;
the maximum over $a$ is monotone, which gives $P(h)$.
\end{proof}

\textbf{Lean remark.} The induction is \texttt{Fin.reverseInduction} on the Lean step
\texttt{h : Fin (H + 1)} (the values at \texttt{h} are \texttt{optZ r β δ hist (H - h)}), with the statement $P$
about $Z$ only, in the form $Z^*_h(s) \in [\Zl^t_h(s), \Zt^t_h(s)]$
(\texttt{optExpValue\_mem\_Icc\_optZ\_of\_pos}); the $U$-inequalities are separate lemmas taking
$P(h+1)$ as a hypothesis (\texttt{optExpQ\_le\_stepUAt\_fst\_of\_pos} for the upper bound,
\texttt{stepUAt\_snd\_le\_optExpQ\_of\_pos} for the lower bound), and the unconditional $U$-statement
is \texttt{optExpQ\_mem\_Icc\_stepUAt\_of\_pos}. The maxima over $a$ are LML's \texttt{Function.max}
(\texttt{optZ\_eq}), handled with \texttt{Function.le\_max} and \texttt{argmax\_spec}. The positivity of
$\phat$ is \texttt{vecExp\_mono}.

\section{The ring value and the certificate}
\label{sec:pos_certificate}

The pessimistic value $\Zl^t$ is below $Z^*$ but need not be comparable to $Z^{\pi}$
($\pi = \pi^{t+1}$ is greedy for $\Ut^t$, not for $\Ul^t$). The analysis-only value $\Zr^t$ of
Definition~\ref{def:ring} follows the policy $\pi$ under the \emph{true} kernel, clipped by the
pessimistic empirical backup; it lies below both $\Zl^t$ and $Z^\pi$, and the certificate
$\pi\cert^t$ dominates $\Zt^t - \Zr^t$. We first record the trivial lower bound on $\Zr^t$,
which the paper uses implicitly.

\begin{lemma}[The ring value is at least one]\label{lem:pos_ring_range}
\lean{Essakine2026Tight.one_le_ringZ_of_pos, Essakine2026Tight.one_le_ringU_of_pos, Essakine2026Tight.one_le_ringUAux_of_pos, Essakine2026Tight.one_le_ringZAt_of_pos, Essakine2026Tight.one_le_ringUAt_of_pos}\leanok
\uses{def:ring, def:entropic_bpi}
For all $t \ge 0$, $h \in \{1, \dots, H + 1\}$, $s$ and $a$: $\Zr_h^t(s) \ge 1$ and, for $h \le H$,
$\Ur_h^t(s,a) \ge 1$.
\end{lemma}

\begin{proof}
\uses{def:ring, lem:exp_value_monotone}
\leanok
Backward induction on $h$; $\Zr_{H+1}^t = 1$. If $\Zr_{h+1}^t \ge 1$ pointwise then
$e^{\beta r_h(s,a)}(p_h\Zr^t_{h+1})(s,a) \ge (p_h \Zr^t_{h+1})(s,a) \ge (p_h 1)(s,a) = 1$ ($\beta r_h \ge 0$;
Lemma~\ref{lem:exp_value_monotone}: $p_h$ is positive and $p_h 1 = 1$). For a visited pair,
$\Ur^t_{h,\mathrm{pes}}(s,a) = \max\{1, \cdot\} \ge 1$, so $\Ur_h^t(s,a)$, the minimum of two numbers
$\ge 1$, is $\ge 1$; for a never-visited pair $\Ur_h^t(s,a) = \min\{e^{\beta r_h}(p_h\Zr^t_{h+1})(s,a), 1\} \ge 1$.
Finally $\Zr_h^t(s) = \Ur_h^t(s, \pi_h(s)) \ge 1$.
\end{proof}

\begin{lemma}[The ring value is a lower bound; Lemma~10 of the paper]\label{lem:ring_pos}
\lean{Essakine2026Tight.ringZ_le_optZ_snd_and_expValue_of_pos, Essakine2026Tight.ringU_le_stepUAt_snd_and_expQ_of_pos, Essakine2026Tight.ringU_le_stepUAt_snd_of_pos, Essakine2026Tight.ringU_le_expQ_of_pos, Essakine2026Tight.ringZAt_le_ZlowerAt_and_expValue_of_pos, Essakine2026Tight.ringUAt_le_UlowerAt_and_expQ_of_pos}\leanok
\uses{def:ring, def:backups, def:greedy, def:exp_value, def:entropic_bpi}
For all $t \ge 0$ (no event is needed), $h \in \{1, \dots, H + 1\}$, $s$ and $a$, with $\pi = \pi^{t+1}$,
\[
  \Ur_h^t(s,a) \le \min\{\Ul_h^t(s,a), U^\pi_h(s,a)\} \quad (h \le H)
  \qquad\text{and}\qquad
  \Zr_h^t(s) \le \min\{\Zl_h^t(s), Z^\pi_h(s)\} .
\]
\end{lemma}

\begin{proof}
\uses{def:ring, def:backups, def:greedy, lem:exp_value_monotone, lem:exp_bellman, def:exp_value, lem:emp_trans_simplex}
\leanok
Backward induction on $h$ with the statement $P(h)$: ``for all $s$,
$\Zr_h^t(s) \le \min\{\Zl_h^t(s), Z^\pi_h(s)\}$''. Base case: all three values are $1$ at $h = H + 1$.
Let $h \le H$, assume $P(h+1)$, fix $(s,a)$.

\emph{Comparison with $U^\pi$.} In both the visited and the never-visited case
$\Ur_h^t(s,a) \le e^{\beta r_h(s,a)}(p_h\Zr^t_{h+1})(s,a)$ (it is a minimum with this quantity), and
$(p_h\Zr^t_{h+1})(s,a) \le (p_h Z^\pi_{h+1})(s,a)$ by $P(h+1)$ and the monotonicity of $p_h$
(Lemma~\ref{lem:exp_value_monotone}); by Definition~\ref{def:exp_value},
$e^{\beta r_h(s,a)}(p_h Z^\pi_{h+1})(s,a) = U^\pi_h(s,a)$.

\emph{Comparison with $\Ul$.} If $n = 0$, $\Ur_h^t(s,a) \le 1 = \Ul_h^t(s,a)$. If $n \ge 1$,
$\Ur_h^t(s,a) \le \Ur^t_{h,\mathrm{pes}}(s,a) = \max\{1, Y'\}$ and $\Ul_h^t(s,a) = \max\{1, Y\}$ with
\[
  Y = e^{\beta r_h}\big[\phat\Zl^t_{h+1} - b - \tfrac1H\phat(\Zt^t_{h+1} - \Zl^t_{h+1})\big], \qquad
  Y' = e^{\beta r_h}\big[\phat\Zr^t_{h+1} - b - \tfrac1H\phat(\Zt^t_{h+1} - \Zr^t_{h+1})\big],
\]
so $Y - Y' = e^{\beta r_h}\big[\phat(\Zl^t_{h+1} - \Zr^t_{h+1}) + \tfrac1H\phat(\Zl^t_{h+1} - \Zr^t_{h+1})\big]
= e^{\beta r_h}\big(1 + \tfrac1H\big)\phat(\Zl^t_{h+1} - \Zr^t_{h+1}) \ge 0$ by $P(h+1)$ and the
positivity of $\phat$ (Lemma~\ref{lem:emp_trans_simplex}); $y \mapsto \max\{1, y\}$ is monotone,
hence $\Ur^t_{h,\mathrm{pes}}(s,a) \le \Ul_h^t(s,a)$. (The paper writes the factor as $1 - \frac1H$; the
correct factor is $1 + \frac1H$, and only its sign matters.)

\emph{From $U$ to $Z$.} With $a = \pi_h(s)$: $\Zr_h^t(s) = \Ur_h^t(s, a) \le \Ul_h^t(s, a) \le \max_{a'}\Ul_h^t(s,a') = \Zl_h^t(s)$
(Definitions~\ref{def:ring}, \ref{def:backups}) and
$\Zr_h^t(s) \le U^\pi_h(s, \pi_h(s)) = Z^\pi_h(s)$ (Lemma~\ref{lem:exp_bellman}). This is $P(h)$.
\end{proof}

\begin{lemma}[One-step certificate bound; Lemma~11 of the paper]\label{lem:certificate_pos}
\lean{Essakine2026Tight.stepUAt_fst_sub_ringU_le_of_pos, Essakine2026Tight.UtildeAt_sub_ringUAt_le_of_pos}\leanok
\uses{def:good_event, def:ring, def:backups, def:bonus, def:entropic_bpi}
On $\evGood^+$, for all $t \ge 0$, $h \in \{1, \dots, H\}$ and $(s,a)$ with $n = n_h^t(s,a) \ge 1$,
\[
  \Ut_h^t(s,a) - \Ur_h^t(s,a)
  \le e^{\beta r_h(s,a)}\Big[3 b + \Big(1 + \frac3H\Big)\phat\big(\Zt^t_{h+1} - \Zr^t_{h+1}\big)\Big] .
\]
\end{lemma}

\begin{proof}
\uses{lem:concentration_pos, lem:optimism_pos, lem:ring_pos, lem:pos_ring_range, lem:kl_bernstein, lem:kl_transport_var, lem:var_le_range_mean, lem:sqrt_add_le, lem:sqrt_mul_le, lem:emp_trans_simplex, lem:backups_range, lem:opt_exp_value_range, def:event_kl, def:bonus, lem:pos_core_third_term}
\leanok
Fix $\omega \in \evGood^+$, $t$, $h$, $(s,a)$ with $n \ge 1$, and write $\Delta := \phat(\Zt^t_{h+1} - \Zr^t_{h+1}) \ge 0$
(Lemmas~\ref{lem:ring_pos}, \ref{lem:optimism_pos}: $\Zr^t_{h+1} \le \Zl^t_{h+1} \le Z^*_{h+1} \le \Zt^t_{h+1}$
pointwise). Throughout, $\Ut_h^t(s,a) \le X := e^{\beta r_h}[\phat\Zt^t_{h+1} + b + \frac1H\phat(\Zt^t_{h+1} - \Zl^t_{h+1})]$
(drop the clip), and $\phat(\Zt^t_{h+1} - \Zl^t_{h+1}) \le \Delta$ (positivity of $\phat$,
Lemma~\ref{lem:emp_trans_simplex}). Recall $\Ur_h^t(s,a) = \min\{e^{\beta r_h}(p_h\Zr^t_{h+1})(s,a), \Ur^t_{h,\mathrm{pes}}(s,a)\}$.

\emph{Case 1: $\Ur_h^t(s,a) = e^{\beta r_h}(p_h\Zr^t_{h+1})(s,a)$.} Then
$\Ut_h^t(s,a) - \Ur_h^t(s,a) \le e^{\beta r_h}\big[b + \frac1H\Delta + \phat\Zt^t_{h+1} - p_h\Zr^t_{h+1}\big]$ and
\[
  \phat\Zt^t_{h+1} - p_h\Zr^t_{h+1}
  = \Delta + (\phat - p_h)Z^*_{h+1}(s,a) + (p_h - \phat)\big(Z^*_{h+1} - \Zr^t_{h+1}\big)(s,a) .
\]
\emph{Second term.} By Lemma~\ref{lem:concentration_pos} and $\abs{\Zt^t_{h+1} - Z^*_{h+1}} = \Zt^t_{h+1} - Z^*_{h+1} \le \Zt^t_{h+1} - \Zr^t_{h+1}$
pointwise, $(\phat - p_h)Z^*_{h+1}(s,a) \le b + \frac1H\Delta$.
\emph{Third term.} Let $f := Z^*_{h+1} - \Zr^t_{h+1}$; then $0 \le f \le B_h - 1 \le B_h$ pointwise
(Lemma~\ref{lem:opt_exp_value_range} and Lemma~\ref{lem:pos_ring_range}). Since $\omega \in \evKL$,
$\KL(\phat\,\|\,p_h(\cdot\mid s,a)) \le \alpha/n$, and Lemma~\ref{lem:kl_bernstein} (with $p = \phat$,
$q = p_h(\cdot \mid s,a)$, range $B_h$) gives
\[
  (p_h - \phat) f(s,a) \le \sqrt{2\Var_{p_h}(f)(s,a)\frac{\alpha}{n}} + \frac23 B_h\frac{\alpha}{n} .
\]
The variance is under the \emph{true} kernel (the second argument of the divergence); the
paper writes it under $\phat$, which is the wrong side. We transport it back:
Lemma~\ref{lem:kl_transport_var} (first inequality) gives $\Var_{p_h}(f) \le 2\Var_{\phat}(f) + 4B_h^2\alpha/n$
and Lemma~\ref{lem:var_le_range_mean} gives $\Var_{\phat}(f) \le B_h\,\phat f$. Hence, by
Lemma~\ref{lem:sqrt_add_le} and Lemma~\ref{lem:sqrt_mul_le} (with $x = \phat f$, $y = 4B_h\alpha/n$,
$\eta = 2/H$),
\[
  \sqrt{2\Var_{p_h}(f)\frac{\alpha}{n}}
  \le \sqrt{\phat f \cdot 4 B_h \frac{\alpha}{n}} + \sqrt{8 B_h^2\frac{\alpha^2}{n^2}}
  \le \frac1H \phat f + H B_h\frac{\alpha}{n} + 2\sqrt2 B_h \frac{\alpha}{n} ,
\]
and $\phat f \le \Delta$. Therefore
\[
  (p_h - \phat)f(s,a) \le \frac1H\Delta + \Big(H + 2\sqrt2 + \frac23\Big) B_h\frac{\alpha}{n}
  \le \frac1H\Delta + (4H + 5) B_h \frac{\alpha}{n} \le \frac1H \Delta + b ,
\]
since $H + 2\sqrt2 + \frac23 \le H + 3.5 \le 4H + 5$ and $b \ge 5B_h\alpha/n + 4HB_h\alpha/n$
(Definition~\ref{def:bonus} with the third term in $\alpha$, Remark~\ref{rem:pos_constants_bonus};
the square-root term of $b$ is nonnegative).
\emph{Sum.} $\phat\Zt^t_{h+1} - p_h\Zr^t_{h+1} \le 2b + (1 + \frac2H)\Delta$, hence
$\Ut_h^t(s,a) - \Ur_h^t(s,a) \le e^{\beta r_h}[3b + (1 + \frac3H)\Delta]$.

\emph{Case 2: $\Ur_h^t(s,a) = \Ur^t_{h,\mathrm{pes}}(s,a)$.} Then
$\Ur_h^t(s,a) \ge e^{\beta r_h}[\phat\Zr^t_{h+1} - b - \frac1H\Delta]$ (the maximum dominates its
second argument), so
\[
  \Ut_h^t(s,a) - \Ur_h^t(s,a) \le e^{\beta r_h}\Big[2b + \Delta + \tfrac1H\phat(\Zt^t_{h+1} - \Zl^t_{h+1}) + \tfrac1H\Delta\Big]
  \le e^{\beta r_h}\Big[2b + \big(1 + \tfrac2H\big)\Delta\Big] ,
\]
which is at most the claimed bound as $b \ge 0$ (Lemma~\ref{lem:backups_range}) and $\Delta \ge 0$.
No event is used in this case.
\end{proof}

\begin{remark}[The constants $3$ and $3/H$]\label{rem:pos_constants_certificate}
\uses{lem:certificate_pos}
The paper's proof of Lemma~11 bounds the third term by
$\frac1H\Delta + (2H + \frac23)B_h\alpha/n$ and then claims that this is at most $\frac1H\Delta + b$.
With the bonus of equation (bonus positive) of the paper this would require
$(2H + \frac23) B_h\alpha/n \le 5B_h\alpha/n + 4HB_h\alpha^*/n$, which is false in general
($\alpha^*/\alpha$ can be as small as $1/S$). With the third term of the bonus in $\alpha$
(Remark~\ref{rem:pos_constants_bonus}) the inequality $(H + 2\sqrt2 + \frac23) B_h \alpha/n \le b$
holds for every $H \ge 1$, and the constants $3$ and $3/H$ of the certificate
(Definition~\ref{def:certificate}) are correct. The coefficient $\frac1H$ in front of $\Delta$ is
what makes the certificate recursion $(1 + \frac3H)$-contracting up to the constant $e^3$
(Lemma~\ref{lem:cert_unrolled_pos}); it is bought at the price of the $HB_h\alpha/n$ term, which is
why the bonus needs an $H\alpha/n$ term.
\end{remark}

\begin{lemma}[The certificate dominates the optimistic--ring gap]\label{lem:cert_dominates_gap_pos}
\lean{Essakine2026Tight.optZ_fst_sub_ringZ_le_cert_of_pos, Essakine2026Tight.ZtildeAt_sub_ringZAt_le_certAt_of_pos}\leanok
\uses{def:good_event, def:certificate, def:backups, def:ring, def:greedy, def:entropic_bpi}
On $\evGood^+$, for all $t \ge 0$, $h \in \{1, \dots, H + 1\}$ and $s \in \Ss$, with $\pi = \pi^{t+1}$,
\[
  \Zt_h^t(s) - \Zr_h^t(s) \le (\pi_h\cert_h^t)(s) .
\]
\end{lemma}

\begin{proof}
\uses{lem:certificate_pos, lem:backups_range, lem:pos_ring_range, lem:emp_trans_simplex, def:certificate, def:greedy, def:ring, lem:certificate_range}
\leanok
Fix $\omega \in \evGood^+$ and $t$; backward induction on $h$. Base case $h = H + 1$: both sides
are $0$ ($\Zt_{H+1}^t = \Zr_{H+1}^t = 1$, $\cert_{H+1}^t = 0$). Induction step: let $h \le H$, $s \in \Ss$
and $a := \pi_h(s)$; by Definitions~\ref{def:greedy} and~\ref{def:ring},
$\Zt_h^t(s) = \Ut_h^t(s,a)$ and $\Zr_h^t(s) = \Ur_h^t(s,a)$. In all cases
$\Zt_h^t(s) - \Zr_h^t(s) \le c_h - 1 \le c_h$ by Lemma~\ref{lem:backups_range} ($\Zt_h^t \le c_h$) and
Lemma~\ref{lem:pos_ring_range} ($\Zr_h^t \ge 1$).

\emph{Never-visited pair} ($n_h^t(s,a) = 0$): $(\pi_h\cert_h^t)(s) = c_h$ (Definition~\ref{def:certificate}),
and we are done.

\emph{Visited pair} ($n \ge 1$): $(\pi_h\cert_h^t)(s) = \min\{c_h, e^{\beta r_h(s,a)}[3b + (1 + \frac3H)\phat(\pi_{h+1}\cert_{h+1}^t)]\}$;
the first argument was handled above, and for the second, Lemma~\ref{lem:certificate_pos} and the
induction hypothesis $\Zt^t_{h+1} - \Zr^t_{h+1} \le \pi_{h+1}\cert^t_{h+1}$ pointwise, integrated
against the probability vector $\phat$ (Lemma~\ref{lem:emp_trans_simplex}), give
$\Zt_h^t(s) - \Zr_h^t(s) \le e^{\beta r_h(s,a)}[3b + (1 + \frac3H)\phat(\pi_{h+1}\cert^t_{h+1})]$.
\end{proof}

\begin{lemma}[The certificate dominates the suboptimality gap; Lemma~9 of the paper]\label{lem:stopping_pos}
\lean{Essakine2026Tight.optExpValue_sub_expValue_le_cert_of_pos, Essakine2026Tight.optExpValue_sub_expValue_le_certAt_of_pos}\leanok
\uses{def:good_event, def:certificate, def:opt_value, def:exp_value, def:greedy, def:entropic_bpi}
On $\evGood^+$, for all $t \ge 0$, $h \in \{1, \dots, H + 1\}$ and $s \in \Ss$, with $\pi = \pi^{t+1}$,
\[
  Z^*_h(s) - Z^\pi_h(s) \le (\pi_h\cert_h^t)(s) .
\]
\end{lemma}

\begin{proof}
\uses{lem:optimism_pos, lem:ring_pos, lem:cert_dominates_gap_pos}
\leanok
$Z^*_h(s) \le \Zt_h^t(s)$ (Lemma~\ref{lem:optimism_pos}) and $\Zr_h^t(s) \le Z^\pi_h(s)$
(Lemma~\ref{lem:ring_pos}), so $Z^*_h(s) - Z^\pi_h(s) \le \Zt_h^t(s) - \Zr_h^t(s) \le (\pi_h\cert_h^t)(s)$
(Lemma~\ref{lem:cert_dominates_gap_pos}).
\end{proof}

\begin{lemma}[Soundness of the stopping rule]\label{lem:pac_at_stopping_pos}
\lean{Essakine2026Tight.optEntropicValue_sub_entropicValue_le_of_stopCond_of_pos, Essakine2026Tight.optEntropicValue_sub_entropicValue_greedyAt_le_of_pos}\leanok
\uses{def:good_event, def:stopping_rule, def:entropic_value, def:opt_value, def:greedy, def:entropic_bpi}
On $\evGood^+$, for every $t \ge 0$ such that the stopping condition of
Definition~\ref{def:stopping_rule} holds after $t$ episodes, i.e.\
$(\pi_1\cert_1^t)(s_1) \le \frac{e^{\beta\eps} - 1}{e^{\beta\eps}}\Zt_1^t(s_1)$ with $\pi = \pi^{t+1}$,
one has $V^*_1(s_1) - V^\pi_1(s_1) \le \eps$.
\end{lemma}

\begin{proof}
\uses{lem:cert_dominates_gap_pos, lem:ring_pos, lem:stopping_pos, lem:value_gap_log, lem:exp_value_range}
\leanok
Write $g := (\pi_1\cert_1^t)(s_1) \ge 0$. Multiplying the stopping condition by $e^{\beta\eps} > 0$
gives $e^{\beta\eps} g \le (e^{\beta\eps} - 1)\Zt_1^t(s_1)$, i.e.\ $g \le (e^{\beta\eps} - 1)(\Zt_1^t(s_1) - g)$.
By Lemma~\ref{lem:cert_dominates_gap_pos} and Lemma~\ref{lem:ring_pos},
$\Zt_1^t(s_1) - g \le \Zr_1^t(s_1) \le Z^\pi_1(s_1)$, and $e^{\beta\eps} - 1 > 0$, so
$g \le (e^{\beta\eps} - 1) Z^\pi_1(s_1)$. By Lemma~\ref{lem:stopping_pos},
$Z^*_1(s_1) - Z^\pi_1(s_1) \le g \le (e^{\beta\eps} - 1)Z^\pi_1(s_1)$, and $Z^\pi_1(s_1) > 0$
(Lemma~\ref{lem:exp_value_range}); Lemma~\ref{lem:value_gap_log} ($\beta > 0$) concludes.
\end{proof}

\textbf{Lean remark.} The stopping condition is \texttt{stopCond r β δ ε s₁ (histAt X Y t ω)};
this lemma is about the history of $t$ episodes for a fixed $t$ and does not mention $\tau$. The
translation from the exponential gap to the entropic gap is entirely in
Lemma~\ref{lem:value_gap_log}; here only real arithmetic is used.

\section{The certificate under the true kernel}
\label{sec:pos_true_kernel}

\begin{lemma}[Certificate recursion under the true kernel]\label{lem:cert_true_recursion_pos}
\lean{Essakine2026Tight.cert_le_of_pos, Essakine2026Tight.certAt_le_of_pos, Essakine2026Tight.expValue_mem_Icc_of_pos}\leanok
\uses{def:good_event, def:certificate, def:rate_min, def:exp_value, def:greedy, def:entropic_bpi}
On $\evGood^+$, for all $t \ge 0$, $h \in \{1, \dots, H\}$ and $s \in \Ss$, with $\pi = \pi^{t+1}$ and
$a := \pi_h(s)$,
\[
  (\pi_h\cert_h^t)(s) \le e^{\beta r_h(s,a)}\Big[
    36\sqrt{\Var_{p_h}(Z^\pi_{h+1})(s,a)\,\rho^{*t}_h(s,a)}
    + \Big(1 + \frac{13}{H}\Big)\big(p_h\,\pi_{h+1}\cert^t_{h+1}\big)(s,a)
    + 84\,H\,c_h\,\rho^t_h(s,a)\Big] ,
\]
where $\rho^t_h(s,a) = \min\{\alpha(n,\delta)/n, 1\}$ and $\rho^{*t}_h(s,a) = \min\{\alpha^*(n,\delta)/n, 1\}$
(both equal to $1$ if $n = n_h^t(s,a) = 0$; Definition~\ref{def:rate_min}).
\end{lemma}

\begin{proof}
\uses{lem:certificate_range, def:certificate, def:bonus, def:event_kl, lem:kl_bernstein, lem:var_le_range_mean, lem:kl_transport_var, lem:transport_var, lem:sqrt_add_le, lem:sqrt_mul_le, lem:rates_props, lem:opt_exp_value_eq, lem:optimism_pos, lem:ring_pos, lem:cert_dominates_gap_pos, lem:backups_range, lem:exp_value_range, lem:exp_value_monotone, lem:emp_trans_simplex, lem:pos_core_recursion}
\leanok
Fix $\omega \in \evGood^+$, $t$, $h$, $s$, $a = \pi_h(s)$; write $g := \pi_{h+1}\cert^t_{h+1} : \Ss \to [0, B_h]$
(Lemma~\ref{lem:certificate_range}, the clip at step $h + 1$ being $c_{h+1} = B_h$),
$\Delta_p := (p_h g)(s,a) \ge 0$ and $\rho, \rho^*$ as in the statement. All three terms in the
bracket are nonnegative and $e^{\beta r_h} \ge 1$.

\emph{Trivial case: $n = 0$ or $\alpha/n \ge 1$.} Then $\rho = 1$, and $(\pi_h\cert_h^t)(s) \le c_h \le 84 H c_h \rho$
(Lemma~\ref{lem:certificate_range}), which is at most the right-hand side.

\emph{Main case: $n \ge 1$ and $\alpha/n < 1$.} Then $\rho = \alpha/n$ and, as $\alpha^* \le \alpha$
(Lemma~\ref{lem:rates_props}), $\rho^* = \alpha^*/n \le \rho$. Dropping the clip in
Definition~\ref{def:certificate}, $(\pi_h\cert_h^t)(s) \le e^{\beta r_h(s,a)}[3b + (1 + \frac3H)\phat g]$.

\emph{Step (i): the empirical expectation of $g$.} Since $\omega \in \evKL$,
Lemma~\ref{lem:kl_bernstein} (range $B_h$) gives
$\abs{(\phat - p_h) g(s,a)} \le \sqrt{2\Var_{p_h}(g)(s,a)\rho} + \frac23 B_h\rho$, and
$\Var_{p_h}(g)(s,a) \le B_h\Delta_p$ (Lemma~\ref{lem:var_le_range_mean}). By
Lemma~\ref{lem:sqrt_mul_le} with $x = \Delta_p$, $y = 2B_h\rho$, $\eta = 2/H$:
$\sqrt{2B_h\Delta_p\rho} \le \frac1H\Delta_p + \frac{H}{2}B_h\rho$. Hence, as $\frac H2 + \frac23 \le 2H$,
\[
  \phat g \le \Big(1 + \frac1H\Big)\Delta_p + 2 H B_h \rho .
\]

\emph{Step (ii): the bonus under the true kernel.} Recall
$b = 2\sqrt2\sqrt{\Var_{\phat}(\Zt^t_{h+1})\rho^*} + 5B_h\rho + 4HB_h\rho$ (Definition~\ref{def:bonus},
Remark~\ref{rem:pos_constants_bonus}). Lemma~\ref{lem:kl_transport_var} (second inequality, with
$f = \Zt^t_{h+1} - 1 \in [0, B_h]$ by Lemma~\ref{lem:backups_range}) gives
$\Var_{\phat}(\Zt^t_{h+1}) \le 2\Var_{p_h}(\Zt^t_{h+1})(s,a) + 4B_h^2\rho$.
Lemma~\ref{lem:transport_var} (first part, with $f = \Zt^t_{h+1} - 1$ and $g' = Z^\pi_{h+1} - 1$, both in
$[0, B_h]$ by Lemmas~\ref{lem:backups_range} and~\ref{lem:exp_value_range}) gives
$\Var_{p_h}(\Zt^t_{h+1}) \le 2\Var_{p_h}(Z^\pi_{h+1}) + 2B_h\,p_h\abs{\Zt^t_{h+1} - Z^\pi_{h+1}}$. Now
$Z^\pi_{h+1} \le Z^*_{h+1} \le \Zt^t_{h+1}$ pointwise (Lemma~\ref{lem:opt_exp_value_eq} for $\beta > 0$ and
Lemma~\ref{lem:optimism_pos}), and
$\Zt^t_{h+1} - Z^\pi_{h+1} \le \Zt^t_{h+1} - \Zr^t_{h+1} \le g$ pointwise
(Lemmas~\ref{lem:ring_pos} and~\ref{lem:cert_dominates_gap_pos}); by the monotonicity of $p_h$
(Lemma~\ref{lem:exp_value_monotone}), $p_h\abs{\Zt^t_{h+1} - Z^\pi_{h+1}}(s,a) \le \Delta_p$. Therefore
\[
  \Var_{\phat}(\Zt^t_{h+1})\rho^* \le 4\Var_{p_h}(Z^\pi_{h+1})(s,a)\rho^* + 4B_h\Delta_p\rho^* + 4B_h^2\rho\rho^* ,
\]
and by Lemma~\ref{lem:sqrt_add_le} (twice), Lemma~\ref{lem:sqrt_mul_le} with $x = \Delta_p$,
$y = 4B_h\rho^*$ and $\eta = 1/(\sqrt2 H)$ (so that $\sqrt{xy} \le \frac{x}{2\sqrt2 H} + \sqrt2 H\cdot 2B_h\rho^*$),
and $\sqrt{\rho\rho^*} \le \rho$,
\[
  2\sqrt2\sqrt{\Var_{\phat}(\Zt^t_{h+1})\rho^*}
  \le 4\sqrt2\sqrt{\Var_{p_h}(Z^\pi_{h+1})(s,a)\rho^*} + \frac1H\Delta_p + 8HB_h\rho^* + 4\sqrt2 B_h\rho .
\]
With $\rho^* \le \rho$, $4\sqrt2 \le 6$ and $8H + 4\sqrt2 + 5 + 4H \le (12 + 4\sqrt2 + 5)H \le 23H$,
\[
  b \le 6\sqrt{\Var_{p_h}(Z^\pi_{h+1})(s,a)\rho^*} + \frac1H\Delta_p + 23 H B_h\rho .
\]

\emph{Step (iii): combination.} Using (i), (ii), $(1 + \frac3H)(1 + \frac1H) = 1 + \frac4H + \frac3{H^2} \le 1 + \frac7H$
and $(1 + \frac3H)\cdot 2H = 2H + 6 \le 8H$ (all for $H \ge 1$),
\begin{align*}
  3b + \Big(1 + \frac3H\Big)\phat g
  &\le 18\sqrt{\Var_{p_h}(Z^\pi_{h+1})(s,a)\rho^*} + \frac3H\Delta_p + 69HB_h\rho
     + \Big(1 + \frac7H\Big)\Delta_p + 8HB_h\rho \\
  &\le 36\sqrt{\Var_{p_h}(Z^\pi_{h+1})(s,a)\rho^*} + \Big(1 + \frac{10}{H}\Big)\Delta_p + 77 H B_h\rho ,
\end{align*}
which is at most the bracket of the statement since $B_h \le c_h$, $10 \le 13$ and $77 \le 84$.
\end{proof}

\begin{remark}[The constants $36$, $13$, $84$ and the additive term]\label{rem:pos_constants_recursion}
\uses{lem:cert_true_recursion_pos}
The paper's computation (Appendix~B.1, ``Sample complexity'') is looser than the one above: it
keeps the coefficient $\frac{2\sqrt2}{H}$ in front of $\Delta_p$ in the bound on $b$, which gives the
coefficient $1 + \frac{4 + 6\sqrt2}{H} + \frac{3}{H^2} \le 1 + \frac{15.5}{H}$ (not $1 + \frac{13}{H}$) after
multiplication by $3$, and it drops a $+9$ when simplifying $81H + (1 + \frac3H)3H$ to $84H$.
Choosing $\eta = 1/(\sqrt2 H)$ instead of $2/H$ in the AM--GM step of (ii), and the sharper
form of Lemma~\ref{lem:sqrt_mul_le} in (i), gives $1 + \frac{10}{H}$ and $77H$, so the paper's
constants $36$, $13$, $84$ are correct with room to spare. The additive term is written with
$c_h = e^{\beta(H+1-h)}$ rather than the paper's $e^{\beta(H-h)}$ (deviation~11): it must dominate
the certificate of a never-visited pair, which is its clip $c_h$, and this is the only place
where the clip of Definition~\ref{def:certificate} enters the sample complexity. The rates are
the truncated rates $\rho = \min\{\alpha/n, 1\}$ (the paper's $\alpha(n)/n \wedge 1$), with the value
$1$ for never-visited pairs, so that the statement has no $\alpha(0)/0$.
\end{remark}

\textbf{Lean remark.} For the Lean step \texttt{h : Fin H} (the paper's $h + 1$) the statement is
about \texttt{cert r β δ hist (H - h) s} (\texttt{certAt X Y r β δ t ω h.castSucc s}) and the certificate
at the next step \texttt{cert r β δ hist (H - 1 - h)} (\texttt{certAt X Y r β δ t ω h.succ}); the clip
$c_h$ is \texttt{Real.exp (β * (H - h : ℕ))}, the variance is
\texttt{vecVar (M.transVec h s a) (expValue M β π h.succ)} and the rate terms are
\texttt{rateMin (alphaStar …) n} and \texttt{rateMin (alphaKL …) n} (\texttt{starRateMinAt},
\texttt{klRateMinAt} for a run). The proof follows the three steps above: step (i) is
Lemma~\ref{lem:pos_core_recursion}(i) with range $B_h$, step (ii) is Lemma~\ref{lem:pos_core_recursion}(ii)
with $G = \Zt^t_{h+1}$, $Z = Z^\pi_{h+1}$ and $c = 1$, and step (iii) is Lemma~\ref{lem:pos_core_recursion}(iii)
with $B = B' = B_h$ and $C = c_h$ (the proof gives the constant $77 B_h \le 84 c_h$).

\begin{lemma}[Unrolling with the factor $1 + c/H$; the terms of the unrolled bound]\label{lem:pos_core_unroll}
\lean{Learning.MDP.Episodic.le_exp_mul_sum_integral_of_backward, Learning.MDP.Episodic.sum_integral_exp_mul_sqrt_vecVar_mul_le, Learning.MDP.Episodic.sum_integral_exp_mul_mul_le, Learning.MDP.Episodic.integrable_comp_obs_action}\leanok
\uses{def:exp_value, def:step_law, def:occupancy}
Let $\pi$ be a policy, $s_1 \in \Ss$, $\beta \in \R$, $r : [H] \times \Ss \times \As \to \R$,
$\E^\pi = \E^\pi_{(s_1, 1)}$ and $w_h := e^{\beta\sum_{i \le h} r_i(S_i, A_i)}$ for $h \in [H]$.
\begin{enumerate}
  \item[(i)] Let $c \ge 0$, $u_h \ge 0$ and $g_h : \Ss \to \R$ ($h \le H + 1$) with $g_{H+1} = 0$ and
  $g_h(s) \le e^{\beta r_h(s,\pi_h(s))}[u_h(s, \pi_h(s)) + (1 + \frac cH)(p_h g_{h+1})(s, \pi_h(s))]$
  for all $h \le H$ and $s \in \Ss$. Then $g_1(s_1) \le e^c\sum_{h=1}^H\E^\pi[w_h\,u_h(S_h, A_h)]$.
  \item[(ii)] If $\mathcal M$ has the reward function $r$ and $\rho_h \ge 0$, then
  \[
    \sum_{h=1}^H\E^\pi\Big[w_h\sqrt{\Var_{p_h}(Z^\pi_{h+1})(S_h, A_h)\,\rho_h(S_h, A_h)}\Big]
    \le \sqrt{\Var_{\Prob^\pi}(e^{\beta R_1^\pi})}\,\sqrt{\sum_{h=1}^H\sum_{s,a}p^\pi_h(s,a)\rho_h(s,a)} .
  \]
  \item[(iii)] If $\rho_h \ge 0$ and $w_h\,c_h(S_h, A_h) \le C$ on every trajectory, then
  $\sum_h\E^\pi[w_h\,c_h(S_h, A_h)\,\rho_h(S_h, A_h)] \le C\sum_h\sum_{s,a}p^\pi_h(s,a)\rho_h(s,a)$.
\end{enumerate}
\end{lemma}

\begin{proof}
\uses{lem:unroll_recursion, lem:one_add_div_pow_le_exp, lem:cauchy_schwarz_traj, lem:entropic_variance_sum, lem:occupancy_expectation}
\leanok
(i) Lemma~\ref{lem:unroll_recursion} with $\kappa = 1 + c/H \ge 1$, and $\kappa^{h-1} \le e^c$ for
$1 \le h \le H$ (Lemma~\ref{lem:one_add_div_pow_le_exp}); the expectations are nonnegative.
(ii) Lemma~\ref{lem:cauchy_schwarz_traj} with the weights $w_h$, $v_h := \Var_{p_h}(Z^\pi_{h+1})(S_h, A_h) \ge 0$
and $u_h := \rho_h(S_h, A_h) \ge 0$; since $w_h^2 = e^{2\beta\sum_{i \le h} r_i(S_i,A_i)}$, the first factor is
$\Var_{\Prob^\pi}(e^{\beta R_1^\pi})$ by Lemma~\ref{lem:entropic_variance_sum}, and the second is the occupancy
sum by Lemma~\ref{lem:occupancy_expectation}. (iii) Integrate the pointwise bound
$w_h\,c_h\,\rho_h \le C\rho_h$ and apply Lemma~\ref{lem:occupancy_expectation}.
\end{proof}

\textbf{Lean remark.} These statements are MDP-general
(\texttt{LeanMachineLearning/ReinforcementLearning/MDP/UnrollBounds.lean}); they are recorded here
because they are the parts of Lemmas~\ref{lem:cert_unrolled_pos}, \ref{lem:ratio_bound_pos},
\ref{lem:cert_unrolled_neg} and~\ref{lem:ratio_bound_neg} that do not depend on the algorithm, and
would fit in Chapter~\ref{chap:pre_mdp}. The integrability side conditions are
\texttt{integrable\_exp\_sum\_mul} and \texttt{integrable\_comp\_obs\_action} (bounded functions of finitely
many rounds).

\begin{lemma}[Unrolled certificate]\label{lem:cert_unrolled_pos}
\lean{Essakine2026Tight.cert_le_exp_mul_sum_integral_of_pos, Essakine2026Tight.certAt_le_exp_mul_sum_integral_of_pos}\leanok
\uses{def:good_event, def:certificate, def:rate_min, def:exp_value, def:step_law, def:greedy, def:entropic_bpi}
On $\evGood^+$, for all $t \ge 0$, with $\pi = \pi^{t+1}$ and $\E^\pi = \E^\pi_{(s_1, 1)}$,
\[
  (\pi_1\cert_1^t)(s_1) \le e^{13}\,\E^\pi\Bigg[\sum_{h=1}^H e^{\beta\sum_{i \le h} r_i(S_i, A_i)}
  \Big(36\sqrt{\Var_{p_h}(Z^\pi_{h+1})(S_h, A_h)\,\rho^{*t}_h(S_h, A_h)}
  + 84\,H\,c_h\,\rho^t_h(S_h, A_h)\Big)\Bigg] .
\]
\end{lemma}

\begin{proof}
\uses{lem:cert_true_recursion_pos, lem:unroll_recursion, lem:one_add_div_pow_le_exp, def:certificate, lem:pos_core_unroll}
\leanok
Fix $\omega \in \evGood^+$ and $t$. Apply Lemma~\ref{lem:unroll_recursion} to the policy $\pi$, the
functions $g_h := \pi_h\cert_h^t$ ($h \le H + 1$; $g_{H+1} = 0$ by Definition~\ref{def:certificate}), the
nonnegative functions
$u_h(s,a) := 36\sqrt{\Var_{p_h}(Z^\pi_{h+1})(s,a)\rho^{*t}_h(s,a)} + 84Hc_h\rho^t_h(s,a)$ and
$\kappa := 1 + \frac{13}{H} \ge 1$: the hypothesis
$g_h(s) \le e^{\beta r_h(s, \pi_h(s))}[u_h(s, \pi_h(s)) + \kappa(p_hg_{h+1})(s, \pi_h(s))]$ for all
$h \le H$ and $s$ is Lemma~\ref{lem:cert_true_recursion_pos}, and the rewards are nonnegative. The
conclusion is $g_1(s_1) \le \sum_{h=1}^H\kappa^{h-1}\E^\pi[e^{\beta\sum_{i \le h} r_i(S_i,A_i)}u_h(S_h,A_h)]$, and
$\kappa^{h-1} \le \kappa^H \le e^{13}$ by Lemma~\ref{lem:one_add_div_pow_le_exp} (with $a = 13$,
$0 \le h - 1 \le H$); all the summands are nonnegative.
\end{proof}

\textbf{Lean remark.} The expectation is an integral against
$\mathrm{stepLaw}\ M\ \pi\ (s_1, 0)$ of a bounded measurable function of the first $H$ rounds;
the rate terms $\rho^t_h(S_h, A_h)$ are, for the fixed $\omega$, deterministic functions of $(h, s, a)$
(the history is frozen), so no measurability issue in $\omega$ arises inside the expectation.
The proof is Lemma~\ref{lem:pos_core_unroll}(i) with $c = 13$ and $g_h$ the certificate. For the
$0$-indexed Lean step $h$ the clip is \texttt{Real.exp (β * (H - h : ℕ))} ($= c_{h+1}$ in the
$1$-indexed notation above) and the weight is $e^{\beta\sum_{i \in \mathrm{Iic}\,h} r_i(S_i, A_i)}$.

\begin{lemma}[Normalized certificate bound]\label{lem:ratio_bound_pos}
\lean{Essakine2026Tight.cert_div_optZ_fst_le_of_pos, Essakine2026Tight.certAt_div_ZtildeAt_le_of_pos}\leanok
\uses{def:good_event, def:certificate, def:backups, def:rate_min, def:occupancy, def:gmax, def:upper_bound_const, def:greedy, def:entropic_bpi}
Let $V_G := (e^{\beta\Gmax} - 1)^2/e^{\beta\Gmax}$ with $\Gmax = \Gmax(\mathcal M)$ (Definition~\ref{def:gmax};
this is the $V_G$ of Definition~\ref{def:upper_bound_const}). On $\evGood^+$, for all $t \ge 0$, with
$\pi = \pi^{t+1}$,
\[
  \frac{(\pi_1\cert_1^t)(s_1)}{\Zt_1^t(s_1)} \le e^{13}\Big[36\sqrt{V_G\,x_t} + 84\,H\,e^{\beta(H+1)}\,y_t\Big],
  \quad
  x_t := \sum_{h=1}^H\sum_{s,a} p^\pi_h(s,a)\rho^{*t}_h(s,a), \quad
  y_t := \sum_{h=1}^H\sum_{s,a} p^\pi_h(s,a)\rho^t_h(s,a),
\]
where $p^\pi_h$ is the occupancy measure of $\pi$ from $s_1$ (Definition~\ref{def:occupancy}).
\end{lemma}

\begin{proof}
\uses{lem:cert_unrolled_pos, lem:optimism_pos, lem:opt_exp_value_eq, lem:exp_value_range, lem:cauchy_schwarz_traj, lem:entropic_variance_sum, lem:variance_ratio_le, lem:gmax_le_horizon, lem:occupancy_expectation, def:entropic_variance, lem:pos_core_unroll}
\leanok
Fix $\omega \in \evGood^+$ and $t$. \emph{Denominator.} $\Zt_1^t(s_1) \ge Z^*_1(s_1) \ge Z^\pi_1(s_1) \ge 1 > 0$
(Lemma~\ref{lem:optimism_pos}, Lemma~\ref{lem:opt_exp_value_eq} for $\beta > 0$,
Lemma~\ref{lem:exp_value_range}), so the left-hand side is at most $(\pi_1\cert_1^t)(s_1)/Z^\pi_1(s_1)$,
which we bound by dividing Lemma~\ref{lem:cert_unrolled_pos} by $Z^\pi_1(s_1)$.

\emph{Variance term.} Put $w_h := e^{\beta\sum_{i \le h} r_i(S_i,A_i)}$ (a nonnegative function of the
trajectory up to step $h$), $v_h := \Var_{p_h}(Z^\pi_{h+1})(S_h,A_h)/Z^\pi_1(s_1)^2$ and
$u_h := \rho^{*t}_h(S_h,A_h)$. Lemma~\ref{lem:cauchy_schwarz_traj} (Cauchy--Schwarz for $\sum_h\E^\pi$) gives
\[
  \E^\pi\Big[\sum_h w_h\sqrt{v_hu_h}\Big] \le \sqrt{\E^\pi\Big[\sum_h w_h^2v_h\Big]}\sqrt{\E^\pi\Big[\sum_h u_h\Big]} .
\]
By Lemma~\ref{lem:entropic_variance_sum},
$\E^\pi[\sum_h w_h^2 v_h] = \Var_{\Prob^\pi}(e^{\beta R_1^\pi})/Z^\pi_1(s_1)^2$, and by
Lemma~\ref{lem:variance_ratio_le} with $G = \Gmax$ (admissible since $R_1^\pi \in [0, \Gmax]$ a.s.\ by
Lemma~\ref{lem:gmax_le_horizon}), this ratio is at most $V_G$. By Lemma~\ref{lem:occupancy_expectation},
$\E^\pi[\sum_h u_h] = x_t$.

\emph{Rate term.} Since $0 \le r_i \le 1$ and $\beta > 0$, $w_h \le e^{\beta h}$, and $e^{\beta h}c_h = e^{\beta(H+1)}$;
with $Z^\pi_1(s_1) \ge 1$ and Lemma~\ref{lem:occupancy_expectation},
$\E^\pi[\sum_h w_h\,84Hc_h\rho^t_h(S_h,A_h)]/Z^\pi_1(s_1) \le 84He^{\beta(H+1)}\,y_t$.
\end{proof}

\textbf{Lean remark.} Lemma~\ref{lem:cauchy_schwarz_traj} allows the weight $w_h$ to depend on
the whole trajectory up to step $h$ (here on the past rewards), not only on $(S_h, A_h)$, which
is needed here. The identity of Lemma~\ref{lem:entropic_variance_sum} is where the entropic
variance $\sigma V^\pi_1(s_1)$ of Definition~\ref{def:entropic_variance} is identified with the
variance of $e^{\beta R_1^\pi}$ under $\Prob^\pi_{(s_1,1)}$. The variance and rate terms are
Lemma~\ref{lem:pos_core_unroll}(ii) and~(iii) (with $c_h = 84He^{\beta(H+1-h)}$ and
$C = 84He^{\beta(H+1)}$). The Lean statement writes $e^{|\beta|(H+1)}$ (equal to $e^{\beta(H+1)}$ here), so
that both signs produce the same hypothesis for Lemma~\ref{lem:pos_core_stopping_time}; the
ratio uses the greedy policy of the history, and the occupancy sums are
$\sum_{h, s, a}\texttt{occupancy M π s₁ h s a} \cdot \rho$.

\section{Bounding the stopping time}
\label{sec:pos_stopping_time}

\begin{lemma}[Progress at a non-stopping episode]\label{lem:progress_pos}
\lean{Essakine2026Tight.progress_le_cert_div_optZ_fst_of_pos, Essakine2026Tight.progress_le_certAt_div_ZtildeAt_of_pos}\leanok
\uses{def:good_event, def:stopping_rule, def:certificate, def:backups, def:upper_bound_const, def:entropic_bpi}
Let $\kappa(\beta, \eps) := (e^{|\beta|\eps} - 1)e^{-2|\beta|\eps}/2$ (Definition~\ref{def:upper_bound_const}).
On $\evGood^+$, for every $t \ge 0$ such that the stopping condition of Definition~\ref{def:stopping_rule}
\emph{fails} after $t$ episodes,
\[
  \kappa(\beta, \eps) \le \frac{(\pi_1\cert_1^t)(s_1)}{\Zt_1^t(s_1)} .
\]
\end{lemma}

\begin{proof}
\uses{def:stopping_rule, lem:backups_range}
\leanok
Failure of the condition means $(\pi_1\cert_1^t)(s_1) > \frac{e^{\beta\eps} - 1}{e^{\beta\eps}}\Zt_1^t(s_1)$, and
$\Zt_1^t(s_1) \ge 1 > 0$ (Lemma~\ref{lem:backups_range}), so the ratio exceeds
$(e^{\beta\eps} - 1)e^{-\beta\eps} \ge (e^{\beta\eps} - 1)e^{-2\beta\eps}/2 = \kappa(\beta,\eps)$ (as $e^{-\beta\eps} \le 1$ and $\frac12 \le 1$;
here $|\beta| = \beta$). No event is actually needed.
\end{proof}

\textbf{Lean remark.} No event is assumed; $\eps \ge 0$ and $\delta \in (0, 1]$ suffice ($\delta$ enters
through the range $\Zt_1^t(s_1) \ge 1$ of Lemma~\ref{lem:backups_range}).

\begin{lemma}[Core of the bound on the stopping time]\label{lem:pos_core_stopping_time}
\lean{Essakine2026Tight.natCast_le_upperBound, Essakine2026Tight.natCast_le_upperBound_of_mul_progress_le, ENat.toENNReal_le_ofReal_of_forall_natCast_le, Essakine2026Tight.one_le_upperBound, Essakine2026Tight.one_le_coeffLin, Essakine2026Tight.progress_pos, Essakine2026Tight.progress_le_one, Essakine2026Tight.progress_mul_coeffSqrt, Essakine2026Tight.progress_mul_coeffLin, Essakine2026Tight.varianceFactor_nonneg, Essakine2026Tight.exp_thirteen_le}\leanok
\uses{def:upper_bound_const, def:rates}
Let $\beta \ne 0$, $\eps > 0$, $\delta \in (0, 1]$, integers $S, A \ge 1$, $H \ge 0$, $G \in \R$, $V_G$, $\kappa = \kappa(\beta, \eps)$ as in
Definition~\ref{def:upper_bound_const}.
\begin{enumerate}
  \item[(i)] Let $T \in \N$ and $x_t \ge 0$, $y_t \in \R$ ($t < T$) be such that
  $\kappa \le e^{13}[36\sqrt{V_Gx_t} + 84He^{|\beta|(H+1)}y_t]$ for every $t < T$,
  $\sum_{t<T}x_t \le 16SAH\,\alpha^*(T-1,\delta)\log(T+1)$ and $\sum_{t<T}y_t \le 16SAH\,\alpha(T-1,\delta)\log(T+1)$.
  Then $T \le \mathrm{upperBound}(S, A, H, \beta, \eps, \delta, G)$.
  \item[(ii)] If $\tau \in \N \cup \{\infty\}$, $U \in \R$ and every integer $T \le \tau$ satisfies $T \le U$, then
  $\tau \le U$; in particular $\tau < \infty$.
\end{enumerate}
\end{lemma}

\begin{proof}
\uses{lem:self_bounding, lem:rates_props}
\leanok
(i) This is Steps~1--3 of the proof of Lemma~\ref{lem:stopping_time_bound_pos}, with the sequences
$x_t$, $y_t$ abstracted. If $T = 0$ the claim is $0 \le \mathrm{upperBound}$, and $\mathrm{upperBound} \ge 1$
because $A_0 \ge 0$ ($3SAH/\delta \ge 3$, or $\log 0 = 0$ if $H = 0$), $C \ge 0$ and $D \ge 0$. If $T \ge 1$ and
$H = 0$, the counting hypotheses give $x_0 \le 0$ and the hypothesis at $t = 0$ gives $\kappa \le 0$,
contradicting $\kappa > 0$. If $T, H \ge 1$: with $L = \log(8eT)$, Steps~1--3 give
$T\kappa \le 144e^{13}\sqrt{V_GSAH}\sqrt{T(A_0L + L^2)} + 1344e^{13}e^{|\beta|(H+1)}H^2SA(A_0L + SL^2)$
(the nonnegativity of $y_t$ is not needed), and $e^{13} \le 3^{13} = 1594323$ gives
$144e^{13} \le 10^9$ and $1344e^{13} \le 10^{10}$. Since $\kappa C = 10^9\sqrt{V_GSAH}$ and
$\kappa D = 10^{10}e^{|\beta|(H+1)}H^2SA$, dividing by $\kappa > 0$ gives the hypothesis of
Lemma~\ref{lem:self_bounding} ($A = A_0 \ge 0$, $B = 1$, $E = S \ge 1$, $\alpha = 8e$, $C \ge 0$, $D \ge 1$ because
$\kappa \le 1$ and $H^2SA \ge 1$), whose conclusion is $T \le \mathrm{upperBound}$.
(ii) If $\tau = \infty$, the integer $T = \lfloor U\rfloor_+ + 1 \le \tau$ satisfies $T > U$; otherwise take $T = \tau$.
\end{proof}

\begin{lemma}[Bound on the stopping time]\label{lem:stopping_time_bound_pos}
\lean{Essakine2026Tight.stoppingTime_le_upperBound_of_pos}\leanok
\uses{def:good_event, def:entropic_bpi, def:upper_bound_const, def:gmax}
At every $\omega \in \evGood^+$ at which the policies played are the greedy policies
($X_t(\omega) = \pi^{t+1}$ for all $t$, which holds almost surely by Lemma~\ref{lem:entropic_bpi_run}),
the stopping time of the run satisfies
$\tau \le \mathrm{upperBound}(S, A, H, \beta, \eps, \delta, \Gmax(\mathcal M))$ (Definition~\ref{def:upper_bound_const});
in particular $\tau < \infty$.
\end{lemma}

\begin{proof}
\uses{lem:progress_pos, lem:ratio_bound_pos, lem:sum_counts_bound, lem:entropic_bpi_run, lem:self_bounding, def:rates, def:upper_bound_const, lem:rates_props, def:rate_min, lem:pos_core_stopping_time}
\leanok
Fix $\omega \in \evGood^+$ (in particular $\omega \in \evCnt$) and write $\kappa := \kappa(\beta,\eps) > 0$,
$A_0 := \log(3SAH/\delta) > 0$, $C := 10^9\sqrt{V_G SAH}/\kappa$ and $D := 10^{10}e^{\beta(H+1)}H^2SA/\kappa$ (the
constants of Definition~\ref{def:upper_bound_const}, with $|\beta| = \beta$).

\emph{Step 1: an inequality for every $T \le \tau$.} Let $T \ge 1$ be an integer with $T \le \tau$
(i.e.\ the stopping condition fails after $t$ episodes for every $t < T$: by
Lemma~\ref{lem:entropic_bpi_run}, $\tau$ is the first $t$ at which it holds). Summing
Lemma~\ref{lem:progress_pos} and Lemma~\ref{lem:ratio_bound_pos} over $t = 0, \dots, T - 1$,
\[
  T\kappa \le \sum_{t < T}\frac{(\pi_1\cert_1^t)(s_1)}{\Zt_1^t(s_1)}
  \le 36e^{13}\sqrt{V_G}\sum_{t < T}\sqrt{x_t} + 84e^{13}He^{\beta(H+1)}\sum_{t < T}y_t .
\]
By the Cauchy--Schwarz inequality, $\sum_{t < T}\sqrt{x_t} \le \sqrt{T}\sqrt{\sum_{t < T}x_t}$.

\emph{Step 2: the counting bounds.} $x_t$ and $y_t$ are the inner sums of
Lemma~\ref{lem:sum_counts_bound} (with the policy $\pi^{t+1}$, the counts $n_h^t$ of the run and the
rates $\alpha^*$, resp.\ $\alpha$), so on $\evCnt$,
$\sum_{t < T}x_t \le 16SAH\,\alpha^*(T-1,\delta)\log(T+1)$ and $\sum_{t < T}y_t \le 16SAH\,\alpha(T-1,\delta)\log(T+1)$.
Put $L := \log(8eT)$; for $T \ge 1$, $\log(T+1) \le L$ and, by Definition~\ref{def:rates},
$\alpha^*(T-1,\delta) = A_0 + L$ and $\alpha(T-1,\delta) = A_0 + SL$. Hence
\[
  \sum_{t < T}x_t \le 16SAH\,(A_0L + L^2), \qquad \sum_{t < T}y_t \le 16SAH\,(A_0L + SL^2) .
\]

\emph{Step 3: the self-bounding inequality.} Combining,
\[
  T \le \frac{144\,e^{13}\sqrt{V_G SAH}}{\kappa}\sqrt{T(A_0L + L^2)}
     + \frac{1344\,e^{13}e^{\beta(H+1)}H^2SA}{\kappa}\,(A_0L + SL^2) .
\]
Numerically $e^{13} \le 4.43\cdot 10^5$, so $144\,e^{13} \le 6.4\cdot 10^7 \le 10^9$ and
$1344\,e^{13} \le 6.0\cdot 10^8 \le 10^{10}$; both brackets are nonnegative, so
\[
  T \le C\sqrt{T(A_0L + L^2)} + D\,(A_0L + SL^2), \qquad L = \log(8eT) .
\]
This is the hypothesis of Lemma~\ref{lem:self_bounding} with $\tau := T$, $A := A_0$, $B := 1$,
$E := S$, $\alpha := 8e \ge e$ (so that $\log(\alpha\tau) = L$) and the constants $C \ge 0$, $D$: indeed
$E \ge B$, $A_0 > 0$, and $D \ge 1$ (because $\kappa \le 1/8$: $(u - 1)/u^2 \le 1/4$ for $u > 1$).
Its conclusion is
\[
  T \le C^2(A_0 + 1)C_1^2 + \big(D + 2\sqrt{D}C\big)(A_0 + S)C_1^2 + 1 = \mathrm{upperBound}(S,A,H,\beta,\eps,\delta,\Gmax),
\]
with the $C_1$ of Lemma~\ref{lem:self_bounding} for $\alpha = 8e$, $A + E = A_0 + S$ and the constants
$C$, $D$, which is the $C_1$ of Definition~\ref{def:upper_bound_const} (Remark~\ref{rem:pos_c1_constant}).

\emph{Step 4: conclusion.} If $\tau = \infty$, every integer $T \ge 1$ satisfies $T \le \tau$, so every
$T \ge 1$ is at most the finite number $\mathrm{upperBound}$: contradiction. Hence $\tau < \infty$; if
$\tau \ge 1$, take $T = \tau$ in Step~3; if $\tau = 0$ the bound holds since $\mathrm{upperBound} \ge 1$.
\end{proof}

\textbf{Lean remark.} $\tau$ is \texttt{IdentAlg.stoppingTime}, with values in \texttt{ℕ∞}; the
statement ``for every $t < \tau$ the stopping condition fails'' is the characterization of the
hitting time (\texttt{hittingAfter}, Lemma~\ref{lem:entropic_bpi_run}), valid also when $\tau = \infty$.
The proof is organized as: (a) a lemma ``for all $T \ge 1$ with $(T : \texttt{ℕ∞}) \le \tau$,
$T \le \mathrm{upperBound}$'' proved on $\evGood^+$ by Steps~1--3; (b) the deduction $\tau \le \mathrm{upperBound}$
by a case split on $\tau = \top$. The step $T\kappa \le \sum_t(\cdot)$ is a sum of $T$ inequalities each
valid at an $\omega$ of the event, the sums $\sum_{t<T}x_t$ are \texttt{Finset.sum (Finset.range T)},
and the final numeric inequalities only need $e^{13} \le 3^{13}$ (from \texttt{Real.exp\_one\_lt\_three}:
$144 \cdot 3^{13} \approx 2.3\cdot 10^8$ and $1344 \cdot 3^{13} \approx 2.1\cdot 10^9$). When
$\Gmax(\mathcal M) = 0$ (all returns vanish a.s.), $V_G = 0$ and $C = 0$, which Lemma~\ref{lem:self_bounding}
allows. Steps~1--3 and~4 are Lemma~\ref{lem:pos_core_stopping_time}(i) and~(ii), with
$x_t = \sum_{h,s,a}p^{X_t}_h(s,a)\rho^{*t}_h(s,a)$ and $y_t$ likewise: the counting bound of
Lemma~\ref{lem:sum_counts_bound} is about the policies $X_t$ played, the normalized certificate bound about
the greedy policies $\pi^{t+1}$, hence the hypothesis $X_t(\omega) = \pi^{t+1}$ in the statement (Lean:
\texttt{hX : ∀ t, X t ω = greedyAt X Y r β δ t ω}). The Lean hypotheses are $\delta \in (0, 1]$ and
$\eps > 0$; the conclusion is \texttt{(τ ω : ℝ≥0∞) ≤ ENNReal.ofReal (upperBound …)}, and no hypothesis that
$(X, Y)$ is a run is needed (the characterization of $\tau$ as a hitting time is pointwise).

\begin{remark}[The constant $C_1$]\label{rem:pos_c1_constant}
\uses{lem:self_bounding, def:upper_bound_const}
The paper's Lemma~29 states the self-bounding inequality with
$C_1 = \frac85\log(11\alpha^2(A + E)(C + D))$, which is false
(Remark~\ref{rem:pana_self_bounding_paper}). Chapter~\ref{chap:pre_analysis} proves it with
$C_1 = \frac85\log(11\alpha^2(A + E)(C + \sqrt D)^2)$, and with $(C + D)^2$ in place of
$(C + \sqrt D)^2$ when $D \ge 1$ (which holds here); Definition~\ref{def:upper_bound_const}
uses this last form with $\alpha = 8e$, $A + E = A_0 + S$, so that the two agree. For
$C + D \ge 1$ this $C_1$ is larger than the paper's by $\frac85\log(C + D)$. Nothing else in
this chapter depends on the form of $C_1$; the asymptotic form $\tilde O(\kappa^{-2}V_GSAH)$ of
the bound is the same in both cases.
\end{remark}

\begin{lemma}[From an event of probability at least $1 - \delta$ to the probability statements]\label{lem:pos_core_pac}
\lean{Essakine2026Tight.one_sub_le_measureReal_stoppingTime_le, Essakine2026Tight.measureReal_bad_le, Essakine2026Tight.stoppingTime_ne_top_of_le}\leanok
\uses{def:entropic_bpi, def:stopping_rule}
Let $(X, Y, \mathrm{out})$ be a run of Entropic-BPI (with any parameters) in any environment, $E \subseteq \Omega$ with
$\Prob(E^c) \le \delta$, $U \in \R$, and call $\omega$ \emph{greedy} if $X_t(\omega) = \pi^{t+1}$ for all $t$.
\begin{enumerate}
  \item[(i)] If $\tau(\omega) \le U$ for every greedy $\omega \in E$, then $\Prob(\tau \le U) \ge 1 - \delta$.
  \item[(ii)] If $\tau(\omega) < \infty$ for every greedy $\omega \in E$, and for every $\omega \in E$ and $t$ such that the
  stopping condition holds after $t$ episodes $V^*_1(s_1) - V^{\pi^{t+1}}_1(s_1) \le \eps$, then
  $\Prob(V^*_1(s_1) - V^{\mathrm{out}}_1(s_1) > \eps) \le \delta$.
\end{enumerate}
\end{lemma}

\begin{proof}
\uses{lem:entropic_bpi_run}
\leanok
Let $N$ be the null set of Lemma~\ref{lem:entropic_bpi_run}. (i) $\{\tau \le U\}^c \setminus N \subseteq E^c$, so
$\Prob(\{\tau \le U\}^c) \le \delta$ and $\Prob(\tau \le U) \ge 1 - \Prob(\{\tau \le U\}^c)$ (subadditivity, no
measurability needed). (ii) For $\omega \in E \setminus N$, $\tau(\omega) < \infty$, the stopping condition holds after
$\tau(\omega)$ episodes and $\mathrm{out}(\omega) = \pi^{\tau(\omega)+1}$, so $\omega$ is not in the bad event; hence the bad
event is contained in $E^c \cup N$.
\end{proof}

\begin{lemma}[PAC guarantee and sample complexity of Entropic-BPI for $\beta > 0$]\label{lem:pac_pos}
\lean{Essakine2026Tight.measureReal_bad_le_of_pos, Essakine2026Tight.one_sub_le_measureReal_stoppingTime_le_of_pos}\leanok
\uses{def:entropic_bpi, def:entropic_value, def:opt_value, def:upper_bound_const, def:gmax, def:good_event}
In the standing setting,
\[
  \Prob\big(V^*_1(s_1) - V^{\mathrm{out}}_1(s_1) > \eps\big) \le \delta
  \qquad\text{and}\qquad
  \Prob\big(\tau \le \mathrm{upperBound}(S,A,H,\beta,\eps,\delta,\Gmax(\mathcal M))\big) \ge 1 - \delta ,
\]
(on the event $\{\tau = \infty\}$, of probability at most $\delta$ by the second statement, the
output $\mathrm{out}$ is whatever policy the run assigns).
\end{lemma}

\begin{proof}
\uses{lem:stopping_time_bound_pos, lem:entropic_bpi_run, lem:pac_at_stopping_pos, lem:good_event, lem:pos_core_pac}
\leanok
By Lemma~\ref{lem:good_event}, $\Prob(\evGood^+) \ge 1 - \delta$, so it suffices to show that, up to a
$\Prob$-null set $N$, $\evGood^+ \setminus N$ is contained in both events. Let $N$ be the null set
outside of which the conclusions of Lemma~\ref{lem:entropic_bpi_run} hold (the actions are the
greedy policies of the histories, the stopping condition holds at $\tau$ when $\tau < \infty$, and
$\mathrm{out} = \pi^{\tau + 1}$ on $\{\tau < \infty\}$). For $\omega \in \evGood^+ \setminus N$:
Lemma~\ref{lem:stopping_time_bound_pos} gives $\tau(\omega) \le \mathrm{upperBound} < \infty$; then the
stopping condition holds after $t := \tau(\omega)$ episodes and $\mathrm{out}(\omega) = \pi^{t+1}$, so
Lemma~\ref{lem:pac_at_stopping_pos} gives $V^*_1(s_1) - V^{\mathrm{out}(\omega)}_1(s_1) \le \eps$.
\end{proof}

\textbf{Lean remark.} This is the run-level statement (\texttt{IdentAlg.IsRun}) that
Chapter~\ref{chap:sample_complexity} converts into the PAC property of the algorithm
(\texttt{IdentAlg.IsPAC} is about the law of the output, Definition~\ref{def:entropic_pac}) and into
Theorem~\ref{thm:upper_bound_complexity}. The measurability of the events $\{V^*_1 - V^{\mathrm{out}}_1 > \eps\}$
and $\{\tau \le \mathrm{upperBound}\}$ is immediate: the policy type is finite and $\tau$ is a stopping time.
The proof needs $\evGood^+$ only through its three defining properties, and the inclusion
$\evGood^+ \setminus N \subseteq \{\tau \le \mathrm{upperBound}\} \cap \{V^*_1 - V^{\mathrm{out}}_1 \le \eps\}$ is proved
pointwise. In Lean the two statements are two lemmas for a run in \texttt{statesEnv M s₁}, with $\delta \in (0, 1]$
and $\eps > 0$, each an instance of Lemma~\ref{lem:pos_core_pac} with $E = \evGood^+$ and
$\Prob((\evGood^+)^c) \le \delta$ (\texttt{measure\_compl\_goodEvent\_le}, Lemma~\ref{lem:good_event}); in fact no
measurability of the events is used (outer measure and complements).

\chapter{Analysis of Entropic-BPI for \texorpdfstring{$\beta < 0$}{beta < 0}}
\label{chap:analysis_neg}

This chapter reproduces Appendix~B.2 of \cite{essakine2026tight} (Lemmas~12--16 and the proof
of the sample complexity for $\beta < 0$). It is the mirror image of
Chapter~\ref{chap:analysis_pos}, with the same lemma structure and proofs: for $\beta < 0$ the map
$x \mapsto e^{\beta x}$ is decreasing, so the optimal exponential value is a \emph{minimum} over
policies and actions (Lemmas~\ref{lem:opt_exp_value_eq}, \ref{lem:opt_bellman}), the values live in
$[e^{\beta(H+1-h)}, 1]$, the optimistic backup $\Ut^t$ is clipped at $1$ and the pessimistic one at
$e^{\beta(H+1-h)}$, the greedy policy minimizes $\Ul^t$ (deviation~2), and the ring value $\Zr^t$ is an
\emph{upper} bound on both $\Zt^t$ and $Z^{\pi^{t+1}}$, the certificate dominating $\Zr^t - \Zl^t$.
The reader is referred to Chapter~\ref{chap:analysis_pos} for the discussion of the constants;
the differences are pointed out in remarks.

\section{Standing setting and notation}
\label{sec:neg_setting}

Throughout this chapter: $\beta < 0$; $\mathcal M$ is a finite-horizon MDP
(Definition~\ref{def:episodic_mdp}) with reward function $r$ with values in $[0, 1]$
(Definition~\ref{def:reward_fn}); $\delta \in (0, 1)$ and $\eps > 0$; $(X, Y, \mathrm{out})$ is a run
of Entropic-BPI (Definition~\ref{def:entropic_bpi}) with parameters $(r, \beta, \delta, \eps, s_1)$
in the episode environment of $\mathcal M$ from $s_1$ (Definition~\ref{def:episode_env}), on a
probability space $(\Omega, \mathcal F, \Prob)$, with stopping time $\tau$. The quantities $n_h^t$,
$\phat_h^t$, $b_h^t$, $\Ut_h^t$, $\Ul_h^t$, $\Zt_h^t$, $\Zl_h^t$, $\pi^{t+1}$, $\cert_h^t$, $\Ur_h^t$, $\Zr_h^t$,
$\rho_h^t$, $\rho^{*t}_h$ are those of Definitions~\ref{def:counts}, \ref{def:emp_trans}, \ref{def:bonus},
\ref{def:backups}, \ref{def:greedy}, \ref{def:certificate}, \ref{def:ring}, \ref{def:rate_min}
(their $\beta < 0$ versions) for the history of the first $t$ episodes. The abbreviations
$\pi := \pi^{t+1}$, $n$, $\phat$, $b$, $\alpha$, $\alpha^*$, $\rho$, $\rho^*$, $\phat f$, $(p_hf)(s,a)$,
$\Var_{\phat}$, $\Var_{p_h}$ are as in Section~\ref{sec:pos_setting}. We write
\[
  \ell_h := e^{\beta(H + 1 - h)} \in (0, 1] \quad\text{(the lower clip of step $h$)}, \qquad
  B_h := 1 - \ell_{h+1} = 1 - e^{\beta(H-h)} \in [0, 1) \quad\text{(the range parameter of step $h$)},
\]
so that $\ell_{H+1} = 1$, $B_H = 0$, $\ell_h \le \ell_{h+1}$ and $B_h \le 1$; the upper clip of $\Ut^t_h$ and
of the certificate at every step is $1$. The ranges used all the time are: $U^*_h \in [\ell_h, 1]$ and
$Z^*_{h+1} \in [\ell_{h+1}, 1]$ (Lemma~\ref{lem:opt_exp_value_range}); $Z^\pi_{h+1} \in [\ell_{h+1}, 1]$ for every
policy $\pi$ (Lemma~\ref{lem:exp_value_range}); $\ell_{h+1} \le \Zl^t_{h+1} \le \Zt^t_{h+1} \le 1$ and
$\ell_h \le \Ul^t_h \le \Ut^t_h \le 1$ (Lemma~\ref{lem:backups_range}); $0 \le \pi_{h+1}\cert^t_{h+1} \le 1$
(Lemma~\ref{lem:certificate_range}). The good event is
$\evGood^- = \evKL \cap \evBern(Z^*, B) \cap \evCnt$ with $B = (B_h)_{h \le H}$
(Definition~\ref{def:good_event}), and ``on $\evGood^-$'' means ``for every $\omega \in \evGood^-$''.
For $\beta < 0$, $e^{\beta r_h(s,a)} \in (0, 1]$: the exponential Bellman backups are contractions. This
is harmless everywhere except in Lemma~\ref{lem:cert_true_recursion_neg}, where the factor
$e^{\beta r_h}$ cannot absorb the clip $1$ of a never-visited pair (Remark~\ref{rem:neg_additive_term}).

\textbf{Lean remark.} The same Lean definitions serve both signs (\texttt{if 0 < β then \ldots else \ldots}),
so the lemmas of this chapter are the \texttt{β < 0} branches of the corresponding declarations, with
the same names suffixed differently; the induction on $j = H + 1 - h$ and the treatment of the
never-visited pairs are as in Chapter~\ref{chap:analysis_pos}. The clips of a never-visited pair
are $\Ut = 1$, $\Ul = \ell_h$, the certificate is $1$ and the rate terms are $\rho = \rho^* = 1$. The
ranges $\ell_h$, $B_h$ are \texttt{Real.exp (β * j)} and \texttt{1 - Real.exp (β * (j - 1))}.

\section{Concentration and optimism}
\label{sec:neg_optimism}

\begin{lemma}[Concentration of the optimal exponential value; Lemma~12 of the paper]\label{lem:concentration_neg}
\uses{def:good_event, def:backups, def:emp_trans, def:counts, def:rates, def:opt_value, def:entropic_bpi}
On $\evGood^-$, for all $t \ge 0$, $h \in \{1, \dots, H\}$ and $(s,a)$ with $n = n_h^t(s,a) \ge 1$,
\[
  \abs{(p_h - \phat) Z^*_{h+1}(s,a)}
  \le 2\sqrt2 \sqrt{\Var_{\phat}(\Zl^t_{h+1}) \frac{\alpha^*}{n}}
  + 5 B_h \frac{\alpha}{n} + 4 H B_h \frac{\alpha}{n}
  + \frac1H \phat \abs{Z^*_{h+1} - \Zl^t_{h+1}} ,
\]
i.e.\ $\abs{(p_h - \phat) Z^*_{h+1}(s,a)} \le b + \frac1H \phat\abs{Z^*_{h+1} - \Zl^t_{h+1}}$ with the bonus
$b = b_h^t(s,a)$ of Definition~\ref{def:bonus} for $\beta < 0$ (variance of the pessimistic value,
factor $B_h = 1 - e^{\beta(H-h)}$, third term in $\alpha$; Remark~\ref{rem:neg_constants}).
\end{lemma}

\begin{proof}
\uses{def:good_event, def:event_bern, def:event_kl, lem:opt_exp_value_range, lem:backups_range, lem:kl_transport_var, lem:transport_var, lem:sqrt_add_le, lem:sqrt_mul_le, lem:rates_props, lem:emp_trans_simplex}
Identical to the proof of Lemma~\ref{lem:concentration_pos} with $f := Z^*_{h+1} - \ell_{h+1}$ and
$g := \Zl^t_{h+1} - \ell_{h+1}$, both with values in $[0, B_h]$ (Lemmas~\ref{lem:opt_exp_value_range}
and~\ref{lem:backups_range}), and the range $B_h$: the Bernstein event gives
$\abs{(p_h - \phat)Z^*_{h+1}(s,a)} \le \sqrt{2\Var_{p_h}(Z^*_{h+1})(s,a)\alpha^*/n} + 3B_h\alpha^*/n$
(Definitions~\ref{def:event_bern}, \ref{def:good_event}); the KL event and
Lemma~\ref{lem:kl_transport_var} give $\Var_{p_h}(Z^*_{h+1}) \le 2\Var_{\phat}(Z^*_{h+1}) + 4B_h^2\alpha/n$;
Lemma~\ref{lem:transport_var} gives $\Var_{\phat}(Z^*_{h+1}) \le 2\Var_{\phat}(\Zl^t_{h+1}) + 2B_h\phat\abs{Z^*_{h+1} - \Zl^t_{h+1}}$;
Lemma~\ref{lem:sqrt_add_le}, Lemma~\ref{lem:sqrt_mul_le} with $\eta = 2/H$ for the middle term and
with $\eta = 1$ for $\sqrt{\alpha\alpha^*} \le (\alpha + \alpha^*)/2$, and $\alpha^* \le \alpha$
(Lemma~\ref{lem:rates_props}) give the bound with the coefficients
$(\sqrt2 + 3)B_h\alpha/n \le 5B_h\alpha/n$ and $(2H + \sqrt2)B_h\alpha/n \le 4HB_h\alpha/n$. When $h = H$ all
the quantities involved vanish ($Z^*_{H+1} = \Zl^t_{H+1} = 1$, $B_H = 0$) and the statement is $0 \le 0$.
\end{proof}

\begin{remark}[Constants and the bonus for $\beta < 0$]\label{rem:neg_constants}
\uses{def:bonus, lem:concentration_neg}
The remarks of Chapter~\ref{chap:analysis_pos} apply verbatim: the constants $2\sqrt2$, $5$, $4H$
hold (Remark~\ref{rem:pos_constants_bonus}), the third term of the bonus must use $\alpha(n, \delta)$,
as in the statement of Lemma~12 of the paper (Remark~\ref{rem:pos_constants_certificate}), and
the statement of the paper's Lemma~12 has two typos corrected here: its second term lacks the
division by $n$ (deviation~4) and its exponent reads $H - hH$. The variance in the bonus is that of
the \emph{pessimistic} value $\Zl^t_{h+1}$: for $\beta < 0$ the greedy policy is greedy for $\Ul^t$, and the
bridge $Z^\pi_{h+1} - \Zl^t_{h+1} \le \Zr^t_{h+1} - \Zl^t_{h+1} \le \pi_{h+1}\cert^t_{h+1}$ of
Lemma~\ref{lem:cert_true_recursion_neg} controls $\abs{\Zl^t_{h+1} - Z^\pi_{h+1}}$, whereas
$\Zt^t_{h+1} - Z^\pi_{h+1}$ has no sign. The last term of the lemma is $\frac1H\phat\abs{Z^*_{h+1} - \Zl^t_{h+1}}$
(no sign assumption, deviation~13); Lemma~\ref{lem:optimism_neg} removes the absolute value.
\end{remark}

\begin{lemma}[Optimism and pessimism; Lemma~13 of the paper]\label{lem:optimism_neg}
\uses{def:good_event, def:backups, def:opt_value, def:entropic_bpi}
On $\evGood^-$, for all $t \ge 0$, $h \in \{1, \dots, H + 1\}$, $s \in \Ss$ and $a \in \As$,
\[
  \Ul_h^t(s,a) \le U^*_h(s,a) \le \Ut_h^t(s,a) \quad (h \le H)
  \qquad\text{and}\qquad
  \Zl_h^t(s) \le Z^*_h(s) \le \Zt_h^t(s) .
\]
\end{lemma}

\begin{proof}
\uses{lem:concentration_neg, lem:opt_exp_value_range, lem:opt_bellman, lem:emp_trans_simplex, def:backups, lem:backups_range}
Fix $\omega \in \evGood^-$ and $t$. Backward induction on $h$ with the statement $P(h)$: ``for all $s$,
$\Zl_h^t(s) \le Z^*_h(s) \le \Zt_h^t(s)$''; base case $h = H+1$: all three are $1$.

\emph{Induction step.} Let $h \le H$, assume $P(h+1)$, fix $(s,a)$. By Lemma~\ref{lem:opt_bellman},
$U^*_h(s,a) = e^{\beta r_h(s,a)}(p_hZ^*_{h+1})(s,a) \in [\ell_h, 1]$ (Lemma~\ref{lem:opt_exp_value_range}).

\emph{Never-visited pair} ($n = 0$): $\Ut_h^t(s,a) = 1 \ge U^*_h(s,a) \ge \ell_h = \Ul_h^t(s,a)$.

\emph{Visited pair} ($n \ge 1$). Upper bound: $\Ut_h^t(s,a) = \min\{1, X\}$ with
$X := e^{\beta r_h(s,a)}[\phat\Zt^t_{h+1} + b + \frac1H\phat(\Zt^t_{h+1} - \Zl^t_{h+1})]$; $U^*_h(s,a) \le 1$, and
\[
  X - U^*_h(s,a) = e^{\beta r_h(s,a)}\Big[\phat(\Zt^t_{h+1} - Z^*_{h+1}) + (\phat - p_h)Z^*_{h+1}(s,a) + b
  + \tfrac1H\phat(\Zt^t_{h+1} - \Zl^t_{h+1})\Big] .
\]
By $P(h+1)$, $\abs{Z^*_{h+1} - \Zl^t_{h+1}} = Z^*_{h+1} - \Zl^t_{h+1}$ pointwise, so
Lemma~\ref{lem:concentration_neg} gives $(\phat - p_h)Z^*_{h+1}(s,a) \ge -b - \frac1H\phat(Z^*_{h+1} - \Zl^t_{h+1})$
and
\[
  X - U^*_h(s,a) \ge e^{\beta r_h(s,a)}\Big[\phat(\Zt^t_{h+1} - Z^*_{h+1}) + \tfrac1H\phat(\Zt^t_{h+1} - Z^*_{h+1})\Big]
  = e^{\beta r_h(s,a)}\big(1 + \tfrac1H\big)\phat(\Zt^t_{h+1} - Z^*_{h+1}) \ge 0
\]
by $P(h+1)$ and the positivity of $\phat$ (Lemma~\ref{lem:emp_trans_simplex}). Lower bound:
$\Ul_h^t(s,a) = \max\{\ell_h, Y\}$ with $Y := e^{\beta r_h(s,a)}[\phat\Zl^t_{h+1} - b - \frac1H\phat(\Zt^t_{h+1} - \Zl^t_{h+1})]$;
$U^*_h(s,a) \ge \ell_h$, and with Lemma~\ref{lem:concentration_neg} in the form
$(p_h - \phat)Z^*_{h+1}(s,a) \ge -b - \frac1H\phat(Z^*_{h+1} - \Zl^t_{h+1})$,
\[
  U^*_h(s,a) - Y = e^{\beta r_h(s,a)}\Big[\phat(Z^*_{h+1} - \Zl^t_{h+1}) + (p_h - \phat)Z^*_{h+1}(s,a) + b + \tfrac1H\phat(\Zt^t_{h+1} - \Zl^t_{h+1})\Big]
  \ge e^{\beta r_h(s,a)}\Big[\big(1 - \tfrac1H\big)\phat(Z^*_{h+1} - \Zl^t_{h+1}) + \tfrac1H\phat(\Zt^t_{h+1} - \Zl^t_{h+1})\Big] \ge 0
\]
($H \ge 1$ and $P(h+1)$). \emph{From $U$ to $Z$:} by Lemma~\ref{lem:opt_bellman} ($\beta < 0$),
$Z^*_h(s) = \min_aU^*_h(s,a)$, and $\Zt_h^t(s) = \min_a\Ut_h^t(s,a)$, $\Zl_h^t(s) = \min_a\Ul_h^t(s,a)$
(Definition~\ref{def:backups}); the minimum over $a$ is monotone.
\end{proof}

\textbf{Lean remark.} Compared with Lemma~\ref{lem:optimism_pos}, the roles of the factors
$1 - \frac1H$ and $1 + \frac1H$ are exchanged between the upper and the lower bound; both are
nonnegative, and the Lean proof is the same term with \texttt{Finset.inf'} in place of \texttt{Finset.sup'}.

\section{The ring value and the certificate}
\label{sec:neg_certificate}

For $\beta < 0$ the value $Z^{\pi}$ of the greedy policy $\pi = \pi^{t+1}$ is \emph{above} $Z^*$
(Lemma~\ref{lem:opt_exp_value_eq}), and the gap to control is $Z^\pi - Z^*$. The ring value $\Zr^t$ of
Definition~\ref{def:ring} ($\beta < 0$) follows $\pi$ under the true kernel, clipped by the
\emph{optimistic} empirical backup; it lies above both $\Zt^t$ and $Z^\pi$, and the certificate
dominates $\Zr^t - \Zl^t$.

\begin{lemma}[The ring value is at most one]\label{lem:neg_ring_range}
\uses{def:ring, def:entropic_bpi}
For all $t \ge 0$, $h \in \{1, \dots, H + 1\}$, $s$ and $a$: $\Zr_h^t(s) \le 1$ and, for $h \le H$,
$\Ur_h^t(s,a) \le 1$.
\end{lemma}

\begin{proof}
\uses{def:ring, lem:exp_value_monotone}
Backward induction on $h$; $\Zr_{H+1}^t = 1$. If $\Zr_{h+1}^t \le 1$ pointwise then
$e^{\beta r_h(s,a)}(p_h\Zr^t_{h+1})(s,a) \le (p_h\Zr^t_{h+1})(s,a) \le 1$ ($e^{\beta r_h} \le 1$ as $\beta < 0 \le r_h$; $p_h$
positive with $p_h1 = 1$, Lemma~\ref{lem:exp_value_monotone}). For a visited pair
$\Ur^t_{h,\mathrm{opt}}(s,a) = \min\{1, \cdot\} \le 1$, so $\Ur_h^t(s,a)$, a maximum of two numbers $\le 1$,
is $\le 1$; for a never-visited pair $\Ur_h^t(s,a) = \max\{e^{\beta r_h}(p_h\Zr^t_{h+1})(s,a), 1\} = 1$.
Finally $\Zr_h^t(s) = \Ur_h^t(s, \pi_h(s)) \le 1$.
\end{proof}

\begin{lemma}[The ring value is an upper bound; Lemma~15 of the paper]\label{lem:ring_neg}
\uses{def:ring, def:backups, def:greedy, def:exp_value, def:entropic_bpi}
For all $t \ge 0$ (no event is needed), $h \in \{1, \dots, H + 1\}$, $s$ and $a$, with $\pi = \pi^{t+1}$,
\[
  \Ur_h^t(s,a) \ge \max\{\Ut_h^t(s,a), U^\pi_h(s,a)\} \quad (h \le H)
  \qquad\text{and}\qquad
  \Zr_h^t(s) \ge \max\{\Zt_h^t(s), Z^\pi_h(s)\} .
\]
\end{lemma}

\begin{proof}
\uses{def:ring, def:backups, def:greedy, lem:exp_value_monotone, lem:exp_bellman, def:exp_value, lem:emp_trans_simplex}
Backward induction on $h$ with $P(h)$: ``for all $s$, $\Zr_h^t(s) \ge \max\{\Zt_h^t(s), Z^\pi_h(s)\}$'';
base case: all three values are $1$ at $h = H + 1$. Let $h \le H$, assume $P(h+1)$, fix $(s,a)$.

\emph{Comparison with $U^\pi$.} In both cases $\Ur_h^t(s,a) \ge e^{\beta r_h(s,a)}(p_h\Zr^t_{h+1})(s,a)$ (it is a
maximum with this quantity), and $(p_h\Zr^t_{h+1})(s,a) \ge (p_hZ^\pi_{h+1})(s,a)$ by $P(h+1)$ and
Lemma~\ref{lem:exp_value_monotone}; $e^{\beta r_h(s,a)}(p_hZ^\pi_{h+1})(s,a) = U^\pi_h(s,a)$
(Definition~\ref{def:exp_value}).

\emph{Comparison with $\Ut$.} If $n = 0$, $\Ur_h^t(s,a) \ge 1 = \Ut_h^t(s,a)$. If $n \ge 1$,
$\Ur_h^t(s,a) \ge \Ur^t_{h,\mathrm{opt}}(s,a) = \min\{1, X'\}$ and $\Ut_h^t(s,a) = \min\{1, X\}$ with
\[
  X = e^{\beta r_h}\big[\phat\Zt^t_{h+1} + b + \tfrac1H\phat(\Zt^t_{h+1} - \Zl^t_{h+1})\big], \qquad
  X' = e^{\beta r_h}\big[\phat\Zr^t_{h+1} + b + \tfrac1H\phat(\Zr^t_{h+1} - \Zl^t_{h+1})\big],
\]
so $X' - X = e^{\beta r_h}(1 + \frac1H)\phat(\Zr^t_{h+1} - \Zt^t_{h+1}) \ge 0$ by $P(h+1)$ and the positivity
of $\phat$ (Lemma~\ref{lem:emp_trans_simplex}); $x \mapsto \min\{1, x\}$ is monotone.

\emph{From $U$ to $Z$.} With $a = \pi_h(s)$: $\Zr_h^t(s) = \Ur_h^t(s,a) \ge \Ut_h^t(s,a) \ge \min_{a'}\Ut_h^t(s,a') = \Zt_h^t(s)$
and $\Zr_h^t(s) \ge U^\pi_h(s, \pi_h(s)) = Z^\pi_h(s)$ (Lemma~\ref{lem:exp_bellman}). This is $P(h)$.
\end{proof}

\begin{lemma}[One-step certificate bound; Lemma~16 of the paper]\label{lem:certificate_neg}
\uses{def:good_event, def:ring, def:backups, def:bonus, def:entropic_bpi}
On $\evGood^-$, for all $t \ge 0$, $h \in \{1, \dots, H\}$ and $(s,a)$ with $n = n_h^t(s,a) \ge 1$,
\[
  \Ur_h^t(s,a) - \Ul_h^t(s,a)
  \le e^{\beta r_h(s,a)}\Big[3 b + \Big(1 + \frac3H\Big)\phat\big(\Zr^t_{h+1} - \Zl^t_{h+1}\big)\Big] .
\]
\end{lemma}

\begin{proof}
\uses{lem:concentration_neg, lem:optimism_neg, lem:ring_neg, lem:neg_ring_range, lem:kl_bernstein, lem:kl_transport_var, lem:var_le_range_mean, lem:sqrt_add_le, lem:sqrt_mul_le, lem:emp_trans_simplex, lem:backups_range, lem:opt_exp_value_range, def:event_kl, def:bonus}
Fix $\omega \in \evGood^-$, $t$, $h$, $(s,a)$ with $n \ge 1$; write $\Delta := \phat(\Zr^t_{h+1} - \Zl^t_{h+1}) \ge 0$
(Lemmas~\ref{lem:optimism_neg}, \ref{lem:ring_neg}: $\Zl^t_{h+1} \le Z^*_{h+1} \le \Zt^t_{h+1} \le \Zr^t_{h+1}$
pointwise). In both cases below, $\Ul_h^t(s,a) \ge Y := e^{\beta r_h}[\phat\Zl^t_{h+1} - b - \frac1H\phat(\Zt^t_{h+1} - \Zl^t_{h+1})]$
(the maximum dominates its second argument), and $\phat(\Zt^t_{h+1} - \Zl^t_{h+1}) \le \Delta$
(Lemma~\ref{lem:emp_trans_simplex}). Recall $\Ur_h^t(s,a) = \max\{e^{\beta r_h}(p_h\Zr^t_{h+1})(s,a), \Ur^t_{h,\mathrm{opt}}(s,a)\}$.

\emph{Case 1: $\Ur_h^t(s,a) = e^{\beta r_h}(p_h\Zr^t_{h+1})(s,a)$.} Then
$\Ur_h^t(s,a) - \Ul_h^t(s,a) \le e^{\beta r_h}[b + \frac1H\Delta + p_h\Zr^t_{h+1} - \phat\Zl^t_{h+1}]$ and
\[
  p_h\Zr^t_{h+1} - \phat\Zl^t_{h+1}
  = \Delta + (p_h - \phat)Z^*_{h+1}(s,a) + (p_h - \phat)\big(\Zr^t_{h+1} - Z^*_{h+1}\big)(s,a) .
\]
\emph{Second term.} By Lemma~\ref{lem:concentration_neg} and $\abs{Z^*_{h+1} - \Zl^t_{h+1}} = Z^*_{h+1} - \Zl^t_{h+1} \le \Zr^t_{h+1} - \Zl^t_{h+1}$
pointwise, $(p_h - \phat)Z^*_{h+1}(s,a) \le b + \frac1H\Delta$. (The paper's proof bounds this term
by $b + \frac1H\phat(\Zt^t_{h+1} - Z^*_{h+1})$, mixing the two chapters; the bound above is the one
that Lemma~\ref{lem:concentration_neg} gives.)
\emph{Third term.} Let $f := \Zr^t_{h+1} - Z^*_{h+1}$; then $0 \le f \le 1 - \ell_{h+1} = B_h$ pointwise
(Lemma~\ref{lem:neg_ring_range} and Lemma~\ref{lem:opt_exp_value_range}). Since $\omega \in \evKL$,
Lemma~\ref{lem:kl_bernstein} (with $p = \phat$, $q = p_h(\cdot \mid s,a)$, range $B_h$) gives
$(p_h - \phat)f(s,a) \le \sqrt{2\Var_{p_h}(f)(s,a)\alpha/n} + \frac23B_h\alpha/n$, with the variance under the
true kernel; Lemma~\ref{lem:kl_transport_var} (first inequality) and
Lemma~\ref{lem:var_le_range_mean} give $\Var_{p_h}(f) \le 2\Var_{\phat}(f) + 4B_h^2\alpha/n \le 2B_h\phat f + 4B_h^2\alpha/n$,
hence by Lemma~\ref{lem:sqrt_add_le} and Lemma~\ref{lem:sqrt_mul_le} ($x = \phat f$, $y = 4B_h\alpha/n$,
$\eta = 2/H$), and $\phat f \le \Delta$,
\[
  (p_h - \phat)f(s,a) \le \frac1H\Delta + \Big(H + 2\sqrt2 + \frac23\Big)B_h\frac{\alpha}{n}
  \le \frac1H\Delta + (4H + 5)B_h\frac{\alpha}{n} \le \frac1H\Delta + b ,
\]
by Definition~\ref{def:bonus} with the third term in $\alpha$ (Remark~\ref{rem:neg_constants}).
\emph{Sum.} $p_h\Zr^t_{h+1} - \phat\Zl^t_{h+1} \le 2b + (1 + \frac2H)\Delta$ and
$\Ur_h^t(s,a) - \Ul_h^t(s,a) \le e^{\beta r_h}[3b + (1 + \frac3H)\Delta]$.

\emph{Case 2: $\Ur_h^t(s,a) = \Ur^t_{h,\mathrm{opt}}(s,a)$.} Then
$\Ur_h^t(s,a) \le e^{\beta r_h}[\phat\Zr^t_{h+1} + b + \frac1H\Delta]$ (the minimum is below its second
argument), so
\[
  \Ur_h^t(s,a) - \Ul_h^t(s,a) \le e^{\beta r_h}\Big[2b + \Delta + \tfrac1H\Delta + \tfrac1H\phat(\Zt^t_{h+1} - \Zl^t_{h+1})\Big]
  \le e^{\beta r_h}\Big[2b + \big(1 + \tfrac2H\big)\Delta\Big] ,
\]
at most the claimed bound since $b \ge 0$ (Lemma~\ref{lem:backups_range}) and $\Delta \ge 0$.
\end{proof}

\begin{lemma}[The certificate dominates the ring--pessimistic gap]\label{lem:cert_dominates_gap_neg}
\uses{def:good_event, def:certificate, def:backups, def:ring, def:greedy, def:entropic_bpi}
On $\evGood^-$, for all $t \ge 0$, $h \in \{1, \dots, H + 1\}$ and $s \in \Ss$, with $\pi = \pi^{t+1}$,
\[
  \Zr_h^t(s) - \Zl_h^t(s) \le (\pi_h\cert_h^t)(s) .
\]
\end{lemma}

\begin{proof}
\uses{lem:certificate_neg, lem:backups_range, lem:neg_ring_range, lem:emp_trans_simplex, def:certificate, def:greedy, def:ring, lem:certificate_range}
Fix $\omega \in \evGood^-$ and $t$; backward induction on $h$. Base case $h = H + 1$: both sides are $0$.
Induction step: let $h \le H$, $s \in \Ss$ and $a := \pi_h(s)$; by Definition~\ref{def:greedy} ($\beta < 0$:
$\pi_h(s) = \argmin_a\Ul_h^t(s,a)$, deviation~2) and Definition~\ref{def:ring},
$\Zl_h^t(s) = \Ul_h^t(s,a)$ and $\Zr_h^t(s) = \Ur_h^t(s,a)$. In all cases
$\Zr_h^t(s) - \Zl_h^t(s) \le 1 - \ell_h \le 1$ by Lemma~\ref{lem:neg_ring_range} and
Lemma~\ref{lem:backups_range} ($\Zl_h^t \ge \ell_h$). \emph{Never-visited pair:} $(\pi_h\cert_h^t)(s) = 1$
(Definition~\ref{def:certificate}). \emph{Visited pair:}
$(\pi_h\cert_h^t)(s) = \min\{1, e^{\beta r_h(s,a)}[3b + (1 + \frac3H)\phat(\pi_{h+1}\cert^t_{h+1})]\}$, and
Lemma~\ref{lem:certificate_neg} with the induction hypothesis integrated against $\phat$
(Lemma~\ref{lem:emp_trans_simplex}) bounds $\Zr_h^t(s) - \Zl_h^t(s)$ by the second argument.
\end{proof}

\textbf{Lean remark.} This is the one place where deviation~2 is essential: with the greedy
policy $\argmin_a\Ut^t_h(s,a)$ of the paper's algorithm box, $\Zl_h^t(s)$ would not equal
$\Ul_h^t(s, \pi_h(s))$ and the induction would not close.

\begin{lemma}[The certificate dominates the suboptimality gap; Lemma~14 of the paper]\label{lem:stopping_neg}
\uses{def:good_event, def:certificate, def:opt_value, def:exp_value, def:greedy, def:entropic_bpi}
On $\evGood^-$, for all $t \ge 0$, $h \in \{1, \dots, H + 1\}$ and $s \in \Ss$, with $\pi = \pi^{t+1}$,
\[
  Z^\pi_h(s) - Z^*_h(s) \le (\pi_h\cert_h^t)(s) .
\]
\end{lemma}

\begin{proof}
\uses{lem:optimism_neg, lem:ring_neg, lem:cert_dominates_gap_neg}
$Z^\pi_h(s) \le \Zr_h^t(s)$ (Lemma~\ref{lem:ring_neg}) and $\Zl_h^t(s) \le Z^*_h(s)$
(Lemma~\ref{lem:optimism_neg}), so $Z^\pi_h(s) - Z^*_h(s) \le \Zr_h^t(s) - \Zl_h^t(s) \le (\pi_h\cert_h^t)(s)$
(Lemma~\ref{lem:cert_dominates_gap_neg}). The statement is for the exponential values $Z$ at every
step (the paper's Lemma~14 mixes $V$ and $Z$ in its last sentence, deviation~4).
\end{proof}

\begin{lemma}[Soundness of the stopping rule]\label{lem:pac_at_stopping_neg}
\uses{def:good_event, def:stopping_rule, def:entropic_value, def:opt_value, def:greedy, def:entropic_bpi}
On $\evGood^-$, for every $t \ge 0$ such that the stopping condition of
Definition~\ref{def:stopping_rule} holds after $t$ episodes, i.e.\
$(\pi_1\cert_1^t)(s_1) \le (1 - e^{\beta\eps})\Zl_1^t(s_1)$ with $\pi = \pi^{t+1}$, one has
$V^*_1(s_1) - V^\pi_1(s_1) \le \eps$.
\end{lemma}

\begin{proof}
\uses{lem:stopping_neg, lem:optimism_neg, lem:opt_exp_value_eq, lem:value_gap_log, lem:exp_value_range, lem:two_sub_le_inv}
Write $g := (\pi_1\cert_1^t)(s_1)$. By Lemma~\ref{lem:stopping_neg}, Lemma~\ref{lem:optimism_neg}
($\Zl_1^t(s_1) \le Z^*_1(s_1)$) and $1 - e^{\beta\eps} > 0$,
\[
  Z^\pi_1(s_1) - Z^*_1(s_1) \le g \le (1 - e^{\beta\eps})\Zl_1^t(s_1) \le (1 - e^{\beta\eps})Z^*_1(s_1) .
\]
Two ways to conclude. \emph{Directly:} $Z^*_1(s_1) \le Z^\pi_1(s_1)$ (Lemma~\ref{lem:opt_exp_value_eq},
$\beta < 0$), so $(1 - e^{\beta\eps})Z^*_1(s_1) \le (1 - e^{\beta\eps})Z^\pi_1(s_1)$ and
Lemma~\ref{lem:value_gap_log} ($\beta < 0$, with $Z^\pi_1(s_1) > 0$ by Lemma~\ref{lem:exp_value_range})
gives $V^*_1(s_1) - V^\pi_1(s_1) \le \eps$. \emph{Through the ratio:} $Z^*_1(s_1) > 0$
(Lemma~\ref{lem:exp_value_range}), so $Z^\pi_1(s_1)/Z^*_1(s_1) \le 2 - e^{\beta\eps} \le e^{-\beta\eps} = e^{|\beta|\eps}$
by Lemma~\ref{lem:two_sub_le_inv} applied to $x = e^{\beta\eps} \in (0, 1]$, and
$V^*_1(s_1) - V^\pi_1(s_1) = \beta^{-1}\log(Z^*_1(s_1)/Z^\pi_1(s_1)) = |\beta|^{-1}\log(Z^\pi_1(s_1)/Z^*_1(s_1)) \le \eps$
(Lemma~\ref{lem:value_gap_log}).
\end{proof}

\begin{remark}[The paper's argument at stopping for $\beta < 0$]\label{rem:neg_stopping}
\uses{lem:pac_at_stopping_neg}
The end of Appendix~B.2 of the paper argues ``since $Z^\pi_1 \ge \Zl_1^t$, the stopping condition is
stronger than $\pi\cert \le (1 - e^{\beta\eps})Z^\pi_1$'' and then writes
$V^*_1 - V^\pi_1 = \beta^{-1}\log(1 + (Z^*_1 - Z^\pi_1)/Z^\pi_1) \le \beta^{-1}\log(1 + \pi\cert/Z^\pi_1)$. The
inequality $Z^\pi_1 \ge \Zl_1^t$ is not among the paper's lemmas (its Lemma~15 gives an upper
bound $Z^\pi_1 \le \Zr_1^t$); it does hold, through $Z^\pi_1 \ge Z^*_1 \ge \Zl_1^t$
(Lemmas~\ref{lem:opt_exp_value_eq} and~\ref{lem:optimism_neg}), which is the direct route above. In
the last display, since $\beta^{-1} < 0$ the map $x \mapsto \beta^{-1}\log(1 + x)$ is decreasing and the
lower bound $Z^*_1 - Z^\pi_1 \ge -\pi\cert$ of Lemma~\ref{lem:stopping_neg} is what is needed, so the
display should read $\le \beta^{-1}\log(1 - \pi\cert/Z^\pi_1) \le \beta^{-1}\log(e^{\beta\eps}) = \eps$. The
alternative route uses only $\Zl_1^t \le Z^*_1$ and the elementary inequality
$2 - e^{\beta\eps} \le e^{|\beta|\eps}$, which is the form recorded in the overview.
\end{remark}

\section{The certificate under the true kernel}
\label{sec:neg_true_kernel}

\begin{lemma}[Certificate recursion under the true kernel]\label{lem:cert_true_recursion_neg}
\uses{def:good_event, def:certificate, def:rate_min, def:exp_value, def:greedy, def:entropic_bpi}
On $\evGood^-$, for all $t \ge 0$, $h \in \{1, \dots, H\}$ and $s \in \Ss$, with $\pi = \pi^{t+1}$ and
$a := \pi_h(s)$,
\[
  (\pi_h\cert_h^t)(s) \le e^{\beta r_h(s,a)}\Big[
    36\sqrt{\Var_{p_h}(Z^\pi_{h+1})(s,a)\,\rho^{*t}_h(s,a)}
    + \Big(1 + \frac{13}{H}\Big)\big(p_h\,\pi_{h+1}\cert^t_{h+1}\big)(s,a)\Big]
    + 84\,H\,\rho^t_h(s,a) ,
\]
with $\rho^t_h(s,a) = \min\{\alpha(n,\delta)/n, 1\}$, $\rho^{*t}_h(s,a) = \min\{\alpha^*(n,\delta)/n, 1\}$
(Definition~\ref{def:rate_min}; both $1$ for a never-visited pair). Note that the additive term
is \emph{not} multiplied by $e^{\beta r_h(s,a)}$ (Remark~\ref{rem:neg_additive_term}); equivalently,
the bracket may contain $84He^{-\beta r_h(s,a)}\rho^t_h(s,a)$, which is at most $84He^{|\beta|}\rho^t_h(s,a)$.
\end{lemma}

\begin{proof}
\uses{lem:certificate_range, def:certificate, def:bonus, def:event_kl, lem:kl_bernstein, lem:var_le_range_mean, lem:kl_transport_var, lem:transport_var, lem:sqrt_add_le, lem:sqrt_mul_le, lem:rates_props, lem:opt_exp_value_eq, lem:optimism_neg, lem:ring_neg, lem:cert_dominates_gap_neg, lem:backups_range, lem:exp_value_range, lem:exp_value_monotone, lem:emp_trans_simplex}
Fix $\omega \in \evGood^-$, $t$, $h$, $s$, $a = \pi_h(s)$; write $g := \pi_{h+1}\cert^t_{h+1} : \Ss \to [0, 1]$
(Lemma~\ref{lem:certificate_range}; the clip of the certificate is $1$ at every step) and
$\Delta_p := (p_hg)(s,a) \ge 0$. The bracket is nonnegative and $e^{\beta r_h(s,a)} \in (0, 1]$.

\emph{Trivial case: $n = 0$ or $\alpha/n \ge 1$.} Then $\rho = 1$ and
$(\pi_h\cert_h^t)(s) \le 1 \le 84H = 84H\rho$ (Lemma~\ref{lem:certificate_range}), which is at most the
right-hand side. (This is where the additive term must stand outside the factor $e^{\beta r_h(s,a)}$:
for a never-visited pair the certificate equals $1$, while $e^{\beta r_h(s,a)}\cdot 84H$ can be smaller
than $1$ when $|\beta| > \log(84H)$.)

\emph{Main case: $n \ge 1$ and $\alpha/n < 1$.} Then $\rho = \alpha/n$, $\rho^* = \alpha^*/n \le \rho$
(Lemma~\ref{lem:rates_props}), and $(\pi_h\cert_h^t)(s) \le e^{\beta r_h(s,a)}[3b + (1 + \frac3H)\phat g]$
(drop the clip in Definition~\ref{def:certificate}).

\emph{Step (i): the empirical expectation of $g$} (range $1$). Since $\omega \in \evKL$,
Lemma~\ref{lem:kl_bernstein} gives $\abs{(\phat - p_h)g(s,a)} \le \sqrt{2\Var_{p_h}(g)(s,a)\rho} + \frac23\rho$,
$\Var_{p_h}(g)(s,a) \le \Delta_p$ (Lemma~\ref{lem:var_le_range_mean}), and Lemma~\ref{lem:sqrt_mul_le}
($x = \Delta_p$, $y = 2\rho$, $\eta = 2/H$) gives $\sqrt{2\Delta_p\rho} \le \frac1H\Delta_p + \frac H2\rho$; hence
$\phat g \le (1 + \frac1H)\Delta_p + 2H\rho$. (The paper uses the range $1 - e^{\beta(H-h)}$ for $g$, which
the definition of the certificate does not provide; the range $1$ costs nothing since
$1 - e^{\beta(H-h)} \le 1$, and it is the reason why the additive term has no factor $B_h$,
deviation~11: the paper's factor vanishes at $h = H$.)

\emph{Step (ii): the bonus under the true kernel.} $b = 2\sqrt2\sqrt{\Var_{\phat}(\Zl^t_{h+1})\rho^*} + 5B_h\rho + 4HB_h\rho$.
Lemma~\ref{lem:kl_transport_var} (second inequality, $f = \Zl^t_{h+1} - \ell_{h+1} \in [0, B_h]$ by
Lemma~\ref{lem:backups_range}) gives $\Var_{\phat}(\Zl^t_{h+1}) \le 2\Var_{p_h}(\Zl^t_{h+1})(s,a) + 4B_h^2\rho$;
Lemma~\ref{lem:transport_var} (with $f = \Zl^t_{h+1} - \ell_{h+1}$, $g' = Z^\pi_{h+1} - \ell_{h+1} \in [0, B_h]$,
Lemma~\ref{lem:exp_value_range}) gives $\Var_{p_h}(\Zl^t_{h+1}) \le 2\Var_{p_h}(Z^\pi_{h+1}) + 2B_h\,p_h\abs{\Zl^t_{h+1} - Z^\pi_{h+1}}$.
Now $\Zl^t_{h+1} \le Z^*_{h+1} \le Z^\pi_{h+1} \le \Zr^t_{h+1}$ pointwise (Lemma~\ref{lem:optimism_neg},
Lemma~\ref{lem:opt_exp_value_eq} for $\beta < 0$, Lemma~\ref{lem:ring_neg}), so
$\abs{\Zl^t_{h+1} - Z^\pi_{h+1}} = Z^\pi_{h+1} - \Zl^t_{h+1} \le \Zr^t_{h+1} - \Zl^t_{h+1} \le g$
(Lemma~\ref{lem:cert_dominates_gap_neg}) and $p_h\abs{\Zl^t_{h+1} - Z^\pi_{h+1}}(s,a) \le \Delta_p$
(Lemma~\ref{lem:exp_value_monotone}). Exactly as in Lemma~\ref{lem:cert_true_recursion_pos}
(Lemma~\ref{lem:sqrt_add_le}, Lemma~\ref{lem:sqrt_mul_le} with $\eta = 1/(\sqrt2H)$, $\sqrt{\rho\rho^*} \le \rho$,
$\rho^* \le \rho$, $4\sqrt2 \le 6$, $12H + 4\sqrt2 + 5 \le 23H$), and using $B_h \le 1$,
\[
  b \le 6\sqrt{\Var_{p_h}(Z^\pi_{h+1})(s,a)\rho^*} + \frac1H\Delta_p + 23H\rho .
\]

\emph{Step (iii): combination.} As in Lemma~\ref{lem:cert_true_recursion_pos},
$3b + (1 + \frac3H)\phat g \le 36\sqrt{\Var_{p_h}(Z^\pi_{h+1})(s,a)\rho^*} + (1 + \frac{10}{H})\Delta_p + 77H\rho$,
and multiplying by $e^{\beta r_h(s,a)} \in (0, 1]$ gives the claim, since
$e^{\beta r_h(s,a)}\cdot 77H\rho \le 84H\rho$ and $10 \le 13$.
\end{proof}

\begin{remark}[The additive term for $\beta < 0$]\label{rem:neg_additive_term}
\uses{lem:cert_true_recursion_neg}
For $\beta > 0$ the factor $e^{\beta r_h} \ge 1$ absorbs the clip of a never-visited pair; for $\beta < 0$
it does not ($e^{\beta r_h} \le 1$), so the additive term must be placed \emph{outside} the factor
$e^{\beta r_h(s,a)}$: the recursion proved is
\[
  (\pi_h\cert_h^t)(s) \le e^{\beta r_h(s,a)}\Big[36\sqrt{\Var_{p_h}(Z^\pi_{h+1})(s,a)\rho^*} + \Big(1 + \frac{13}{H}\Big)(p_h\pi_{h+1}\cert^t_{h+1})(s,a)\Big] + 84H\rho^t_h(s,a) .
\]
This is weaker than the paper's form (which has the factor $e^{\beta r_h} \le 1$ on the additive term
as well), but the paper's form fails for a never-visited pair when $|\beta| > \log(84H)$; in the
main case the proof gives the paper's form with $77H$. It is the form that the unrolling of Lemma~\ref{lem:cert_unrolled_neg} consumes
(Lemma~\ref{lem:unroll_recursion} accepts $u_h$ of the form $e^{-\beta r_h}\cdot 84H\rho$, i.e.\ the
recursion with $u_h(s,a) := 36\sqrt{\cdot} + 84He^{-\beta r_h(s,a)}\rho^t_h(s,a)$, and since
$e^{\beta\sum_{i \le h}r_i}e^{-\beta r_h} = e^{\beta\sum_{i < h}r_i} \le 1$ for $\beta < 0$, the additive term
contributes $84H\rho$ per step to the unrolled bound). Both signs are dominated
by the uniform term $84He^{|\beta|(H+1)}\rho$ (deviation~11), which the constant $D$ of
Definition~\ref{def:upper_bound_const} reflects.
\end{remark}

\begin{lemma}[Unrolled certificate]\label{lem:cert_unrolled_neg}
\uses{def:good_event, def:certificate, def:rate_min, def:exp_value, def:step_law, def:greedy, def:entropic_bpi}
On $\evGood^-$, for all $t \ge 0$, with $\pi = \pi^{t+1}$ and $\E^\pi = \E^\pi_{(s_1, 1)}$,
\[
  (\pi_1\cert_1^t)(s_1) \le e^{13}\,\E^\pi\Bigg[\sum_{h=1}^H
  \Big(36\, e^{\beta\sum_{i \le h} r_i(S_i, A_i)}\sqrt{\Var_{p_h}(Z^\pi_{h+1})(S_h, A_h)\,\rho^{*t}_h(S_h, A_h)}
  + 84\,H\,\rho^t_h(S_h, A_h)\Big)\Bigg] .
\]
\end{lemma}

\begin{proof}
\uses{lem:cert_true_recursion_neg, lem:unroll_recursion, lem:one_add_div_pow_le_exp, def:certificate}
Fix $\omega \in \evGood^-$ and $t$. Apply Lemma~\ref{lem:unroll_recursion} (valid for $\beta$ of either
sign) to $\pi$, $g_h := \pi_h\cert_h^t$ ($g_{H+1} = 0$), $\kappa := 1 + \frac{13}{H}$ and
$u_h(s,a) := 36\sqrt{\Var_{p_h}(Z^\pi_{h+1})(s,a)\rho^{*t}_h(s,a)} + 84He^{-\beta r_h(s,a)}\rho^t_h(s,a) \ge 0$: the
hypothesis is Lemma~\ref{lem:cert_true_recursion_neg} in the form of
Remark~\ref{rem:neg_additive_term}. The conclusion is
$g_1(s_1) \le \sum_h\kappa^{h-1}\E^\pi[e^{\beta\sum_{i \le h}r_i(S_i,A_i)}u_h(S_h,A_h)]$; $\kappa^{h-1} \le e^{13}$
(Lemma~\ref{lem:one_add_div_pow_le_exp}), and
$e^{\beta\sum_{i \le h}r_i}\cdot 84He^{-\beta r_h}\rho = 84He^{\beta\sum_{i < h}r_i}\rho \le 84H\rho$ since $\beta < 0$ and $r \ge 0$.
\end{proof}

\begin{lemma}[Normalized certificate bound]\label{lem:ratio_bound_neg}
\uses{def:good_event, def:certificate, def:backups, def:rate_min, def:occupancy, def:gmax, def:upper_bound_const, def:greedy, def:entropic_bpi}
Let $V_G := (e^{|\beta|\Gmax} - 1)^2/e^{|\beta|\Gmax}$ with $\Gmax = \Gmax(\mathcal M)$ (the $V_G$ of
Definition~\ref{def:upper_bound_const}). On $\evGood^-$, for all $t \ge 0$, with $\pi = \pi^{t+1}$,
\[
  \frac{(\pi_1\cert_1^t)(s_1)}{\Zl_1^t(s_1) + (\pi_1\cert_1^t)(s_1)}
  \le \frac{(\pi_1\cert_1^t)(s_1)}{Z^\pi_1(s_1)}
  \le e^{13}\Big[36\sqrt{V_G\,x_t} + 84\,H\,e^{|\beta|H}\,y_t\Big] ,
\]
with $x_t := \sum_{h}\sum_{s,a}p^\pi_h(s,a)\rho^{*t}_h(s,a)$ and $y_t := \sum_h\sum_{s,a}p^\pi_h(s,a)\rho^t_h(s,a)$
(Definition~\ref{def:occupancy}); in particular the right-hand side is at most
$e^{13}[36\sqrt{V_Gx_t} + 84He^{|\beta|(H+1)}y_t]$.
\end{lemma}

\begin{proof}
\uses{lem:cert_unrolled_neg, lem:ring_neg, lem:cert_dominates_gap_neg, lem:backups_range, lem:exp_value_range, lem:cauchy_schwarz_traj, lem:entropic_variance_sum, lem:variance_ratio_le, lem:gmax_le_horizon, lem:occupancy_expectation, def:entropic_variance}
Fix $\omega \in \evGood^-$ and $t$; write $g := (\pi_1\cert_1^t)(s_1) \ge 0$.

\emph{Denominators.} $0 < Z^\pi_1(s_1) \le \Zr_1^t(s_1) \le \Zl_1^t(s_1) + g$ (Lemma~\ref{lem:exp_value_range},
Lemma~\ref{lem:ring_neg}, Lemma~\ref{lem:cert_dominates_gap_neg}), and $\Zl_1^t(s_1) \ge \ell_1 > 0$
(Lemma~\ref{lem:backups_range}); hence the first inequality. Also $Z^\pi_1(s_1) \ge \ell_1 = e^{\beta H}$, i.e.\
$1/Z^\pi_1(s_1) \le e^{|\beta|H}$.

\emph{Variance term.} Divide Lemma~\ref{lem:cert_unrolled_neg} by $Z^\pi_1(s_1)$. With
$w_h := e^{\beta\sum_{i \le h}r_i(S_i,A_i)}$, $v_h := \Var_{p_h}(Z^\pi_{h+1})(S_h,A_h)/Z^\pi_1(s_1)^2$ and
$u_h := \rho^{*t}_h(S_h,A_h)$, Lemma~\ref{lem:cauchy_schwarz_traj} gives
$\E^\pi[\sum_hw_h\sqrt{v_hu_h}] \le \sqrt{\E^\pi[\sum_hw_h^2v_h]}\sqrt{\E^\pi[\sum_hu_h]}$;
Lemma~\ref{lem:entropic_variance_sum} identifies the first factor squared with
$\Var_{\Prob^\pi}(e^{\beta R_1^\pi})/Z^\pi_1(s_1)^2$, which is at most $V_G$ by Lemma~\ref{lem:variance_ratio_le}
(with $G = \Gmax$, admissible by Lemma~\ref{lem:gmax_le_horizon}; the bound is in terms of $|\beta|$), and
Lemma~\ref{lem:occupancy_expectation} gives $\E^\pi[\sum_hu_h] = x_t$.

\emph{Rate term.} $\E^\pi[\sum_h 84H\rho^t_h(S_h,A_h)]/Z^\pi_1(s_1) \le 84He^{|\beta|H}y_t$
(Lemma~\ref{lem:occupancy_expectation}).
\end{proof}

\textbf{Lean remark.} For $\beta < 0$ the natural computable ratio is
$g/(\Zl_1^t(s_1) + g)$ rather than $g/\Zt_1^t(s_1)$: the stopping rule compares $g$ with
$\Zl_1^t(s_1)$, and the only available lower bound on $Z^\pi_1(s_1)$ in terms of computable quantities is
$Z^\pi_1 \ge \Zr_1^t - g \ge \Zt^t_1 - g$, while the upper bound $Z^\pi_1 \le \Zl_1^t + g$ is what the
normalization needs. The lemma is stated with both ratios so that Lemma~\ref{lem:progress_neg} can
use the computable one.

\section{Bounding the stopping time}
\label{sec:neg_stopping_time}

\begin{lemma}[Progress at a non-stopping episode]\label{lem:progress_neg}
\uses{def:good_event, def:stopping_rule, def:certificate, def:backups, def:upper_bound_const, def:entropic_bpi}
Let $\kappa(\beta, \eps) := (e^{|\beta|\eps} - 1)e^{-2|\beta|\eps}/2$ (Definition~\ref{def:upper_bound_const}).
On $\evGood^-$, for every $t \ge 0$ such that the stopping condition of Definition~\ref{def:stopping_rule}
\emph{fails} after $t$ episodes, with $g := (\pi_1\cert_1^t)(s_1)$,
\[
  \kappa(\beta, \eps) \le \frac{g}{\Zl_1^t(s_1) + g} .
\]
\end{lemma}

\begin{proof}
\uses{def:stopping_rule, lem:backups_range, lem:certificate_range}
Put $c := \Zl_1^t(s_1) \ge \ell_1 > 0$ (Lemma~\ref{lem:backups_range}) and $u := e^{|\beta|\eps} = e^{-\beta\eps} > 1$,
so that $e^{\beta\eps} = 1/u$. Failure of the condition means $g > (1 - e^{\beta\eps})c =: y \ge 0$. The map
$x \mapsto x/(c + x)$ is increasing on $[0, \infty)$, hence
\[
  \frac{g}{c + g} \ge \frac{y}{c + y} = \frac{1 - e^{\beta\eps}}{2 - e^{\beta\eps}} = \frac{u - 1}{2u - 1}
  \ge \frac{u - 1}{2u^2} = \kappa(\beta, \eps),
\]
the last inequality because $2u - 1 \le 2u^2$ (as $2u^2 - 2u + 1 = 2(u - \tfrac12)^2 + \tfrac12 > 0$). No event
is actually needed ($g \ge 0$ by Lemma~\ref{lem:certificate_range}).
\end{proof}

\begin{lemma}[Bound on the stopping time]\label{lem:stopping_time_bound_neg}
\uses{def:good_event, def:entropic_bpi, def:upper_bound_const, def:gmax}
On $\evGood^-$, the stopping time of the run satisfies
$\tau \le \mathrm{upperBound}(S, A, H, \beta, \eps, \delta, \Gmax(\mathcal M))$ (Definition~\ref{def:upper_bound_const});
in particular $\tau < \infty$.
\end{lemma}

\begin{proof}
\uses{lem:progress_neg, lem:ratio_bound_neg, lem:sum_counts_bound, lem:entropic_bpi_run, lem:self_bounding, def:rates, def:upper_bound_const, lem:rates_props, def:rate_min}
Identical to the proof of Lemma~\ref{lem:stopping_time_bound_pos}, with the ratio
$g_t/(\Zl_1^t(s_1) + g_t)$, $g_t := (\pi_1\cert_1^t)(s_1)$, in place of $g_t/\Zt_1^t(s_1)$. Fix
$\omega \in \evGood^-$, $\kappa := \kappa(\beta,\eps)$, $A_0 := \log(3SAH/\delta)$, and the constants
$C := 10^9\sqrt{V_GSAH}/\kappa$, $D := 10^{10}e^{|\beta|(H+1)}H^2SA/\kappa$ of Definition~\ref{def:upper_bound_const}.

\emph{Step 1.} For an integer $T \ge 1$ with $T \le \tau$ (the stopping condition fails after $t$
episodes for all $t < T$, Lemma~\ref{lem:entropic_bpi_run}), Lemmas~\ref{lem:progress_neg}
and~\ref{lem:ratio_bound_neg} summed over $t < T$ give, with the Cauchy--Schwarz inequality
$\sum_{t<T}\sqrt{x_t} \le \sqrt T\sqrt{\sum_{t<T}x_t}$,
\[
  T\kappa \le 36e^{13}\sqrt{V_G}\sqrt{T}\sqrt{\sum_{t<T}x_t} + 84e^{13}He^{|\beta|(H+1)}\sum_{t<T}y_t
\]
(we used $e^{|\beta|H} \le e^{|\beta|(H+1)}$).

\emph{Step 2.} On $\evCnt$, Lemma~\ref{lem:sum_counts_bound} with $L := \log(8eT)$, $\log(T+1) \le L$,
$\alpha^*(T-1,\delta) = A_0 + L$ and $\alpha(T-1,\delta) = A_0 + SL$ (Definition~\ref{def:rates}) gives
$\sum_{t<T}x_t \le 16SAH(A_0L + L^2)$ and $\sum_{t<T}y_t \le 16SAH(A_0L + SL^2)$.

\emph{Step 3.} Hence
$T \le \frac{144\,e^{13}\sqrt{V_GSAH}}{\kappa}\sqrt{T(A_0L + L^2)} + \frac{1344\,e^{13}e^{|\beta|(H+1)}H^2SA}{\kappa}(A_0L + SL^2)$,
and $144e^{13} \le 10^9$, $1344e^{13} \le 10^{10}$ give $T \le C\sqrt{T(A_0L + L^2)} + D(A_0L + SL^2)$.
Lemma~\ref{lem:self_bounding} with $A = A_0$, $B = 1$, $E = S$, $\alpha = 8e$ ($C \ge 0$, $D \ge 1$) yields
$T \le C^2(A_0 + 1)C_1^2 + (D + 2\sqrt DC)(A_0 + S)C_1^2 + 1 = \mathrm{upperBound}$ with the $C_1$ of that
lemma, which is the one of Definition~\ref{def:upper_bound_const} (Remark~\ref{rem:pos_c1_constant}).

\emph{Step 4.} As in Lemma~\ref{lem:stopping_time_bound_pos}: $\tau = \infty$ is impossible, and
$T = \tau$ gives the bound.
\end{proof}

\begin{remark}[The two signs share the constants]\label{rem:neg_constants_final}
\uses{lem:stopping_time_bound_neg, lem:stopping_time_bound_pos}
The paper's proof for $\beta < 0$ ends with the constant $81$ instead of $84$ in the additive term
and with $D = 81e^{13}e^{\beta H}H^2SA/(e^{|\beta|\eps} - 1)$, where $e^{\beta H} < 1$ is presumably a typo
for the $e^{|\beta|H}$ of its computation. Here both signs use the same $\kappa(\beta,\eps)$, the same
constants $36$, $13$, $84$, the same additive factor $e^{|\beta|(H+1)}$ and the same $C$, $D$, so that
Definition~\ref{def:upper_bound_const} is sign-independent. The progress per non-stopping
episode is $(e^{\beta\eps} - 1)e^{-\beta\eps}$ for $\beta > 0$ (Lemma~\ref{lem:progress_pos}) and
$(u - 1)/(2u - 1)$ with $u = e^{|\beta|\eps}$ for $\beta < 0$ (Lemma~\ref{lem:progress_neg}); both are at
least $\kappa(\beta, \eps) = (e^{|\beta|\eps} - 1)e^{-2|\beta|\eps}/2$ and of order $|\beta|\eps$ for small $\eps$, so
the common $\kappa$ costs at most a factor $2e^{2|\beta|\eps}$ in $1/\kappa$ relative to the paper's
sign-dependent constants (its Theorem~4 displays $e^{2\max\{0,\beta\}\eps}$; its $\beta < 0$ proof carries
$e^{2|\beta|\eps}$), which is absorbed in the generous constants $10^9$, $10^{10}$.
\end{remark}

\begin{lemma}[PAC guarantee and sample complexity of Entropic-BPI for $\beta < 0$]\label{lem:pac_neg}
\uses{def:entropic_bpi, def:entropic_value, def:opt_value, def:upper_bound_const, def:gmax, def:good_event}
In the standing setting,
\[
  \Prob\big(V^*_1(s_1) - V^{\mathrm{out}}_1(s_1) > \eps\big) \le \delta
  \qquad\text{and}\qquad
  \Prob\big(\tau \le \mathrm{upperBound}(S,A,H,\beta,\eps,\delta,\Gmax(\mathcal M))\big) \ge 1 - \delta .
\]
\end{lemma}

\begin{proof}
\uses{lem:stopping_time_bound_neg, lem:entropic_bpi_run, lem:pac_at_stopping_neg, lem:good_event}
As for Lemma~\ref{lem:pac_pos}: $\Prob(\evGood^-) \ge 1 - \delta$ (Lemma~\ref{lem:good_event}); outside the
null set of Lemma~\ref{lem:entropic_bpi_run}, every $\omega \in \evGood^-$ has $\tau(\omega) \le \mathrm{upperBound} < \infty$
(Lemma~\ref{lem:stopping_time_bound_neg}), the stopping condition holds after $\tau(\omega)$ episodes
and $\mathrm{out}(\omega) = \pi^{\tau(\omega)+1}$, so Lemma~\ref{lem:pac_at_stopping_neg} gives
$V^*_1(s_1) - V^{\mathrm{out}(\omega)}_1(s_1) \le \eps$.
\end{proof}

\textbf{Lean remark.} Chapter~\ref{chap:sample_complexity} combines Lemmas~\ref{lem:pac_pos}
and~\ref{lem:pac_neg} by a case split on the sign of $\beta \ne 0$ into
Theorems~\ref{thm:upper_bound_pac} and~\ref{thm:upper_bound_complexity}; the two lemmas have the
same Lean statement with the hypotheses \texttt{0 < β} and \texttt{β < 0}.

\chapter{Sample complexity of Entropic-BPI}
\label{chap:sample_complexity}

This chapter proves Theorem~4 of the paper (Section~4, restated with explicit constants in
Appendix~B), in the two-part form of deviation~6: the PAC guarantee
(Theorem~\ref{thm:upper_bound_pac}) and the high-probability bound on the number of episodes
(Theorem~\ref{thm:upper_bound_complexity}). The sign-specific work is done in
Chapters~\ref{chap:analysis_pos} and~\ref{chap:analysis_neg}; what remains here is common to
both signs: the counting argument of Appendix~B.1 (the paper's Lemmas~27 and~28), which
converts a sum of rate terms $\alpha(n)/n \wedge 1$ along the episodes into a logarithmic
bound (Lemma~\ref{lem:sum_counts_bound}), the definition of the constants of the bound
(Definition~\ref{def:upper_bound_const}), and the final assembly.

Standing setting. $\beta \ne 0$; $\mathcal M$ is a finite-horizon MDP (Definition~\ref{def:episodic_mdp})
with the reward function $r$ (Definition~\ref{def:reward_fn}) with values in $[0,1]$;
$\delta \in (0,1)$, $\eps > 0$, $s_1 \in \Ss$; $(X, Y, \mathrm{out})$ is a run of an algorithm
in the episode environment of $\mathcal M$ from $s_1$ (Definition~\ref{def:episode_env}) on a
probability space $(\Omega, \Prob)$, with $X_t = \pi^{t+1}$ the policy of episode $t + 1$ and
$Y_t$ its state sequence; $n_h^t(s,a)$ are the counts (Definition~\ref{def:counts}),
$\nbar_h^t(s,a) = \sum_{i \le t} p^{\pi^i}_h(s,a)$ the (random) pseudo-counts
(Definition~\ref{def:pseudo_counts}, deviation~10), $p^\pi_h(s,a)$ the occupancy measure
(Definition~\ref{def:occupancy}), and $\alpha(\cdot,\delta)$, $\alpha^*(\cdot,\delta)$,
$\alpha^{\mathrm{cnt}}(\delta)$ the rates of Definition~\ref{def:rates}. We write $S := \abs{\Ss}$,
$A := \abs{\As}$. When the algorithm is Entropic-BPI (Definition~\ref{def:entropic_bpi}), $\tau$
denotes its stopping time (Definition~\ref{def:bpi_alg}).

\section{The counting argument}
\label{sec:sc_counting}

\begin{definition}[Rate terms]\label{def:rate_min}
\lean{Essakine2026Tight.rateMin, Essakine2026Tight.klRateMinAt, Essakine2026Tight.starRateMinAt}\leanok
\uses{def:counts, def:rates}
For a function $\alpha : [0, \infty) \to \R$ and $n \in \N$ let
\[
  \rho_\alpha(n) := \min\Big\{\frac{\alpha(n)}{n},\, 1\Big\} \quad (n \ge 1), \qquad \rho_\alpha(0) := 1 .
\]
The \emph{rate terms} of a run are, for $t \in \N$, $h \in [H]$ and $(s,a) \in \Ss\times\As$,
\[
  \rho^t_h(s,a) := \rho_{\alpha(\cdot,\delta)}\big(n_h^t(s,a)\big), \qquad
  \rho^{*t}_h(s,a) := \rho_{\alpha^*(\cdot,\delta)}\big(n_h^t(s,a)\big) .
\]
(The paper writes $\alpha(n)/n \wedge 1$ with the convention $\alpha(0)/0 = \infty$; the value
$1$ for an unvisited pair is deviation~11.)
\end{definition}

\textbf{Lean remark.} \texttt{rateMin (α : ℝ → ℝ) (n : ℕ) : ℝ := if n = 0 then 1 else min (α n / n) 1};
on a run, \texttt{klRateMinAt X Y δ h s a t ω := rateMin (alphaKL S A H δ) (visitCountAt X Y h s a t ω)}
and \texttt{starRateMinAt} with \texttt{alphaStar}. The rates of Definition~\ref{def:rates} are
defined as functions of a \emph{real} argument $x \ge 0$ by the same formulas
($\alpha(x,\delta) = \log(3SAH/\delta) + S\log(8e(x+1))$, etc.): Lemma~\ref{lem:pseudo_counts_bound}
evaluates them at the real-valued pseudo-counts. Lean's $\alpha(0)/0 = 0$ never appears in a
statement thanks to the case $n = 0$ of \texttt{rateMin}.

\begin{lemma}[The rates as real functions]\label{lem:sc_rates_real}
\uses{def:rates}
Let $\delta \in (0, 1]$ and let $\alpha$ be one of $\alpha(\cdot,\delta)$, $\alpha^*(\cdot,\delta)$,
viewed as a function of a real argument $x \ge 0$. Then
\begin{enumerate}
  \item[(i)] $\alpha$ is nondecreasing on $[0,\infty)$;
  \item[(ii)] $x \mapsto \alpha(x)/x$ is nonincreasing on $[1, \infty)$;
  \item[(iii)] $\alpha(x) \ge \alpha^{\mathrm{cnt}}(\delta) = \log(3SAH/\delta) \ge \log 3 > 1$ for all $x \ge 0$.
\end{enumerate}
\end{lemma}

\begin{proof}
\uses{def:rates}
Write $a_0 := \log(3SAH/\delta)$ and $\alpha(x) = a_0 + k\log(8e(x+1))$ with $k = S$ or $k = 1$.
(i) $x \mapsto \log(8e(x+1))$ is increasing. (ii) $\alpha(x)/x = a_0/x + k\,\log(8e(x+1))/x$;
$a_0/x$ is nonincreasing on $(0,\infty)$ since $a_0 \ge 0$, and $f(x) := \log(8e(x+1))/x$ has
derivative $f'(x) = \big(\tfrac{x}{x+1} - \log(8e(x+1))\big)/x^2 < 0$ on $(0,\infty)$ because
$\tfrac{x}{x+1} < 1 < 1 + \log 8 = \log(8e) \le \log(8e(x+1))$. (iii) $k\log(8e(x+1)) \ge 0$ since
$8e(x+1) \ge 8e > 1$, and $a_0 \ge \log 3 > 1$ because $3SAH/\delta \ge 3$ ($S, A, H \ge 1$,
$\delta \le 1$) and $\log 3 > 1.09$.
\end{proof}

\begin{lemma}[Rate terms against pseudo-counts: the arithmetic]\label{lem:sc_rate_min_arith}
\uses{def:rate_min}
Let $a_0 \ge 1$ and let $\alpha : [0,\infty) \to \R$ satisfy
(i) $\alpha$ nondecreasing, (ii) $x \mapsto \alpha(x)/x$ nonincreasing on $[1,\infty)$ and
(iii) $\alpha(x) \ge a_0$ for all $x \ge 0$. Let $n \in \N$ and $\nbar \ge 0$ be such that
$n \ge \nbar/2 - a_0$. Then
\[
  \rho_\alpha(n) \le \frac{4\,\alpha(\nbar)}{\max\{\nbar, 1\}} .
\]
\end{lemma}

\begin{proof}
\uses{def:rate_min}
This is the case analysis of \cite[Lemma~7]{menard2021fast}.

\emph{Case $\nbar \ge 4a_0$.} Then $n \ge \nbar/2 - a_0 \ge \nbar/2 - \nbar/4 = \nbar/4 \ge a_0 \ge 1$;
in particular $n \ge 1$, so $\rho_\alpha(n) \le \alpha(n)/n$. By (ii) applied to
$1 \le \nbar/4 \le n$ and then (i),
\[
  \frac{\alpha(n)}{n} \le \frac{\alpha(\nbar/4)}{\nbar/4} = \frac{4\,\alpha(\nbar/4)}{\nbar} \le \frac{4\,\alpha(\nbar)}{\nbar}
  = \frac{4\,\alpha(\nbar)}{\max\{\nbar,1\}},
\]
the last equality because $\nbar \ge 4a_0 \ge 4$.

\emph{Case $\nbar < 4a_0$.} Here $\rho_\alpha(n) \le 1$ in all cases (by definition when $n = 0$,
by the minimum when $n \ge 1$), and by (iii) $4\alpha(\nbar)/\max\{\nbar,1\} \ge 4a_0/\max\{\nbar,1\}$.
If $\nbar \le 1$ the right-hand side is $4a_0 \ge 4 \ge 1$; if $\nbar > 1$ it is
$4a_0/\nbar > 4a_0/(4a_0) = 1$. Hence $\rho_\alpha(n) \le 1 \le 4\alpha(\nbar)/\max\{\nbar,1\}$.
\end{proof}

\begin{remark}[On the paper's Lemma~27]\label{rem:sc_lemma27}
The paper states Lemma~27 (from \cite[Lemma~7]{menard2021fast}) for every rate $\alpha$ that is
nondecreasing with $\alpha(x)/x$ nonincreasing on $x \ge 1$, and for $t \ge 1$. These hypotheses
do not suffice: the unvisited case $n = 0$, where the left-hand side is $1$, needs
$\alpha(\nbar) \ge \max\{\nbar,1\}/4$, which is hypothesis (iii) of
Lemma~\ref{lem:sc_rate_min_arith} together with $a_0 \ge 1$ (for $\alpha \equiv 1/8$, $n = 0$ and
$\nbar = 1/2$ the event inequality $0 \ge 1/4 - \alpha^{\mathrm{cnt}}$ holds while $1 > 4\cdot\frac18$).
The proof of M\'enard et al.\ uses, exactly as above, that the rate dominates the count rate
$\alpha^{\mathrm{cnt}}(\delta) = \log(3SAH/\delta) \ge \log 3 > 1$ of the event
$\evCnt$ (Definition~\ref{def:event_cnt}), which holds for both rates of the paper
(Lemma~\ref{lem:sc_rates_real}) and for $\delta \le 1$; this is how the lemma is stated below.
The restriction $t \ge 1$ is not needed (at $t = 0$, $n = \nbar = 0$ and the inequality reads
$1 \le 4\alpha(0)$).
\end{remark}

\begin{lemma}[Lemma~27: rate terms against pseudo-counts]\label{lem:pseudo_counts_bound}
\uses{def:rate_min, def:pseudo_counts, def:rates, def:event_cnt}
Let $\delta \in (0, 1]$ and let $\alpha$ be one of $\alpha(\cdot,\delta)$, $\alpha^*(\cdot,\delta)$.
For every run of any algorithm in the episode environment of $\mathcal M$, on the event
$\evCnt$, for all $t \in \N$, $h \in [H]$, $(s,a) \in \Ss\times\As$, with
$n := n_h^t(s,a)$ and $\nbar := \nbar_h^t(s,a)$,
\[
  \rho_\alpha(n) \le \frac{4\,\alpha(\nbar)}{\max\{\nbar, 1\}} \le \frac{4\,\alpha(t)}{\max\{\nbar,1\}} .
\]
In particular $\rho^t_h(s,a) \le 4\alpha(t,\delta)/\max\{\nbar_h^t(s,a),1\}$ and
$\rho^{*t}_h(s,a) \le 4\alpha^*(t,\delta)/\max\{\nbar_h^t(s,a),1\}$ on $\evCnt$.
\end{lemma}

\begin{proof}
\uses{lem:sc_rate_min_arith, lem:sc_rates_real, lem:pseudo_counts_increments, def:event_cnt}
On $\evCnt$, $n \ge \nbar/2 - \alpha^{\mathrm{cnt}}(\delta)$ (Definition~\ref{def:event_cnt}, with the
random pseudo-counts of deviation~10). Lemma~\ref{lem:sc_rates_real} gives the hypotheses
(i)--(iii) of Lemma~\ref{lem:sc_rate_min_arith} with $a_0 := \alpha^{\mathrm{cnt}}(\delta) \ge 1$,
which yields the first inequality. For the second, $\nbar_h^t(s,a) \le t$
(Lemma~\ref{lem:pseudo_counts_increments}: the increments are probabilities) and $\alpha$ is
nondecreasing.
\end{proof}

\begin{lemma}[Counting argument]\label{lem:sum_counts_bound}
\uses{def:rate_min, def:pseudo_counts, def:occupancy, def:event_cnt, def:rates, def:episode_env}
Let $\delta \in (0,1]$ and let $\alpha$ be one of $\alpha(\cdot,\delta)$, $\alpha^*(\cdot,\delta)$.
For every run of any algorithm in the episode environment of $\mathcal M$, with
$\pi^{t+1} := X_t$, on the event $\evCnt$, for every integer $T \ge 1$,
\[
  \sum_{t=0}^{T-1}\sum_{h=1}^H\sum_{(s,a)\in\Ss\times\As} p^{\pi^{t+1}}_h(s,a)\,\rho_\alpha\big(n_h^t(s,a)\big)
  \;\le\; 16\, S A H\, \alpha(T-1)\log(T+1) \;\le\; 16\, SAH\,\alpha(T)\log(T+1) .
\]
\end{lemma}

\begin{proof}
\uses{lem:pseudo_counts_bound, lem:pseudo_counts_increments, lem:log_sum_bound, lem:sc_rates_real}
Fix $h$ and $(s,a)$ and put $u_{t+1} := p^{\pi^{t+1}}_h(s,a) \in [0,1]$ and
$U_t := \sum_{i=1}^t u_i$ for $t \ge 0$. By Lemma~\ref{lem:pseudo_counts_increments},
$U_t = \nbar_h^t(s,a)$ (both vanish at $t = 0$ and have the same increments). By
Lemma~\ref{lem:pseudo_counts_bound}, on $\evCnt$ and for $0 \le t \le T-1$,
$\rho_\alpha(n_h^t(s,a)) \le 4\alpha(t)/\max\{U_t,1\} \le 4\alpha(T-1)/\max\{U_t,1\}$ (the rate
is nondecreasing, Lemma~\ref{lem:sc_rates_real}). Hence
\[
  \sum_{t=0}^{T-1} u_{t+1}\,\rho_\alpha\big(n_h^t(s,a)\big) \le 4\alpha(T-1)\sum_{t=0}^{T-1}\frac{u_{t+1}}{\max\{U_t,1\}}
  \le 4\alpha(T-1)\cdot 4\log(U_T + 1) \le 16\,\alpha(T-1)\log(T+1),
\]
by Lemma~\ref{lem:log_sum_bound} (applied with its $T$ equal to $T - 1 \ge 0$) and
$U_T = \nbar_h^T(s,a) \le T$. Summing over the $SAH$ triples $(h, s, a)$ gives the first bound;
the second follows from $\alpha(T-1) \le \alpha(T)$.
\end{proof}

\begin{remark}[On the paper's constant]\label{rem:sc_counting_constant}
The paper's Appendix~B.1 states this bound with the constant $3$ (``$\le 3SAH\alpha(\tau-1,\delta)\log(\tau+1)$'').
Its chain of inequalities drops the factor $4$ of Lemma~27 (it writes
$\alpha(\tilde n \vee 1)$ for $4\alpha(\tilde n)/(\tilde n \vee 1)$) and the factor $4$ of
Lemma~28; the product of the two factors is $16$. The value $\alpha(T-1)$ (rather than
$\alpha(T)$) is the one used by Lemmas~\ref{lem:stopping_time_bound_pos}
and~\ref{lem:stopping_time_bound_neg}: $\alpha(T-1,\delta) = \log(3SAH/\delta) + S\log(8eT)$
and $\alpha^*(T-1,\delta) = \log(3SAH/\delta) + \log(8eT)$ exactly.
\end{remark}

\textbf{Lean remark.} In a run, $p^{\pi^{t+1}}_h(s,a)$ is \texttt{occupancy M (X t ω) s₁ h s a}
and the pseudo-count is the finite sum of these over the past episodes, so the per-$(h,s,a)$
inequality is a statement about the real sequence $t \mapsto \texttt{occupancy M (X t ω) s₁ h s a}$
on the event; the lemma is a \texttt{Finset.sum\_le\_sum} over \texttt{Finset.range T} and over
$(h, s, a)$ on top of Lemma~\ref{lem:log_sum_bound}. Nothing about the algorithm is used
(for Entropic-BPI, $X_t = \pi^{t+1}$ is the greedy policy, Lemma~\ref{lem:entropic_bpi_run}).

\section{The sample complexity theorem}
\label{sec:sc_theorem}

\begin{definition}[Constants of the upper bound]\label{def:upper_bound_const}
\lean{Essakine2026Tight.progress, Essakine2026Tight.varianceFactor, Essakine2026Tight.coeffSqrt, Essakine2026Tight.coeffLin, Essakine2026Tight.logFactor, Essakine2026Tight.upperBound}\leanok
\uses{def:rates, def:gmax}
For $\beta \ne 0$, $\eps > 0$, $\delta \in (0,1)$, integers $S, A, H \ge 1$ and $G \ge 0$ let
\begin{align*}
  \kappa(\beta,\eps) &:= \frac{(e^{|\beta|\eps} - 1)\,e^{-2|\beta|\eps}}{2}, &
  V_G &:= \frac{(e^{|\beta| G} - 1)^2}{e^{|\beta| G}}, &
  A_0 &:= \log\frac{3SAH}{\delta}, \\
  C &:= \frac{10^9\sqrt{V_G\, S A H}}{\kappa(\beta,\eps)}, &
  D &:= \frac{10^{10}\, e^{|\beta|(H+1)} H^2 S A}{\kappa(\beta,\eps)}, &
  C_1 &:= \frac85\log\big(11\,(8e)^2\,(A_0 + S)\,(C + D)^2\big),
\end{align*}
and
\[
  \mathrm{upperBound}(S, A, H, \beta, \eps, \delta, G) := C^2 (A_0 + 1)\, C_1^2 + \big(D + 2\sqrt{D}\, C\big)(A_0 + S)\, C_1^2 + 1 .
\]
The bound is applied with $G = \Gmax(\mathcal M)$ (Definition~\ref{def:gmax}).
\end{definition}

\textbf{Lean remark.} \texttt{upperBound (S A H : ℕ) (β ε δ G : ℝ) : ℝ} in
\texttt{Essakine2026Tight/EV2026/Constants.lean}, with $S$, $A$ the cardinalities of the
state and action types. All quantities are real numbers; $\kappa > 0$ for $\eps > 0$, $\beta \ne 0$,
and $A_0 > 0$ for $\delta < 1$.

\begin{remark}[Where the constants come from; the asymptotic form]\label{rem:sc_constants}
The constants are chosen so that the recursive inequality proved in
Lemmas~\ref{lem:stopping_time_bound_pos} and~\ref{lem:stopping_time_bound_neg} is exactly the
hypothesis of the self-bounding Lemma~\ref{lem:self_bounding} (the paper's Lemma~29). We
record the computation. Let $T \ge 1$ with $T \le \tau$ and write $L := \log(8eT)$, so that
$\log(T+1) \le L$, $\alpha^*(T-1,\delta) = A_0 + L$ and $\alpha(T-1,\delta) = A_0 + SL$. On the
good event, the stopping condition fails at every $t < T$, so by
Lemma~\ref{lem:progress_pos} (resp.~\ref{lem:progress_neg}) and
Lemma~\ref{lem:ratio_bound_pos} (resp.~\ref{lem:ratio_bound_neg}, whose additive constant is
smaller), with $x_t := \sum_{h,s,a} p^{\pi^{t+1}}_h(s,a)\rho^{*t}_h(s,a)$ and
$y_t := \sum_{h,s,a} p^{\pi^{t+1}}_h(s,a)\rho^t_h(s,a)$,
\[
  T\kappa \le e^{13}\Big[36\sqrt{V_{\Gmax}}\sum_{t<T}\sqrt{x_t} + 84 H e^{|\beta|(H+1)}\sum_{t<T} y_t\Big]
  \le e^{13}\Big[36\sqrt{V_{\Gmax}\, T \textstyle\sum_{t<T} x_t} + 84 H e^{|\beta|(H+1)}\sum_{t<T} y_t\Big]
\]
by the Cauchy--Schwarz inequality, and Lemma~\ref{lem:sum_counts_bound} (on $\evCnt$) gives
$\sum_{t<T} x_t \le 16 SAH(A_0 L + L^2)$ and $\sum_{t<T} y_t \le 16SAH(A_0 L + S L^2)$. Hence
\[
  T \le C_0\sqrt{T(A_0 L + L^2)} + D_0 (A_0 L + S L^2), \qquad
  C_0 := \frac{144 e^{13}\sqrt{V_{\Gmax} SAH}}{\kappa}, \quad
  D_0 := \frac{1344 e^{13} e^{|\beta|(H+1)} H^2 SA}{\kappa}.
\]
Since $e^{13} < 4.5\cdot 10^5$, $144 e^{13} < 6.5\cdot 10^7 \le 10^9$ and
$1344 e^{13} < 6.0\cdot 10^8 \le 10^{10}$, so $C_0 \le C$ and $D_0 \le D$ and the inequality holds
with $C$, $D$. Lemma~\ref{lem:self_bounding} with $A = A_0$, $B = 1$, $E = S$, $\alpha = 8e$
(so $L = \log(\alpha T)$, $\alpha \ge e$, $1 \le B \le E$) then gives
$T \le C^2(A_0 + 1)C_1^2 + (D + 2\sqrt D C)(A_0+S)C_1^2 + 1 = \mathrm{upperBound}$.
The slack (a factor $\ge 15$ in each of $C$, $D$) is deliberate (deviation~6): the constants of
the analysis chapters may be adjusted without touching this definition.

\emph{Degenerate case.} Lemma~\ref{lem:self_bounding} is stated for $C > 0$, while
$C = 0$ when $\Gmax(\mathcal M) = 0$ (all returns vanish, $V_{\Gmax} = 0$). The recursive
inequality with $C = 0$ implies the one with any $C' > 0$, and the conclusion of
Lemma~\ref{lem:self_bounding} is continuous in $C'$ at $0$, so the bound holds by letting
$C' \downarrow 0$; the formalization should rather state Lemma~\ref{lem:self_bounding} for
$C \ge 0$ (the proof of \cite[Lemma~13]{menard2021fast} only uses $C \ge 0$).

\emph{Asymptotic form.} Since $1/\kappa^2 = 4e^{4|\beta|\eps}/(e^{|\beta|\eps}-1)^2$, the first
term is
\[
  4\cdot 10^{18}\,\frac{e^{4|\beta|\eps}}{(e^{|\beta|\eps} - 1)^2}\,
  \frac{(e^{|\beta|\Gmax(\mathcal M)} - 1)^2}{e^{|\beta|\Gmax(\mathcal M)}}\, SAH\,(A_0 + 1)\,C_1^2 ,
\]
and $C_1 = O(\log(SAH e^{|\beta| H}/(\kappa\delta)))$, so it is
$\tilde O\big(e^{4|\beta|\eps}(e^{|\beta|\eps}-1)^{-2} V_{\Gmax} SAH\log(1/\delta)\big)$: the
paper's form (Theorem~4: $e^{2\max\{\beta,0\}\eps}(e^{|\beta|\eps}-1)^{-2}V_{\Gmax}SAH$, up to
logarithms), the exponent $4|\beta|\eps$ instead of $2\max\{\beta,0\}\eps$ coming from the
common progress constant $\kappa$ of the two signs (deviation~11); under the hypothesis
$\eps \le 2/(|\beta|HS)$ of Theorem~\ref{thm:upper_bound_complexity}, $e^{4|\beta|\eps} \le e^8$
is a constant. The second term $(D + 2\sqrt D C)(A_0+S)C_1^2$ is of lower order in
$1/(e^{|\beta|\eps}-1)$ but carries $e^{|\beta|(H+1)}$ instead of $e^{|\beta|\Gmax}$; the
paper's ``in particular'' single-term bound assumes that the first term dominates, which the
hypothesis on $\eps$ does not guarantee (deviation~6). The paper's displayed two-term bound
omits the absolute constants of its own $C$ and $D$ (e.g.\ the factor $(36e^{13})^2$ of $C^2$),
and its $C_1$ is not the $C_1$ of its Lemma~29 evaluated at its $C$, $D$; we use
Lemma~\ref{lem:self_bounding} verbatim.
\end{remark}

\begin{theorem}[Theorem~4, PAC guarantee]\label{thm:upper_bound_pac}
\lean{Essakine2026Tight.isEntropicPAC_entropicBPI}\leanok
\uses{def:entropic_bpi, def:entropic_pac, def:reward_fn, def:episode_env}
Let $\beta \ne 0$, $H \ge 1$, $\Ss$ and $\As$ finite, $r : [H]\times\Ss\times\As \to [0,1]$ a
reward function, $\delta \in (0,1)$, $\eps \in (0, 2/(|\beta| H S)]$ and $s_1 \in \Ss$. Then
Entropic-BPI with parameters $(r, \beta, \delta, \eps, s_1)$ (Definition~\ref{def:entropic_bpi})
is $(\eps,\delta)$-PAC for entropic best-policy identification with reward function $r$,
parameter $\beta$ and initial state $s_1$ (Definition~\ref{def:entropic_pac}): for every
$\mathcal M \in \mathcal M_r$, the output $\hat\pi$ of Entropic-BPI in the episode environment
of $\mathcal M$ from $s_1$ satisfies
\[
  \Prob\big(V^*_1(s_1) - V^{\hat\pi}_1(s_1) > \eps\big) \le \delta .
\]
\end{theorem}

\begin{proof}
\uses{lem:pac_pos, lem:pac_neg, def:entropic_pac, def:bpi_alg, def:episode_env, lem:entropic_bpi_run}
Fix $\mathcal M \in \mathcal M_r$. The PAC property is a statement about the law of the output
of the algorithm in the episode environment of $\mathcal M$, which is the same for every run
(Definition~\ref{def:bpi_alg}); a run exists (the canonical run of LML). Let
$(X, Y, \mathrm{out})$ be a run on $(\Omega, \Prob)$. If $\beta > 0$, Lemma~\ref{lem:pac_pos}
gives $\Prob(V^*_1(s_1) - V^{\mathrm{out}}_1(s_1) > \eps) \le \delta$; if $\beta < 0$,
Lemma~\ref{lem:pac_neg} gives the same. (Both lemmas hold for every $\delta \in (0,1)$ and
$\eps > 0$; the hypothesis $\eps \le 2/(|\beta|HS)$ is not used, see
Remark~\ref{rem:sc_hypotheses}.) Since the law of $\mathrm{out}$ is the output law, the
output law gives mass at most $\delta$ to the bad set $\{\pi : V^*_1(s_1) - V^\pi_1(s_1) > \eps\}$
of $\mathcal M$; this for every $\mathcal M \in \mathcal M_r$ is the PAC property.
\end{proof}

\textbf{Lean remark.} \texttt{Essakine2026Tight.isEntropicPAC\_entropicBPI}. LML's
\texttt{IdentAlg.IsPAC} is a property of \texttt{IdentAlg.outputMeasure alg env}: the law of the
output, defined as the composition of the law of the stopped history with the output kernel;
\texttt{IsEntropicPAC 𝒜 r β ε δ s₁ := 𝒜.IsPAC (fun M : \{M // M.HasRewardFn r\} ↦ statesEnv M.1 s₁) (fun M π ↦ ε < optEntropicValue M.1 β 0 s₁ - entropicValue M.1 β π 0 s₁) δ}.
Lemmas~\ref{lem:pac_pos} and~\ref{lem:pac_neg} are statements about a run
(\texttt{IdentAlg.IsRun}) on an arbitrary probability space; the transfer to
\texttt{outputMeasure} is \texttt{IdentAlg.IsRun.hasLaw\_output} (the law of the output of a
run is \texttt{outputMeasure}) applied to the canonical run
\texttt{IdentAlg.runMeasure}/\texttt{IdentAlg.isRun\_runMeasure} on the Ionescu-Tulcea space of
trajectories, or directly the criterion \texttt{IdentAlg.IsPAC.of\_forall\_isRun}, which
quantifies over probability spaces in the universe of the trajectory space (as in the sister
project \texttt{colt-2026-83}, \texttt{LinearBandit.isPAC\_of\_forall\_isRun}). Conversely
\texttt{IsPAC.measureReal\_bad\_of\_isRun} gives $\Prob(\mathrm{bad}) \le \delta$ in every run; this
direction is the one used by the lower bound (Lemma~\ref{lem:hard_pac_event}). The event
$\{V^*_1(s_1) - V^{\mathrm{out}}_1(s_1) > \eps\}$ is measurable since the policy type is finite
(discrete $\sigma$-algebra).

\begin{theorem}[Theorem~4, sample complexity]\label{thm:upper_bound_complexity}
\lean{Essakine2026Tight.probReal_stoppingTime_entropicBPI_le_ge}\leanok
\uses{def:entropic_bpi, def:upper_bound_const, def:gmax, def:bpi_alg, def:reward_fn, def:episode_env}
Under the hypotheses of Theorem~\ref{thm:upper_bound_pac}, for every $\mathcal M \in \mathcal M_r$
and every run $(X, Y, \mathrm{out})$ of Entropic-BPI in the episode environment of $\mathcal M$
from $s_1$ on a probability space $(\Omega, \Prob)$, the number of episodes $\tau$ satisfies
\[
  \Prob\big(\tau \le \mathrm{upperBound}(S, A, H, \beta, \eps, \delta, \Gmax(\mathcal M))\big) \ge 1 - \delta ;
\]
in particular $\Prob(\tau < \infty) \ge 1 - \delta$.
\end{theorem}

\begin{proof}
\uses{lem:stopping_time_bound_pos, lem:stopping_time_bound_neg, lem:good_event, def:good_event}
Fix $\mathcal M \in \mathcal M_r$ and a run. If $\beta > 0$, Lemma~\ref{lem:stopping_time_bound_pos}
gives $\tau \le \mathrm{upperBound}(S,A,H,\beta,\eps,\delta,\Gmax(\mathcal M))$ (in particular
$\tau < \infty$) at every $\omega \in \evGood^+$, and $\Prob(\evGood^+) \ge 1 - \delta$ by
Lemma~\ref{lem:good_event}. If $\beta < 0$, the same holds with
Lemma~\ref{lem:stopping_time_bound_neg} and $\evGood^-$. Hence
$\{\tau \le \mathrm{upperBound}\} \supseteq \evGood^\pm$ has probability at least $1 - \delta$.
\end{proof}

\textbf{Lean remark.} \texttt{Essakine2026Tight.probReal\_stoppingTime\_entropicBPI\_le\_ge}.
The stopping time \texttt{IdentAlg.stoppingTime} of a run takes values in $\N \cup \{\infty\}$
(\texttt{ℕ∞}); the event is $\{\omega : (\tau(\omega) : \texttt{ℝ≥0∞}) \le \texttt{ENNReal.ofReal}\ \mathrm{upperBound}\}$,
which contains $\tau < \infty$. The run hypothesis (\texttt{IdentAlg.IsRun}) makes
$\tau$ the hitting time of the stopping set by the history process, so that the good event and
$\{\tau \le B\}$ are events of the same probability space; no law transfer is needed here, the
statement being about an arbitrary run (deviation~6 separates the two conclusions of Theorem~4
because this one is naturally a statement about runs while the PAC one is about the output law).

\begin{remark}[On the hypotheses]\label{rem:sc_hypotheses}
The paper states Theorem~4 for $\delta \in [0,1]$ and $\eps > 0$ ``small enough'', made precise
in the appendix as $\eps \in (0, 2/(|\beta|HS)]$. We require $\delta \in (0,1)$ (for $\delta = 0$
the rates are infinite, for $\delta = 1$ the statements are empty). The hypothesis on $\eps$ is
kept for fidelity to the paper, but it is not used by the proof as decomposed in
Chapters~\ref{chap:analysis_pos}--\ref{chap:analysis_neg} and above: every lemma holds for all
$\eps > 0$, the bound $\mathrm{upperBound}$ being valid for every $\eps > 0$ (it only becomes
weak, through $1/\kappa$, when $|\beta|\eps$ is large); its only role is to make
$e^{4|\beta|\eps} \le e^8$ a constant in Remark~\ref{rem:sc_constants}. Finally, the theorem
is stated for the class $\mathcal M_r$ of MDPs with the reward function $r$ known to the
algorithm (deviation~7), which is the paper's setting (``rewards are known and
deterministic'').
\end{remark}

\chapter{Lower bound}
\label{chap:lower_bound}

This chapter proves Theorem~3 of the paper (Section~3, proof in Appendix~C): on a hard
instance $\mathcal M_0$, every $(\eps,\delta)$-PAC algorithm uses in expectation at least
$\mathrm{lowerBound}(S,A,H,\beta,\eps,\delta,\Gmax(\mathcal M_0))$ episodes
(Theorem~\ref{thm:lower_bound}, Definition~\ref{def:lower_bound_const}). The hard instances are
those of \cite{domingues2021episodic}, with the leaf transition probabilities of the paper: a
waiting state $s_w$, an $A$-ary tree of $S - 3$ states, and two absorbing states $s_g$
(rewarding from the step $\tilde H$ on) and $s_b$. An episode of a deterministic policy
realizes exactly one \emph{triple} $u = (\text{exit step}, \text{leaf}, \text{leaf action})$, and
the instance $\mathcal M_u$ differs from $\mathcal M_0$ only in the probability of reaching $s_g$
after the triple $u$ ($p_+$ instead of $p_-$). The family thus behaves like a bandit with
$|\mathcal U| = \bar H L A = \Theta(SAH)$ arms, and the change of measure of
\cite[Lemma~1]{kaufmann2016complexity} (Theorem~\ref{thm:change_of_measure}, the paper's
Lemma~17) gives the bound.

What differs from the paper (deviations~5, 7 and~8, and the corrections recorded in
Remarks~\ref{rem:lb_condition_a}, \ref{rem:lb_params}, \ref{rem:lb_kl_paper}
and~\ref{rem:lb_constant}):
\begin{enumerate}
  \item[(i)] the tree is not required to be full: its nodes are the first $S - 3$ nodes of the
        infinite $A$-ary tree in breadth-first order, its leaves the childless nodes, and an
        action whose child is missing leads to the first child (Definition~\ref{def:hard_tree},
        Lemma~\ref{lem:tree_leaves}); the leaf action is played on arrival at the leaf, at a step
        that depends on the depth of the leaf;
  \item[(ii)] the PAC hypothesis is the one of Definition~\ref{def:entropic_pac} for the reward
        function $r_0$ of the hard instances only (deviation~7), and the bad event of the PAC
        property is the strict inequality $V^* - V^{\hat\pi} > \eps$, so the gap $\Delta = p_+ - p_-$
        is chosen so that a policy missing the triple is \emph{strictly} more than
        $\eps$-suboptimal (Definition~\ref{def:hard_params});
  \item[(iii)] Condition~A is stated with the effective horizon $H' = H - \lfloor H/3\rfloor - d$
        of the construction and with a margin ($p_+ \le (1 + p_-)/2$): the paper's condition,
        stated with $H$, does not imply admissibility with $H'$, and its bound on the binary
        divergence $\klbin(p_-, p_+)$ is false near the boundary of its condition
        (Remark~\ref{rem:lb_kl_paper});
  \item[(iv)] the change of measure is applied in LML generality to the two stationary
        environments $\mathrm{statesEnv}(\mathcal M_0)$, $\mathrm{statesEnv}(\mathcal M_u)$,
        whose actions are the policies (Lemma~\ref{lem:hard_change_of_measure});
  \item[(v)] the constant $1/1650$ becomes $1/5000$ (the proof gives $1/3456$).
\end{enumerate}
Throughout, $S := \abs{\Ss} \ge 6$ and $A := \abs{\As} \ge 2$ are the numbers of states and
actions, $\beta \ne 0$, $\eps > 0$, and $q := e^{|\beta|\eps} > 1$.

\section{The tree}
\label{sec:lb_tree}

\begin{definition}[Depth of the tree; Assumption~1]\label{def:tree_depth}
\lean{Essakine2026Tight.treeDepth}\leanok
For integers $S \ge 6$ and $A \ge 2$ let $N := (S - 3)(A - 1) + 1$ and
\[
  d := \lceil \log_A N \rceil = \big\lceil \log_A\big((S-3)(A-1) + 1\big)\big\rceil .
\]
\emph{Assumption~1} of the paper is $H \ge 3d$.
\end{definition}

\textbf{Lean remark.} \texttt{treeDepth (S A : ℕ) : ℕ := Nat.clog A ((S - 3) * (A - 1) + 1)}
(Mathlib's \texttt{Nat.clog b n} is $\lceil\log_b n\rceil$).

\begin{lemma}[Properties of the depth]\label{lem:tree_depth_props}
\uses{def:tree_depth}
For $S \ge 6$ and $A \ge 2$: $d \ge 2$, $A^{d-1} < N \le A^d$, and consequently
\[
  \frac{A^{d-1} - 1}{A - 1} < S - 3 \le \frac{A^d - 1}{A - 1} .
\]
Since $\sum_{j=0}^{k} A^j = (A^{k+1} - 1)/(A - 1)$ is the number of nodes of depth at most $k$ of
the infinite $A$-ary tree, this says that the full tree of depth $d - 2$ has fewer than $S - 3$
nodes and that the full tree of depth $d - 1$ has at least $S - 3$ nodes.
\end{lemma}

\begin{proof}
\uses{def:tree_depth}
$N \ge 3(A - 1) + 1 = 3A - 2 > A$ because $A \ge 2$, so $\log_A N > 1$ and $d \ge 2$. The
ceiling gives $d - 1 < \log_A N \le d$, i.e.\ $A^{d-1} < N \le A^d$ (Mathlib:
\texttt{Nat.pow\_pred\_clog\_lt\_self} for $N \ge 2$ and \texttt{Nat.le\_pow\_clog}). Subtracting
$1$ and dividing by $A - 1 \ge 1$, using $N - 1 = (S-3)(A-1)$, gives the displayed inequalities.
\end{proof}

\begin{definition}[The tree]\label{def:hard_tree}
\uses{def:tree_depth, lem:tree_depth_props}
The tree states are $x_0, x_1, \dots, x_{S-4}$ ($S - 3$ nodes), $x_0$ being the root, numbered in
the breadth-first order of the infinite $A$-ary tree. For $j \ge 1$ the \emph{parent} of $x_j$ is
$x_{\mathrm{par}(j)}$ with $\mathrm{par}(j) := \lfloor (j - 1)/A \rfloor$, and the \emph{depth} is
$\mathrm{dep}(x_0) := 0$, $\mathrm{dep}(x_j) := \mathrm{dep}(x_{\mathrm{par}(j)}) + 1$; equivalently
$\mathrm{dep}(x_j) = \lfloor \log_A((A-1)j + 1)\rfloor$ is the $k$ with
$(A^k - 1)/(A-1) \le j < (A^{k+1} - 1)/(A-1)$. The \emph{children} of $x_i$ are the nodes
$x_{Ai + a}$, $a \in [A] = \{1, \dots, A\}$, with $Ai + a \le S - 4$; note that
$\mathrm{par}(Ai + a) = i$ for $a \in [A]$. The node $x_i$ is \emph{internal} if it has at least
one child, i.e.\ $Ai + 1 \le S - 4$, and a \emph{leaf} otherwise; $\mathcal L$ denotes the set of
leaves and $L := \abs{\mathcal L}$. For an internal node $x_i$ and $a \in [A]$ the \emph{child
function} is
\[
  \mathrm{child}(x_i, a) := \begin{cases} x_{Ai + a} & \text{if } Ai + a \le S - 4, \\ x_{Ai + 1} & \text{otherwise (the first child),}\end{cases}
\]
so that $\mathrm{dep}(\mathrm{child}(x, a)) = \mathrm{dep}(x) + 1$ for every internal $x$ and every $a$.
\end{definition}

\textbf{Lean remark.} The state type of the hard instances has $S$ elements: an inductive type
\texttt{HardState S} with constructors \texttt{w}, \texttt{g}, \texttt{b} and
\texttt{node (j : Fin (S - 3))} (or \texttt{Fin S} with $s_w, s_g, s_b = 0, 1, 2$ and
$x_j = j + 3$), with \texttt{Fintype.card (HardState S) = S}. Actions are \texttt{Fin A}, the
paper's $a \in [A]$ being \texttt{a.val + 1}. The depth is best \emph{defined} by the closed
form \texttt{Nat.log A ((A - 1) * j + 1)} and the recursion proved as a lemma; the child
function and leaf predicate are index arithmetic (\texttt{A * i + a + 1 ≤ S - 4}).

\begin{lemma}[Leaves]\label{lem:tree_leaves}
\uses{def:hard_tree, lem:tree_depth_props}
For $S \ge 6$ and $A \ge 2$:
\begin{enumerate}
  \item[(i)] the internal nodes are $x_0, \dots, x_{I-1}$ with $I := \lceil (S-4)/A \rceil$, every
        internal node except possibly $x_{I-1}$ has all its $A$ children, and
        $L = (S - 3) - \lceil (S-4)/A \rceil$;
  \item[(ii)] $L \ge (S-3)/2 \ge S/4$ (in particular $L \ge S/6$), and $L \ge 2$;
  \item[(iii)] $\mathrm{dep}(x_j) \le d - 1$ for every $j \le S - 4$;
  \item[(iv)] for every $j \ge 1$, $a_j := j - A\,\mathrm{par}(j) \in [A]$ and
        $\mathrm{child}(x_{\mathrm{par}(j)}, a_j) = x_j$; hence for every node $x$ there is a
        sequence of $\mathrm{dep}(x)$ actions leading from $x_0$ to $x$ through the child function.
\end{enumerate}
\end{lemma}

\begin{proof}
\uses{def:hard_tree, lem:tree_depth_props}
(i) $x_i$ is internal iff $Ai + 1 \le S - 4$ iff $i \le (S-5)/A$ iff
$i \le \lfloor (S-5)/A \rfloor = \lceil (S-4)/A \rceil - 1$ (for integers $m \ge 1$,
$\lfloor (m-1)/A \rfloor + 1 = \lceil m/A \rceil$). It has all its children iff $Ai + A \le S - 4$
iff $i \le (S-4)/A - 1$, which holds for $i \le I - 2 < (S-4)/A - 1$. The leaves are the
remaining $S - 3 - I$ nodes.

(ii) $\lceil (S-4)/A \rceil \le \lceil (S-4)/2 \rceil \le (S-3)/2$ (for $S$ even
$\lceil (S-4)/2\rceil = (S-4)/2$, for $S$ odd it is $(S-3)/2$), so
$L \ge (S-3) - (S-3)/2 = (S-3)/2$; and $(S-3)/2 \ge S/4$ iff $S \ge 6$. As $L$ is an integer
and $(S-3)/2 \ge 3/2$, $L \ge 2$.

(iii) By strong induction on $j$, $\mathrm{dep}(x_j) = k$ iff
$(A^k - 1)/(A-1) \le j < (A^{k+1}-1)/(A-1)$: for $j = 0$ both sides say $k = 0$; for $j \ge 1$ with
$i = \mathrm{par}(j)$, i.e.\ $Ai + 1 \le j \le Ai + A$, if $(A^{k-1}-1)/(A-1) \le i < (A^k-1)/(A-1)$
then $j \ge A(A^{k-1}-1)/(A-1) + 1 = (A^k - 1)/(A-1)$ and
$j \le A i + A \le A\big((A^k-1)/(A-1) - 1\big) + A = (A^{k+1} - A)/(A-1) < (A^{k+1}-1)/(A-1)$.
Since $j \le S - 4 < S - 3 \le (A^d - 1)/(A-1)$ (Lemma~\ref{lem:tree_depth_props}), the depth of
$x_j$ is at most $d - 1$.

(iv) $j - 1 - A\lfloor (j-1)/A \rfloor \in \{0, \dots, A-1\}$, so $a_j \in [A]$, and
$A\,\mathrm{par}(j) + a_j = j \le S - 4$, so $x_{\mathrm{par}(j)}$ is internal and
$\mathrm{child}(x_{\mathrm{par}(j)}, a_j) = x_j$. Iterating along the ancestors
$x_j, x_{\mathrm{par}(j)}, \dots, x_0$ (a chain of length $\mathrm{dep}(x_j)$) gives the action
sequence.
\end{proof}

One checks in the same way that every leaf has depth $d - 2$ or $d - 1$; this is not needed.

\section{Parameters of the hard instances}
\label{sec:lb_params}

\begin{definition}[Condition~A]\label{def:condition_a}
\lean{Essakine2026Tight.ConditionA, Essakine2026Tight.effHorizon}\leanok
Let $\beta \ne 0$, $\eps > 0$, let $H' \ge 1$ be an integer (the \emph{effective horizon},
$H' = H - \lfloor H/3\rfloor - d$ in Definition~\ref{def:hard_params}), and put
$c := e^{|\beta| H'} - 1 > 0$ and $q := e^{|\beta|\eps} > 1$. \emph{Condition~A} for $(\beta, H', \eps)$
is
\[
  q - 1 \le \frac{c\,(2c+1)}{4\,(3c+2)} \quad (\beta > 0), \qquad\qquad
  q - 1 \le \frac{c}{10c + 8} \;\Big(\text{i.e.\ } q \le \frac{11c + 8}{10c + 8}\Big) \quad (\beta < 0).
\]
\end{definition}

\textbf{Lean remark.} \texttt{ConditionA (β : ℝ) (H' : ℕ) (ε : ℝ) : Prop}, with the two
cases by \texttt{if 0 < β then \dots\ else \dots}; Theorem~\ref{thm:lower_bound} assumes it for
\texttt{H' := H - H / 3 - treeDepth S A} (natural-number division).

\begin{definition}[Parameters of the hard instances]\label{def:hard_params}
\uses{def:tree_depth}
Let $S \ge 6$, $A \ge 2$, $H \ge 3d$ (Definition~\ref{def:tree_depth}), $\beta \ne 0$ and $\eps > 0$.
Define
\[
  \bar H := \lfloor H/3 \rfloor, \qquad \tilde H := \bar H + d + 1, \qquad H' := H - \bar H - d, \qquad
  c := e^{|\beta| H'} - 1, \qquad q := e^{|\beta|\eps}, \qquad m := \frac{3c + 2}{2(c+1)} = 1 + \frac{c}{2(c+1)},
\]
and the leaf probabilities
\begin{align*}
  \beta > 0:\qquad & p_- := \frac{1}{2(c+1)}, & \Delta &:= \frac{(q - 1)(3c+2)}{c(c+1)}, & p_+ &:= p_- + \Delta ; \\
  \beta < 0:\qquad & p_- := 1 - \frac{1}{2(c+1)}, & \Delta &:= \frac{(q-1)(3c+2)}{c(c+1)(2q-1)}, & p_+ &:= p_- + \Delta .
\end{align*}
These satisfy the identities (direct computation, using $c\Delta = 2(q-1)m$ for $\beta > 0$ and
$c\Delta = 2(q-1)m/(2q-1)$ for $\beta < 0$)
\[
  \beta > 0:\quad 1 + c\,p_- = m, \quad 1 + c\,p_+ = (2q - 1)\,m; \qquad\qquad
  \beta < 0:\quad 1 + c\,(1 - p_-) = m, \quad 1 + c\,(1 - p_+) = \frac{m}{2q-1} .
\]
Under the standing hypotheses: $d \ge 2$ (Lemma~\ref{lem:tree_depth_props}), hence $H \ge 6$,
$\bar H \ge 2$, and $\bar H \le H/3$, $d \le H/3$ give $H' \ge H/3 \ge 2$; thus $c > 0$,
$\tilde H = H - H' + 1 \le H - 1$ and $\bar H + d = H - H' \le H - 2$. Only $(\beta, \eps, H')$ enter
$c, q, m, p_\pm, \Delta$.
\end{definition}

\textbf{Lean remark.} \texttt{hardPMinus β H' : ℝ}, \texttt{hardDelta β H' ε : ℝ},
\texttt{hardPPlus β H' ε : ℝ} as functions of $(\beta, H', \eps)$ (case split on the sign of $\beta$),
and \texttt{effHorizon S A H := H - H / 3 - treeDepth S A}, \texttt{waitHorizon H := H / 3},
\texttt{rewardStart S A H := H / 3 + treeDepth S A + 1}.

\begin{remark}[Design of the parameters; comparison with the paper]\label{rem:lb_params}
By Lemma~\ref{lem:hard_exp_value} the exponential value of a policy $\pi$ in an instance is
$1 + c\,p(\pi)$ ($\beta > 0$) or $e^{\beta H'}(1 + c(1 - p(\pi)))$ ($\beta < 0$), where
$p(\pi) \in \{p_-, p_+\}$ is its probability of reaching $s_g$. The identities of
Definition~\ref{def:hard_params} say that the ratio of the exponential values of a policy with
$p_+$ and of one with $p_-$ is $2q - 1$ (resp.\ $1/(2q-1)$), so that the entropic gap between them
is $|\beta|^{-1}\log(2q - 1) > |\beta|^{-1}\log q = \eps$ (Lemma~\ref{lem:hard_eps_suboptimal}).
The paper chooses $\Delta$ so that a policy with success probability $p_- + \Delta/2$ (a
\emph{randomized} policy realizing the triple with probability $1/2$) is exactly
$\eps$-suboptimal: $\beta > 0$: $(1 + c(p_- + \Delta))/(1 + c(p_- + \Delta/2)) = q$, i.e.\
$\Delta = (q-1)(3c+2)/(c(c+1)(2 - q))$, which is our $\Delta$ times $1/(2-q)$; $\beta < 0$: the
analogous equation has the solution $(q-1)(3c+2)/(c(c+1)(2q-1))$, which is our $\Delta$, while
the paper writes $(q-1)(3c+2)/(c(c+1)(q - 1/2))$, twice the solution of its own equation.
Policies being deterministic here (Definition~\ref{def:policy}), only the gap between $p_+$ and
$p_-$ matters; it must exceed $\eps$ strictly since the bad event of
Definition~\ref{def:entropic_pac} is $V^* - V^{\hat\pi} > \eps$, and our choice is the
paper's design with the midpoint policy $p_- + \Delta/2$ made exactly $\eps$-suboptimal with
respect to the baseline $p_-$ (indeed $1 + c(p_- + \Delta/2) = q\,m$ for $\beta > 0$).
\end{remark}

\begin{lemma}[Admissibility of the parameters]\label{lem:hard_params_valid}
\uses{def:hard_params, def:condition_a}
Let $\beta \ne 0$, $\eps > 0$, $H' \ge 1$, and let $c, q, m, p_-, \Delta, p_+$ be as in
Definition~\ref{def:hard_params}. Then $c > 0$, $q > 1$, $\Delta > 0$, $0 < p_- < 1$,
$1 - p_- = (2c+1)/(2(c+1))$ if $\beta > 0$ and $1 - p_- = 1/(2(c+1))$ if $\beta < 0$, and
Condition~A for $(\beta, H', \eps)$ is equivalent to $\Delta \le (1 - p_-)/2$. Under Condition~A:
\[
  0 < p_- < p_+ \le \frac{1 + p_-}{2} < 1 \qquad\text{and}\qquad p_+(1 - p_+) \ge \frac{1}{8(c+1)} .
\]
\end{lemma}

\begin{proof}
\uses{def:hard_params, def:condition_a}
$c = e^{|\beta|H'} - 1 > 0$ since $|\beta| H' > 0$, and $q > 1$ since $|\beta|\eps > 0$; hence
$\Delta > 0$ (all factors are positive; $2q - 1 > 1$). $p_- = 1/(2(c+1)) \in (0, 1/2]$ for
$\beta > 0$ and $p_- = 1 - 1/(2(c+1)) \in [1/2, 1)$ for $\beta < 0$; the values of $1 - p_-$ follow.

\emph{Equivalence.} For $\beta > 0$, $(1-p_-)/2 = (2c+1)/(4(c+1))$ and
\[
  \Delta \le \frac{2c+1}{4(c+1)} \iff \frac{(q-1)(3c+2)}{c(c+1)} \le \frac{2c+1}{4(c+1)}
  \iff 4(q-1)(3c+2) \le c(2c+1) \iff q - 1 \le \frac{c(2c+1)}{4(3c+2)} .
\]
For $\beta < 0$, $(1-p_-)/2 = 1/(4(c+1))$ and
\[
  \Delta \le \frac{1}{4(c+1)} \iff 4(q-1)(3c+2) \le c(2q-1) \iff q(12c + 8) - 12c - 8 \le 2cq - c
  \iff q(10c + 8) \le 11c + 8 \iff q - 1 \le \frac{c}{10c+8} .
\]

\emph{Consequences.} $p_+ > p_-$ since $\Delta > 0$, and $p_+ = p_- + \Delta \le p_- + (1-p_-)/2 = (1+p_-)/2 < 1$
since $p_- < 1$. For the product: $p_+ \ge p_-$ and $1 - p_+ \ge (1 - p_-)/2$. If $\beta > 0$,
$p_- = 1/(2(c+1))$ and $(1-p_-)/2 = (2c+1)/(4(c+1)) \ge 1/4$ (as $2c + 1 \ge c + 1$), so
$p_+(1-p_+) \ge 1/(8(c+1))$. If $\beta < 0$, $p_- \ge 1/2$ and $(1-p_-)/2 = 1/(4(c+1))$, so again
$p_+(1-p_+) \ge 1/(8(c+1))$.
\end{proof}

\begin{remark}[On Condition~A]\label{rem:lb_condition_a}
\emph{The effective horizon.} The paper states Condition~A with $c = e^{|\beta|H} - 1$ while its
construction uses $c' = e^{|\beta|H'} - 1$ with $H' \le H$. The admissibility thresholds
$\varphi_+(c) := c(2c+1)/(4(3c+2))$ and $\varphi_-(c) := c/(10c+8)$ of
Definition~\ref{def:condition_a} are increasing in $c$ (the derivative of $c(2c+1)/(3c+2)$ is
$(6c^2 + 8c + 2)/(3c+2)^2 > 0$, that of $c/(10c+8)$ is $8/(10c+8)^2 > 0$), and so are the
paper's ($4(c+1)^2/(2c^2+7c+4)$ has derivative of the sign of $12c^2 + 16c + 4 > 0$; the other
one is $\varphi_-$). Since $c' \le c$, a condition stated with $H$ is \emph{weaker} than the same
condition with $H'$ and does not imply admissibility ($p_+ < 1$) for the construction: this is
the monotonicity issue of the outline, and the reason why Condition~A is stated with $H'$. As
$H' \ge H/3$ and the thresholds are increasing, Condition~A for $(\beta, \lceil H/3 \rceil, \eps)$,
or the simple sufficient condition $e^{|\beta|\eps} - 1 \le \min\{c, 1\}/20$ (both signs; for
$c \le 1$, $\varphi_+(c) \ge c/20$ and $\varphi_-(c) \ge c/18$, and for $c \ge 1$,
$\varphi_+(c) \ge 3/20$, $\varphi_-(c) \ge 1/18$), implies it.

\emph{The margin.} The paper's Condition~A is exactly the condition $p_+ < 1$ for its own
parameters (for $\beta > 0$, $p_+ < 1$ iff $q < 4(c+1)^2/(2c^2+7c+4)$, and the cap $16/13$ is used
in its bound on $\klbin(p_-,p_+)$; for $\beta < 0$, $p_+ < 1$ iff $q < (11c+8)/(10c+8)$, and its
cap $11/10$ is redundant since $(11c+8)/(10c+8) < 11/10$). Ours requires the margin
$p_+ \le (1 + p_-)/2$, which bounds $p_+(1 - p_+)$ from below (Lemma~\ref{lem:hard_params_valid})
and is what the divergence bound of Lemma~\ref{lem:hard_kl_params} needs; without a margin that
bound is false (Remark~\ref{rem:lb_kl_paper}). For $\beta < 0$ our condition coincides with the
paper's (with $H'$ for $H$); for $\beta > 0$ neither implies the other in general, ours being
more demanding for $c \le 3.1$ and less demanding for $c \ge 3.2$ (no cap at $16/13$).
\end{remark}

\section{The hard instances}
\label{sec:lb_instances}

\begin{definition}[The hard instances]\label{def:hard_mdp}
\uses{def:hard_tree, def:hard_params, def:episodic_mdp, def:reward_fn, lem:tree_leaves}
Let $S \ge 6$, $A \ge 2$, $H \ge 3d$, $\beta \ne 0$, $\eps > 0$ and the parameters of
Definition~\ref{def:hard_params}. The state space is
$\Ss := \{s_w, s_g, s_b\} \cup \{x_0, \dots, x_{S-4}\}$ ($S$ states: the waiting state, the good
and the bad absorbing states, and the tree), the action space is $\As = [A]$ with the fixed
\emph{waiting action} $a_w := 1$, the horizon is $H$ and the initial state is $s_1 := s_w$. The set
of \emph{triples} is
\[
  \mathcal U := [\bar H] \times \mathcal L \times \As, \qquad \abs{\mathcal U} = \bar H\, L\, A,
\]
and the \emph{arrival step} of $u = (h_e, \ell^*, a^*) \in \mathcal U$ is
$h_u := h_e + 1 + \mathrm{dep}(\ell^*) \le \bar H + d \le H - 2$ (Lemma~\ref{lem:tree_leaves}~(iii),
Definition~\ref{def:hard_params}). The map $u \mapsto (h_u, \ell^*, a^*)$ is injective on
$\mathcal U$ ($h_e = h_u - 1 - \mathrm{dep}(\ell^*)$). The \emph{leaf success probabilities} are,
for $h \in [H]$, $\ell \in \mathcal L$, $a \in \As$:
\[
  p^{0}_h(\ell, a) := p_- ; \qquad
  p^{u}_h(\ell, a) := \begin{cases} p_+ & \text{if } (h, \ell, a) = (h_u, \ell^*, a^*), \\ p_- & \text{otherwise.}\end{cases}
\]
For $\mathcal M \in \{\mathcal M_0\} \cup \{\mathcal M_u : u \in \mathcal U\}$, with the
corresponding $p^{\mathcal M}_h$, the transition kernel at step $h \in [H]$ is
\begin{align*}
  p_h(\cdot \mid s_w, a) &:= \delta_{s_w} \text{ if } a = a_w \text{ and } h < \bar H, \qquad \delta_{x_0} \text{ otherwise}; \\
  p_h(\cdot \mid x, a) &:= \delta_{\mathrm{child}(x, a)} \quad \text{for an internal node } x; \\
  p_h(\cdot \mid \ell, a) &:= p^{\mathcal M}_h(\ell,a)\,\delta_{s_g} + \big(1 - p^{\mathcal M}_h(\ell,a)\big)\,\delta_{s_b} \quad \text{for } \ell \in \mathcal L; \\
  p_h(\cdot \mid s, a) &:= \delta_s \quad \text{for } s \in \{s_g, s_b\} .
\end{align*}
The reward function (the same for all instances) is
\[
  r_0 : \quad r_h(s, a) := \indic\{s = s_g\}\,\indic\{h \ge \tilde H\} \in \{0, 1\},
\]
deterministic and with values in $[0,1]$, so that $\mathcal M_0, \mathcal M_u \in \mathcal M_{r_0}$
(Definition~\ref{def:reward_fn}). Thus $\mathcal M_u$ differs from $\mathcal M_0$ at exactly one
(step, state, action), namely $(h_u, \ell^*, a^*)$, where the probability of $s_g$ is $p_+$
instead of $p_-$. The definition requires $p_\pm \in [0,1]$, which holds under Condition~A
(Lemma~\ref{lem:hard_params_valid}).
\end{definition}

\textbf{Lean remark.} \texttt{hardMDP S A H β ε (u : Option (HardTriple S A H)) : EpisodicMDP (HardState S) (Fin A) H}
with $\mathcal M_0 = $ \texttt{hardMDP \dots\ none}, built from a transition function
\texttt{Fin H → HardState S → Fin A → Measure (HardState S)} by \texttt{Kernel.ofFunOfCountable}
(finite discrete types) and the deterministic reward \texttt{hardReward S A H} through
\texttt{EpisodicMDP.ofDet}. The leaf transition is the Bernoulli measure of the parameter
\emph{clamped} to $[0,1]$ (\texttt{Set.projIcc}), so that the definition makes sense for all
parameters and agrees with the above under Condition~A. The triple type is
\texttt{HardTriple S A H := Fin (H / 3) × Leaf S A × Fin A}, the paper's $h_e \in [\bar H]$ being
\texttt{h.val + 1}.

\begin{definition}[Exit step, leaf, leaf action; the triple of a policy]\label{def:hard_triple}
\uses{def:hard_mdp, def:hard_tree, def:policy, lem:tree_leaves}
Let $\pi = (\pi_h)_{h \le H}$ be a policy (Definition~\ref{def:policy}). Its \emph{exit step} is
\[
  h_e(\pi) := \min\big(\{h \in [\bar H] : \pi_h(s_w) \ne a_w\} \cup \{\bar H\}\big) \in [\bar H] .
\]
Its \emph{tree walk} is $y_0 := x_0$ and, as long as $y_k$ is internal,
$y_{k+1} := \mathrm{child}\big(y_k, \pi_{h_e(\pi) + 1 + k}(y_k)\big)$; since
$\mathrm{dep}(y_k) = k$ and all depths are at most $d - 1$ (Lemma~\ref{lem:tree_leaves}~(iii)),
the walk reaches a leaf after $K(\pi) := \min\{k : y_k \in \mathcal L\} \le d - 1$ steps. The
\emph{leaf} of $\pi$ is $\ell(\pi) := y_{K(\pi)}$, its \emph{arrival step} is
\[
  h'(\pi) := h_e(\pi) + 1 + K(\pi) = h_e(\pi) + 1 + \mathrm{dep}(\ell(\pi)) \le \bar H + d \le H - 2,
\]
its \emph{leaf action} is $a(\pi) := \pi_{h'(\pi)}(\ell(\pi))$, and its \emph{triple} is
\[
  u(\pi) := \big(h_e(\pi), \ell(\pi), a(\pi)\big) \in \mathcal U .
\]
We say that $\pi$ \emph{realizes} $u \in \mathcal U$ if $u(\pi) = u$; every policy realizes
exactly one triple. The \emph{deterministic prefix} of $\pi$ is the sequence of states
$x^\pi_h := s_w$ for $1 \le h \le h_e(\pi)$ and $x^\pi_{h_e(\pi) + 1 + k} := y_k$ for
$0 \le k \le K(\pi)$ (so $x^\pi_{h'(\pi)} = \ell(\pi)$), and
$\Phi_\pi : \{0, 1\} \to \Ss^{H+1}$ is the injective map
\[
  \Phi_\pi(b)_h := x^\pi_h \ (h \le h'(\pi)), \qquad \Phi_\pi(b)_h := s_g \text{ if } b = 1,\ s_b \text{ if } b = 0 \ (h'(\pi) < h \le H + 1).
\]
In a run $(X, Y, \mathrm{out})$ of an algorithm in the episode environment of an instance
(Definition~\ref{def:episode_env}), with $\pi^{t+1} = X_t$, the \emph{number of episodes with
triple $u$} among the first $t$ is $N_u(t) := \sum_{i < t}\indic\{u(\pi^{i+1}) = u\}$, and
$N_u(\tau)$ is its value at the stopping time on $\{\tau < \infty\}$; $E_u := \{u(\mathrm{out}) = u\}$
is the event that the output realizes $u$.
\end{definition}

\textbf{Lean remark.} \texttt{exitStep π : Fin (H / 3)} by \texttt{Nat.find} on a decidable
predicate over a finite range (or a fold over \texttt{List.range}); the tree walk by recursion on
$k \le d$ (\texttt{treeWalk π k : HardState S}), \texttt{arrivalDepth π := Nat.find} of the first
leaf (bounded by $d - 1$ by Lemma~\ref{lem:tree_leaves}), \texttt{hardTriple π : HardTriple S A H},
\texttt{Realizes π u := hardTriple π = u}; \texttt{hardTraj π : Bool → Traj (HardState S) H}
is $\Phi_\pi$ (\texttt{Traj S H := Fin (H + 1) → S}). $N_u(t)$ is a \texttt{Finset.card} of a
filtered range, a measurable function of the history.

\begin{lemma}[A realizing policy exists]\label{lem:lb_realizing_exists}
\uses{def:hard_triple, lem:tree_leaves}
For every $u = (h_e, \ell^*, a^*) \in \mathcal U$ there is a policy $\pi$ with $u(\pi) = u$.
\end{lemma}

\begin{proof}
\uses{def:hard_triple, lem:tree_leaves}
Let $a_1, \dots, a_{k}$ ($k = \mathrm{dep}(\ell^*)$) be the action sequence of
Lemma~\ref{lem:tree_leaves}~(iv) leading from $x_0$ to $\ell^*$ through the nodes
$x_0 = z_0, z_1, \dots, z_k = \ell^*$. Define $\pi_h(s_w) := a_w$ for $h < h_e$, $\pi_{h_e}(s_w) := $
any action different from $a_w$ (there is one since $A \ge 2$), $\pi_{h_e + 1 + i}(z_i) := a_{i+1}$
for $0 \le i < k$, $\pi_{h_e + 1 + k}(\ell^*) := a^*$, and $\pi$ arbitrary elsewhere. Then
$h_e(\pi) = h_e$ (if $h_e < \bar H$ because $\pi_h(s_w) = a_w$ for $h < h_e$ and
$\pi_{h_e}(s_w) \ne a_w$; if $h_e = \bar H$ by the cap), the tree walk is
$y_i = z_i$, which is internal for $i < k$ and a leaf for $i = k$, so $K(\pi) = k$,
$\ell(\pi) = \ell^*$, $h'(\pi) = h_e + 1 + k$ and $a(\pi) = \pi_{h_e+1+k}(\ell^*) = a^*$.
\end{proof}

\begin{lemma}[Trajectories of the hard instances]\label{lem:hard_trajectory}
\uses{def:hard_mdp, def:hard_triple, def:step_law, def:episode_return, lem:hard_params_valid}
Let $\mathcal M \in \{\mathcal M_0\} \cup \{\mathcal M_u : u \in \mathcal U\}$ (under
Condition~A) and let $\pi$ be a policy, $h' := h'(\pi)$ and
$p := p^{\mathcal M}_{h'}(\ell(\pi), a(\pi))$. Under $\Prob^{\pi}_{(s_w, 1)}$ of $\mathcal M$
(Definition~\ref{def:step_law}), almost surely,
\begin{enumerate}
  \item[(i)] $S_h = x^\pi_h$ for $1 \le h \le h'$ (in particular $S_{h'} = \ell(\pi)$ and $A_{h'} = a(\pi)$);
  \item[(ii)] $S_{h'+1} \in \{s_g, s_b\}$, with $\Prob(S_{h'+1} = s_g) = p$, and $S_h = S_{h'+1}$ for $h' + 1 \le h \le H + 1$.
\end{enumerate}
Consequently the state sequence $(S_1, \dots, S_{H+1})$ (Definition~\ref{def:episode_return}) has
the law $\Phi_\pi \# \mathrm{Ber}(p)$, the image of the Bernoulli law of parameter $p$ by
$\Phi_\pi$; $S_{\tilde H} \in \{s_g, s_b\}$ a.s.; and the return satisfies
$R^\pi_1 = H'\,\indic\{S_{\tilde H} = s_g\}$ a.s.
\end{lemma}

\begin{proof}
\uses{lem:step_law_markov, def:hard_mdp, def:hard_triple, def:episode_return}
(i) By induction on $h$ with the Markov property (Lemma~\ref{lem:step_law_markov}), which says
that given $S_h = s$ the action is $\pi_h(s)$ and $S_{h+1} \sim p_h(\cdot \mid s, \pi_h(s))$.
$S_1 = s_w = x^\pi_1$. For $h < h_e(\pi)$: $S_h = s_w$, $\pi_h(s_w) = a_w$ (by minimality of
$h_e(\pi)$) and $h < h_e(\pi) \le \bar H$, so $S_{h+1} = s_w = x^\pi_{h+1}$. At $h = h_e(\pi)$:
$S_h = s_w$ and either $\pi_h(s_w) \ne a_w$ or $h = \bar H$; in both cases
$S_{h+1} = x_0 = y_0 = x^\pi_{h+1}$. For $h = h_e(\pi) + 1 + k$ with $k < K(\pi)$: $S_h = y_k$ is
internal and $S_{h+1} = \mathrm{child}(y_k, \pi_h(y_k)) = y_{k+1} = x^\pi_{h+1}$ (all transitions
so far are Dirac masses). (ii) At $h = h'$: $S_{h'} = \ell(\pi)$ is a leaf, the action is
$a(\pi)$ and $S_{h'+1} \sim p\,\delta_{s_g} + (1-p)\,\delta_{s_b}$; afterwards $s_g$ and $s_b$
are absorbing, so $S_h = S_{h'+1}$ for $h > h'$ (induction with the Dirac transitions).

Since $\Phi_\pi(b)$ is exactly the sequence described by (i)--(ii) with $b = \indic\{S_{h'+1} = s_g\}$,
the state sequence is $\Phi_\pi(B)$ a.s.\ with $B \sim \mathrm{Ber}(p)$, whence the law. As
$h' + 1 \le \bar H + d + 1 = \tilde H$, $S_{\tilde H} = S_{h'+1} \in \{s_g, s_b\}$. Finally
$R^\pi_1 = \sum_{h=1}^H r_h(S_h, A_h) = \sum_{h=\tilde H}^{H}\indic\{S_h = s_g\}
= (H - \tilde H + 1)\,\indic\{S_{\tilde H} = s_g\} = H'\,\indic\{S_{\tilde H} = s_g\}$, since
$S_h = S_{\tilde H}$ for $h \ge \tilde H$ and $H - \tilde H + 1 = H'$.
\end{proof}

\textbf{Lean remark.} The useful form is the law statement: for the episode kernel
\texttt{statesKernel (hardMDP \dots\ u) w : Kernel (Policy \dots) (Traj \dots)}
(Definition~\ref{def:episode_env}),
\texttt{statesKernel (hardMDP \dots) w π = (bernoulliMeasure p).map (hardTraj π)}. It follows
from Lemma~\ref{lem:step_law_markov} by induction on the steps of the deterministic prefix
(each step being a \texttt{Kernel.deterministic}); this is the one place where the concrete
kernels of Definition~\ref{def:hard_mdp} are unfolded.

\begin{lemma}[Success probabilities]\label{lem:hard_success}
\uses{lem:hard_trajectory, def:hard_triple, def:hard_mdp}
For an instance $\mathcal M$ and a policy $\pi$ let
$p^{\mathcal M}(\pi) := \Prob^{\pi}_{(s_w,1)}(S_{\tilde H} = s_g)$ under $\mathcal M$. Then
$p^{\mathcal M}(\pi) = p^{\mathcal M}_{h'(\pi)}(\ell(\pi), a(\pi))$; hence
$p^{\mathcal M_0}(\pi) = p_-$ for every $\pi$, and for $u \in \mathcal U$,
\[
  p^{\mathcal M_u}(\pi) = p_+ \text{ if } \pi \text{ realizes } u, \qquad p^{\mathcal M_u}(\pi) = p_- \text{ otherwise.}
\]
\end{lemma}

\begin{proof}
\uses{lem:hard_trajectory, def:hard_triple, def:hard_mdp}
The first claim is Lemma~\ref{lem:hard_trajectory}~(ii) ($S_{\tilde H} = S_{h'+1}$). In
$\mathcal M_0$ all leaf probabilities are $p_-$. In $\mathcal M_u$ with $u = (h_e, \ell^*, a^*)$,
$p^{u}_{h'(\pi)}(\ell(\pi), a(\pi)) = p_+$ iff $(h'(\pi), \ell(\pi), a(\pi)) = (h_u, \ell^*, a^*)$ iff
$\ell(\pi) = \ell^*$, $a(\pi) = a^*$ and $h_e(\pi) + 1 + \mathrm{dep}(\ell^*) = h_e + 1 + \mathrm{dep}(\ell^*)$,
i.e.\ iff $u(\pi) = u$ (Definition~\ref{def:hard_triple}: $h'(\pi) = h_e(\pi) + 1 + \mathrm{dep}(\ell(\pi))$).
\end{proof}

\begin{lemma}[Exponential values]\label{lem:hard_exp_value}
\uses{def:exp_value, lem:hard_trajectory, lem:hard_success, def:hard_params}
For every instance $\mathcal M$ and policy $\pi$, with $p := p^{\mathcal M}(\pi)$,
\[
  Z^\pi_1(s_w) = 1 + p\,(e^{\beta H'} - 1) = \begin{cases} 1 + c\,p & (\beta > 0), \\ e^{\beta H'}\big(1 + c\,(1 - p)\big) & (\beta < 0). \end{cases}
\]
\end{lemma}

\begin{proof}
\uses{def:exp_value, lem:hard_trajectory, lem:hard_success}
By Lemma~\ref{lem:hard_trajectory}, $R^\pi_1 = H' B$ a.s.\ with $B \sim \mathrm{Ber}(p)$, so
$Z^\pi_1(s_w) = \E[e^{\beta R^\pi_1}] = (1 - p) + p\,e^{\beta H'}$. For $\beta > 0$ this is
$1 + c p$ with $c = e^{\beta H'} - 1$. For $\beta < 0$, $e^{\beta H'} = 1/(c+1)$ and
$(1-p) + p/(c+1) = \frac{(1-p)(c+1) + p}{c+1} = e^{\beta H'}\big(1 + c(1-p)\big)$.
\end{proof}

\begin{lemma}[Optimal exponential values]\label{lem:hard_opt_value}
\uses{def:opt_value, lem:opt_exp_value_eq, lem:hard_exp_value, lem:hard_success, lem:lb_realizing_exists, def:hard_params, lem:hard_params_valid}
Under Condition~A, for every $u \in \mathcal U$, in $\mathcal M_u$:
\[
  Z^*_1(s_w) = 1 + c\,p_+ = (2q-1)\,m \quad (\beta > 0), \qquad
  Z^*_1(s_w) = e^{\beta H'}\big(1 + c(1 - p_+)\big) = e^{\beta H'}\frac{m}{2q-1} \quad (\beta < 0),
\]
attained exactly by the policies realizing $u$; and in $\mathcal M_0$, $Z^*_1(s_w) = 1 + c\,p_- = m$
($\beta > 0$), $= e^{\beta H'}(1 + c(1-p_-)) = e^{\beta H'} m$ ($\beta < 0$).
\end{lemma}

\begin{proof}
\uses{lem:opt_exp_value_eq, lem:hard_exp_value, lem:hard_success, lem:lb_realizing_exists, lem:hard_params_valid}
By Lemma~\ref{lem:opt_exp_value_eq}, $Z^*_1(s_w) = \max_\pi Z^\pi_1(s_w)$ if $\beta > 0$ and
$= \min_\pi Z^\pi_1(s_w)$ if $\beta < 0$. By Lemma~\ref{lem:hard_exp_value}, $Z^\pi_1(s_w)$ is
$1 + c\,p^{\mathcal M}(\pi)$, increasing in $p^{\mathcal M}(\pi)$ ($c > 0$), resp.\
$e^{\beta H'}(1 + c(1 - p^{\mathcal M}(\pi)))$, decreasing in $p^{\mathcal M}(\pi)$. In
$\mathcal M_u$, $p^{\mathcal M_u}(\pi) \in \{p_-, p_+\}$ with $p_+ > p_-$
(Lemma~\ref{lem:hard_params_valid}) and $p_+$ is attained exactly by the realizing policies
(Lemma~\ref{lem:hard_success}), which exist (Lemma~\ref{lem:lb_realizing_exists}); so the
extremum is the value at $p_+$. In $\mathcal M_0$ all policies have $p_-$. The closed forms
are the identities of Definition~\ref{def:hard_params}.
\end{proof}

\begin{lemma}[Missing the triple is more than $\eps$-suboptimal]\label{lem:hard_eps_suboptimal}
\uses{def:entropic_value, def:opt_value, lem:hard_opt_value, lem:hard_exp_value, lem:hard_success, lem:hard_params_valid, def:condition_a}
Under Condition~A, for every $u \in \mathcal U$ and every policy $\pi$ that does not realize
$u$, in $\mathcal M_u$,
\[
  V^*_1(s_w) - V^\pi_1(s_w) = \frac{1}{|\beta|}\log\big(2e^{|\beta|\eps} - 1\big) > \eps .
\]
Equivalently: if $V^*_1(s_w) - V^\pi_1(s_w) \le \eps$ in $\mathcal M_u$ then $\pi$ realizes $u$
(the paper's equation (hardness)).
\end{lemma}

\begin{proof}
\uses{def:entropic_value, lem:hard_opt_value, lem:hard_exp_value, lem:hard_success, lem:hard_params_valid}
By Lemma~\ref{lem:hard_success}, $p^{\mathcal M_u}(\pi) = p_-$. If $\beta > 0$,
Lemmas~\ref{lem:hard_opt_value} and~\ref{lem:hard_exp_value} give
$V^*_1(s_w) - V^\pi_1(s_w) = \beta^{-1}\log\frac{1 + c\,p_+}{1 + c\,p_-} = \beta^{-1}\log\frac{(2q-1)m}{m} = \beta^{-1}\log(2q-1)$.
If $\beta < 0$, $V^*_1(s_w) - V^\pi_1(s_w) = \beta^{-1}\log\frac{Z^*_1(s_w)}{Z^\pi_1(s_w)}
= -|\beta|^{-1}\log\frac{e^{\beta H'} m/(2q-1)}{e^{\beta H'} m} = |\beta|^{-1}\log(2q-1)$.
(All quantities are positive: $m > 0$, $2q - 1 > 1$.) Finally $2q - 1 > q$ since $q > 1$, so
$|\beta|^{-1}\log(2q-1) > |\beta|^{-1}\log q = \eps$.
\end{proof}

\begin{lemma}[Maximal return of the hard instances]\label{lem:hard_gmax}
\uses{def:gmax, lem:hard_trajectory, lem:hard_success, lem:hard_params_valid}
Under Condition~A, $\Gmax(\mathcal M_0) = H'$ (and $\Gmax(\mathcal M_u) = H'$ for every $u$);
hence $c = e^{|\beta|\Gmax(\mathcal M_0)} - 1$ and
$c^2/(c+1) = (e^{|\beta|\Gmax(\mathcal M_0)} - 1)^2/e^{|\beta|\Gmax(\mathcal M_0)}$.
\end{lemma}

\begin{proof}
\uses{def:gmax, lem:hard_trajectory, lem:hard_success, lem:hard_params_valid}
For every policy $\pi$, $R^\pi_1 = H' B$ a.s.\ with $B \sim \mathrm{Ber}(p)$ and
$p = p^{\mathcal M}(\pi) \in \{p_-, p_+\} \subset (0, 1)$ (Lemmas~\ref{lem:hard_trajectory},
\ref{lem:hard_success}, \ref{lem:hard_params_valid}). Both values $0$ and $H'$ of $R^\pi_1$ have
positive probability, so the essential supremum of $R^\pi_1$ under $\Prob^\pi_{(s_w,1)}$ is
$H'$; the supremum over $\pi$ (Definition~\ref{def:gmax}) is $H'$.
\end{proof}

\section{Change of measure}
\label{sec:lb_change_of_measure}

For a BPI algorithm $\mathcal A$ (Definition~\ref{def:bpi_alg}) and an instance $\mathcal M$, the
law of the pair (stopped history, output) $(\mathcal H_\tau, \mathrm{out})$ of $\mathcal A$ in the
episode environment of $\mathcal M$ from $s_w$ is the same for all runs (Definition~\ref{def:bpi_alg});
we write $\Prob_0$, $\Prob_u$ for these laws under $\mathcal M_0$, $\mathcal M_u$, and $\E_0$ for
the expectation under $\Prob_0$. The number of episodes $\tau$ and the counts $N_u(\tau)$ are
functions of $\mathcal H_\tau$, so $\E_0[\tau]$ and $\E_0[N_u(\tau)]$ (in $[0, \infty]$) are defined.

\begin{lemma}[PAC algorithms find the triple]\label{lem:hard_pac_event}
\uses{def:entropic_pac, def:bpi_alg, def:hard_mdp, def:hard_triple, lem:hard_eps_suboptimal, def:reward_fn}
Assume Condition~A and let $\mathcal A$ be $(\eps,\delta)$-PAC for entropic BPI with reward
function $r_0$, parameter $\beta$ and initial state $s_w$ (Definition~\ref{def:entropic_pac}).
Then for every $u \in \mathcal U$,
\[
  \Prob_u(E_u) \ge 1 - \delta, \qquad E_u = \{u(\mathrm{out}) = u\} .
\]
Moreover, in every instance, the events $(E_u)_{u \in \mathcal U}$ are pairwise disjoint; in
particular $\sum_{u \in \mathcal U}\Prob_0(E_u) \le 1$.
\end{lemma}

\begin{proof}
\uses{def:entropic_pac, lem:hard_eps_suboptimal, def:hard_triple, def:hard_mdp}
$\mathcal M_u \in \mathcal M_{r_0}$ (Definition~\ref{def:hard_mdp}), so the PAC property gives
$\Prob_u(V^*_1(s_w) - V^{\mathrm{out}}_1(s_w) > \eps) \le \delta$ (the values being those of
$\mathcal M_u$). By Lemma~\ref{lem:hard_eps_suboptimal},
$\{u(\mathrm{out}) \ne u\} \subseteq \{V^*_1(s_w) - V^{\mathrm{out}}_1(s_w) > \eps\}$, hence
$\Prob_u(E_u^c) \le \delta$. Disjointness: $u(\mathrm{out})$ is a single element of $\mathcal U$
(Definition~\ref{def:hard_triple}), so at most one $E_u$ occurs; a sum of probabilities of
pairwise disjoint events is at most $1$.
\end{proof}

\textbf{Lean remark.} The first claim is \texttt{IdentAlg.IsPAC.measureReal\_bad\_of\_isRun}
applied to the member $\mathcal M_u$ of the class \texttt{\{M // M.HasRewardFn r₀\}}, in any run
(or to \texttt{outputMeasure} directly); the second is \texttt{measure\_biUnion\_finset} for the
pairwise disjoint preimages of the singletons $\{u\}$ by \texttt{hardTriple ∘ out}.

\begin{lemma}[Divergence of one episode]\label{lem:hard_kl_episode}
\uses{def:episode_env, lem:hard_trajectory, lem:hard_success, def:klbin, lem:klbin_bernoulli, def:hard_triple}
Assume Condition~A. For an instance $\mathcal M$ and a policy $\pi$ let
$\nu_{\mathcal M}(\pi)$ be the law of the state sequence of an episode of $\pi$ in $\mathcal M$
from $s_w$, i.e.\ the feedback kernel of the episode environment
$\mathrm{statesEnv}(\mathcal M, s_w)$ evaluated at the action $\pi$
(Definition~\ref{def:episode_env}). Then for every $u \in \mathcal U$ and every $\pi$,
\[
  \KL\big(\nu_{\mathcal M_0}(\pi) \,\big\|\, \nu_{\mathcal M_u}(\pi)\big) \le \klbin(p_-, p_+) \text{ if } \pi \text{ realizes } u,
  \qquad = 0 \text{ otherwise.}
\]
\end{lemma}

\begin{proof}
\uses{lem:hard_trajectory, lem:hard_success, lem:klbin_bernoulli, def:hard_triple}
By Lemma~\ref{lem:hard_trajectory}, $\nu_{\mathcal M}(\pi) = \Phi_\pi \# \mathrm{Ber}(p^{\mathcal M}(\pi))$
with the \emph{same} map $\Phi_\pi$ for all instances (it depends only on $\pi$ and the tree).
If $\pi$ does not realize $u$, $p^{\mathcal M_0}(\pi) = p^{\mathcal M_u}(\pi) = p_-$
(Lemma~\ref{lem:hard_success}), the two laws coincide and the divergence is $0$. If $\pi$
realizes $u$, the parameters are $p_-$ and $p_+$, and the data-processing inequality for the
measurable map $\Phi_\pi$ gives
$\KL(\Phi_\pi\#\mathrm{Ber}(p_-) \| \Phi_\pi\#\mathrm{Ber}(p_+)) \le \KL(\mathrm{Ber}(p_-)\|\mathrm{Ber}(p_+)) = \klbin(p_-, p_+)$
(Lemma~\ref{lem:klbin_bernoulli}; $p_+ \in (0,1)$ by Lemma~\ref{lem:hard_params_valid}).
\end{proof}

\textbf{Lean remark.} Data processing is Mathlib's \texttt{klDiv\_map\_le}
(\texttt{ProbabilityTheory.klDiv} of the images by a measurable map). Equality holds since
$\Phi_\pi$ is injective (apply data processing to a left inverse), but only the inequality is
used.

\begin{lemma}[Change of measure on the hard instances]\label{lem:hard_change_of_measure}
\uses{thm:change_of_measure, lem:hard_kl_episode, def:episode_env, def:bpi_alg, def:hard_triple, lem:kl_output_pair, def:klbin}
Assume Condition~A. Let $\mathcal A$ be a BPI algorithm, $u \in \mathcal U$, and assume
$\E_0[\tau] < \infty$. Then
\[
  \klbin\big(\Prob_0(E_u), \Prob_u(E_u)\big) \le \E_0\big[N_u(\tau)\big]\,\klbin(p_-, p_+) .
\]
\end{lemma}

\begin{proof}
\uses{thm:change_of_measure, lem:hard_kl_episode, def:episode_env, lem:kl_output_pair, def:hard_triple}
In LML terms (deviation~8), the episode environments $\mathrm{statesEnv}(\mathcal M_0, s_w)$ and
$\mathrm{statesEnv}(\mathcal M_u, s_w)$ (Definition~\ref{def:episode_env}) are the stationary
environments $\mathrm{stationaryEnv}(\nu_{\mathcal M_0})$ and $\mathrm{stationaryEnv}(\nu_{\mathcal M_u})$
with the finite action set $\mathcal P := \{\text{policies}\}$ and feedback kernels
$\pi \mapsto \nu_{\mathcal M_0}(\pi)$, $\pi \mapsto \nu_{\mathcal M_u}(\pi)$ (the laws of the state
sequences, Lemma~\ref{lem:hard_kl_episode}): a ``bandit'' whose arms are the policies and whose
observations are the state sequences. The sampling rule of $\mathcal A$ is an LML algorithm on
this action set, its stopping time $\tau$ is a stopping time of the history filtration, a.s.\
finite under $\Prob_0$ since $\E_0[\tau] < \infty$, and its output kernel is the same under both
environments (Lemma~\ref{lem:kl_output_pair}). For $\pi \in \mathcal P$ let
$N_\pi(\tau) := \sum_{t < \tau}\indic\{X_t = \pi\}$ be the number of pulls of the arm $\pi$
before $\tau$. Theorem~\ref{thm:change_of_measure} applied to these two environments, to the
stopping time $\tau$ and to the event $E_u = \{u(\mathrm{out}) = u\}$ of $(\mathcal H_\tau, \mathrm{out})$
gives
\[
  \klbin\big(\Prob_0(E_u), \Prob_u(E_u)\big) \le \sum_{\pi \in \mathcal P}\E_0[N_\pi(\tau)]\,\KL\big(\nu_{\mathcal M_0}(\pi)\,\big\|\,\nu_{\mathcal M_u}(\pi)\big) .
\]
By Lemma~\ref{lem:hard_kl_episode} the divergence is at most $\klbin(p_-,p_+)$ when $\pi$
realizes $u$ and $0$ otherwise, so the right-hand side is at most
$\klbin(p_-,p_+)\sum_{\pi : u(\pi) = u}\E_0[N_\pi(\tau)] = \klbin(p_-,p_+)\,\E_0[N_u(\tau)]$,
since $N_u(\tau) = \sum_{\pi : u(\pi) = u} N_\pi(\tau)$ (Definition~\ref{def:hard_triple}) and
all terms are finite ($N_\pi(\tau) \le \tau$).
\end{proof}

\textbf{Lean remark.} \texttt{statesEnv M s₁ := stationaryEnv (statesKernel M s₁)} by
definition, so Theorem~\ref{thm:change_of_measure}
(\texttt{Essakine2026Tight/LeanMachineLearning/SequentialLearning/ChangeOfMeasure.lean})
applies verbatim with the action type \texttt{Policy S A H} (a \texttt{Fintype}) and the
feedback type \texttt{Traj S H}. The pull counts $N_\pi(\tau)$ are LML's counts of the rounds
with action $\pi$ before the stopping time; the sum over $\mathcal P$ is a \texttt{Finset.sum}
over \texttt{Finset.univ}, split by \texttt{Finset.sum\_filter} according to
\texttt{hardTriple π = u}. Under $\Prob_u$ the stopping time need not be finite; whatever
convention LML uses for the output on $\{\tau = \infty\}$ is inherited by $\Prob_u(E_u)$, and
Theorem~\ref{thm:change_of_measure} is stated for it.

\begin{lemma}[Divergence of the leaf parameters]\label{lem:hard_kl_params}
\uses{def:hard_params, def:condition_a, lem:hard_params_valid, lem:klbin_upper, def:klbin}
Under Condition~A, for both signs of $\beta$,
\[
  0 < \klbin(p_-, p_+) \le 72\,\frac{c+1}{c^2}\,\big(e^{|\beta|\eps} - 1\big)^2 < \infty .
\]
\end{lemma}

\begin{proof}
\uses{lem:hard_params_valid, lem:klbin_upper, lem:klbin_bernoulli}
By Lemma~\ref{lem:hard_params_valid}, $p_+ \in (0,1)$, $p_- \in (0,1)$, $p_- \ne p_+$ and
$p_+(1-p_+) \ge 1/(8(c+1))$. Positivity and finiteness follow ($\klbin(x,y) > 0$ for
$x \ne y$, $y \in (0,1)$: Gibbs' inequality for the Bernoulli laws through
Lemma~\ref{lem:klbin_bernoulli}, Mathlib \texttt{klDiv\_eq\_zero\_iff}; and it is finite since $y \in (0,1)$). By Lemma~\ref{lem:klbin_upper}
(the bound $\klbin(x,y) \le (x-y)^2/(y(1-y))$, with $y = p_+$ the \emph{second} argument),
\[
  \klbin(p_-, p_+) \le \frac{\Delta^2}{p_+(1-p_+)} \le 8(c+1)\,\Delta^2 .
\]
Next, $(3c+2)/(c+1) \le 3$, so for $\beta > 0$, $\Delta = (q-1)(3c+2)/(c(c+1)) \le 3(q-1)/c$, and for
$\beta < 0$, $\Delta = (q-1)(3c+2)/(c(c+1)(2q-1)) \le 3(q-1)/c$ as well since $2q - 1 \ge 1$. Hence
$\klbin(p_-,p_+) \le 8(c+1)\cdot 9(q-1)^2/c^2 = 72\,(c+1)(q-1)^2/c^2$.
\end{proof}

\begin{remark}[The paper's bound on the divergence]\label{rem:lb_kl_paper}
The paper bounds $\klbin(p_-,p_+)$ by $(p_+ - p_-)^2/(p_-(1-p_-))$. The elementary bound
$\klbin(x,y) \le (x-y)^2/(y(1-y))$ (Lemma~\ref{lem:klbin_upper}, from $\log t \le t - 1$) has the
\emph{second} argument in the denominator, so the paper's expression bounds $\klbin(p_+, p_-)$,
not $\klbin(p_-, p_+)$, and no bound of the form $(x - y)^2/(x(1-x))$ can hold in general since
$\klbin(x, y) \to \infty$ as $y \to 1$ for fixed $x < 1$. This matters: the paper's Condition~A
only ensures $p_+ < 1$ without margin, and its final bound
$\klbin(p_-,p_+) \le \frac{1521}{100}\frac{c+1}{c^2}(q-1)^2$ ($\beta > 0$) is false near the
boundary. For instance with $c = 1$ and $q = 1.23$ (admissible: $q < 16/13$ and
$4(c+1)^2/(2c^2+7c+4) = 16/13$), the paper's parameters are $p_- = 1/4$,
$\Delta = 5\cdot 0.23/(2\cdot 0.77) \approx 0.7468$, $p_+ \approx 0.9968$, and
$\klbin(0.25, 0.99675) \approx 0.25\log(0.2508) + 0.75\log(230.8) \approx 3.7$, while the claimed
bound is $15.21\cdot 2\cdot 0.23^2 \approx 1.61$. Even granting the denominator, the paper drops
a factor $4$ ($p_-(1-p_-) = (2c+1)/(4(c+1)^2) \ge 1/(4(c+1))$, not $\ge 1/(c+1)$). For $\beta < 0$
the paper bounds $1/(q - 1/2)^2$ by $9e^{2|\beta|\eps}$ where $4$ suffices ($q \ge 1$); this
is the only source of the factor $e^{2\min\{\beta,0\}\eps}$ in its Theorem~3, which is therefore
unnecessary (it is $\le 1$ and is kept in Definition~\ref{def:lower_bound_const} only to match
the paper's statement). Our Condition~A includes the margin $p_+ \le (1+p_-)/2$, which makes
the bound $\klbin(p_-,p_+) \le \Delta^2/(p_+(1-p_+)) \le 8(c+1)\Delta^2$ valid.
\end{remark}

\begin{lemma}[The counts sum to the stopping time]\label{lem:hard_count_sum}
\uses{def:hard_triple, def:bpi_alg}
In every run of a BPI algorithm in the episode environment of an instance, on $\{\tau < \infty\}$,
$\sum_{u \in \mathcal U} N_u(\tau) = \tau$. Consequently, if $\tau < \infty$ $\Prob_0$-a.s.,
$\E_0[\tau] = \sum_{u \in \mathcal U}\E_0[N_u(\tau)]$ (in $[0,\infty]$).
\end{lemma}

\begin{proof}
\uses{def:hard_triple}
Each of the $\tau$ episodes $t < \tau$ has exactly one triple $u(\pi^{t+1})$
(Definition~\ref{def:hard_triple}), so $\sum_u N_u(\tau) = \sum_{t<\tau}\sum_u\indic\{u(\pi^{t+1}) = u\} = \tau$.
Take expectations (a finite sum of nonnegative terms).
\end{proof}

\begin{lemma}[Lower bound on the sum of divergences]\label{lem:hard_count_lower}
\uses{def:klbin, lem:klbin_lower}
Let $\mathcal U$ be a finite set with $\abs{\mathcal U} \ge 4$, $\delta \in (0, 1/16]$,
$(x_u)_{u \in \mathcal U}$ in $[0,1]$ with $\sum_u x_u \le 1$ and $(y_u)_{u \in \mathcal U}$ in
$[1 - \delta, 1]$. Then, in $[0, \infty]$,
\[
  \sum_{u \in \mathcal U}\klbin(x_u, y_u) \;\ge\; \sum_{u \in \mathcal U}\Big[(1 - x_u)\log\frac1\delta - \log 2\Big] \;\ge\; \frac{\abs{\mathcal U}}{2}\log\frac1\delta .
\]
\end{lemma}

\begin{proof}
\uses{lem:klbin_lower}
\emph{Termwise.} If $y_u < 1$, Lemma~\ref{lem:klbin_lower} gives
$\klbin(x_u,y_u) \ge (1-x_u)\log\frac{1}{1-y_u} - \log 2 \ge (1-x_u)\log\frac1\delta - \log 2$
since $1 - y_u \le \delta$ and $1 - x_u \ge 0$. If $y_u = 1$ and $x_u < 1$, $\klbin(x_u, 1) = \infty$
(Definition~\ref{def:klbin}). If $x_u = y_u = 1$, $\klbin(1,1) = 0 \ge -\log 2$.

\emph{Summing.} $\sum_u[(1-x_u)\log\frac1\delta - \log 2] = (\abs{\mathcal U} - \sum_u x_u)\log\frac1\delta - \abs{\mathcal U}\log 2
\ge (\abs{\mathcal U} - 1)\log\frac1\delta - \abs{\mathcal U}\log 2$. This is at least
$\frac{\abs{\mathcal U}}{2}\log\frac1\delta$ iff $(\frac{\abs{\mathcal U}}{2} - 1)\log\frac1\delta \ge \abs{\mathcal U}\log 2$,
which holds because $\frac{\abs{\mathcal U}}{2} - 1 \ge \frac{\abs{\mathcal U}}{4}$ (as
$\abs{\mathcal U} \ge 4$) and $\frac14\log\frac1\delta \ge \frac14\log 16 = \log 2$ (as $\delta \le 1/16$).
\end{proof}

\textbf{Lean remark.} $\klbin$ is the $\texttt{ℝ≥0∞}$-valued \texttt{klBin} of
Definition~\ref{def:klbin} (as \texttt{klBer} in the sister project, infinite when the second
argument is $1$ and the first is not); the sum is a \texttt{Finset.sum} in \texttt{ℝ≥0∞} and the
right-hand side is \texttt{ENNReal.ofReal} of a real number. The outline's version
($\abs{\mathcal U} \ge 2$, constant $1/4$) is weaker; the paper's constant $1/2$ is correct
since $\abs{\mathcal U} \ge 4$ in the application.

\section{The lower bound}
\label{sec:lb_theorem}

\begin{definition}[The lower bound]\label{def:lower_bound_const}
\lean{Essakine2026Tight.lowerBound}\leanok
\uses{def:gmax}
For integers $S, A, H \ge 1$, $\beta \ne 0$, $\eps > 0$, $\delta \in (0,1)$ and $G \ge 0$,
\[
  \mathrm{lowerBound}(S, A, H, \beta, \eps, \delta, G) := \frac{1}{5000}\,\frac{(e^{|\beta| G} - 1)^2}{e^{|\beta| G}}\,
  \frac{e^{2\min\{\beta, 0\}\eps}\, S A H}{(e^{|\beta|\eps} - 1)^2}\,\log\frac1\delta .
\]
\end{definition}

\textbf{Lean remark.} \texttt{lowerBound (S A H : ℕ) (β ε δ G : ℝ) : ℝ} in
\texttt{Essakine2026Tight/EV2026/Constants.lean}.

\begin{theorem}[Theorem~3, lower bound]\label{thm:lower_bound}
\lean{Essakine2026Tight.exists_hardMDP_lowerBound_le_lintegral_stoppingTime}\leanok
\uses{def:tree_depth, def:condition_a, def:hard_params, def:hard_mdp, def:entropic_pac, def:bpi_alg, def:episode_env, def:gmax, def:lower_bound_const, def:reward_fn}
Let $S \ge 6$, $A \ge 2$, $H \ge 3d$ with $d$ as in Definition~\ref{def:tree_depth}, $\beta \ne 0$,
$\delta \in (0, 1/16]$ and $\eps > 0$ satisfying Condition~A for $(\beta, H', \eps)$ with
$H' := H - \lfloor H/3 \rfloor - d$ (Definition~\ref{def:condition_a}). There are a reward
function $r_0$ with values in $[0,1]$, a finite-horizon MDP $\mathcal M_0$ with $S$ states, $A$
actions, horizon $H$ and reward function $r_0$, and an initial state $s_1$, such that every BPI
algorithm $\mathcal A$ which is $(\eps,\delta)$-PAC for entropic BPI with reward function $r_0$,
parameter $\beta$ and initial state $s_1$ (Definition~\ref{def:entropic_pac}) satisfies, in every
run in the episode environment of $\mathcal M_0$ from $s_1$,
\[
  \E_0[\tau] \;\ge\; \mathrm{lowerBound}\big(S, A, H, \beta, \eps, \delta, \Gmax(\mathcal M_0)\big) .
\]
More precisely the proof gives $\E_0[\tau] \ge \frac{1}{3456}\frac{(e^{|\beta|\Gmax(\mathcal M_0)} - 1)^2}{e^{|\beta|\Gmax(\mathcal M_0)}}\frac{SAH}{(e^{|\beta|\eps}-1)^2}\log\frac1\delta$.
\end{theorem}

\begin{proof}
\uses{lem:hard_params_valid, lem:hard_gmax, lem:hard_pac_event, lem:hard_change_of_measure, lem:hard_kl_params, lem:hard_count_sum, lem:hard_count_lower, lem:tree_leaves, lem:tree_depth_props, def:hard_triple}
Take $r_0$, $\mathcal M_0$, the family $(\mathcal M_u)_{u \in \mathcal U}$ and $s_1 = s_w$ of
Definition~\ref{def:hard_mdp}, which is admissible by Lemma~\ref{lem:hard_params_valid}. Let
$\mathcal A$ be $(\eps,\delta)$-PAC for $r_0$, fix a run in the episode environment of
$\mathcal M_0$, and write $B := \mathrm{lowerBound}(S,A,H,\beta,\eps,\delta,\Gmax(\mathcal M_0)) < \infty$.

\emph{Step 0.} If $\E_0[\tau] = \infty$ there is nothing to prove. Assume $\E_0[\tau] < \infty$;
then $\tau < \infty$ $\Prob_0$-a.s.

\emph{Step 1: change of measure.} By Lemma~\ref{lem:hard_change_of_measure}, for every
$u \in \mathcal U$, $\klbin(\Prob_0(E_u), \Prob_u(E_u)) \le \E_0[N_u(\tau)]\,\klbin(p_-,p_+)$. Summing
over $u$ and using Lemma~\ref{lem:hard_count_sum},
\[
  \sum_{u \in \mathcal U}\klbin\big(\Prob_0(E_u), \Prob_u(E_u)\big) \le \klbin(p_-,p_+)\sum_{u}\E_0[N_u(\tau)] = \klbin(p_-,p_+)\,\E_0[\tau] .
\]

\emph{Step 2: the PAC property.} By Lemma~\ref{lem:hard_pac_event}, $x_u := \Prob_0(E_u)$
satisfies $\sum_u x_u \le 1$ and $y_u := \Prob_u(E_u) \ge 1 - \delta$. Moreover
$\abs{\mathcal U} = \bar H L A \ge 1\cdot 2\cdot 2 = 4$ ($\bar H \ge 1$, $L \ge 2$ by
Lemma~\ref{lem:tree_leaves}, $A \ge 2$), so Lemma~\ref{lem:hard_count_lower} gives
$\sum_u \klbin(x_u, y_u) \ge \frac{\abs{\mathcal U}}{2}\log\frac1\delta$. With Step~1 and
$0 < \klbin(p_-,p_+) < \infty$ (Lemma~\ref{lem:hard_kl_params}),
\[
  \E_0[\tau] \ge \frac{\abs{\mathcal U}\log(1/\delta)}{2\,\klbin(p_-,p_+)} \ge \frac{\abs{\mathcal U}\log(1/\delta)}{2\cdot 72}\,\frac{c^2}{(c+1)(q-1)^2}
  = \frac{\abs{\mathcal U}}{144}\,\frac{c^2}{c+1}\,\frac{\log(1/\delta)}{(e^{|\beta|\eps}-1)^2} .
\]

\emph{Step 3: the number of triples.} $\bar H = \lfloor H/3 \rfloor \ge H/6$ for $H \ge 3$ (for
$H \ge 4$, $\lfloor H/3\rfloor \ge (H-2)/3 \ge H/6$; for $H = 3$, $1 \ge 1/2$), $L \ge S/4$
(Lemma~\ref{lem:tree_leaves}~(ii)), so $\abs{\mathcal U} = \bar H L A \ge SAH/24$.

\emph{Step 4: the constant.} By Lemma~\ref{lem:hard_gmax},
$c^2/(c+1) = (e^{|\beta|\Gmax(\mathcal M_0)}-1)^2/e^{|\beta|\Gmax(\mathcal M_0)}$. Hence
\[
  \E_0[\tau] \ge \frac{1}{24\cdot 144}\,\frac{(e^{|\beta|\Gmax(\mathcal M_0)}-1)^2}{e^{|\beta|\Gmax(\mathcal M_0)}}\,\frac{SAH}{(e^{|\beta|\eps}-1)^2}\log\frac1\delta ,
\]
$24 \cdot 144 = 3456 \le 5000$, and $e^{2\min\{\beta,0\}\eps} \le 1$; this is $\E_0[\tau] \ge B$.
\end{proof}

\textbf{Lean remark.} \texttt{Essakine2026Tight.exists\_hardMDP\_lowerBound\_le\_lintegral\_stoppingTime}:
for the state type \texttt{HardState S} (cardinality $S$) and action type \texttt{Fin A}, the
statement is
$\exists\, r_0\ M_0\ s_1,\ M_0.\texttt{HasRewardFn}\ r_0 \wedge (\forall a\, h\, s, r_0\, h\, s\, a \in [0,1]) \wedge
\forall\, \mathcal A,\ \texttt{IsEntropicPAC}\ \mathcal A\ r_0\ \beta\ \eps\ \delta\ s_1 \to \forall \text{ run},\
\texttt{ENNReal.ofReal}\ (\mathrm{lowerBound}\ \dots\ (\Gmax\ M_0\ s_1)) \le \int \tau\, d\Prob$,
with $\tau : \Omega \to \texttt{ℕ∞}$ cast to \texttt{ℝ≥0∞} (so $\E_0[\tau] = \infty$ is allowed and
Step~0 is the case \texttt{lintegral = ⊤}). The hard MDP $M_0$ is \texttt{hardMDP S A H β ε none}
and the family $\mathcal M_u$ only appears in the proof. Condition~A is a hypothesis on
$(\beta, \texttt{effHorizon S A H}, \eps)$.

\begin{remark}[On the constant and the hypotheses]\label{rem:lb_constant}
The proof gives the constant $1/3456 = 1/(2\cdot 72\cdot 24)$: $2$ from
Lemma~\ref{lem:hard_count_lower}, $72$ from Lemma~\ref{lem:hard_kl_params} and $24 = 6\cdot 4$
from $\bar H \ge H/6$, $L \ge S/4$. Under Assumption~1 one even has $H \ge 6$ and
$\lfloor H/3\rfloor \ge H/4$, giving $1/2304$; we state $1/5000$ (deviation~5) to leave slack.
The paper's $1/1650$ rests on its erroneous divergence bound (Remark~\ref{rem:lb_kl_paper}),
on $L \ge S/4$ for a full tree (which our Lemma~\ref{lem:tree_leaves} recovers for the
breadth-first tree) and on $\bar H = H/3$ taken as an integer. The hypotheses $\delta \le 1/16$
(used only in Lemma~\ref{lem:hard_count_lower}) and $H \ge 3d$ (used to fit the waiting phase,
the tree and a rewarding phase of length $H' \ge H/3$ into the horizon) are the paper's; the
paper's Theorem~3 writes $\delta \in (0, 1/16)$ in the main text and $\delta \le 1/16$ in the
appendix, and the proof only needs $\log(1/\delta) \ge 4\log 2$. The lower bound is stated for
algorithms that are PAC for the single reward function $r_0$ (deviation~7), a weaker hypothesis
than PAC for every reward function, hence a stronger conclusion.
\end{remark}


\part{Prerequisites to be added to Mathlib and LML}
\label{part:prereq}

\textbf{Overview}

The proofs of Part~\ref{part:main} rely on results of probability theory and information theory
that are standard but not yet available in Mathlib or in LML. This part develops them, at the
level of generality of a library: each chapter is a self-contained area that can be formalized
(and contributed upstream) independently of the rest.
\begin{itemize}
  \item Chapter~\ref{chap:pre_mdp}: finite-horizon MDPs: occupancy measures, unrolling of
        recursive inequalities along trajectories, entropic variances and their Bellman
        recursion, sampling at the visits of a state-action pair.
  \item Chapter~\ref{chap:pre_concentration}: Sanov's high-probability bound, a maximal
        inequality for Bernoulli sequences, the self-normalized Bernstein inequality, the
        KL-Bernstein inequality and the transportation of variances.
  \item Chapter~\ref{chap:pre_info}: the binary relative entropy and the change of measure at
        a stopping time for algorithm-environment sequences.
  \item Chapter~\ref{chap:pre_analysis}: elementary inequalities (Bhatia--Davis, a normalized
        variance bound for exponential transforms, a logarithmic sum bound, a self-bounding
        inequality).
\end{itemize}

\chapter{Finite-horizon MDPs: trajectories, occupancy measures and variances}
\label{chap:pre_mdp}

This chapter develops the finite-horizon MDP theory that the analysis of Entropic-BPI uses
beyond the model of Chapter~\ref{chap:mdp}: the iterated Markov property of trajectory laws,
occupancy measures, the unrolling of a backward recursive inequality along the trajectory
(the step of Appendix~B.1 of \cite{essakine2026tight} that turns the certificate recursion
into an expectation over the episode), a Cauchy--Schwarz inequality along trajectories, the
entropic variances and their Bellman recursion (the paper's Lemma~23), the telescoping
identity that expresses the variance of the exponentiated return as a sum of one-step
variances, and the sampling structure at the visits of a state-action pair in a run of an
arbitrary algorithm (what the concentration events of Chapter~\ref{chap:empirical} consume).
Everything is stated for an arbitrary finite-horizon MDP and policy, in library generality
(namespace \texttt{Learning.MDP}, files \texttt{Episodic.lean}, \texttt{Entropic.lean} and
\texttt{Empirical.lean} of
\texttt{Essakine2026Tight/LeanMachineLearning/ReinforcementLearning/MDP/}).

\textbf{Standing notation.} $\mathcal M = (\Ss, \As, H, (p_h)_{h \le H}, (R_h)_{h \le H})$ is
a finite-horizon MDP (Definition~\ref{def:episodic_mdp}), $\pi$ a deterministic policy
(Definition~\ref{def:policy}), $s_1 \in \Ss$ an initial state and $\beta \in \R$. For
$1 \le h \le H + 1$ and $s \in \Ss$, $\Prob^\pi_{(s,h)}$ is the law of the trajectory
$(S_i, A_i, R_i)_{i \ge h}$ of $\pi$ started at $s$ at step $h$ (Definition~\ref{def:step_law}),
$\E^\pi_{(s,h)}$ its expectation, $R^\pi_h := \sum_{i = h}^H R_i$ the return from step $h$
(Definition~\ref{def:episode_return}, $R^\pi_{H+1} = 0$), and $\Prob^\pi := \Prob^\pi_{(s_1, 1)}$,
$\E^\pi := \E^\pi_{(s_1,1)}$. For $f : \Ss \to \R$ we write
\[
  (p_h f)(s, a) := \sum_{s' \in \Ss} p_h(s' \mid s, a)\, f(s'), \qquad
  \Var_{p_h}(f)(s, a) := \sum_{s'} p_h(s' \mid s, a)\big(f(s') - (p_h f)(s,a)\big)^2
  = (p_h f^2)(s,a) - \big((p_h f)(s,a)\big)^2 ,
\]
the mean and the variance of $f$ under the transition law $p_h(\cdot \mid s, a)$ (Lean:
\texttt{vecExp}, \texttt{vecVar} in \texttt{Vec.lean}). All functions on the finite sets
$\Ss$, $\As$, $\Ss \times \As$ are bounded, hence integrable against any probability measure.

\textbf{The Markov property.} We use Lemma~\ref{lem:step_law_markov} in the following form.
For $1 \le h \le H$ and $s \in \Ss$: under $\Prob^\pi_{(s,h)}$, $S_h = s$ and $A_h = \pi_h(s)$
almost surely, $(R_h, S_{h+1}) \sim R_h(\cdot \mid s, \pi_h(s)) \otimes p_h(\cdot \mid s, \pi_h(s))$,
and conditionally on the first round $(S_h, A_h, R_h, S_{h+1})$ the shifted trajectory
$\theta := (S_i, A_i, R_i)_{i \ge h+1}$ has law $\Prob^\pi_{(S_{h+1}, h+1)}$; in integrated
form, for bounded measurable $\psi : \As \times \R \times \Ss \to \R$ and $F$ on trajectories,
\begin{equation}\label{eq:pmdp_markov_integrated}
  \E^\pi_{(s,h)}\big[\psi(A_h, R_h, S_{h+1})\, F(\theta)\big]
  = \E_{(R, s') \sim (R_h \otimes p_h)(s, \pi_h(s))}\Big[\psi(\pi_h(s), R, s')\;\E^\pi_{(s',h+1)}[F]\Big] .
\end{equation}
For $h = H + 1$ the trajectory law $\Prob^\pi_{(s, H+1)}$ is the law of the absorbing
terminal layer (all rewards $0$), and $R^\pi_{H+1} = 0$.

\textbf{What Mathlib/LML provides.} LML's \texttt{MDP.policyMeasure} and its Ionescu-Tulcea
structure (\texttt{trajMeasure}, \texttt{IsAlgEnvSeq}, \texttt{hasCondDistrib\_feedback},
\texttt{hasCondDistrib\_action}), the stationary environment \texttt{stationaryEnv} with
\texttt{feedback\_stationaryEnv}, and, for bandits, the law of the reward of the $k$-th pull
of an arm (\texttt{Online/Bandit/RewardByCountMeasure.lean}: \texttt{hasLaw\_rewardByCount},
\texttt{iIndepFun\_rewardByCount}). Mathlib: \texttt{ProbabilityTheory.variance},
\texttt{variance\_def'}, \texttt{MeasureTheory.integral\_finset\_sum},
\texttt{integral\_indicator}, \texttt{Nat.nth} for the enumeration of the visits, and the
Cauchy--Schwarz inequalities \texttt{Finset.sum\_mul\_sq\_le\_sq\_mul\_sq} and
\texttt{integral\_mul\_le\_Lp\_mul\_Lq\_of\_nonneg}.

\textbf{What is missing.} Everything below: the iterated Markov property for finite-horizon
trajectory laws, occupancy measures, the unrolling lemma, entropic variances and the
sampling-at-visits lemma (the episodic analogue of \texttt{RewardByCountMeasure.lean}, a
candidate for LML).

\section{The iterated Markov property}
\label{sec:pre_mdp_markov}

\begin{lemma}[Iterated Markov property]\label{lem:pmdp_markov_iter}
\uses{def:episodic_mdp, def:policy, def:step_law}
Let $1 \le h' \le h \le H + 1$ and $s \in \Ss$. For every bounded measurable function $\Psi$ of
the rounds $(S_i, A_i, R_i)_{h' \le i < h}$ and every bounded measurable function $F$ of the
trajectory $(S_i, A_i, R_i)_{i \ge h}$,
\[
  \E^\pi_{(s,h')}\big[\Psi\cdot F\big] = \E^\pi_{(s,h')}\big[\Psi\cdot \Phi_F(S_h)\big],
  \qquad \Phi_F(s') := \E^\pi_{(s', h)}[F] ,
\]
i.e.\ conditionally on the rounds before $h$, the trajectory from $h$ has law
$\Prob^\pi_{(S_h, h)}$. Consequently, for $h \le H$, bounded measurable
$\psi : \As \times \R \times \Ss \to \R$ and every $f : \Ss \to \R$,
\begin{align*}
  \E^\pi_{(s,h')}\big[\Psi\cdot\psi(A_h, R_h, S_{h+1})\big]
  &= \E^\pi_{(s,h')}\Big[\Psi\cdot\E_{(R, s'') \sim (R_h \otimes p_h)(S_h, \pi_h(S_h))}\big[\psi(\pi_h(S_h), R, s'')\big]\Big], \\
  \E^\pi_{(s,h')}\big[\Psi\cdot f(S_{h+1})\big]
  &= \E^\pi_{(s,h')}\big[\Psi\cdot (p_h f)(S_h, \pi_h(S_h))\big] .
\end{align*}
\end{lemma}

\begin{proof}
\uses{lem:step_law_markov}
Induction on $k := h - h' \ge 0$, for all $h'$ and $s$ simultaneously. For $k = 0$ there are
no rounds before $h = h'$, so $\Psi$ is a constant $c$; since $S_{h'} = s$ a.s.\ under
$\Prob^\pi_{(s,h')}$ (Lemma~\ref{lem:step_law_markov}), $\Phi_F(S_{h'}) = \Phi_F(s) =
\E^\pi_{(s,h')}[F]$ a.s., and both sides equal $c\,\E^\pi_{(s,h')}[F]$. For the step, let
$h = h' + k + 1 \le H + 1$ (so $h' \le H$). Write the rounds before $h$ as the first round
$(S_{h'}, A_{h'}, R_{h'}, S_{h'+1})$ followed by the rounds $h' + 1 \le i < h$ of the shifted
trajectory $\theta$, and apply~\eqref{eq:pmdp_markov_integrated} to the function
$(A_{h'}, R_{h'}, S_{h'+1}, \theta) \mapsto \Psi\cdot F$:
\[
  \E^\pi_{(s,h')}[\Psi F]
  = \E_{(R,s')}\Big[\E^\pi_{(s',h'+1)}\big[\Psi_{R,s'}\cdot F\big]\Big],
\]
where $(R, s') \sim (R_{h'} \otimes p_{h'})(s, \pi_{h'}(s))$ and $\Psi_{R,s'}$ denotes $\Psi$
with its first round frozen at $(s, \pi_{h'}(s), R, s')$, a bounded measurable function of the
rounds $h' + 1 \le i < h$. By the induction hypothesis (with $h' + 1$ in place of $h'$,
$h - (h' + 1) = k$, and the starting state $s'$),
$\E^\pi_{(s',h'+1)}[\Psi_{R,s'} F] = \E^\pi_{(s',h'+1)}[\Psi_{R,s'}\Phi_F(S_h)]$. Applying
\eqref{eq:pmdp_markov_integrated} again, now to the function
$(A_{h'}, R_{h'}, S_{h'+1}, \theta) \mapsto \Psi\cdot\Phi_F(S_h)$, gives
$\E_{(R,s')}[\E^\pi_{(s',h'+1)}[\Psi_{R,s'}\Phi_F(S_h)]] = \E^\pi_{(s,h')}[\Psi\,\Phi_F(S_h)]$,
which proves the claim. For the consequences, take $F := \psi(A_h, R_h, S_{h+1})$, a function
of the first round of the trajectory from $h$: by Lemma~\ref{lem:step_law_markov} under
$\Prob^\pi_{(s',h)}$, $\Phi_F(s') = \E_{(R,s'') \sim (R_h\otimes p_h)(s', \pi_h(s'))}[\psi(\pi_h(s'), R, s'')]$;
with $\psi(a, r, s'') := f(s'')$ this is $\sum_{s''} p_h(s'' \mid s', \pi_h(s')) f(s'') = (p_h f)(s', \pi_h(s'))$.
\end{proof}

\textbf{Lean remark.} In \texttt{stepLaw M π (s, h)} the trajectory is a sequence of rounds
indexed by $\N$ and the paper's step $i$ is the round $i - h'$ (with the $0$-indexed
translation of Chapter~\ref{chap:mdp}); ``the trajectory from $h$'' is the shift by $h - h'$
rounds, and ``the rounds before $h$'' are the first $h - h'$ rounds. State the lemma as a
\texttt{HasCondDistrib} statement (the conditional law of the shifted trajectory given the
first $h - h'$ rounds is \texttt{stepLaw M π (S\_h, h)}), or directly in the integral form
above, by induction on $h - h'$ from LML's one-step statement. The layer component of the
state at round $k$ is deterministic ($h' + k$); this is part of the description of
\texttt{layerMDP} and is what allows writing $\pi_h(S_h)$ for the action.

\begin{lemma}[Actions and rewards along a trajectory]\label{lem:pmdp_action_eq_policy}
\uses{def:episodic_mdp, def:policy, def:step_law, def:reward_fn}
Let $1 \le h' \le h \le H$ and $s \in \Ss$. Under $\Prob^\pi_{(s,h')}$: $S_{h'} = s$ and
$A_h = \pi_h(S_h)$ almost surely; and if $\mathcal M$ has the reward function $r$
(Definition~\ref{def:reward_fn}) then $R_h = r_h(S_h, A_h)$ almost surely.
\end{lemma}

\begin{proof}
\uses{lem:step_law_markov, lem:pmdp_markov_iter}
$S_{h'} = s$ a.s.\ is part of Lemma~\ref{lem:step_law_markov}. Apply
Lemma~\ref{lem:pmdp_markov_iter} with $\Psi = 1$ and $F := \indic\{A_h \ne \pi_h(S_h)\}$, a
function of the first round of the trajectory from $h$: $\Phi_F(s') = \Prob^\pi_{(s',h)}(A_h \ne \pi_h(S_h)) = 0$
for every $s'$, because $S_h = s'$ and $A_h = \pi_h(s')$ a.s.\ under $\Prob^\pi_{(s',h)}$
(Lemma~\ref{lem:step_law_markov}). Hence $\Prob^\pi_{(s,h')}(A_h \ne \pi_h(S_h)) = 0$. For the
rewards, take $F := \indic\{R_h \ne r_h(S_h, A_h)\}$: under $\Prob^\pi_{(s',h)}$, $R_h$ has law
$R_h(\cdot \mid s', \pi_h(s')) = \delta_{r_h(s', \pi_h(s'))}$, so again $\Phi_F \equiv 0$.
\end{proof}

\section{Occupancy measures}
\label{sec:pre_mdp_occupancy}

\begin{definition}[Occupancy measure]\label{def:occupancy}
\lean{Learning.MDP.Episodic.occupancy}\leanok
\uses{def:episodic_mdp, def:policy, def:step_law}
For $1 \le h \le H$, $s \in \Ss$ and $a \in \As$, the \emph{occupancy measure} of $\pi$ from
$s_1$ at step $h$ is
\[
  p^\pi_h(s, a) := \Prob^\pi(S_h = s,\ A_h = a) = \Prob^\pi_{(s_1,1)}(S_h = s,\ A_h = a) .
\]
\end{definition}

\textbf{Lean remark.} \texttt{occupancy M π s₁ h s a := (stepLaw M π (s₁, 0)).real \{ω | (ω h).obs.1 = s ∧ (ω h).action = a\}}
with the $0$-indexed step \texttt{h : Fin H} (the paper's $h$ is Lean's $h - 1$), a real number
(\texttt{Measure.real}). The dependence on $s_1$ is kept explicit.

\begin{lemma}[Occupancy measures are probability vectors]\label{lem:occupancy_sum_one}
\uses{def:occupancy}
For $1 \le h \le H$: $p^\pi_h(s,a) \ge 0$ for all $(s, a)$, $\sum_{s,a} p^\pi_h(s,a) = 1$, and
\[
  p^\pi_h(s, a) = \indic\{a = \pi_h(s)\}\,\Prob^\pi(S_h = s) .
\]
\end{lemma}

\begin{proof}
\uses{lem:pmdp_action_eq_policy}
Nonnegativity is clear. The events $\{S_h = s, A_h = a\}$, $(s, a) \in \Ss \times \As$, are
pairwise disjoint and cover $\Omega$, so their probabilities sum to $1$ (finite additivity,
Mathlib \texttt{measureReal\_biUnion\_finset} or \texttt{sum\_measureReal\_preimage\_singleton}
for the map $\omega \mapsto (S_h, A_h)$). By Lemma~\ref{lem:pmdp_action_eq_policy},
$A_h = \pi_h(S_h)$ a.s., so up to a null set $\{S_h = s, A_h = a\} = \{S_h = s, \pi_h(s) = a\}$,
which is $\{S_h = s\}$ if $a = \pi_h(s)$ and empty otherwise.
\end{proof}

\begin{lemma}[Expectations of sums along the trajectory]\label{lem:occupancy_expectation}
\uses{def:occupancy}
For functions $f_h : \Ss \times \As \to \R$, $1 \le h \le H$,
\[
  \E^\pi\Big[\sum_{h=1}^H f_h(S_h, A_h)\Big] = \sum_{h=1}^H\sum_{s,a} p^\pi_h(s,a)\, f_h(s,a) .
\]
\end{lemma}

\begin{proof}
\uses{def:occupancy}
Pointwise, $f_h(S_h, A_h) = \sum_{s,a} f_h(s,a)\,\indic\{S_h = s, A_h = a\}$ (exactly one term
of the sum is nonzero). All functions are bounded, so the integral commutes with the finite
sums (\texttt{integral\_finset\_sum}, \texttt{integral\_const\_mul}), and
$\E^\pi[\indic\{S_h = s, A_h = a\}] = p^\pi_h(s,a)$ (\texttt{integral\_indicator\_one} or
\texttt{integral\_indicator} with \texttt{Measure.real}).
\end{proof}

\begin{lemma}[Forward recursion of the occupancy measures]\label{lem:occupancy_recursion}
\uses{def:occupancy}
For $1 \le h \le H - 1$ and $(s', a') \in \Ss \times \As$,
\[
  p^\pi_{h+1}(s', a') = \indic\{a' = \pi_{h+1}(s')\}\sum_{s,a} p^\pi_h(s,a)\, p_h(s' \mid s, a) .
\]
\end{lemma}

\begin{proof}
\uses{lem:occupancy_sum_one, lem:pmdp_markov_iter}
By Lemma~\ref{lem:occupancy_sum_one} at step $h + 1$ it suffices to compute
$\Prob^\pi(S_{h+1} = s')$. Decomposing over the value of $S_h$,
\[
  \Prob^\pi(S_{h+1} = s') = \sum_{s}\E^\pi\big[\indic\{S_h = s\}\,\indic\{S_{h+1} = s'\}\big]
  = \sum_{s}\E^\pi\big[\indic\{S_h = s\}\,p_h(s' \mid S_h, \pi_h(S_h))\big]
  = \sum_{s}\Prob^\pi(S_h = s)\,p_h(s' \mid s, \pi_h(s)) ,
\]
where the second equality is Lemma~\ref{lem:pmdp_markov_iter} (with $h' = 1$,
$\Psi = \indic\{S_h = s\}$ and $f = \indic\{\cdot = s'\}$, so that $(p_h f)(s, \pi_h(s)) = p_h(s' \mid s, \pi_h(s))$).
By Lemma~\ref{lem:occupancy_sum_one}, $\Prob^\pi(S_h = s)\,p_h(s'\mid s, \pi_h(s)) = \sum_a p^\pi_h(s,a)\,p_h(s' \mid s, a)$
(the terms with $a \ne \pi_h(s)$ vanish), which gives the claim.
\end{proof}

\section{Unrolling a backward recursion along the trajectory}
\label{sec:pre_mdp_unroll}

The certificate of Entropic-BPI satisfies, on the good event, a backward inequality of the
form $g_h(s) \le e^{\beta r_h(s,a)}[u_h(s,a) + \kappa\,(p_h g_{h+1})(s,a)]$ at the action
$a = \pi_h(s)$ of the greedy policy (Lemma~\ref{lem:cert_true_recursion_pos}). The next lemma
unrolls such an inequality into an expectation along the trajectory of $\pi$; it is the
unrolling step of Appendix~B.1 of \cite{essakine2026tight} (``Sample complexity''), stated for
arbitrary functions.

\begin{lemma}[Unrolling a backward recursive inequality]\label{lem:unroll_recursion}
\uses{def:episodic_mdp, def:policy, def:step_law}
Let $\beta \in \R$, $\kappa \ge 0$, and let $r_h, u_h : \Ss \times \As \to \R$ ($1 \le h \le H$)
and $g_h : \Ss \to \R$ ($1 \le h \le H + 1$) be arbitrary functions with $g_{H+1} = 0$. Assume
that for all $1 \le h \le H$ and $s \in \Ss$, with $a := \pi_h(s)$,
\[
  g_h(s) \le e^{\beta r_h(s,a)}\Big[u_h(s,a) + \kappa\,(p_h g_{h+1})(s,a)\Big] .
\]
Then for every $1 \le h \le H + 1$,
\begin{equation}\label{eq:pmdp_unroll_ih}
  \E^\pi\Big[e^{\beta\sum_{i < h} r_i(S_i, A_i)}\,g_h(S_h)\Big]
  \le \sum_{h' = h}^{H}\kappa^{h' - h}\,\E^\pi\Big[e^{\beta\sum_{i \le h'} r_i(S_i,A_i)}\,u_{h'}(S_{h'}, A_{h'})\Big] ,
\end{equation}
and in particular ($h = 1$, using $S_1 = s_1$ a.s.)
\[
  g_1(s_1) \le \sum_{h = 1}^{H}\kappa^{h - 1}\,\E^\pi\Big[e^{\beta\sum_{i \le h} r_i(S_i,A_i)}\,u_h(S_h, A_h)\Big] .
\]
\end{lemma}

\begin{proof}
\uses{lem:pmdp_action_eq_policy, lem:pmdp_markov_iter}
Write $W_h := e^{\beta\sum_{i < h} r_i(S_i, A_i)}$ for $1 \le h \le H + 1$, a positive function
of the rounds before $h$, with $W_1 = 1$ and $W_{h+1} = W_h\,e^{\beta r_h(S_h, A_h)}$. We prove
\eqref{eq:pmdp_unroll_ih} by backward induction on $h$.

\emph{Base case $h = H + 1$.} The left-hand side is $\E^\pi[W_{H+1}\,g_{H+1}(S_{H+1})] = 0$
since $g_{H+1} = 0$, and the right-hand side is an empty sum, equal to $0$.

\emph{Induction step.} Let $1 \le h \le H$ and assume \eqref{eq:pmdp_unroll_ih} for $h + 1$.
By Lemma~\ref{lem:pmdp_action_eq_policy}, $A_h = \pi_h(S_h)$ a.s., so the hypothesis applied at
$s = S_h$ gives, almost surely,
\[
  g_h(S_h) \le e^{\beta r_h(S_h, A_h)}\Big[u_h(S_h, A_h) + \kappa\,(p_h g_{h+1})(S_h, A_h)\Big] .
\]
Multiplying by $W_h > 0$, using $W_h e^{\beta r_h(S_h,A_h)} = W_{h+1}$, and taking expectations
(all functions are bounded):
\[
  \E^\pi[W_h\,g_h(S_h)] \le \E^\pi\big[W_{h+1}\,u_h(S_h, A_h)\big]
  + \kappa\,\E^\pi\big[W_{h+1}\,(p_h g_{h+1})(S_h, A_h)\big] .
\]
The first term is the summand $h' = h$ of the right-hand side of~\eqref{eq:pmdp_unroll_ih}
(with $\kappa^0 = 1$). For the second, $W_{h+1}$ is a function of the rounds before $h + 1$
and $A_h = \pi_h(S_h)$ a.s., so Lemma~\ref{lem:pmdp_markov_iter} (with $\Psi = W_{h+1}$ and
$f = g_{h+1}$) gives $\E^\pi[W_{h+1}\,(p_h g_{h+1})(S_h, A_h)] = \E^\pi[W_{h+1}\,g_{h+1}(S_{h+1})]$,
which by the induction hypothesis is at most
$\sum_{h' = h+1}^{H}\kappa^{h' - h - 1}\E^\pi[e^{\beta\sum_{i \le h'} r_i}u_{h'}(S_{h'},A_{h'})]$.
Multiplying by $\kappa \ge 0$ turns $\kappa^{h' - h - 1}$ into $\kappa^{h' - h}$, and adding the
summand $h' = h$ gives \eqref{eq:pmdp_unroll_ih} for $h$.

For the last claim, take $h = 1$: $W_1 = 1$ and $S_1 = s_1$ a.s.\ (Lemma~\ref{lem:pmdp_action_eq_policy}),
so the left-hand side is $g_1(s_1)$.
\end{proof}

\begin{remark}[Hypotheses]\label{rem:pmdp_unroll_hypotheses}
The paper applies the unrolling with $\kappa = 1 + 13/H$, nonnegative $u_h$ and rewards in
$[0, 1]$; none of $\kappa \ge 1$, $u_h \ge 0$, $r_h \ge 0$ is needed, only $\kappa \ge 0$ (to
multiply an inequality) and $e^{\beta r} > 0$ (to keep the direction of the inequality when
multiplying by $W_h$). The functions $r_h$ are arbitrary: the lemma is about the state-action
process only, and the reward kernels of $\mathcal M$ play no role. \textbf{Lean remark.} The
induction is on the number $j = H + 1 - h$ of steps to the end (the convention of the backward
recursions of Chapter~\ref{chap:algorithm}); $\sum_{i \le h} r_i(S_i, A_i)$ is a
\texttt{Finset.sum} over \texttt{Finset.range h} of the rounds of the trajectory, and
\eqref{eq:pmdp_unroll_ih} is the statement to prove by induction (the ``in particular'' is its
instance $h = 1$).
\end{remark}

\begin{lemma}[Cauchy--Schwarz along a trajectory]\label{lem:cauchy_schwarz_traj}
Let $(\Omega, \Prob)$ be a probability space and, for $1 \le h \le H$, let
$W_h, V_h, U_h : \Omega \to [0, \infty)$ be bounded random variables. Then
\[
  \E\Big[\sum_{h = 1}^H W_h\sqrt{V_h U_h}\Big]
  \le \sqrt{\E\Big[\sum_{h=1}^H W_h^2 V_h\Big]}\;\sqrt{\E\Big[\sum_{h=1}^H U_h\Big]} .
\]
In particular this applies to $W_h = w_h((S_i, A_i)_{i \le h})$, $V_h = v_h(S_h, A_h)$,
$U_h = u_h(S_h, A_h)$ under $\Prob^\pi$ for nonnegative functions $w_h, v_h, u_h$.
\end{lemma}

\begin{proof}
Let $X_h := W_h\sqrt{V_h} \ge 0$ and $Y_h := \sqrt{U_h} \ge 0$, so that
$X_h Y_h = W_h\sqrt{V_h U_h}$ (\texttt{Real.sqrt\_mul}), $X_h^2 = W_h^2 V_h$ and $Y_h^2 = U_h$
(\texttt{Real.sq\_sqrt}). By the Cauchy--Schwarz inequality for integrals (Mathlib
\texttt{integral\_mul\_le\_Lp\_mul\_Lq\_of\_nonneg} with $p = q = 2$; the variables are bounded,
hence in $L^2$), $\E[X_h Y_h] \le \sqrt{\E X_h^2}\sqrt{\E Y_h^2}$ for each $h$. Summing over
$h$ and applying the Cauchy--Schwarz inequality for finite sums to the vectors
$(\sqrt{\E X_h^2})_h$ and $(\sqrt{\E Y_h^2})_h$ (\texttt{Finset.sum\_mul\_sq\_le\_sq\_mul\_sq},
then \texttt{Real.sqrt\_le\_sqrt} and \texttt{Real.sqrt\_mul}),
\[
  \E\Big[\sum_h X_h Y_h\Big] = \sum_h\E[X_h Y_h] \le \sum_h\sqrt{\E X_h^2}\sqrt{\E Y_h^2}
  \le \sqrt{\sum_h\E X_h^2}\,\sqrt{\sum_h\E Y_h^2}
  = \sqrt{\E\Big[\sum_h X_h^2\Big]}\,\sqrt{\E\Big[\sum_h Y_h^2\Big]} ,
\]
using the linearity of the integral for the finite sums (\texttt{integral\_finset\_sum}).
\end{proof}

\textbf{Lean remark.} The lemma is pure probability (no MDP); it is used in
Lemma~\ref{lem:ratio_bound_pos} with $W_h = e^{\beta\sum_{i \le h} r_i(S_i,A_i)}$, which is
a function of the whole prefix of the trajectory, not of $(S_h, A_h)$ alone: this is why the
statement is for general random variables and not through the occupancy measures.

\section{Entropic variances}
\label{sec:pre_mdp_entropic_variance}

Throughout this section the reward kernels of $\mathcal M$ are supported in a bounded
interval (the paper's rewards are deterministic with values in $[0, 1]$), so that every
exponentiated return $e^{\beta R^\pi_h}$ is bounded and square-integrable, and for the
recursion (Lemma~\ref{lem:entropic_variance_recursion} onwards) $\mathcal M$ has a
deterministic reward function $r$ (Definition~\ref{def:reward_fn}).

\begin{definition}[Entropic variances]\label{def:entropic_variance}
\lean{Learning.MDP.Episodic.entropicVarQ, Learning.MDP.Episodic.entropicVarV}\leanok
\uses{def:episodic_mdp, def:policy, def:step_law, def:episode_return, def:exp_value}
For $1 \le h \le H$ and $(s, a) \in \Ss \times \As$, let $\Prob^\pi_{(s,a,h)}$ be the law of
$(R, (S_i, A_i, R_i)_{i \ge h+1})$ where $(R, S') \sim (R_h \otimes p_h)(s, a)$ and, given
$(R, S')$, the trajectory from step $h + 1$ has law $\Prob^\pi_{(S', h+1)}$: the law of the
rest of the episode when the action $a$ is played at $s$ at step $h$ and $\pi$ is followed
afterwards. Under it, $R + R^\pi_{h+1}$ is the return from $(s, a)$, and by
Definition~\ref{def:exp_value} and Fubini,
$\E^\pi_{(s,a,h)}\big[e^{\beta(R + R^\pi_{h+1})}\big] = \E\big[e^{\beta R}Z^\pi_{h+1}(S')\big] = U^\pi_h(s,a)$.
The \emph{entropic variances} of $\pi$ are
\[
  \sigma Q^\pi_h(s, a) := \Var_{\Prob^\pi_{(s,a,h)}}\big(e^{\beta(R + R^\pi_{h+1})}\big)
  = \E^\pi_{(s,a,h)}\Big[\big(e^{\beta(R + R^\pi_{h+1})} - U^\pi_h(s,a)\big)^2\Big],
  \qquad
  \sigma V^\pi_h(s) := \sigma Q^\pi_h(s, \pi_h(s)),
\]
for $1 \le h \le H$, and $\sigma V^\pi_{H+1} := 0$.
\end{definition}

\textbf{Lean remark.} \texttt{entropicVarQ M β π h s a} is
\texttt{variance (fun x ↦ exp (β * (x.1 + episodeReturn x.2))) (saLaw M π h s a)} where
\texttt{saLaw M π h s a} is the composition of the step kernel $(R_h \otimes p_h)(s,a)$ with
\texttt{fun (R, s') ↦ (dirac R).prod (stepLaw M π (s', h + 1))} (``\texttt{stepLaw} from
\texttt{(x.2, h+1)} after the first transition''); \texttt{entropicVarV M β π h s :=
entropicVarQ M β π h s (π h s)} for $h \le H$ and $0$ at the terminal step. This matches the
paper's ``$\E^\pi[\,\cdot \mid S_h = s, A_h = a]$'' (Lemma~23) without conditioning on an event
of probability zero when $a \ne \pi_h(s)$.

\begin{lemma}[The entropic variance of a state is a variance under the trajectory law]\label{lem:pmdp_entropic_variance_V}
\uses{def:entropic_variance, def:exp_value}
For $1 \le h \le H$ and $s \in \Ss$,
\[
  \sigma V^\pi_h(s) = \Var_{\Prob^\pi_{(s,h)}}\big(e^{\beta R^\pi_h}\big)
  = \E^\pi_{(s,h)}\Big[\big(e^{\beta R^\pi_h} - Z^\pi_h(s)\big)^2\Big] ,
\]
and both sides are $0$ for $h = H + 1$.
\end{lemma}

\begin{proof}
\uses{lem:step_law_markov, lem:exp_bellman}
Under $\Prob^\pi_{(s,h)}$, $A_h = \pi_h(s)$ a.s.\ and the pair $(R_h, \theta)$, with $\theta$
the trajectory from $h + 1$, has law $\Prob^\pi_{(s, \pi_h(s), h)}$ by
Lemma~\ref{lem:step_law_markov} (the law of $(R_h, S_{h+1})$ is $(R_h \otimes p_h)(s, \pi_h(s))$
and, given the first round, $\theta \sim \Prob^\pi_{(S_{h+1}, h+1)}$), and
$R^\pi_h = R_h + R^\pi_{h+1}(\theta)$. So $e^{\beta R^\pi_h}$ under $\Prob^\pi_{(s,h)}$ and
$e^{\beta(R + R^\pi_{h+1})}$ under $\Prob^\pi_{(s, \pi_h(s), h)}$ have the same law, hence the same
variance (\texttt{variance} depends only on the law: \texttt{variance\_map} /
\texttt{IdentDistrib.variance\_eq}). The mean of $e^{\beta R^\pi_h}$ is $Z^\pi_h(s)$
(Definition~\ref{def:exp_value}), which equals $U^\pi_h(s, \pi_h(s))$ by
Lemma~\ref{lem:exp_bellman}, and for a square-integrable variable the variance is the mean
square deviation from the mean (\texttt{variance\_def'} or the definition of
\texttt{variance} through \texttt{evariance}). For $h = H + 1$, $R^\pi_{H+1} = 0$, the variable
is the constant $1$ and $Z^\pi_{H+1} = 1$, so both sides are $0 = \sigma V^\pi_{H+1}$.
\end{proof}

\begin{lemma}[Entropic variance recursion]\label{lem:entropic_variance_recursion}
\uses{def:entropic_variance, def:exp_value, def:reward_fn}
Assume $\mathcal M$ has the reward function $r$. For $1 \le h \le H$ and $(s,a) \in \Ss\times\As$,
\[
  \sigma Q^\pi_h(s, a) = e^{2\beta r_h(s,a)}\,\Var_{p_h}\big(Z^\pi_{h+1}\big)(s,a)
  + e^{2\beta r_h(s,a)}\,\big(p_h\,\sigma V^\pi_{h+1}\big)(s,a) .
\]
\end{lemma}

\begin{proof}
\uses{lem:pmdp_entropic_variance_V, def:exp_value}
Write $r := r_h(s,a)$, $Z := Z^\pi_{h+1}$ and $m := (p_h Z)(s,a)$. Under $\Prob^\pi_{(s,a,h)}$ the
reward kernel is $\delta_{r}$, so $R = r$ a.s.\ and $e^{\beta(R + R^\pi_{h+1})} = e^{\beta r}X$ with
$X := e^{\beta R^\pi_{h+1}}$; hence $\sigma Q^\pi_h(s,a) = e^{2\beta r}\Var(X)$
(\texttt{variance\_mul\_const}). The law of $(S', \theta)$ under $\Prob^\pi_{(s,a,h)}$ is the
mixture $\sum_{s'} p_h(s' \mid s,a)\,\delta_{s'} \otimes \Prob^\pi_{(s',h+1)}$ and $X$ is a
function of $\theta$, so for every bounded measurable $\phi$,
$\E^\pi_{(s,a,h)}[\phi(X)] = \sum_{s'} p_h(s' \mid s,a)\,\E^\pi_{(s',h+1)}[\phi(X)]$. With
$\phi = \mathrm{id}$ this gives $\E[X] = \sum_{s'} p_h(s'\mid s,a) Z(s') = m$
(Definition~\ref{def:exp_value}). We now make the law of total variance (conditioning on
$S' = s'$) explicit. For each $s'$, since $\E^\pi_{(s',h+1)}[X] = Z(s')$,
\[
  \E^\pi_{(s',h+1)}\big[(X - m)^2\big]
  = \E^\pi_{(s',h+1)}\big[(X - Z(s'))^2\big] + 2\,(Z(s') - m)\,\E^\pi_{(s',h+1)}\big[X - Z(s')\big] + (Z(s') - m)^2
  = \sigma V^\pi_{h+1}(s') + (Z(s') - m)^2 ,
\]
by Lemma~\ref{lem:pmdp_entropic_variance_V} at step $h + 1$ (for $h = H$: $X = 1 = Z(s')$ and
$\sigma V^\pi_{H+1} = 0$, consistent). Therefore
\[
  \Var(X) = \E^\pi_{(s,a,h)}\big[(X - m)^2\big]
  = \sum_{s'} p_h(s' \mid s,a)\Big(\sigma V^\pi_{h+1}(s') + (Z(s') - m)^2\Big)
  = \big(p_h\,\sigma V^\pi_{h+1}\big)(s,a) + \Var_{p_h}(Z)(s,a) ,
\]
and multiplying by $e^{2\beta r}$ gives the claim. This is the paper's Lemma~23.
\end{proof}

\begin{lemma}[Telescoping of the entropic variances]\label{lem:entropic_variance_sum}
\uses{def:entropic_variance, def:exp_value, def:reward_fn}
Assume $\mathcal M$ has the reward function $r$. Then
\[
  \sum_{h=1}^{H}\E^\pi\Big[e^{2\beta\sum_{i \le h} r_i(S_i, A_i)}\,\Var_{p_h}\big(Z^\pi_{h+1}\big)(S_h, A_h)\Big]
  = \sigma V^\pi_1(s_1) = \Var_{\Prob^\pi}\big(e^{\beta R^\pi_1}\big) .
\]
\end{lemma}

\begin{proof}
\uses{lem:entropic_variance_recursion, lem:pmdp_entropic_variance_V, lem:pmdp_action_eq_policy, lem:pmdp_markov_iter}
Let $W_h := e^{2\beta\sum_{i < h} r_i(S_i, A_i)}$ for $1 \le h \le H + 1$ ($W_1 = 1$,
$W_{h+1} = W_h e^{2\beta r_h(S_h, A_h)}$), a bounded function of the rounds before $h$. Fix
$1 \le h \le H$. By Lemma~\ref{lem:pmdp_action_eq_policy}, $A_h = \pi_h(S_h)$ a.s., so
$\sigma V^\pi_h(S_h) = \sigma Q^\pi_h(S_h, A_h)$ a.s., and Lemma~\ref{lem:entropic_variance_recursion}
at $(S_h, A_h)$ gives, after multiplication by $W_h$ and integration,
\[
  \E^\pi\big[W_h\,\sigma V^\pi_h(S_h)\big]
  = \E^\pi\big[W_{h+1}\Var_{p_h}(Z^\pi_{h+1})(S_h, A_h)\big]
  + \E^\pi\big[W_{h+1}\,(p_h\,\sigma V^\pi_{h+1})(S_h, A_h)\big] .
\]
By Lemma~\ref{lem:pmdp_markov_iter} ($\Psi = W_{h+1}$, a function of the rounds before $h + 1$,
and $f = \sigma V^\pi_{h+1}$, with $A_h = \pi_h(S_h)$ a.s.), the last term equals
$\E^\pi[W_{h+1}\,\sigma V^\pi_{h+1}(S_{h+1})]$. Hence, with $D_h := \E^\pi[W_h\,\sigma V^\pi_h(S_h)]$
for $1 \le h \le H + 1$,
\[
  \E^\pi\big[W_{h+1}\Var_{p_h}(Z^\pi_{h+1})(S_h, A_h)\big] = D_h - D_{h+1} ,
\]
and $W_{h+1} = e^{2\beta\sum_{i \le h} r_i(S_i,A_i)}$. Summing over $h = 1, \dots, H$ telescopes
(\texttt{Finset.sum\_range\_sub'}) to $D_1 - D_{H+1}$. Now $D_{H+1} = 0$ because
$\sigma V^\pi_{H+1} = 0$, and $D_1 = \sigma V^\pi_1(s_1)$ because $W_1 = 1$ and $S_1 = s_1$ a.s.
Finally $\sigma V^\pi_1(s_1) = \Var_{\Prob^\pi}(e^{\beta R^\pi_1})$ by
Lemma~\ref{lem:pmdp_entropic_variance_V}.
\end{proof}

\begin{lemma}[Normalized variance of the exponentiated return]\label{lem:variance_ratio_le}
\uses{def:exp_value, def:episode_return, def:gmax}
Let $G \ge 0$ be such that $R^\pi_1 \in [0, G]$ almost surely under $\Prob^\pi$. Then
$Z^\pi_1(s_1) > 0$ and
\[
  \frac{\Var_{\Prob^\pi}\big(e^{\beta R^\pi_1}\big)}{Z^\pi_1(s_1)^2}
  \le \frac{(e^{\abs\beta G} - 1)^2}{4\,e^{\abs\beta G}}
  \le \frac{(e^{\abs\beta G} - 1)^2}{e^{\abs\beta G}} .
\]
In particular, if $\mathcal M$ has a reward function with values in $[0, 1]$, this holds with
$G = \Gmax(\mathcal M)$.
\end{lemma}

\begin{proof}
\uses{lem:normalized_variance_bound, lem:gmax_le_horizon, lem:pmdp_action_eq_policy}
Apply Lemma~\ref{lem:normalized_variance_bound} on $(\Omega, \Prob^\pi)$ with $f := R^\pi_1$ and
$R := G$: the mean of $Y = e^{\beta R^\pi_1}$ is $Z^\pi_1(s_1)$ (Definition~\ref{def:exp_value}),
which is positive, and the first inequality follows; the second holds because
$4e^{\abs\beta G} \ge e^{\abs\beta G} > 0$. For the last claim: by
Lemma~\ref{lem:gmax_le_horizon}, $\Gmax(\mathcal M) \ge 0$ and $R^\pi_1 \le \Gmax(\mathcal M)$
a.s.\ under $\Prob^\pi$; and $R^\pi_1 = \sum_{i \le H} r_i(S_i, A_i) \ge 0$ a.s.\ by
Lemma~\ref{lem:pmdp_action_eq_policy} since $r \ge 0$.
\end{proof}

\textbf{Lean remark.} This is where the paper's Lemma~24 enters (Lemma~\ref{lem:ratio_bound_pos}).
The bound is used with the paper's constant $V_G = (e^{\abs\beta G} - 1)^2/e^{\abs\beta G}$
(Definition~\ref{def:upper_bound_const}); the factor $4$ is slack.

\section{Sampling at the visits of a state-action pair}
\label{sec:pre_mdp_visits}

The concentration events of Chapter~\ref{chap:empirical} are about the transitions observed
at the visits of a fixed pair $(s, a)$ at a fixed step $h$, across the episodes of a run of an
arbitrary algorithm. The visits are chosen by the algorithm (a policy is played, the pair may
or may not be visited), so the observed next states are not literally an i.i.d.\ sample;
what is true is that each observed next state is drawn from $p_h(\cdot \mid s, a)$
independently of everything before it. We state this in the three forms that the events use.

\textbf{Setting and notation.} $(X, Y)$ is a run of an arbitrary algorithm in the episode
environment of $\mathcal M$ from $s_1$ (Definition~\ref{def:episode_env}), on $(\Omega, \Prob)$:
$X_t = \pi^{t+1}$ is the policy of episode $t + 1$ and $Y_t = (s^{t+1}_1, \dots, s^{t+1}_{H+1})$
its state sequence; $a^{t+1}_h := \pi^{t+1}_h(s^{t+1}_h)$; $\mathrm{hist}_t := (X_i, Y_i)_{i < t}$
is the history of the first $t$ episodes, and $n^t_h(s,a)$, $n^t_h(s,a,s')$, $\phat^t_h$ are the
counts and empirical transitions of Definitions~\ref{def:counts} and~\ref{def:emp_trans}. Fix
$1 \le h \le H$ and $(s, a) \in \Ss \times \As$. For $i \ge 1$ let
$w_i := \indic\{(s^i_h, a^i_h) = (s, a)\}$ be the indicator that episode $i$ visits $(s,a)$ at
step $h$, so that $n^t_h(s,a) = \sum_{i \le t} w_i$. The \emph{$k$-th visit time} ($k \ge 1$) is
\[
  T_k := \min\Big\{i \ge 1 : \sum_{j \le i} w_j = k\Big\} \in \{1, 2, \dots\} \cup \{\infty\},
\]
the index of the $k$-th episode with $w_i = 1$ ($T_k = \infty$ if there are fewer than $k$ such
episodes), and $W_k := s^{T_k}_{h+1}$ on $\{T_k < \infty\}$ (and $W_k := s$, say, elsewhere) is
the \emph{$k$-th observed next state}. Then $T_1 < T_2 < \cdots$ on the event where they are
finite, $\{T_k = i\}$ is determined by $\mathrm{hist}_{i-1}$, $X_{i-1}$ and $s^i_h$
($T_k = i$ iff $\sum_{j < i} w_j = k - 1$ and $w_i = 1$), and $n^t_h(s,a) \ge k$ iff $T_k \le t$.
For $t \ge 0$ let $\mathcal G_t$ be the $\sigma$-algebra generated by
$(\mathrm{hist}_t, X_t, s^{t+1}_1, \dots, s^{t+1}_h)$: the information available just before the
transition of episode $t + 1$ at step $h$; $w_{t+1}$ is $\mathcal G_t$-measurable.

\begin{lemma}[Counts and episodes]\label{lem:count_le_episodes}
\uses{def:counts}
For every $t \ge 0$ and $h \le H$: $0 \le n^t_h(s,a) \le t$ for every $(s,a)$, and
$\sum_{s,a} n^t_h(s,a) = t$.
\end{lemma}

\begin{proof}
\uses{def:counts}
$n^t_h(s,a)$ is a sum of $t$ indicators, hence between $0$ and $t$
(\texttt{Finset.sum\_le\_card\_nsmul}). Exchanging the sums,
$\sum_{s,a} n^t_h(s,a) = \sum_{i \le t}\sum_{s,a}\indic\{(s^i_h, a^i_h) = (s,a)\} = \sum_{i \le t} 1 = t$,
since for each episode $i$ exactly one pair equals $(s^i_h, a^i_h)$ (\texttt{Finset.sum\_comm},
\texttt{Finset.sum\_ite\_eq}).
\end{proof}

\begin{lemma}[The empirical transitions are the empirical distribution of the observed next states]\label{lem:emp_trans_reindex}
\uses{def:counts, def:emp_trans}
For every $t \ge 0$, let $n := n^t_h(s,a)$. Then $\{i \le t : w_i = 1\} = \{T_1, \dots, T_n\}$
with $T_1 < \dots < T_n \le t < T_{n+1}$, and for every $s' \in \Ss$,
\[
  n^t_h(s, a, s') = \sum_{k = 1}^{n}\indic\{W_k = s'\}, \qquad\text{hence, if } n \ge 1,\qquad
  \phat^t_h(s' \mid s, a) = \frac1n\sum_{k=1}^n\indic\{W_k = s'\} .
\]
In particular $\phat^t_h(\cdot \mid s,a)$ depends on $t$ only through $n^t_h(s,a)$: it is the
empirical distribution of $W_1, \dots, W_{n}$.
\end{lemma}

\begin{proof}
\uses{def:counts, def:emp_trans}
By definition, $T_k$ is the $k$-th element (in increasing order) of the set
$I := \{i \ge 1 : w_i = 1\}$, and $I \cap \{1, \dots, t\}$ has exactly $n = \sum_{i \le t} w_i$
elements; so $I \cap \{1,\dots,t\} = \{T_1, \dots, T_n\}$, $T_n \le t$ and $T_{n+1} > t$. Then
\[
  n^t_h(s,a,s') = \sum_{i \le t} w_i\,\indic\{s^i_{h+1} = s'\}
  = \sum_{i \in I \cap \{1,\dots,t\}}\indic\{s^i_{h+1} = s'\}
  = \sum_{k=1}^n\indic\{s^{T_k}_{h+1} = s'\} = \sum_{k=1}^n\indic\{W_k = s'\}
\]
(reindexing the sum by the bijection $k \mapsto T_k$, \texttt{Finset.sum\_bij} or
\texttt{Finset.sum\_map}), and $\phat^t_h(s'\mid s,a) = n^t_h(s,a,s')/n$ for $n \ge 1$
(Definition~\ref{def:emp_trans}).
\end{proof}

\textbf{Lean remark.} With the episodes $0$-indexed, $T_k - 1 = $ \texttt{Nat.nth (fun i ↦ w i = 1) (k - 1)}
(Mathlib \texttt{Nat.nth}, with \texttt{Nat.count}, \texttt{Nat.nth\_count},
\texttt{Nat.nth\_mem\_of\_lt\_card}, \texttt{Nat.count\_nth\_of\_infinite}/\texttt{count\_nth});
$\{T_k < \infty\}$ is ``$k - 1 <$ the number of visits'' or ``the set of visits is infinite''.
Alternatively, avoid $\infty$ by working with the finite enumeration
\texttt{Finset.orderEmbOfFin} of $\{i < t : w_i = 1\}$ for each $t$; the statement (i) below
about $\{T_n < \infty\}$ is then a statement about $\bigcup_t\{n^t_h(s,a) = n\}$.

\begin{lemma}[Sampling at the visits]\label{lem:visits_iid}
\uses{def:episode_env, def:counts, def:emp_trans, lem:emp_trans_reindex}
In the setting above, with $p := p_h(\cdot \mid s, a)$:
\begin{enumerate}
  \item[(i)] (\emph{Dominated law of the observed next states.}) For every $n \ge 1$ and
        $(w_1, \dots, w_n) \in \Ss^n$,
        \[
          \Prob\big(T_n < \infty,\ W_1 = w_1, \dots, W_n = w_n\big)
          = \Prob\big(T_n < \infty,\ W_1 = w_1, \dots, W_{n-1} = w_{n-1}\big)\,p(w_n)
          \le \prod_{k=1}^n p(w_k) .
        \]
        Consequently, for every $B \subseteq \Ss^n$,
        $\Prob\big(T_n < \infty,\ (W_1, \dots, W_n) \in B\big) \le p^{\otimes n}(B)$, and since
        $\{T_n < \infty\} = \{\exists t \ge 0 : n^t_h(s,a) = n\}$, for every $B$,
        \[
          \Prob\big(\exists t \ge 0 :\ n^t_h(s,a) = n \text{ and } (W_1, \dots, W_n) \in B\big)
          \le p^{\otimes n}(B) ,
        \]
        where $p^{\otimes n}$ is the law of $n$ i.i.d.\ draws from $p$. In words: on the event
        that the $n$-th visit happens, the first $n$ observed next states are dominated by an
        i.i.d.\ sample of $p$; by Lemma~\ref{lem:emp_trans_reindex}, this applies to every
        event about $\phat^t_h(\cdot \mid s,a)$ with $n^t_h(s,a) = n$, uniformly in $t$.
  \item[(ii)] (\emph{Martingale differences along the episodes.}) For every $t \ge 0$, every
        bounded $\mathcal G_t$-measurable $\Psi$ and every $f : \Ss \to \R$,
        \[
          \E\big[\Psi\, f(s^{t+1}_{h+1})\big] = \E\big[\Psi\,(p_h f)(s^{t+1}_h, a^{t+1}_h)\big] ,
          \qquad\text{in particular}\qquad
          \E\big[\Psi\, w_{t+1}\, f(s^{t+1}_{h+1})\big] = \E\big[\Psi\, w_{t+1}\big]\,(pf) ;
        \]
        i.e.\ $\E[f(s^{t+1}_{h+1}) \mid \mathcal G_t] = (p_h f)(s^{t+1}_h, a^{t+1}_h)$ a.s., and
        $w_{t+1}\big(f(s^{t+1}_{h+1}) - pf\big)$ is a martingale difference with respect to
        $(\mathcal G_t)_t$, with $w_{t+1}$ predictable.
  \item[(iii)] (\emph{Martingale differences along the visits.}) Let $k \ge 1$ and let $E$ be an
        event such that $E \cap \{T_k = t + 1\}$ is $\mathcal G_t$-measurable for every $t \ge 0$
        (an event of the information available at the $k$-th visit; for instance
        $E = \{W_1 = w_1, \dots, W_{k-1} = w_{k-1}\}$ or $E = \Omega$). Then for every $f : \Ss \to \R$,
        \[
          \E\big[\indic_{E}\,\indic\{T_k < \infty\}\, f(W_k)\big]
          = \Prob\big(E \cap \{T_k < \infty\}\big)\,(pf) .
        \]
\end{enumerate}
\end{lemma}

\begin{proof}
\uses{lem:episode_env_law, lem:pmdp_markov_iter, lem:emp_trans_reindex}
\emph{(ii).} Lemma~\ref{lem:episode_env_law} states that, conditionally on
$\mathrm{hist}_t$, $X_t$ and $s^{t+1}_1, \dots, s^{t+1}_h$, the next state $s^{t+1}_{h+1}$ has law
$p_h(\cdot \mid s^{t+1}_h, a^{t+1}_h)$: $\E[f(s^{t+1}_{h+1}) \mid \mathcal G_t] = (p_h f)(s^{t+1}_h, a^{t+1}_h)$
a.s. Multiplying by the bounded $\mathcal G_t$-measurable $\Psi$ and integrating gives the
first identity (this is also obtained directly: by Lemma~\ref{lem:episode_env_law} the
conditional law of $Y_t$ given $(\mathrm{hist}_t, X_t)$ is the law of the state sequence under
$\Prob^{X_t}_{(s_1,1)}$, and under that law Lemma~\ref{lem:pmdp_markov_iter} with $h' = 1$,
$\Psi$ a function of the rounds before $h + 1$ and $f$ gives
$\E^{\pi}[\Psi f(S_{h+1})] = \E^\pi[\Psi\,(p_h f)(S_h, \pi_h(S_h))]$; then integrate over
$(\mathrm{hist}_t, X_t)$). For the particular case, $\Psi w_{t+1}$ is bounded and
$\mathcal G_t$-measurable, and on $\{w_{t+1} = 1\}$, $(p_h f)(s^{t+1}_h, a^{t+1}_h) = pf$.

\emph{(iii).} For $t \ge 0$ let $E_t := E \cap \{T_k = t + 1\}$, which is $\mathcal G_t$-measurable
by assumption. On $E_t$, $w_{t+1} = 1$, so $(s^{t+1}_h, a^{t+1}_h) = (s, a)$ and
$W_k = s^{t+1}_{h+1}$. By (ii) with $\Psi = \indic_{E_t}$,
\[
  \E\big[\indic_{E_t}\, f(W_k)\big] = \E\big[\indic_{E_t}\, f(s^{t+1}_{h+1})\big]
  = \E\big[\indic_{E_t}\,(p_h f)(s^{t+1}_h, a^{t+1}_h)\big] = \Prob(E_t)\,(pf) .
\]
The events $E_t$, $t \ge 0$, are pairwise disjoint with union $E \cap \{T_k < \infty\}$; summing
over $t$ (countable additivity of $\Prob$ and of the integral: \texttt{integral\_iUnion} or
\texttt{tsum} of the indicators, all terms being bounded by $\sup\abs f\,\Prob(E_t)$ with
$\sum_t\Prob(E_t) \le 1$) gives the claim.

\emph{(i).} Fix $n \ge 1$ and $w_1, \dots, w_n$. The event $E := \{W_1 = w_1, \dots, W_{n-1} = w_{n-1}\}$
($E = \Omega$ if $n = 1$) satisfies the hypothesis of (iii) with $k = n$: on $\{T_n = t + 1\}$
the visit times $T_1 < \dots < T_{n-1}$ are at most $t$, so $W_1, \dots, W_{n-1}$ are functions
of $\mathrm{hist}_t$, and $\{T_n = t+1\}$ is $\mathcal G_t$-measurable. Applying (iii) with
$f = \indic\{\cdot = w_n\}$ gives the equality
$\Prob(T_n < \infty, W_1 = w_1, \dots, W_n = w_n) = \Prob(T_n < \infty, W_1 = w_1, \dots, W_{n-1} = w_{n-1})\,p(w_n)$.
For the inequality, induct on $n$: for $n = 1$ the right-hand side is $\Prob(T_1 < \infty)p(w_1) \le p(w_1)$;
for $n \ge 2$, $\{T_n < \infty\} \subseteq \{T_{n-1} < \infty\}$, so the first factor is at most
$\Prob(T_{n-1} < \infty, W_1 = w_1, \dots, W_{n-1} = w_{n-1}) \le \prod_{k < n} p(w_k)$ by the
induction hypothesis. For $B \subseteq \Ss^n$ (a finite set), sum over $(w_1,\dots,w_n) \in B$:
$\Prob(T_n < \infty, (W_1,\dots,W_n) \in B) = \sum_{w \in B}\Prob(T_n < \infty, W_{\le n} = w) \le \sum_{w \in B}\prod_k p(w_k) = p^{\otimes n}(B)$.
Finally $t \mapsto n^t_h(s,a)$ is nondecreasing, starts at $0$ and increases by at most $1$ per
episode, so it takes the value $n$ at some $t$ iff it reaches $n$, iff $T_n < \infty$
(indeed $n^{T_n}_h(s,a) = n$ by Lemma~\ref{lem:emp_trans_reindex}).
\end{proof}

\begin{remark}[Use and Lean encoding]\label{rem:pmdp_visits_iid}
Item (i) is what the KL event needs (Lemma~\ref{lem:event_kl_prob}): for the event
$\{\KL(\phat^t_h(\cdot\mid s,a)\,\|\,p) > x\}$ on $\{n^t_h(s,a) = n\}$, which by
Lemma~\ref{lem:emp_trans_reindex} is $\{(W_1,\dots,W_n) \in B\}$ with
$B = \{w \in \Ss^n : \KL(\mathrm{emp}(w)\|p) > x\}$, (i) bounds its probability, uniformly in
$t$, by the i.i.d.\ probability that Sanov's bound (Theorem~\ref{thm:sanov_hp}) controls.
Items (ii) and (iii) are what the Bernstein and count events need
(Lemmas~\ref{lem:event_bern_prob} and~\ref{lem:event_cnt_prob}): (ii) with the filtration
$(\mathcal G_t)$ of the episodes and the predictable weights $w_{t+1} \in \{0, 1\}$, and (iii)
when the time index is the number of visits. \textbf{Lean remark.} This is the episodic
analogue of LML's \texttt{Online/Bandit/RewardByCountMeasure.lean}
(\texttt{stepsUntil}, \texttt{rewardByCount}, \texttt{hasLaw\_rewardByCount},
\texttt{iIndepFun\_rewardByCount}): the ``arm'' is the triple $(h, s, a)$ (episode $i$ ``pulls''
it iff $w_i = 1$) and the ``reward'' is the next state $s^i_{h+1}$. The bandit file obtains the
exact i.i.d.\ law of the rewards by count by extending the probability space with an
independent i.i.d.\ stream (\texttt{P.prod (streamMeasure ν)}) that fills in the pulls that never
happen; the dominated form (i), stated on $\{T_n < \infty\}$ without any extension, is all the
union bound needs and is the statement to contribute to LML (\texttt{Empirical.lean}). Item
(ii) is \texttt{lem:episode\_env\_law} in integral form (\texttt{HasCondDistrib} of
$s^{t+1}_{h+1}$ given $\mathcal G_t$, from LML's \texttt{IsAlgEnvSeq.hasCondDistrib\_feedback}
and \texttt{feedback\_stationaryEnv}); the $\sigma$-algebra $\mathcal G_t$ is the comap of the
map $\omega \mapsto (\mathrm{hist}_t, X_t, (s^{t+1}_i)_{i \le h})$.
\end{remark}

\chapter{Concentration inequalities}
\label{chap:pre_concentration}

This chapter contains the probabilistic and information-theoretic tools of the analysis
(Appendix~E, Lemmas~19--22, and Appendix~F, Lemmas~25--26 of \cite{essakine2026tight}), stated in
Mathlib-shaped generality with no reference to the MDP: the method of types and Sanov's
high-probability bound (Section~\ref{sec:pconc_sanov}), Ville's maximal inequality and the
Bernoulli maximal inequality of \cite{dann2017unifying} (Section~\ref{sec:pconc_ville}), the
self-normalized Bernstein inequality of \cite{menard2021fast} (Section~\ref{sec:pconc_bernstein}),
and the KL--Bernstein inequality of \cite{talebi2018variance} with the variance transport lemmas
(Section~\ref{sec:pconc_kl}). The paper only states Lemmas~20, 21, 22, 25 and~26; complete proofs
are given here. Every constant is verified; the two places where the paper is loose are recorded
in Remarks~\ref{rem:pconc_sanov_paper} and~\ref{rem:pconc_bernstein_constants}.

\textbf{Conventions.} A \emph{probability vector} on a finite set $\mathcal X$ is a map
$p : \mathcal X \to [0, 1]$ with $\sum_x p(x) = 1$; for $f : \mathcal X \to \R$ we write
$pf := \sum_x p(x) f(x)$, $\Var_p(f) := p(f - pf)^2 = pf^2 - (pf)^2$, and
$\norm{p - q}_1 := \sum_x \abs{p(x) - q(x)}$. The Kullback--Leibler divergence is
$\KL(p \| q) := \sum_{x : p(x) > 0} p(x) \log(p(x)/q(x)) \in [0, \infty]$, equal to $+\infty$ when
$p(x) > 0 = q(x)$ for some $x$ (we then say that $p$ is not dominated by $q$). Lean: $p$ is a
\texttt{weightedMeasure} on \texttt{Fin m} (a probability measure), $pf$ its integral,
$\Var_p(f)$ is \texttt{ProbabilityTheory.variance}, and $\KL$ is
\texttt{InformationTheory.klDiv}, valued in $[0, \infty]$ (\texttt{ℝ≥0∞}), with
\texttt{klDiv\_of\_not\_ac} for the infinite case and \texttt{toReal\_klDiv} for the finite-sum
formula. Throughout, $\log$ is the natural logarithm and $e$ its base.

\section{The method of types and Sanov's bound}
\label{sec:pconc_sanov}

Let $\mathcal X$ be a finite set with $m := \abs{\mathcal X} \ge 1$ elements, $p$ a probability
vector on $\mathcal X$, $n \ge 1$, and $X_1, \dots, X_n$ i.i.d.\ random variables with law $p$ on
a probability space $(\Omega, \Prob)$. The \emph{empirical distribution} is
$\phat_n(x) := \frac1n \sum_{k \le n} \indic\{X_k = x\}$, a random probability vector. A
\emph{type} (of denominator $n$) is a probability vector $q$ with $n q(x) \in \N$ for every $x$;
$\phat_n$ is a type.

\begin{lemma}[Probability of a type]\label{lem:type_prob_le}
For every type $q$ of denominator $n$,
$\Prob(\phat_n = q) \le e^{-n \KL(q \| p)}$, with $e^{-\infty} := 0$.
\end{lemma}

\begin{proof}
Write $N(x) := n q(x) \in \N$, so that $\sum_x N(x) = n$. The event $\{\phat_n = q\}$ is the
set $T(q)$ of sequences $(x_1, \dots, x_n) \in \mathcal X^n$ in which each $x$ appears exactly
$N(x)$ times. By independence, each such sequence has probability
$\prod_{k \le n} p(x_k) = \prod_x p(x)^{N(x)}$, so $\Prob(\phat_n = q) = \abs{T(q)} \prod_x p(x)^{N(x)}$.
If $q$ is not dominated by $p$, some $x$ has $N(x) > 0 = p(x)$ and the product vanishes: the
claim is $0 \le 0$. Otherwise $p(x) > 0$ whenever $N(x) > 0$, and
\[
  \prod_x p(x)^{N(x)} = \exp\Big(\sum_{x : q(x) > 0} n\, q(x) \log p(x)\Big)
  = \exp\Big(-n \KL(q \| p) - n H(q)\Big), \qquad H(q) := -\sum_{x : q(x) > 0} q(x) \log q(x),
\]
because $q(x) \log p(x) = -q(x) \log(q(x)/p(x)) - q(x)\log(1/q(x))$ on $\{q > 0\}$. It remains to
show $\abs{T(q)} \le e^{n H(q)}$: under the product law $q^{\otimes n}$ on $\mathcal X^n$, each
sequence of $T(q)$ has probability $\prod_x q(x)^{N(x)} = e^{-n H(q)}$ by the same computation
(with $p := q$, for which $\KL(q \| q) = 0$), hence
$1 \ge q^{\otimes n}(T(q)) = \abs{T(q)}\, e^{-n H(q)}$. This is Theorem~11.1.4 of
\cite{cover2006elements}.
\end{proof}

\begin{lemma}[Number of types]\label{lem:num_types}
The number of types of denominator $n$ on a set with $m \ge 1$ elements is
$\binom{n + m - 1}{m - 1}$, and $\binom{n + m - 1}{m - 1} \le (n + 1)^{m - 1}$.
\end{lemma}

\begin{proof}
Types of denominator $n$ correspond bijectively to families $(N(x))_{x \in \mathcal X}$ of
nonnegative integers with $\sum_x N(x) = n$ (via $N(x) = n q(x)$), i.e.\ to multisets of size
$n$ over $\mathcal X$, whose number is the multichoose number $\binom{n + m - 1}{n} = \binom{n + m - 1}{m - 1}$
(stars and bars; Mathlib \texttt{Sym.card\_sym\_eq\_multichoose} and
\texttt{Nat.multichoose\_eq}). For the bound, $\binom{n + m - 1}{m - 1} = \prod_{j = 1}^{m - 1} \frac{n + j}{j}$
and each factor satisfies $\frac{n + j}{j} = 1 + \frac nj \le 1 + n$ for $j \ge 1$; the product of
$m - 1$ factors, each in $[1, n + 1]$, is at most $(n + 1)^{m - 1}$. (Induction on $m$:
$\binom{n + m}{m} = \binom{n + m - 1}{m - 1} \cdot \frac{n + m}{m}$, Mathlib
\texttt{Nat.succ\_mul\_choose\_eq}.)
\end{proof}

\begin{theorem}[Lemma 19: Sanov's high-probability bound]\label{thm:sanov_hp}
\uses{lem:type_prob_le, lem:num_types}
For every $x \in \R$, $\Prob(\KL(\phat_n \| p) > x) \le (n + 1)^{m - 1} e^{-n x}$. Consequently,
for every $\delta \in (0, 1)$, with probability at least $1 - \delta$,
\[
  \KL(\phat_n \| p) \le \frac{(m - 1)\log(n + 1) + \log(1/\delta)}{n} .
\]
\end{theorem}

\begin{proof}
\uses{lem:type_prob_le, lem:num_types}
Since $\phat_n$ is a type of denominator $n$, the event $\{\KL(\phat_n \| p) > x\}$ is the
disjoint union, over the types $q$ of denominator $n$ with $\KL(q \| p) > x$, of the events
$\{\phat_n = q\}$. By Lemma~\ref{lem:type_prob_le} each has probability at most
$e^{-n \KL(q \| p)} \le e^{-nx}$ (also when $\KL(q\|p) = \infty$), and by Lemma~\ref{lem:num_types}
there are at most $(n + 1)^{m - 1}$ of them; the union bound gives the first claim. For the
second, take $x := ((m - 1)\log(n + 1) + \log(1/\delta))/n$, for which
$(n + 1)^{m - 1} e^{-nx} = (n + 1)^{m - 1} (n + 1)^{-(m - 1)} \delta = \delta$.
This is Theorem~11.4.1 of \cite{cover2006elements} applied to the set
$\{q : \KL(q \| p) > x\}$.
\end{proof}

\begin{remark}[The paper's proof, and where the bound is used]\label{rem:pconc_sanov_paper}
The proof of Lemma~19 in \cite{essakine2026tight} writes $(n + 1)^m$ in the intermediate
bound and $(m - 1)$ in the conclusion (and $D(p \| q)$ for $D(q \| p)$ in the definition of the
set); the version above with $(n + 1)^{m - 1}$ is the correct one (the type class has
$\binom{n + m - 1}{m - 1} \le (n + 1)^{m - 1}$ elements). The theorem is a fixed-$n$ statement.
The KL event of Chapter~\ref{chap:empirical} needs a bound uniform in $n$, which cannot be
obtained from this theorem by a union bound over $n$ with the paper's rate
(Remark~\ref{rem:empirical_sanov_union}); the time-uniform bound
(Lemma~\ref{lem:empirical_kl_time_uniform}) is proved there with Lemmas~\ref{lem:ville}
and~\ref{lem:num_types}, and Theorem~\ref{thm:sanov_hp} is only referenced in remarks.
\end{remark}

\textbf{Lean remark.} The sample is \texttt{X : Fin n → Ω → Fin m} with
\texttt{ProbabilityTheory.iIndepFun} and \texttt{HasLaw (X k) (weightedMeasure p)}; the empirical
distribution is the \texttt{weightedMeasure} of the vector of frequencies. The proof of
Lemma~\ref{lem:type_prob_le} needs the law of the whole vector, \texttt{(fun ω k => X k ω)},
to be the product measure (\texttt{iIndepFun\_iff\_map\_fun\_eq\_pi\_map}), and the statement of
Theorem~\ref{thm:sanov_hp} compares \texttt{klDiv} (in \texttt{ℝ≥0∞}) with a real threshold
through \texttt{ENNReal.ofReal}.

\section{Ville's inequality and the Bernoulli maximal inequality}
\label{sec:pconc_ville}

Throughout this section and the next, $(\Omega, \mathcal F, \Prob)$ is a probability space
with a filtration $(\mathcal F_t)_{t \ge 0}$ (Mathlib \texttt{MeasureTheory.Filtration ℕ}),
and conditional expectations are \texttt{MeasureTheory.condExp}; all (in)equalities between
random variables hold almost surely.

\begin{lemma}[One-step bound for a Bernoulli variable]\label{lem:bernoulli_mgf_step}
Let $\mathcal G \subseteq \mathcal F$ be a sub-$\sigma$-algebra, $X$ a random variable with
values in $\{0, 1\}$ and $P$ a $\mathcal G$-measurable random variable with values in $[0, 1]$
such that $\E[X \mid \mathcal G] = P$ (i.e.\ $P = \Prob(X = 1 \mid \mathcal G)$). Then
$\E[e^{P/2 - X} \mid \mathcal G] \le 1$.
\end{lemma}

\begin{proof}
Since $X \in \{0, 1\}$, $e^{-X} = 1 - (1 - e^{-1}) X$ pointwise, so by linearity
$\E[e^{-X} \mid \mathcal G] = 1 - (1 - e^{-1}) P$. As $e^{P/2}$ is bounded and
$\mathcal G$-measurable, $\E[e^{P/2 - X} \mid \mathcal G] = e^{P/2}\big(1 - (1 - e^{-1})P\big)$
(\texttt{condExp\_mul\_of\_stronglyMeasurable\_left}). Now $e \ge 2$ gives $1 - e^{-1} \ge 1/2$,
so $1 - (1 - e^{-1})P \le 1 - P/2 \le e^{-P/2}$ (using $P \ge 0$ and $1 + y \le e^y$ with
$y = -P/2$, \texttt{Real.add\_one\_le\_exp}). Multiplying by $e^{P/2} > 0$ gives the claim.
\end{proof}

\begin{lemma}[Ville's maximal inequality]\label{lem:ville}
Let $(M_t)_{t \ge 0}$ be a nonnegative supermartingale with respect to $(\mathcal F_t)$ with
$\E[M_0] \le 1$. Then for every $c > 0$, $\Prob(\exists t \ge 0,\ M_t \ge c) \le 1/c$; in
particular $\Prob(\exists t,\ M_t \ge 1/\delta) \le \delta$ for $\delta \in (0, 1]$.
\end{lemma}

\begin{proof}
Let $\tau := \inf\{t : M_t \ge c\}$ (with $\inf \emptyset = \infty$), a stopping time of
$(\mathcal F_t)$ (\texttt{hitting} of the closed set $[c, \infty)$,
\texttt{MeasureTheory.hitting\_isStoppingTime}). Fix $N \ge 0$ and let $\tau_N := \tau \wedge N$,
a bounded stopping time. The optional stopping theorem for supermartingales, in the form
$\E[M_{\tau_N}] \le \E[M_0]$ (Mathlib \texttt{Submartingale.expected\_stoppedValue\_mono}
applied to the submartingale $-M$ with the stopping times $0 \le \tau_N \le N$), gives
$\E[M_{\tau_N}] \le \E[M_0] \le 1$. On $\{\tau \le N\}$ we have $\tau_N = \tau$ and
$M_{\tau_N} = M_\tau \ge c$, and $M_{\tau_N} \ge 0$ everywhere, so
$c\, \Prob(\tau \le N) \le \E[M_{\tau_N} \indic\{\tau \le N\}] \le \E[M_{\tau_N}] \le 1$.
The events $\{\tau \le N\}$ increase to $\{\tau < \infty\} = \{\exists t, M_t \ge c\}$, so
$\Prob(\exists t, M_t \ge c) = \lim_N \Prob(\tau \le N) \le 1/c$
(\texttt{MeasureTheory.tendsto\_measure\_iUnion\_atTop}). The second claim is $c = 1/\delta$.
\end{proof}

\textbf{Lean remark.} Mathlib's \texttt{MeasureTheory.maximal\_ineq} is Doob's inequality
for nonnegative \emph{sub}martingales and does not apply directly; the proof above uses the
stopped process. Alternatively, \texttt{Supermartingale.stoppedProcess} (the stopped
supermartingale is a supermartingale) with \texttt{stoppedValue} at $N$ gives the same
$\E[M_{\tau \wedge N}] \le \E[M_0]$. The lemma is used in the form
$\Prob(\sup_t M_t \ge c) \le \E[M_0]/c$ by Lemma~\ref{lem:empirical_kl_time_uniform},
Theorem~\ref{thm:bernoulli_concentration} and Theorem~\ref{thm:bernstein_self_normalized}.

\begin{theorem}[Lemma 20: Bernoulli maximal inequality]\label{thm:bernoulli_concentration}
\uses{lem:bernoulli_mgf_step, lem:ville}
Let $(X_i)_{i \ge 1}$ be random variables with values in $\{0, 1\}$, $X_i$ being
$\mathcal F_i$-measurable, and $(P_i)_{i \ge 1}$ random variables with values in $[0, 1]$, $P_i$
being $\mathcal F_{i - 1}$-measurable, with $\E[X_i \mid \mathcal F_{i-1}] = P_i$ for every
$i \ge 1$. Then for every $\delta \in (0, 1)$,
\[
  \Prob\Big( \exists n \ge 1,\ \sum_{i \le n} X_i < \frac12 \sum_{i \le n} P_i - \log\frac1\delta \Big) \le \delta .
\]
\end{theorem}

\begin{proof}
\uses{lem:bernoulli_mgf_step, lem:ville}
Let $M_n := \exp\big(\sum_{i \le n} (P_i/2 - X_i)\big)$, $M_0 := 1$. Each $M_n$ is
$\mathcal F_n$-measurable, positive, and bounded by $e^{n/2}$, hence integrable. For $n \ge 1$,
$M_n = M_{n-1}\, e^{P_n/2 - X_n}$ with $M_{n-1}$ bounded and $\mathcal F_{n-1}$-measurable, so
\[
  \E[M_n \mid \mathcal F_{n-1}] = M_{n-1}\, \E[e^{P_n/2 - X_n} \mid \mathcal F_{n-1}] \le M_{n-1}
\]
by Lemma~\ref{lem:bernoulli_mgf_step} (with $\mathcal G = \mathcal F_{n-1}$, $X = X_n$,
$P = P_n$) and $M_{n-1} \ge 0$: $(M_n)$ is a nonnegative supermartingale with $M_0 = 1$.
Lemma~\ref{lem:ville} gives $\Prob(\exists n, M_n \ge 1/\delta) \le \delta$. Finally, if
$\sum_{i \le n} X_i < \frac12 \sum_{i \le n} P_i - \log(1/\delta)$ then
$\sum_{i \le n}(P_i/2 - X_i) > \log(1/\delta)$, i.e.\ $M_n > 1/\delta$; so the event of the
statement is contained in $\{\exists n, M_n \ge 1/\delta\}$.
\end{proof}

\textbf{Lean remark.} This is Lemma~F.4 of \cite{dann2017unifying}. The hypothesis
$\E[X_i \mid \mathcal F_{i-1}] = P_i$ is stated with \texttt{condExp}; in
Chapter~\ref{chap:empirical} it is supplied by a conditional-law statement
(\texttt{HasCondDistrib}, Lemma~\ref{lem:pseudo_counts_increments}), from which the
conditional expectation of the indicator follows. The process $M$ is
\texttt{fun n ω => Real.exp (∑ i ∈ Finset.range n, (P (i+1) ω / 2 - X (i+1) ω))}; its
supermartingale property is \texttt{supermartingale\_of\_setIntegral\_succ\_le} or the
direct \texttt{condExp} computation above.

\section{The self-normalized Bernstein inequality}
\label{sec:pconc_bernstein}

Let $h(x) := (x + 1)\log(x + 1) - x$ for $x \ge 0$ (the Cramér transform of the Poisson
distribution of parameter $1$) and, for $b > 0$, $\varphi_b(\lambda) := e^{\lambda b} - 1 - \lambda b$
for $\lambda \in \R$.

\begin{lemma}[Elementary properties of the Poisson Cramér transform]\label{lem:poisson_cramer}
\begin{enumerate}
\item[(i)] For $x \ge 0$: $h(x) \ge \dfrac{x^2}{2(1 + x/3)} \ge 0$. For $\lambda \ge 0$ and $b > 0$:
  $\varphi_b(\lambda) \ge 0$.
\item[(ii)] For $c > 0$, $s \ge 0$, $b > 0$ and $\hat\lambda := \frac1b \log(1 + s/c) \ge 0$:
  $\hat\lambda\, b s - c\, \varphi_b(\hat\lambda) = c\, h(s/c)$.
\item[(iii)] If $c > 0$, $s \ge 0$, $\ell \ge 0$ and $c\, h(s/c) \le \ell$, then
  $s \le \sqrt{2 c \ell} + \frac23 \ell$.
\end{enumerate}
\end{lemma}

\begin{proof}
\uses{lem:sqrt_add_le}
(i) $\varphi_b(\lambda) \ge 0$ is $1 + u \le e^u$ with $u = \lambda b$ (\texttt{Real.add\_one\_le\_exp}).
For $h$, let $\psi(x) := h(x) - \frac{3x^2}{2(3 + x)}$ on $[0, \infty)$. Then $\psi(0) = 0$,
$\psi'(x) = \log(1 + x) - \frac{3x(6 + x)}{2(3 + x)^2}$ (using $h'(x) = \log(1 + x)$), $\psi'(0) = 0$,
and $\psi''(x) = \frac1{1 + x} - \frac{27}{(3 + x)^3} \ge 0$ because
$(3 + x)^3 = 27 + 27x + 9x^2 + x^3 \ge 27(1 + x)$. Hence $\psi'$ is nondecreasing, so
$\psi' \ge \psi'(0) = 0$, so $\psi$ is nondecreasing and $\psi \ge \psi(0) = 0$
(\texttt{monotoneOn\_of\_deriv\_nonneg} twice). Finally $\frac{x^2}{2(1 + x/3)} = \frac{3x^2}{2(3 + x)} \ge 0$.

(ii) Put $x := s/c \ge 0$, so $e^{\hat\lambda b} = 1 + x$ and $\varphi_b(\hat\lambda) = x - \log(1 + x)$.
Then $\hat\lambda b s - c\varphi_b(\hat\lambda) = c\,[x \log(1 + x) - x + \log(1 + x)] = c\,[(1 + x)\log(1 + x) - x] = c\, h(x)$.
(This value is the maximum of $\lambda \mapsto \lambda b s - c\varphi_b(\lambda)$ over $\lambda \ge 0$,
which is not needed below.)

(iii) By (i) with $x = s/c$, $\ell \ge c\, h(s/c) \ge \frac{c\, (s/c)^2}{2(1 + s/(3c))} = \frac{s^2}{2c + 2s/3}$,
so $s^2 \le 2c\ell + \frac{2\ell}{3} s$, i.e.\ $(s - \ell/3)^2 \le \ell^2/9 + 2c\ell$. Hence
$s \le \ell/3 + \sqrt{\ell^2/9 + 2c\ell} \le \ell/3 + \ell/3 + \sqrt{2c\ell}$ by
Lemma~\ref{lem:sqrt_add_le} (\texttt{Real.sqrt\_le\_sqrt}, \texttt{abs\_le\_of\_sq\_le\_sq'}).
\end{proof}

\begin{lemma}[A Taylor bound for the exponential of a bounded variable]\label{lem:pconc_exp_taylor_bound}
For $b > 0$, $\lambda \ge 0$ and $y \in [-b, b]$:
$e^{\lambda y} \le 1 + \lambda y + \dfrac{y^2}{b^2}\, \varphi_b(\lambda)$.
\end{lemma}

\begin{proof}
By the power series of the exponential (\texttt{Real.exp\_eq\_exp\_ℝ},
\texttt{NormedSpace.exp\_eq\_tsum\_div}), $e^{\lambda y} - 1 - \lambda y = \sum_{k \ge 2} (\lambda y)^k/k!$
and $\varphi_b(\lambda) = \sum_{k \ge 2} (\lambda b)^k/k!$, both series being absolutely convergent.
For $k \ge 2$, since $\lambda \ge 0$ and $\abs y \le b$,
$(\lambda y)^k \le \abs{\lambda y}^k = \lambda^k y^2 \abs y^{k - 2} \le \lambda^k y^2 b^{k - 2} = \frac{y^2}{b^2}(\lambda b)^k$.
Summing (\texttt{Summable.tsum\_le\_tsum}) gives
$e^{\lambda y} - 1 - \lambda y \le \frac{y^2}{b^2} \sum_{k \ge 2} (\lambda b)^k/k! = \frac{y^2}{b^2}\varphi_b(\lambda)$.
\end{proof}

\begin{lemma}[The exponential supermartingale]\label{lem:bernstein_exp_supermartingale}
\uses{lem:pconc_exp_taylor_bound}
Let $b > 0$, $(Y_t)_{t \ge 1}$ random variables with $Y_t$ $\mathcal F_t$-measurable,
$\abs{Y_t} \le b$ and $\E[Y_t \mid \mathcal F_{t-1}] = 0$, and $(w_t)_{t \ge 1}$ random variables
with $w_t$ $\mathcal F_{t-1}$-measurable and $w_t \in [0, 1]$. Let $\sigma_t^2 := \E[Y_t^2 \mid \mathcal F_{t-1}]$
(an $\mathcal F_{t-1}$-measurable version with values in $[0, b^2]$) and, for $t \ge 0$,
\[
  S_t := \sum_{s \le t} w_s Y_s, \qquad V_t := \sum_{s \le t} w_s^2 \sigma_s^2 \qquad (S_0 = V_0 = 0),
\]
so that $\abs{S_t} \le tb$ and $0 \le V_t \le tb^2$. Then for every $\lambda \ge 0$:
\begin{enumerate}
\item[(i)] $\E\big[\exp\big(\lambda w_t Y_t - \varphi_b(\lambda) w_t^2 \sigma_t^2/b^2\big) \mid \mathcal F_{t-1}\big] \le 1$ for $t \ge 1$;
\item[(ii)] $M^\lambda_t := \exp\big(\lambda S_t - \varphi_b(\lambda) V_t/b^2\big)$ is a nonnegative
  supermartingale with respect to $(\mathcal F_t)$, with $M^\lambda_0 = 1$ and $M^\lambda_t \le e^{\lambda b t}$;
\item[(iii)] the same holds for $-Y$ in place of $Y$, i.e.\ for $\exp(-\lambda S_t - \varphi_b(\lambda) V_t/b^2)$.
\end{enumerate}
\end{lemma}

\begin{proof}
\uses{lem:pconc_exp_taylor_bound}
The bounds $\abs{S_t} \le tb$ and $V_t \le tb^2$ follow from $w_s \in [0, 1]$, $\abs{Y_s} \le b$
and $\sigma_s^2 \le b^2$ (conditional expectation of a variable in $[0, b^2]$).

(i) Since $\abs{w_t Y_t} \le b$, Lemma~\ref{lem:pconc_exp_taylor_bound} with $y = w_t Y_t$ gives
$e^{\lambda w_t Y_t} \le 1 + \lambda w_t Y_t + \varphi_b(\lambda)\, w_t^2 Y_t^2/b^2$ pointwise. All terms are
bounded, and $w_t$ is $\mathcal F_{t-1}$-measurable, so taking conditional expectations
(monotonicity, linearity and \texttt{condExp\_mul\_of\_stronglyMeasurable\_left}):
\[
  \E[e^{\lambda w_t Y_t} \mid \mathcal F_{t-1}] \le 1 + \lambda w_t\, \E[Y_t \mid \mathcal F_{t-1}] + \varphi_b(\lambda) \frac{w_t^2 \sigma_t^2}{b^2}
  = 1 + \varphi_b(\lambda) \frac{w_t^2 \sigma_t^2}{b^2} \le \exp\Big(\varphi_b(\lambda) \frac{w_t^2 \sigma_t^2}{b^2}\Big)
\]
by $1 + u \le e^u$. The factor $\exp(-\varphi_b(\lambda) w_t^2 \sigma_t^2/b^2)$ is
$\mathcal F_{t-1}$-measurable with values in $(0, 1]$ ($\varphi_b(\lambda) \ge 0$ by
Lemma~\ref{lem:poisson_cramer}(i)); pulling it out gives (i).

(ii) $M^\lambda_t$ is $\mathcal F_t$-measurable, positive, and
$M^\lambda_t \le e^{\lambda \abs{S_t}} \le e^{\lambda b t}$ since $\varphi_b(\lambda) V_t \ge 0$; so it is
integrable, and $M^\lambda_0 = e^0 = 1$. For $t \ge 1$,
$M^\lambda_t = M^\lambda_{t-1} \exp(\lambda w_t Y_t - \varphi_b(\lambda) w_t^2 \sigma_t^2/b^2)$ with
$M^\lambda_{t-1}$ bounded, nonnegative and $\mathcal F_{t-1}$-measurable, so
$\E[M^\lambda_t \mid \mathcal F_{t-1}] = M^\lambda_{t-1}\, \E[\exp(\cdots) \mid \mathcal F_{t-1}] \le M^\lambda_{t-1}$ by (i).

(iii) $-Y_t$ satisfies the same hypotheses ($\abs{-Y_t} \le b$, $\E[-Y_t \mid \mathcal F_{t-1}] = 0$,
$(-Y_t)^2 = Y_t^2$), and its sums are $-S_t$ and $V_t$.
\end{proof}

\textbf{Lean remark.} The hypothesis $\lambda \ge 0$ is necessary: for $\lambda < 0$ the bound
$\E[e^{\lambda Y}] \le \exp(\E[Y^2]\varphi_b(\lambda)/b^2)$ fails (take $Y = \pm b$ with probability
$1/2$ each and $\lambda b = -2$: $\cosh 2 \approx 3.76 > \exp(e^{-2} + 1) \approx 3.11$); the
outline's ``$\lambda \in \R$'' must be read with $\varphi_b(\abs\lambda)$, which is (iii). In Lean,
$\sigma_t^2$ is \texttt{μ[Y t \^{} 2 | ℱ (t - 1)]} (a fixed version), $S$ and $V$ are
\texttt{Finset.sum} over \texttt{Finset.Icc 1 t}, and the supermartingale property is
proved by \texttt{supermartingale\_of\_setIntegral\_succ\_le} or by the \texttt{condExp}
computation above; \texttt{ProbabilityTheory.variance} does not appear (the conditional
variance is the conditional expectation of the square of a centered variable).

The paper's Lemma~21 is proved in \cite{menard2021fast} by a peeling argument over
$\lambda$. We use instead the method of mixtures: the exponential supermartingales of
Lemma~\ref{lem:bernstein_exp_supermartingale} are integrated against the exponential prior
$e^{-\mu}\,d\mu$ on $\mu = \lambda b \in [0, \infty)$, and a Laplace-type lower bound on the
mixture recovers the Cramér transform up to a factor $8(1 + \sqrt{2t})$, which is smaller than the
paper's $4e(2t + 1)$. Write $\varphi := \varphi_1$, i.e.\ $\varphi(\mu) = e^\mu - 1 - \mu$.

\begin{lemma}[Lower bound on the exponential mixture]\label{lem:pconc_mixture_lower}
\uses{lem:poisson_cramer}
For all $s \ge 0$ and $v \ge 0$, with $T := (v + 1)\, h\big(s/(v + 1)\big)$,
\[
  \int_0^\infty \exp\big(\mu s - \varphi(\mu)\, v - \mu\big)\, d\mu \ \ge\ \frac{e^{T}}{4\,(1 + \sqrt{s + v})}
\]
(the left-hand side is the Lebesgue integral of a nonnegative continuous function, possibly $+\infty$).
\end{lemma}

\begin{proof}
\uses{lem:poisson_cramer}
If $s = v = 0$ then $T = h(0) = 0$ and the integral is $\int_0^\infty e^{-\mu}\,d\mu = 1 \ge 1/4$.
Otherwise $s + v > 0$; put $x := s/(v + 1) \ge 0$, $\hat\mu := \log(1 + x) \ge 0$,
$\theta_0 := 1/\sqrt{s + v} > 0$, $\ell := \log(1 + \theta_0) > 0$ and
$g(\mu) := \mu s - \varphi(\mu) v$. By Lemma~\ref{lem:poisson_cramer}(ii) (with $b = 1$, $c = v + 1$),
$\hat\mu s - (v + 1)\varphi(\hat\mu) = T$, hence $g(\hat\mu) = T + \varphi(\hat\mu)$ with
$\varphi(\hat\mu) = x - \log(1 + x)$.

\emph{Lower bound on $g$ near $\hat\mu$.} Let $\mu \in [\hat\mu, \hat\mu + \ell]$ and
$\theta := e^{\mu - \hat\mu} - 1 \in [0, \theta_0]$, so that $\mu - \hat\mu = \log(1 + \theta) \ge \theta/(1 + \theta) \ge \theta - \theta^2$
(\texttt{Real.one\_sub\_inv\_le\_log\_of\_pos}). Since $e^{\hat\mu} = 1 + x$,
$\varphi(\mu) - \varphi(\hat\mu) = e^{\hat\mu}(e^{\mu - \hat\mu} - 1) - (\mu - \hat\mu) = (1 + x)\theta - (\mu - \hat\mu)$, and therefore
\[
  g(\mu) - g(\hat\mu) = (\mu - \hat\mu)(s + v) - v(1 + x)\theta
  \ge (\theta - \theta^2)(s + v) - v(1 + x)\theta = \theta\, x - (s + v)\theta^2 \ge -(s + v)\theta_0^2 = -1 ,
\]
using $(s + v) - v(1 + x) = s - vx = s/(v + 1) = x \ge 0$.

\emph{Integration.} The integrand is nonnegative, so the integral over $[0, \infty)$ is at least the
integral over $[\hat\mu, \hat\mu + \ell]$, on which $e^{g(\mu)} \ge e^{T + \varphi(\hat\mu) - 1}$:
\[
  \int_0^\infty e^{g(\mu) - \mu}\, d\mu \ge e^{T + \varphi(\hat\mu) - 1} \int_{\hat\mu}^{\hat\mu + \ell} e^{-\mu}\, d\mu
  = e^{T - 1}\, \frac{e^{\varphi(\hat\mu)}}{1 + x}\, \big(1 - e^{-\ell}\big)
  = e^{T - 1}\, \frac{e^{x}}{(1 + x)^2}\, \frac{\theta_0}{1 + \theta_0} ,
\]
since $e^{-\hat\mu} = 1/(1 + x)$, $e^{\varphi(\hat\mu)} = e^x/(1 + x)$ and $1 - e^{-\ell} = \theta_0/(1 + \theta_0)$.
Finally $e^{x}/(1 + x)^2 \ge e/4$ for $x \ge 0$: indeed $e^{(x - 1)/2} \ge 1 + (x - 1)/2 = (1 + x)/2 \ge 0$
(\texttt{Real.add\_one\_le\_exp}), and squaring gives $e^{x - 1} \ge (1 + x)^2/4$
(\texttt{pow\_le\_pow\_left}); and $\theta_0/(1 + \theta_0) = 1/(1 + \sqrt{s + v})$. Hence the
integral is at least $e^{T - 1} \cdot (e/4) \cdot 1/(1 + \sqrt{s + v}) = e^T/(4(1 + \sqrt{s + v}))$.
\end{proof}

\begin{lemma}[The mixture supermartingale]\label{lem:pconc_bernstein_mixture}
\uses{lem:bernstein_exp_supermartingale}
In the setting of Lemma~\ref{lem:bernstein_exp_supermartingale}, define for $t \ge 0$
\[
  M_t := \frac12 \int_0^\infty \Big[ \exp\Big(\frac{\mu S_t}{b} - \varphi(\mu)\frac{V_t}{b^2}\Big) + \exp\Big(-\frac{\mu S_t}{b} - \varphi(\mu)\frac{V_t}{b^2}\Big) \Big] e^{-\mu}\, d\mu .
\]
Then $M_t$ is finite almost surely (we set $M_t := 0$ where it is not), $M_0 = 1$, $(M_t)_{t \ge 0}$ is
a nonnegative supermartingale with respect to $(\mathcal F_t)$, and almost surely, for every $t \ge 0$,
\[
  M_t \ \ge\ \frac{e^{T_t}}{8\,(1 + \sqrt{2t})}, \qquad
  T_t := \Big(\frac{V_t}{b^2} + 1\Big)\, h\Big(\frac{b\, \abs{S_t}}{V_t + b^2}\Big) .
\]
\end{lemma}

\begin{proof}
\uses{lem:bernstein_exp_supermartingale, lem:pconc_mixture_lower}
For $\mu \ge 0$ let $\lambda := \mu/b \ge 0$, so that $\mu S_t/b - \varphi(\mu)V_t/b^2 = \lambda S_t - \varphi_b(\lambda)V_t/b^2$
and the two exponentials are the processes $M^{\lambda}_t$ and their $-Y$ versions of
Lemma~\ref{lem:bernstein_exp_supermartingale}(ii)--(iii), which are nonnegative supermartingales
with value $1$ at $t = 0$ and expectation at most $1$. The integrand is jointly measurable in
$(\omega, \mu)$ (measurable in $\omega$, continuous in $\mu$).

\emph{Finiteness and integrability.} By Tonelli,
$\E[M_t] = \frac12\int_0^\infty (\E[M^{\mu/b}_t] + \E[M^{-,\mu/b}_t]) e^{-\mu}\,d\mu \le \int_0^\infty e^{-\mu}\, d\mu = 1$,
so $M_t < \infty$ almost surely and $M_t$ is integrable; $M_t$ is $\mathcal F_t$-measurable, and
$M_0 = \frac12\int_0^\infty 2 e^{-\mu}\,d\mu = 1$.

\emph{Supermartingale.} For $t \ge 1$ and $A \in \mathcal F_{t-1}$, Tonelli and the supermartingale
property of each $M^{\pm, \lambda}$ in set-integral form ($\int_A M^{\pm,\lambda}_t\, d\Prob \le \int_A M^{\pm,\lambda}_{t-1}\, d\Prob$,
from $\E[M^{\pm,\lambda}_t \mid \mathcal F_{t-1}] \le M^{\pm,\lambda}_{t-1}$) give
$\int_A M_t\, d\Prob = \frac12 \int_0^\infty e^{-\mu} \int_A (M^{\mu/b}_t + M^{-,\mu/b}_t)\, d\Prob\, d\mu \le \int_A M_{t-1}\, d\Prob$.
Since $M_t$ and $M_{t-1}$ are integrable and $M_{t-1}$ is $\mathcal F_{t-1}$-measurable, this is
$\E[M_t \mid \mathcal F_{t-1}] \le M_{t-1}$ (\texttt{supermartingale\_of\_setIntegral\_succ\_le}).

\emph{Lower bound.} Fix $\omega$ where $M_t(\omega) < \infty$ and put $s := \abs{S_t}/b \in [0, t]$,
$v := V_t/b^2 \in [0, t]$ (Lemma~\ref{lem:bernstein_exp_supermartingale}). One of the two
exponentials equals $\exp(\mu s - \varphi(\mu) v)$ (the one whose sign matches $S_t$) and the other
is nonnegative, so $M_t \ge \frac12 \int_0^\infty \exp(\mu s - \varphi(\mu) v - \mu)\, d\mu \ge \frac12 \cdot \frac{e^{T}}{4(1 + \sqrt{s + v})}$
by Lemma~\ref{lem:pconc_mixture_lower}, where $T = (v + 1)h(s/(v + 1)) = T_t$ because
$s/(v + 1) = (\abs{S_t}/b)/((V_t + b^2)/b^2) = b\abs{S_t}/(V_t + b^2)$. Finally $s + v \le 2t$ gives
$1 + \sqrt{s + v} \le 1 + \sqrt{2t}$.
\end{proof}

\begin{theorem}[Lemma 21: self-normalized Bernstein inequality]\label{thm:bernstein_self_normalized}
\uses{lem:bernstein_exp_supermartingale}
Let $b > 0$, $(Y_t)_{t \ge 1}$, $(w_t)_{t \ge 1}$, $S_t$ and $V_t$ be as in
Lemma~\ref{lem:bernstein_exp_supermartingale} ($Y_t$ adapted, $\abs{Y_t} \le b$, centered
conditionally on $\mathcal F_{t-1}$; $w_t$ predictable with values in $[0, 1]$). Then for every
$\delta \in (0, 1)$,
\[
  \Prob\Big( \exists t \ge 1,\ \Big(\frac{V_t}{b^2} + 1\Big) h\Big(\frac{b\abs{S_t}}{V_t + b^2}\Big) \ge \log\frac1\delta + \log\big(8(1 + \sqrt{2t})\big) \Big) \le \delta ,
\]
and, since $8(1 + \sqrt{2t}) \le 4e(2t + 1)$ for $t \ge 1$, also
\[
  \Prob\Big( \exists t \ge 1,\ \Big(\frac{V_t}{b^2} + 1\Big) h\Big(\frac{b\abs{S_t}}{V_t + b^2}\Big) \ge \log\frac1\delta + \log\big(4e(2t + 1)\big) \Big) \le \delta .
\]
\end{theorem}

\begin{proof}
\uses{lem:pconc_bernstein_mixture, lem:ville}
Let $M$ be the mixture supermartingale of Lemma~\ref{lem:pconc_bernstein_mixture}: it is
nonnegative with $M_0 = 1$, so Lemma~\ref{lem:ville} gives $\Prob(\exists t, M_t \ge 1/\delta) \le \delta$.
Almost surely, for every $t \ge 1$, $M_t \ge e^{T_t}/(8(1 + \sqrt{2t}))$; hence if
$T_t \ge \log(1/\delta) + \log(8(1 + \sqrt{2t}))$ then $e^{T_t} \ge 8(1 + \sqrt{2t})/\delta$ and
$M_t \ge 1/\delta$. So the event of the first claim is contained, up to a null set, in
$\{\exists t, M_t \ge 1/\delta\}$. The second claim follows from the first because
$\sqrt{2t} \le 2t$ for $t \ge 1/2$ (\texttt{Real.sqrt\_le\_self}-type: $2t \ge 1$) and
$8 \le 4e$ ($e \ge 2$), so $8(1 + \sqrt{2t}) \le 4e(2t + 1)$ and the second event is contained in
the first.
\end{proof}

\begin{corollary}[Explicit form of the self-normalized Bernstein inequality]\label{cor:bernstein_explicit}
\uses{thm:bernstein_self_normalized}
In the setting of Theorem~\ref{thm:bernstein_self_normalized}, for every $\delta \in (0, 1)$,
with probability at least $1 - \delta$, for all $t \ge 1$,
\[
  \abs{S_t} \le \sqrt{2 V_t L_t} + 3 b L_t, \qquad L_t := \log\frac{4e(2t + 1)}{\delta} ,
\]
and in fact $\abs{S_t} \le \sqrt{2(V_t + b^2) L_t} + \frac23 b L_t$.
\end{corollary}

\begin{proof}
\uses{thm:bernstein_self_normalized, lem:poisson_cramer, lem:sqrt_add_le}
On the complement of the event of Theorem~\ref{thm:bernstein_self_normalized} (second form), which
has probability at least $1 - \delta$, for every $t \ge 1$: $c\, h(s/c) < L_t$ with
$c := V_t/b^2 + 1 > 0$ and $s := \abs{S_t}/b \ge 0$, and $L_t \ge 0$. Lemma~\ref{lem:poisson_cramer}(iii)
gives $s \le \sqrt{2cL_t} + \frac23 L_t$, i.e., multiplying by $b$ and using
$b\sqrt{2(V_t/b^2 + 1)L_t} = \sqrt{2(V_t + b^2)L_t}$, the second claim
$\abs{S_t} \le \sqrt{2(V_t + b^2)L_t} + \frac23 bL_t$. By Lemma~\ref{lem:sqrt_add_le},
$\sqrt{2(V_t + b^2)L_t} \le \sqrt{2V_tL_t} + b\sqrt{2L_t}$. Since $\delta < 1$ and $t \ge 1$,
$L_t \ge \log(12e) > 2$, so $2L_t \le L_t^2$ and $\sqrt{2L_t} \le L_t$; therefore
$\abs{S_t} \le \sqrt{2V_tL_t} + bL_t + \frac23 bL_t = \sqrt{2V_tL_t} + \frac53 bL_t \le \sqrt{2V_tL_t} + 3bL_t$.
\end{proof}

\begin{remark}[Constants]\label{rem:pconc_bernstein_constants}
The paper (and \cite{menard2021fast}, Lemma~9, from the peeling proof of Domingues et al.)
states Lemma~21 with the term $\log(4e(2t + 1))$; the mixture proof above gives the smaller
$\log(8(1 + \sqrt{2t}))$, and the explicit form holds with $\frac53 b L_t$ in place of $3bL_t$. We
keep the paper's constants in the statements used downstream: Chapter~\ref{chap:empirical}
applies Corollary~\ref{cor:bernstein_explicit} with the process indexed by the number $n$ of
visits and bounds $L_n = \log(4e(2n + 1)/\delta') \le \log(8e(n + 1)) + \log(1/\delta')$, which
is exactly the rate $\alpha^*(n, \delta)$ of Definition~\ref{def:rates}; any threshold below
$\log(8e(n + 1)/\delta')$ is absorbed the same way. The theorem is stated for $\delta \in (0, 1)$;
for $\delta \ge 1$ it is trivial.
\end{remark}

\textbf{Lean remark.} The theorem is library material (\texttt{Essakine2026Tight/Mathlib/}),
with the filtration \texttt{ℱ : Filtration ℕ mΩ}, \texttt{Y w : ℕ → Ω → ℝ},
\texttt{Adapted ℱ Y}, \texttt{∀ t, StronglyMeasurable[ℱ (t-1)] (w (t))} (predictability),
\texttt{∀ t, ∀ᵐ ω, |Y t ω| ≤ b}, \texttt{∀ t, μ[Y t | ℱ (t-1)] =ᵐ[μ] 0}. The mixture is
\texttt{fun t ω => (1/2) * ∫ μ in Set.Ioi 0, (...) * Real.exp (-μ)}; its measurability uses
\texttt{StronglyMeasurable.integral\_prod\_right'} (or \texttt{Measurable.lintegral\_prod\_right}
after passing to \texttt{ENNReal.ofReal}), its integrability Tonelli
(\texttt{MeasureTheory.lintegral\_prod}, \texttt{integrable\_prod\_iff}), and the set-integral
inequality \texttt{MeasureTheory.integral\_prod} together with
\texttt{Supermartingale.setIntegral\_le}. The final union over $t$ is the countable union
inside Lemma~\ref{lem:ville}; the bounds $\abs{S_t} \le tb$, $V_t \le tb^2$ are
\texttt{Finset.sum\_le\_card\_nsmul}. The analytic Lemma~\ref{lem:pconc_mixture_lower} needs only
$\int_a^{a + \ell} e^{-\mu}\,d\mu = e^{-a}(1 - e^{-\ell})$
(\texttt{intervalIntegral.integral\_eq\_sub\_of\_hasDerivAt}) and the monotonicity of the
integral in the domain (\texttt{MeasureTheory.setIntegral\_mono\_set}, nonnegative integrand).

\section{KL--Bernstein and variance transport inequalities}
\label{sec:pconc_kl}

In this section $\mathcal X$ is a finite set, $p, q$ are probability vectors on $\mathcal X$ and
$f, g : \mathcal X \to \R$; the conventions are those of the beginning of the chapter.

\begin{lemma}[Gibbs variational inequality on a finite set]\label{lem:donsker_varadhan_finite}
For every $g : \mathcal X \to \R$: $\KL(p \| q) \ge pg - \log\big(q e^{g}\big)$, where
$q e^g := \sum_x q(x) e^{g(x)}$ (the inequality is trivial when $\KL(p\|q) = +\infty$).
\end{lemma}

\begin{proof}
Assume $p$ is dominated by $q$. Let $Z := q e^g > 0$ and $r(x) := q(x) e^{g(x)}/Z$, a probability
vector with $r(x) > 0$ whenever $q(x) > 0$, hence whenever $p(x) > 0$. For $x$ with $p(x) > 0$,
$\log(p(x)/q(x)) = \log(p(x)/r(x)) + g(x) - \log Z$, so summing against $p(x)$,
\[
  \KL(p \| q) = \sum_{x : p(x) > 0} p(x)\log\frac{p(x)}{r(x)} + pg - \log Z = \KL(p \| r) + pg - \log Z .
\]
It remains to see $\KL(p \| r) \ge 0$ (Gibbs' inequality): by $\log u \ge 1 - 1/u$ for $u > 0$
(\texttt{Real.one\_sub\_inv\_le\_log\_of\_pos}) with $u = p(x)/r(x)$,
$\sum_{x : p(x) > 0} p(x)\log\frac{p(x)}{r(x)} \ge \sum_{x : p(x) > 0} (p(x) - r(x)) = 1 - \sum_{x : p(x) > 0} r(x) \ge 0$.
\end{proof}

\textbf{Lean remark.} Mathlib's \texttt{InformationTheory.klDiv} is defined through the
log-likelihood ratio (\texttt{llr}) and \texttt{klFun}; nonnegativity for probability
measures is available (\texttt{klDiv\_eq\_zero\_iff}/\texttt{integral\_llr\_nonneg}-type
lemmas), and the identity above is the chain \texttt{toReal\_klDiv} for
\texttt{weightedMeasure}s (finite sums). The Donsker--Varadhan variational formula is not
needed in its full generality: only this one-sided inequality with the tilted vector $r$.

\begin{lemma}[Bernstein bound on the log-moment generating function]\label{lem:bernstein_mgf_bounded}
\uses{lem:pconc_exp_taylor_bound}
Let $b > 0$, $f$ with values in $[0, b]$, and $\lambda \ge 0$. Then
$\log\big(q\, e^{\lambda(f - qf)}\big) \le \Var_q(f)\, \varphi_b(\lambda)/b^2$,
where $\varphi_b(\lambda) = e^{\lambda b} - 1 - \lambda b$.
\end{lemma}

\begin{proof}
\uses{lem:pconc_exp_taylor_bound}
Since $f \in [0, b]$, $qf \in [0, b]$ and $y(x) := f(x) - qf \in [-b, b]$ for every $x$. By
Lemma~\ref{lem:pconc_exp_taylor_bound}, $e^{\lambda y(x)} \le 1 + \lambda y(x) + y(x)^2 \varphi_b(\lambda)/b^2$;
averaging under $q$, with $q y = 0$ and $q y^2 = \Var_q(f)$,
$q e^{\lambda(f - qf)} \le 1 + \Var_q(f)\varphi_b(\lambda)/b^2 \le \exp(\Var_q(f)\varphi_b(\lambda)/b^2)$
($1 + u \le e^u$). Taking logarithms (the left-hand side is positive) gives the claim.
\end{proof}

\begin{lemma}[Lemma 22: KL--Bernstein inequality]\label{lem:kl_bernstein}
Let $\alpha \ge 0$, $b > 0$ and $f$ with values in $[0, b]$. If $\KL(p \| q) \le \alpha$, then
\[
  \abs{pf - qf} \le \sqrt{2 \Var_q(f)\, \alpha} + \frac23\, b\, \alpha .
\]
\end{lemma}

\begin{proof}
\uses{lem:donsker_varadhan_finite, lem:bernstein_mgf_bounded, lem:poisson_cramer}
Put $d := pf - qf$ and $\sigma^2 := \Var_q(f) \ge 0$. For every $\lambda \ge 0$,
Lemma~\ref{lem:donsker_varadhan_finite} with $g := \lambda(f - qf)$ (note
$pg = \lambda(pf - qf) = \lambda d$, since $p(qf) = qf$) and Lemma~\ref{lem:bernstein_mgf_bounded} give
\[
  \alpha \ge \KL(p \| q) \ge \lambda d - \log\big(q e^{\lambda(f - qf)}\big) \ge \lambda d - \sigma^2 \varphi_b(\lambda)/b^2 .
\]
\emph{Case $d \ge 0$.} If $\sigma^2 = 0$, letting $\lambda \to \infty$ in $\lambda d \le \alpha$ gives
$d = 0$, and the claim holds. If $\sigma^2 > 0$, apply the inequality at
$\hat\lambda := \frac1b \log(1 + s/c)$ with $c := \sigma^2/b^2 > 0$ and $s := d/b \ge 0$:
by Lemma~\ref{lem:poisson_cramer}(ii), $\hat\lambda d - \sigma^2\varphi_b(\hat\lambda)/b^2 = \hat\lambda b s - c\varphi_b(\hat\lambda) = c\, h(s/c)$,
so $c\, h(s/c) \le \alpha$, and Lemma~\ref{lem:poisson_cramer}(iii) gives
$s \le \sqrt{2c\alpha} + \frac23\alpha$, i.e.\ $d \le b\sqrt{2\sigma^2\alpha/b^2} + \frac23 b\alpha = \sqrt{2\sigma^2\alpha} + \frac23 b\alpha$.

\emph{Case $d < 0$.} Apply the first case to $f' := b - f$, which takes values in $[0, b]$ and
satisfies $pf' - qf' = -d > 0$ and $\Var_q(f') = \Var_q(f)$.
\end{proof}

\textbf{Lean remark.} This is Lemma~3 of \cite{talebi2018variance} (with the constant $2/3$ of
the classical Bernstein inequality; the paper's statement has $\alpha > 0$, and the version with
$\alpha \ge 0$ costs nothing). The hypothesis $\KL(p\|q) \le \alpha$ is
\texttt{klDiv (weightedMeasure p) (weightedMeasure q) ≤ ENNReal.ofReal α}, which forces
absolute continuity (\texttt{klDiv\_of\_not\_ac}), and the conclusion uses
\texttt{ProbabilityTheory.variance f (weightedMeasure q)} and the integral $pf$.

\begin{lemma}[Lemma 25: variance transport under a KL constraint]\label{lem:kl_transport_var}
Let $\alpha \ge 0$, $b > 0$ and $f$ with values in $[0, b]$. If $\KL(p \| q) \le \alpha$, then
\[
  \Var_q(f) \le 2\Var_p(f) + 4b^2\alpha \qquad \text{and} \qquad \Var_p(f) \le 2\Var_q(f) + 4b^2\alpha .
\]
\end{lemma}

\begin{proof}
\uses{lem:kl_bernstein, lem:var_le_range_mean, lem:sqrt_mul_le}
Recall that for any constant $m$, $p(f - m)^2 = \Var_p(f) + (pf - m)^2 \ge \Var_p(f)$, and that
$\abs{f(x) - m} \le b$ whenever $m \in [0, b]$, so that $(f - m)^2$ takes values in $[0, b^2]$.

\emph{First inequality.} Let $g := (f - pf)^2 \in [0, b^2]$ ($pf \in [0, b]$), so that $pg = \Var_p(f)$
and $\Var_q(f) \le qg$. Lemma~\ref{lem:kl_bernstein} applied to $g$ with the range $b^2$ gives
$qg \le pg + \sqrt{2\Var_q(g)\alpha} + \frac23 b^2\alpha$, and Lemma~\ref{lem:var_le_range_mean}
(with $g \in [0, b^2]$) gives $\Var_q(g) \le b^2\, qg$. By Lemma~\ref{lem:sqrt_mul_le} with $\eta = 1$,
$\sqrt{2\alpha b^2 \cdot qg} \le \frac12 qg + b^2\alpha$. Hence
$qg \le \Var_p(f) + \frac12 qg + b^2\alpha + \frac23 b^2\alpha$, i.e.\
$qg \le 2\Var_p(f) + \frac{10}{3} b^2\alpha \le 2\Var_p(f) + 4b^2\alpha$, and $\Var_q(f) \le qg$.

\emph{Second inequality.} Let $g' := (f - qf)^2 \in [0, b^2]$, so that $qg' = \Var_q(f)$ and
$\Var_p(f) \le pg'$. Lemma~\ref{lem:kl_bernstein} applied to $g'$ gives
$pg' \le \Var_q(f) + \sqrt{2\Var_q(g')\alpha} + \frac23 b^2\alpha$, Lemma~\ref{lem:var_le_range_mean}
gives $\Var_q(g') \le b^2 qg' = b^2\Var_q(f)$, and Lemma~\ref{lem:sqrt_mul_le} with $\eta = 1$ gives
$\sqrt{2\alpha b^2 \Var_q(f)} \le \frac12\Var_q(f) + b^2\alpha$. Hence
$\Var_p(f) \le \frac32\Var_q(f) + \frac53 b^2\alpha \le 2\Var_q(f) + 4b^2\alpha$.
\end{proof}

\textbf{Lean remark.} This is Lemma~11 of \cite{menard2021fast}; the proof there is the same
(KL--Bernstein applied to the squared centered function). The constants $10/3$ and $5/3$ show
that $4$ is not tight; the analysis chapters use $4$.

\begin{lemma}[Lemma 26: variance transport between functions and between measures]\label{lem:transport_var}
Let $b \ge 0$ and $f, g$ with values in $[0, b]$. Then
\[
  \Var_p(f) \le 2\Var_p(g) + 2b\, p\abs{f - g} \qquad \text{and} \qquad
  \Var_q(f) \le \Var_p(f) + b^2 \norm{p - q}_1 \le \Var_p(f) + 3b^2\norm{p - q}_1 ,
\]
where $\abs{f - g}(x) := \abs{f(x) - g(x)}$.
\end{lemma}

\begin{proof}
\emph{First inequality.} $\Var_p(f) \le p(f - pg)^2$ (variance minimizes the centered second
moment: $p(f - m)^2 = \Var_p(f) + (pf - m)^2$). Pointwise,
$(f - pg)^2 = \big((f - g) + (g - pg)\big)^2 \le 2(f - g)^2 + 2(g - pg)^2$
($(u + v)^2 \le 2u^2 + 2v^2$, from $2uv \le u^2 + v^2$, \texttt{two\_mul\_le\_add\_sq}), and $(f - g)^2 = \abs{f - g}^2 \le b\abs{f - g}$
because $\abs{f - g} \le b$ for $f, g \in [0, b]$. Averaging under $p$:
$\Var_p(f) \le 2b\,p\abs{f - g} + 2\Var_p(g)$.

\emph{Second inequality.} $\Var_q(f) \le q(f - pf)^2 = p(f - pf)^2 + \sum_x (q(x) - p(x))(f(x) - pf)^2$,
and $0 \le (f(x) - pf)^2 \le b^2$ since $f(x), pf \in [0, b]$, so the last sum is at most
$b^2\sum_x \abs{q(x) - p(x)} = b^2\norm{p - q}_1$. Finally $b^2\norm{p - q}_1 \le 3b^2\norm{p - q}_1$.
\end{proof}

\textbf{Lean remark.} This is Lemma~12 of \cite{menard2021fast}, stated there with the constant
$3$ in the second inequality; the proof gives $1$. Only the first inequality is used by the
analysis (Lemmas~\ref{lem:concentration_pos} and~\ref{lem:concentration_neg}); the second is
kept for completeness. In Lean, $p\abs{f - g}$ is the integral of \texttt{fun x => |f x - g x|}
and $\norm{p - q}_1$ is \texttt{∑ x, |p x - q x|} (for \texttt{weightedMeasure}s it is twice the
total variation distance).

\begin{lemma}[Variance bounded by the range times the mean]\label{lem:var_le_range_mean}
Let $b \ge 0$ and $f$ with values in $[0, b]$. Then $\Var_p(f) \le pf^2 \le b\, pf$.
\end{lemma}

\begin{proof}
$\Var_p(f) = pf^2 - (pf)^2 \le pf^2$ since $(pf)^2 \ge 0$ (\texttt{variance\_def'}), and
$f(x)^2 = f(x) \cdot f(x) \le b\, f(x)$ for every $x$ because $0 \le f(x) \le b$
(\texttt{mul\_le\_mul\_of\_nonneg\_right}); averaging under $p$ gives $pf^2 \le b\,pf$.
\end{proof}

\chapter{Change of measure}
\label{chap:pre_info}

This chapter proves the change-of-measure inequality used by the lower bound
(Chapter~\ref{chap:lower_bound}): the paper's Lemma~17 (Appendix~C of
\cite{essakine2026tight}), which is Lemma~1 of \cite{kaufmann2016complexity} for bandit models.
We state and prove it in LML generality: an arbitrary algorithm interacts with one of two
\emph{stationary environments} with feedback kernels $\nu, \nu' : \As \to \Delta(\mathcal Y)$
(the ``arms'' of the bandit model are the actions $a \in \As$, a finite set), it stops at the
stopping time $\tau$ of a stopping rule and outputs a decision drawn from a Markov kernel of
the stopped history. The conclusion is that for every event $E$ of the pair (stopped history,
output),
\[
  \klbin\big(\Prob_\nu(E), \Prob_{\nu'}(E)\big)
  \le \sum_{a \in \As} \E_\nu[N_a(\tau)]\,\KL(\nu_a \,\|\, \nu'_a),
\]
where $N_a(\tau)$ is the number of rounds before $\tau$ at which $a$ was played and $\klbin$
is the binary relative entropy (Section~\ref{sec:pinfo_klbin}). The right-hand side is the
\emph{divergence decomposition} (Section~\ref{sec:pinfo_decomposition}): the divergence between
the laws of the histories of $n$ rounds under the two environments is
$\sum_a \E_\nu[N_a(n)] \KL(\nu_a\|\nu'_a)$ (Lemma~\ref{lem:kl_history_le}, formalized in LML),
and the same holds at a stopping time (Lemma~\ref{lem:kl_stopped_hist_le}). The episode
environment of an MDP (Definition~\ref{def:episode_env}) is such a stationary environment, with
actions the policies and feedback the state sequence of the episode, which is how
Chapter~\ref{chap:lower_bound} applies the theorem (Remark~\ref{rem:pinfo_paper_lemma}).

\textbf{What Mathlib and LML provide.} Mathlib defines the Kullback--Leibler divergence
\texttt{InformationTheory.klDiv μ ν} of two measures as an element of $[0, \infty]$
(\texttt{ℝ≥0∞}), equal to $\int \log\frac{d\mu}{d\nu}\,d\mu + \nu(\Omega) - \mu(\Omega)$ when
$\mu \ll \nu$ and the log-likelihood ratio is $\mu$-integrable and to $\infty$ otherwise
(\texttt{klDiv\_of\_ac\_of\_integrable}, \texttt{klDiv\_of\_not\_ac}), and equal to the
$f$-divergence $\int \mathrm{klFun}(d\mu/d\nu)\,d\nu$ with $\mathrm{klFun}(x) = x\log x + 1 - x$
when $\mu \ll \nu$ (\texttt{klDiv\_eq\_lintegral\_klFun\_of\_ac}). It proves the chain rule
\texttt{klDiv\_compProd\_eq\_add}, $\KL(\mu\otimes\kappa\,\|\,\nu\otimes\eta) = \KL(\mu\|\nu) +
\KL(\mu\otimes\kappa\,\|\,\mu\otimes\eta)$, its special case \texttt{klDiv\_compProd\_left}
($\KL(\mu\otimes\kappa\,\|\,\nu\otimes\kappa) = \KL(\mu\|\nu)$) and the data-processing
inequality \texttt{klDiv\_map\_le} ($\KL(g_*\mu\|g_*\nu) \le \KL(\mu\|\nu)$ for measurable $g$).
Bernoulli measures are \texttt{ProbabilityTheory.bernoulliMeasure x y p} (mass $p$ at $x$,
$1 - p$ at $y$, $p \in [0,1]$). LML adds
(\texttt{LeanMachineLearning/ForMathlib/InformationTheory/KullbackLeibler/}): the invariance
under measurable equivalences (\texttt{klDiv\_map\_measurableEquiv}), the integrated chain rule
$\KL(\mu\otimes\kappa\,\|\,\mu\otimes\eta) = \int^- \KL(\kappa_x\|\eta_x)\,\mu(dx)$ for a countably
generated target (\texttt{klDiv\_compProd\_right\_eq\_lintegral},
\texttt{measurable\_klDiv\_kernel}), the behaviour under restrictions
(\texttt{klDiv\_restrict\_add\_restrict\_compl}, \texttt{klDiv\_restrict\_le},
\texttt{klDiv\_add\_add\_of\_measure\_eq\_zero}, \texttt{klDiv\_eq\_iSup\_restrict}), and the
divergence decomposition for a fixed number of rounds
(\texttt{SequentialLearning/DivergenceDecomposition.lean}:
\texttt{IsAlgEnvSeq.klDiv\_map\_history}, Lemma~\ref{lem:kl_history_le}). Stopping rules,
stopping times and stopped histories of algorithm-environment sequences are in
\texttt{SequentialLearning/StoppedHistory.lean} (\texttt{sigmaHistory},
\texttt{stoppedHistMeasure}) and identification algorithms in
\texttt{SequentialLearning/IdentificationAlg.lean} (\texttt{IdentAlg}, \texttt{stoppingTime},
\texttt{stoppedHist}, \texttt{IsRun}, \texttt{outputMeasure}). The binary divergence
\texttt{klBin} and its relation to Bernoulli measures are in LMLPapers
(\texttt{LeanMachineLearning/ForMathlib/InformationTheory/Pinsker.lean}); the sister project
\texttt{colt-2026-83} (\texttt{Maiti2026Power/Mathlib/InformationTheory/KLStoppedPrefix.lean},
\texttt{Maiti2026Power/LeanMachineLearning/\{StoppedHistory, DivergenceDecomposition,
RunDivergence\}.lean}) contains the divergence decomposition at a stopping time, which
Lemmas~\ref{lem:pinfo_stopped_chain_rule} and~\ref{lem:kl_stopped_hist_le} follow.

\textbf{What is missing.} The elementary inequalities on $\klbin$
(Lemmas~\ref{lem:klbin_lower} and~\ref{lem:klbin_upper}), its convexity and lower semicontinuity
(Lemmas~\ref{lem:pinfo_klbin_convex_right}, \ref{lem:pinfo_klbin_lsc}), and the change of measure
itself (Theorem~\ref{thm:change_of_measure}) in the one-sided form needed by the lower bound
(the stopping time is only assumed almost surely finite under the \emph{first} environment).

\section{The binary relative entropy}
\label{sec:pinfo_klbin}

\begin{definition}[Binary relative entropy]\label{def:klbin}
For $x, y \in [0, 1]$, the \emph{binary relative entropy} is
\[
  \klbin(x, y) := x\log\frac{x}{y} + (1 - x)\log\frac{1 - x}{1 - y} \in [0, \infty],
\]
with the conventions $0\log(0/y) = 0$ for every $y$ (so that $\klbin(0, 0) = \klbin(1, 1) = 0$)
and $x\log(x/0) = +\infty$ for $x > 0$ (so that $\klbin(x, y) = +\infty$ when $y \in \{0, 1\}$
and $x \ne y$). Equivalently, $\klbin(x, y) = \KL(\mathrm{Ber}(x)\,\|\,\mathrm{Ber}(y))$ for the
Bernoulli measures $\mathrm{Ber}(p) := p\,\delta_1 + (1 - p)\,\delta_0$ on $\{0, 1\}$
(Lemma~\ref{lem:klbin_bernoulli}). For $y \in (0, 1)$ the value is finite and
$\klbin(x, y) = x\log x + (1-x)\log(1-x) - x\log y - (1-x)\log(1-y)$ with $0\log 0 = 0$.
\end{definition}

\textbf{Lean remark.} LMLPapers defines the real-valued
\texttt{InformationTheory.klBin (p q : ℝ) : ℝ := p * log (p / q) + (1 - p) * log ((1 - p) / (1 - q))},
which agrees with $\klbin$ for $p \in [0,1]$ and $q \in (0,1)$ (Lean's \texttt{Real.log 0 = 0}
gives the convention $0\log 0 = 0$ for free) but is a junk real number when $q \in \{0, 1\}$.
Statements in which $y$ may be $0$ or $1$ are therefore formalized with the $[0,\infty]$-valued
\texttt{klDiv (bernoulliMeasure true false x) (bernoulliMeasure true false y)}, and the real
inequalities below are stated for $y \in (0, 1)$; the passage between the two forms is
\texttt{klDiv\_bernoulliMeasure} (LMLPapers) together with \texttt{ENNReal.ofReal\_le\_iff}.
The project keeps a copy of the LMLPapers file in \texttt{Essakine2026Tight/Mathlib/}.

\begin{lemma}[Divergence of two Bernoulli measures]\label{lem:klbin_bernoulli}
\uses{def:klbin}
Let $x \in [0, 1]$ and $y \in (0, 1)$. Then
$\KL(\mathrm{Ber}(x)\,\|\,\mathrm{Ber}(y)) = \klbin(x, y) < \infty$. If $y \in \{0, 1\}$ and
$x \ne y$ then $\KL(\mathrm{Ber}(x)\|\mathrm{Ber}(y)) = \infty = \klbin(x, y)$, and if $x = y$
both sides vanish.
\end{lemma}

\begin{proof}
\uses{def:klbin}
For $y \in (0, 1)$, $\mathrm{Ber}(y)$ charges both points, so $\mathrm{Ber}(x) \ll \mathrm{Ber}(y)$
with Radon--Nikodym derivative $x/y$ at $1$ and $(1-x)/(1-y)$ at $0$, and every function on a
two-point space is integrable. By Mathlib \texttt{klDiv\_eq\_lintegral\_klFun\_of\_ac} the
divergence is $y\,\mathrm{klFun}(x/y) + (1-y)\,\mathrm{klFun}((1-x)/(1-y))$ with
$\mathrm{klFun}(u) = u\log u + 1 - u$, which expands to
$x\log(x/y) + (1-x)\log((1-x)/(1-y)) + (y - x) + ((1-y) - (1-x)) = \klbin(x, y)$; the
convention $0\log 0 = 0$ is $\mathrm{klFun}(0) = 1$ (\texttt{klFun\_zero}). If $y = 0$ and
$x > 0$ then $\mathrm{Ber}(x)(\{1\}) > 0 = \mathrm{Ber}(0)(\{1\})$, so
$\mathrm{Ber}(x) \not\ll \mathrm{Ber}(y)$ and the divergence is $\infty$
(\texttt{klDiv\_of\_not\_ac}); $y = 1$ is symmetric. If $x = y$ the two measures coincide and
\texttt{klDiv\_self} gives $0$.
\end{proof}

\textbf{Lean remark.} This is \texttt{InformationTheory.klDiv\_bernoulliMeasure} of LMLPapers
(for \texttt{x ≠ y} points and $0 < q < 1$) and \texttt{klDiv\_bernoulliMeasure\_eq\_klBerReal}
of the sister project; the infinite case follows from \texttt{bernoulliMeasure\_zero},
\texttt{bernoulliMeasure\_one} and \texttt{klDiv\_of\_not\_ac}.

\begin{lemma}[Data processing to an event]\label{lem:klbin_le_kl}
\uses{def:klbin, lem:klbin_bernoulli}
Let $\mu, \nu$ be probability measures on a measurable space and $E$ a measurable set. Then
\[
  \klbin\big(\mu(E), \nu(E)\big) \le \KL(\mu \,\|\, \nu) \qquad\text{in } [0, \infty].
\]
\end{lemma}

\begin{proof}
\uses{lem:klbin_bernoulli}
Let $g := \indic_E$, a measurable map to $\{0, 1\}$. Then $g_*\mu = \mathrm{Ber}(\mu(E))$ and
$g_*\nu = \mathrm{Ber}(\nu(E))$, and the data-processing inequality (Mathlib
\texttt{klDiv\_map\_le}) gives $\KL(\mathrm{Ber}(\mu(E))\,\|\,\mathrm{Ber}(\nu(E))) \le \KL(\mu\|\nu)$.
The left-hand side is $\klbin(\mu(E), \nu(E))$ by Lemma~\ref{lem:klbin_bernoulli}, in all cases
(finite, infinite, or $\mu(E) = \nu(E)$).
\end{proof}

\textbf{Lean remark.} The map is \texttt{fun ω ↦ decide (ω ∈ E)} and its image measure is
\texttt{bernoulliMeasure true false (μ.real E)} (LMLPapers
\texttt{map\_decide\_mem\_eq\_bernoulliMeasure}); this is the first half of the proof of
Pinsker's inequality \texttt{pinsker\_measureReal} there, and of
\texttt{InformationTheory.sq\_sub\_le\_klDiv} in the sister project. The statement to record
is \texttt{klDiv (bernoulliMeasure true false (μ.real E)) (bernoulliMeasure true false (ν.real E))
≤ klDiv μ ν}, from which the real form \texttt{klBin (μ.real E) (ν.real E) ≤ (klDiv μ ν).toReal}
follows when $\nu(E) \in (0,1)$ and $\KL(\mu\|\nu) < \infty$.

\begin{lemma}[Lower bound on the binary relative entropy]\label{lem:klbin_lower}
\uses{def:klbin}
For $x \in [0, 1]$ and $y \in (0, 1)$,
\[
  \klbin(x, y) \ge (1 - x)\log\frac{1}{1 - y} - \log 2 .
\]
\end{lemma}

\begin{proof}
\uses{def:klbin}
Expand $\klbin(x, y) = \big[x\log x + (1-x)\log(1-x)\big] - x\log y - (1 - x)\log(1 - y)$ with
$0 \log 0 = 0$. The bracket is minus the binary entropy of $x$, which is at most $\log 2$
(concavity of $\log$: $x\log(1/x) + (1-x)\log(1/(1-x)) \le \log(x \cdot \tfrac1x + (1-x)\cdot\tfrac1{1-x}) = \log 2$
by Jensen, the terms with $x \in \{0,1\}$ being $0$); so the bracket is $\ge -\log 2$. Next
$-x\log y \ge 0$ since $y \le 1$ and $x \ge 0$. Finally
$-(1-x)\log(1-y) = (1-x)\log\frac{1}{1-y}$. Summing the three bounds gives the claim.
\end{proof}

\begin{lemma}[Upper bound on the binary relative entropy]\label{lem:klbin_upper}
\uses{def:klbin}
For $x \in [0, 1]$ and $y \in (0, 1)$,
\[
  \klbin(x, y) \le \frac{(x - y)^2}{y(1 - y)} .
\]
\end{lemma}

\begin{proof}
\uses{def:klbin}
Use $\log u \le u - 1$ for $u > 0$ (Mathlib \texttt{Real.log\_le\_sub\_one\_of\_pos}). If
$x > 0$, $x\log\frac{x}{y} \le x\big(\frac{x}{y} - 1\big) = \frac{x(x - y)}{y}$, and this also
holds for $x = 0$ (both sides vanish). If $x < 1$,
$(1-x)\log\frac{1-x}{1-y} \le (1-x)\big(\frac{1-x}{1-y} - 1\big) = \frac{(1-x)(y - x)}{1-y}$, and
this also holds for $x = 1$. Summing,
\[
  \klbin(x, y) \le (x - y)\Big[\frac{x}{y} - \frac{1 - x}{1 - y}\Big]
  = (x - y)\,\frac{x(1-y) - y(1-x)}{y(1-y)} = \frac{(x - y)^2}{y(1-y)} .
\]
\end{proof}

\begin{lemma}[Convexity in the second argument]\label{lem:pinfo_klbin_convex_right}
\uses{def:klbin}
Fix $x \in [0, 1]$. The function $y \mapsto \klbin(x, y)$ is convex on $[0, 1]$ (as a function
with values in $[0, \infty]$). Consequently, for $0 \le y_1 \le y \le y_2 \le 1$,
\[
  \klbin(x, y) \le \max\big\{\klbin(x, y_1), \klbin(x, y_2)\big\}.
\]
\end{lemma}

\begin{proof}
\uses{def:klbin}
On $(0, 1)$, $\klbin(x, y) = c_x - x\log y - (1-x)\log(1-y)$ with $c_x := x\log x + (1-x)\log(1-x)$
independent of $y$; $y \mapsto -\log y$ and $y \mapsto -\log(1-y)$ are convex on $(0,1)$
(Mathlib \texttt{strictConcaveOn\_log\_Ioi} and its composition with the affine map
$y \mapsto 1 - y$), and nonnegative combinations of convex functions are convex, so the
function is convex on $(0, 1)$. At the endpoints: if $x \in (0, 1)$ then $\klbin(x, 0) =
\klbin(x, 1) = \infty$ and convexity on $[0,1]$ is the convexity on $(0,1)$ (a chord ending at
an infinite value imposes no constraint). If $x = 1$ then $\klbin(1, y) = -\log y$ on $(0, 1]$
(value $0$ at $y = 1$, $\infty$ at $y = 0$), convex on $[0, 1]$; $x = 0$ is symmetric. The
consequence is the standard property of convex functions on an interval: on $[y_1, y_2]$ a convex
function is bounded by the chord, hence by the larger endpoint value.
\end{proof}

\begin{lemma}[Lower semicontinuity]\label{lem:pinfo_klbin_lsc}
\uses{def:klbin}
Let $(x_n, y_n) \to (x, y)$ in $[0, 1]^2$ and $D \in [0, \infty)$ with $\klbin(x_n, y_n) \le D$
for all $n$. Then $\klbin(x, y) \le D$.
\end{lemma}

\begin{proof}
\uses{def:klbin}
Write $\klbin(u, v) = c(u) + \ell_1(u, v) + \ell_0(u, v)$ with $c(u) := u\log u + (1-u)\log(1-u)$,
$\ell_1(u, v) := u\log(1/v)$ and $\ell_0(u, v) := (1-u)\log(1/(1-v))$, all three with the
conventions of Definition~\ref{def:klbin} ($\ell_1(0, v) = 0$, $\ell_1(u, 0) = \infty$ for
$u > 0$, and similarly for $\ell_0$). The function $c$ is continuous on $[0,1]$ (Mathlib
\texttt{Real.continuous\_mul\_log}). The function $\ell_1 : [0,1]^2 \to [0,\infty]$ is lower
semicontinuous: it is continuous on $\{v > 0\}$ and on $\{u = 0\}$, and at a point $(u, 0)$ with
$u > 0$, $\ell_1(u_n, v_n) = u_n\log(1/v_n) \to \infty$ since $u_n \to u > 0$ and
$\log(1/v_n) \to \infty$. Likewise for $\ell_0$. A sum of lower semicontinuous functions with
values in $[0,\infty]$ plus a continuous function is lower semicontinuous, so
$\klbin(x, y) \le \liminf_n \klbin(x_n, y_n) \le D$.
\end{proof}

\textbf{Lean remark.} Lemma~\ref{lem:pinfo_klbin_lsc} is the lower semicontinuity of
$\KL(\mathrm{Ber}(\cdot)\|\mathrm{Ber}(\cdot))$; in Lean it is easiest to prove directly the
statement used in Theorem~\ref{thm:change_of_measure}: for the $[0,\infty]$-valued
\texttt{klDiv} of Bernoulli measures, \texttt{Tendsto} of the parameters and an eventual bound by
\texttt{D} give the bound at the limit, via \texttt{ENNReal.le\_of\_forall\_pos\_le\_add} and the
continuity of \texttt{Real.log} on \texttt{(0, 1]} together with the divergence to \texttt{⊤} at
\texttt{0} (\texttt{Real.tendsto\_log\_nhdsGT\_zero}).

\section{The divergence decomposition}
\label{sec:pinfo_decomposition}

\textbf{Standing setting.} Let $\As$ be a finite set (the actions; every subset of a finite
type is measurable) and $\mathcal Y$ a countably generated measurable space (the feedback;
finite types, $\R$ and countable products of such spaces qualify). Let $\nu, \nu' : \As \to
\Delta(\mathcal Y)$ be two Markov kernels and $\mathrm{env} := \mathrm{stationaryEnv}(\nu)$,
$\mathrm{env}' := \mathrm{stationaryEnv}(\nu')$ the two stationary environments without
observations (LML \texttt{Learning.stationaryEnv}: at every round, given the past and the
action $a$, the feedback has law $\nu(a)$, \texttt{feedback\_stationaryEnv}). Let
$\mathrm{alg}$ be an algorithm with actions in $\As$ and feedback in $\mathcal Y$ (policy
kernels $\pi_n$ from histories of $n$ rounds to $\As$). We consider two
\emph{algorithm-environment sequences} (LML \texttt{IsAlgEnvSeq}): $(X_t, Y_t)_{t \in \N}$ on
$(\Omega, \Prob)$ for $(\mathrm{alg}, \mathrm{env})$ and $(X'_t, Y'_t)_{t\in\N}$ on
$(\Omega', \Prob')$ for $(\mathrm{alg}, \mathrm{env}')$: the action at round $t$ has
conditional law $\pi_t(\hist_t)$ given the past and the feedback has conditional law
$\nu(X_t)$ (resp.\ $\nu'(X'_t)$) given the past and the action. Here
$\hist_n := (X_t, Y_t)_{t < n}$ is the \emph{history of the first $n$ rounds}, valued in
$\mathcal H_n := (\{0, \dots, n-1\} \to \As \times \mathcal Y)$ (LML \texttt{Learning.history};
$\hist_{n+1} = (\hist_n, (X_n, Y_n))$), and $\langle n, \hist_n\rangle$ the same history in the
space $\mathcal H := \bigsqcup_n \mathcal H_n$ of histories of variable length
(\texttt{sigmaHistory}). The \emph{history filtration} is $\mathcal F_n := \sigma(\hist_n)$.
For $a \in \As$ and $n \in \N \cup \{\infty\}$,
\[
  N_a(n) := \sum_{t < n} \indic\{X_t = a\} \in \N \cup \{\infty\}
\]
is the number of rounds before $n$ at which $a$ is played, and $\E_\nu[N_a(n)] := \int N_a(n)\,d\Prob
\in [0, \infty]$; we write $\E_\nu$, $\Prob_\nu$ for expectations and probabilities under
$\Prob$ and $\E_{\nu'}$, $\Prob_{\nu'}$ under $\Prob'$ (by LML
\texttt{IsAlgEnvSeq.hasLaw\_history}, the law of $\hist_n$ is determined by $(\mathrm{alg},
\mathrm{env})$, so all statements below about laws of histories do not depend on the
probability spaces). Products $c \cdot \infty$ with $c \in [0,\infty]$ follow the convention
$0 \cdot \infty = 0$ of \texttt{ℝ≥0∞}: an action that is never played contributes nothing even
if its one-step divergence is infinite.

A \emph{stopping rule} is a measurable set $S \subseteq \mathcal H$; its \emph{stopping time}
on a sequence is $\tau_S := \inf\{n : \langle n, \hist_n\rangle \in S\} \in \N \cup \{\infty\}$
(LML: \texttt{hittingAfter (sigmaHistory O X Y) S 0}), a stopping time of the history
filtration; conversely every stopping time $\tau$ of the history filtration on one space is
$\tau_S$ for $S := \{\langle n, h\rangle : \{\hist_n = h\} \subseteq \{\tau = n\}\}$ (Doob--Dynkin
for the sets $\{\tau = n\} \in \sigma(\hist_n)$), so we lose nothing by working with stopping
rules, which make sense simultaneously on both spaces. The \emph{stopped history} is
$\hist_\tau := \langle \tau, \hist_\tau\rangle \in \mathcal H$ on $\{\tau < \infty\}$ (Mathlib
\texttt{stoppedValue}; an arbitrary element of $\mathcal H$ on $\{\tau = \infty\}$), and the
history stopped at time $M$ is $\hist_{\tau \wedge M}$ (Mathlib \texttt{stoppedProcess}). We
write $\mathrm{trunc}_M : \mathcal H \to \mathcal H$ for the truncation to the first $M$ rounds,
so that $\mathrm{trunc}_M(\hist_\tau) = \hist_{\tau\wedge M}$ on $\{\tau < \infty\}$.

\begin{lemma}[Divergence decomposition for a fixed number of rounds]\label{lem:kl_history_le}
In the standing setting, for every $n \in \N$,
\[
  \KL\big(\mathrm{law}_\nu(\hist_n)\,\big\|\,\mathrm{law}_{\nu'}(\hist_n)\big)
  = \sum_{t < n} \E_\nu\big[\KL(\nu(X_t)\,\|\,\nu'(X_t))\big]
  = \sum_{a \in \As} \E_\nu[N_a(n)]\,\KL(\nu_a \,\|\, \nu'_a) \qquad\text{in } [0, \infty],
\]
where $\mathrm{law}_\nu(\hist_n)$ is the law of $\hist_n$ under $\Prob$ and
$\mathrm{law}_{\nu'}(\hist_n)$ the law of $\hist'_n := (X'_t, Y'_t)_{t<n}$ under $\Prob'$.
\end{lemma}

\begin{proof}
\uses{lem:pinfo_kl_step}
Induction on $n$. For $n = 0$, $\mathcal H_0$ is a one-point space, both laws are the Dirac
mass at the empty history and both sides vanish (\texttt{klDiv\_self}). For the step, the law
of $\hist_{n+1} = (\hist_n, (X_n, Y_n))$ is (up to the measurable equivalence
$\mathcal H_{n+1} \simeq \mathcal H_n \times (\As\times\mathcal Y)$, which leaves $\KL$ invariant,
LML \texttt{klDiv\_map\_measurableEquiv}) the composition-product
$\mathrm{law}_\nu(\hist_n) \otimes (\pi_n \otimes \tilde\nu)$ where $\tilde\nu(h, a) := \nu(a)$
(LML \texttt{IsAlgEnvSeq.hasCondDistrib\_step}, \texttt{stepKernel\_stationaryEnv}), and
similarly with primes and the \emph{same} policy kernel $\pi_n$. The chain rule
\texttt{klDiv\_compProd\_eq\_add} gives
\[
  \KL(\hist_{n+1}) = \KL(\hist_n)
  + \KL\big(\mathrm{law}_\nu(\hist_n)\otimes(\pi_n\otimes\tilde\nu)\,\big\|\,
      \mathrm{law}_\nu(\hist_n)\otimes(\pi_n\otimes\tilde\nu')\big),
\]
and the last term is $\E_\nu[\KL(\nu(X_n)\|\nu'(X_n))]$ by Lemma~\ref{lem:pinfo_kl_step} (the
law of $X_n$ under $\Prob$ is $\pi_n \circ \mathrm{law}_\nu(\hist_n)$, LML
\texttt{hasCondDistrib\_action}). The induction hypothesis gives the first equality. For the
second, $\As$ being finite,
$\KL(\nu(X_t)\|\nu'(X_t)) = \sum_a \indic\{X_t = a\}\KL(\nu_a\|\nu'_a)$ pointwise, and summing
over $t < n$ and exchanging the finite sums with the Lebesgue integral
(\texttt{lintegral\_finset\_sum}, \texttt{lintegral\_const\_mul}),
$\sum_{t<n}\E_\nu[\KL(\nu(X_t)\|\nu'(X_t))] = \sum_a \KL(\nu_a\|\nu'_a)\,\E_\nu[\sum_{t<n}\indic\{X_t = a\}]$.
\end{proof}

\textbf{Lean remark.} The first equality is exactly LML's
\texttt{Learning.IsAlgEnvSeq.klDiv\_map\_history}: for \texttt{h : IsAlgEnvSeq O X Y alg
(stationaryEnv κ) P} and \texttt{h' : IsAlgEnvSeq O' X' Y' alg (stationaryEnv κ') P'} and
\texttt{[MeasurableSpace.CountablyGenerated 𝓨]}, \texttt{klDiv (P.map (history O X Y M))
(P'.map (history O' X' Y' M)) = ∑ t ∈ range M, ∫⁻ ω, klDiv (κ (X t ω)) (κ' (X t ω)) ∂P}
(already formalized; its composition-product form
\texttt{klDiv\_map\_history\_compProd} needs no countable generation). Only the reindexing by
actions, which uses \texttt{[Fintype 𝓐]} and \texttt{Finset.sum\_comm}, is to be added; the
count $N_a(n)$ is LML's \texttt{Learning.pullCount X a n} (file
\texttt{SequentialLearning/FiniteActions.lean}), valued in \texttt{ℕ}, coerced to
\texttt{ℝ≥0∞} inside the Lebesgue integral.

\begin{lemma}[Restrictions and the divergence]\label{lem:pinfo_kl_restrict}
Let $\mu, \nu$ be finite measures on a measurable space $\Omega$ and $s$ a measurable set.
\begin{enumerate}
  \item[(i)] $\KL(\mu\|\nu) = \KL(\mu|_s\,\|\,\nu|_s) + \KL(\mu|_{s^c}\,\|\,\nu|_{s^c})$; in
        particular $\KL(\mu|_s\|\nu|_s) \le \KL(\mu\|\nu)$.
  \item[(ii)] If $\mu = \mu_1 + \mu_2$ and $\nu = \nu_1 + \nu_2$ with $\mu_1, \nu_1$ carried by
        $s$ and $\mu_2, \nu_2$ carried by $s^c$, then $\KL(\mu\|\nu) = \KL(\mu_1\|\nu_1) + \KL(\mu_2\|\nu_2)$.
  \item[(iii)] If $s_n \uparrow \Omega$ is an increasing sequence of measurable sets, then
        $\KL(\mu\|\nu) = \sup_n \KL(\mu|_{s_n}\,\|\,\nu|_{s_n})$.
\end{enumerate}
\end{lemma}

\begin{proof}
The Radon--Nikodym derivative of $\mu|_s$ with respect to $\nu|_s$ is that of $\mu$ with
respect to $\nu$ (LML \texttt{rnDeriv\_restrict\_restrict}), so when $\mu \ll \nu$ all three
claims are identities for the integral $\int \mathrm{klFun}(d\mu/d\nu)\,d\nu$ over $s$, $s^c$ and
$s_n$: additivity of the integral over a partition for (i) and (ii) (where
$\mu|_s = \mu_1$, $\nu|_s = \nu_1$ since $\mu_2, \nu_2$ vanish on $s$), and monotone convergence
(\texttt{lintegral\_iSup} for the nonnegative integrand $\mathrm{klFun}(d\mu/d\nu)\indic_{s_n}$)
for (iii). When $\mu \not\ll \nu$, pick $A$ with $\nu(A) = 0 < \mu(A)$; then $A \cap s$ or
$A \cap s^c$ has positive $\mu$-measure and zero $\nu$-measure, so the corresponding restricted
divergence is $\infty$ too, and $\mu(A \cap s_n) \uparrow \mu(A) > 0$ gives some $n$ with
$\KL(\mu|_{s_n}\|\nu|_{s_n}) = \infty$.
\end{proof}

\textbf{Lean remark.} Formalized in LML
(\texttt{ForMathlib/InformationTheory/KullbackLeibler/Restrict.lean}):
\texttt{klDiv\_restrict\_add\_restrict\_compl}, \texttt{klDiv\_restrict\_le},
\texttt{klDiv\_add\_add\_of\_measure\_eq\_zero} and \texttt{klDiv\_eq\_iSup\_restrict}.

\begin{lemma}[One step of the algorithm against two environments]\label{lem:pinfo_kl_step}
Let $P$ be a finite measure on a measurable space $\mathcal H_0$, $\pi$ a Markov kernel from
$\mathcal H_0$ to $\As$ and $\tilde\nu, \tilde\nu'$ the kernels $(h, a) \mapsto \nu(a)$,
$(h, a) \mapsto \nu'(a)$ from $\mathcal H_0 \times \As$ to $\mathcal Y$ (Mathlib
\texttt{Kernel.prodMkLeft}). Then
\[
  \KL\big(P\otimes(\pi\otimes\tilde\nu)\,\big\|\,P\otimes(\pi\otimes\tilde\nu')\big)
  = \KL\big((\pi\circ P)\otimes\nu\,\big\|\,(\pi\circ P)\otimes\nu'\big)
  = \int^- \KL(\nu(a)\|\nu'(a))\,(\pi\circ P)(da)
  = \sum_{a\in\As} (\pi\circ P)(\{a\})\,\KL(\nu_a\|\nu'_a),
\]
where $\pi \circ P := \int \pi(h, \cdot)\,P(dh)$ is the law of the action.
\end{lemma}

\begin{proof}
The measure $P \otimes (\pi \otimes \tilde\nu)$ is the image of $(P \otimes \pi) \otimes \tilde\nu$
under the associativity equivalence $(h, (a, y)) \leftrightarrow ((h, a), y)$ (Mathlib
\texttt{Measure.compProd\_assoc}), and $\KL$ is invariant under measurable equivalences. Since
$\tilde\nu$ does not depend on $h$, the chain rule \texttt{klDiv\_compProd\_eq\_add} applied to
$(P\otimes\pi)\otimes\tilde\nu$ versus $(P\otimes\pi)\otimes\tilde\nu'$ gives
$\KL(P\otimes\pi\,\|\,P\otimes\pi) + \KL((P\otimes\pi)\otimes\tilde\nu\,\|\,(P\otimes\pi)\otimes\tilde\nu') $,
the first term being $0$; and the conditional divergence of kernels which depend on the
conditioning variable $(h, a)$ only through $a$ is the conditional divergence given $a$
(LML \texttt{klDiv\_compProd\_comap}; the law of $a$ under $P\otimes\pi$ is
$\pi\circ P$, Mathlib \texttt{Measure.snd\_compProd}). This is the first equality. The second is
the integrated chain rule for the countably generated space $\mathcal Y$ (LML
\texttt{klDiv\_compProd\_right\_eq\_lintegral}), and the third is the integral of a function on
the finite set $\As$ (\texttt{lintegral\_fintype}).
\end{proof}

\textbf{Lean remark.} The first two equalities are LML's
\texttt{Learning.klDiv\_compProd\_compProd\_prodMkLeft\_eq\_klDiv\_comp\_compProd}
(\texttt{SequentialLearning/DivergenceDecomposition.lean}) and
\texttt{klDiv\_compProd\_right\_eq\_lintegral}; the version with an observation kernel is
\texttt{klDiv\_compProd\_compProd\_compProd\_prodMkLeft\_eq\_klDiv\_comp\_compProd}.

\begin{lemma}[Chain rule at a bounded stopping time]\label{lem:pinfo_stopped_chain_rule}
\uses{lem:pinfo_kl_restrict, lem:pinfo_kl_step}
In the standing setting, let $S$ be a stopping rule with stopping times $\tau$ on $\Omega$ and
$\tau'$ on $\Omega'$, and let $M \in \N$. Then
\[
  \KL\big(\mathrm{law}_\nu(\hist_{\tau\wedge M})\,\big\|\,\mathrm{law}_{\nu'}(\hist'_{\tau'\wedge M})\big)
  = \sum_{t < M} \KL\big(\Prob|_{\{t < \tau\}}\circ X_t^{-1}\otimes\nu\,\big\|\,
      \Prob|_{\{t<\tau\}}\circ X_t^{-1}\otimes\nu'\big)
  = \sum_{a\in\As}\E_\nu[N_a(\tau\wedge M)]\,\KL(\nu_a\|\nu'_a),
\]
where $\Prob|_{\{t<\tau\}}\circ X_t^{-1}$ is the law of $X_t$ on the event $\{t < \tau\}$ (a
sub-probability measure on $\As$) and $\hist_{\tau\wedge M} = \langle\tau\wedge M,
\hist_{\tau\wedge M}\rangle \in \mathcal H$ is the history stopped at time $M$. No finiteness
assumption on $\tau$ is needed.
\end{lemma}

\begin{proof}
\uses{lem:pinfo_kl_restrict, lem:pinfo_kl_step}
\emph{First equality, by induction on $M$.} For $M = 0$ both laws are the Dirac mass at the
empty history and both sides vanish. Assume the identity for $M$ and write
$Z_M := \hist_{\tau\wedge M}$, $\mu_M := \mathrm{law}_\nu(Z_M)$, $\mu'_M := \mathrm{law}_{\nu'}(Z'_M)$.
Since $\tau$ is a stopping time of the history filtration, $\{\tau \le M\} = \{\hist_M \in B\}$
for a measurable set $B \subseteq \mathcal H_M$ which depends only on $S$ (namely $B$ is the set of
histories of length $M$ having a prefix, of length $\le M$, in $S$), and the same $B$ works on
$\Omega'$. The law of $Z_{M+1}$ splits according to whether $\tau \le M$:
\[
  \mu_{M+1} = \mathrm{law}_\nu(Z_M ; \tau\le M) + \mathrm{law}_\nu(\hist_{M+1} ; M < \tau),
\]
the first part being carried by $\mathcal H_{\le M} := \{\text{length} < M\} \cup \{\langle M, h\rangle : h \in B\}$
and the second by $\{\text{length} = M + 1\}$, disjoint sets; likewise
$\mu_M = \mathrm{law}_\nu(Z_M ; \tau \le M) + \mathrm{law}_\nu(\hist_M ; M < \tau)$ with parts
carried by $\mathcal H_{\le M}$ and by $\{\langle M, h\rangle : h \notin B\}$, disjoint. The same
decompositions hold under $\Prob'$ with the same sets. By Lemma~\ref{lem:pinfo_kl_restrict}(ii),
\begin{align*}
  \KL(\mu_{M+1}\|\mu'_{M+1}) &= \KL\big(\mathrm{law}_\nu(Z_M;\tau\le M)\,\|\,\mathrm{law}_{\nu'}(Z'_M;\tau'\le M)\big)
     + \KL\big(\mathrm{law}_\nu(\hist_{M+1}; M<\tau)\,\|\,\mathrm{law}_{\nu'}(\hist'_{M+1}; M<\tau')\big),\\
  \KL(\mu_M\|\mu'_M) &= \KL\big(\mathrm{law}_\nu(Z_M;\tau\le M)\,\|\,\mathrm{law}_{\nu'}(Z'_M;\tau'\le M)\big)
     + \KL\big(\mathrm{law}_\nu(\hist_M; M<\tau)\,\|\,\mathrm{law}_{\nu'}(\hist'_M; M<\tau')\big).
\end{align*}
On the event $\{M < \tau\} = \{\hist_M \notin B\}$, determined by $\hist_M$, the conditional
law of $(X_M, Y_M)$ given $\hist_M$ is still the step kernel $\pi_M \otimes \tilde\nu$ (a
conditional law given $\hist_M$ survives the restriction to an event of $\hist_M$; sister
project \texttt{hasCondDistrib\_step\_restrict\_lt\_hittingAfter\_sigmaHistory}), so
$\mathrm{law}_\nu(\hist_{M+1}; M<\tau) = \rho_M \otimes (\pi_M\otimes\tilde\nu)$ with
$\rho_M := \mathrm{law}_\nu(\hist_M; M < \tau)$, and likewise with primes, where
$\rho'_M := \mathrm{law}_{\nu'}(\hist'_M; M<\tau')$. The chain rule \texttt{klDiv\_compProd\_eq\_add}
and Lemma~\ref{lem:pinfo_kl_step} (for the finite measure $\rho_M$) give
\[
  \KL\big(\rho_M\otimes(\pi_M\otimes\tilde\nu)\,\|\,\rho'_M\otimes(\pi_M\otimes\tilde\nu')\big)
  = \KL(\rho_M\|\rho'_M) + \KL\big((\pi_M\circ\rho_M)\otimes\nu\,\|\,(\pi_M\circ\rho_M)\otimes\nu'\big),
\]
and $\pi_M \circ \rho_M = \Prob|_{\{M<\tau\}}\circ X_M^{-1}$ (LML
\texttt{hasCondDistrib\_action}, restricted to the event $\{M < \tau\}$ of $\hist_M$).
Subtracting the two displayed identities (all terms are in $[0,\infty]$; if a common term is
infinite, both $\KL(\mu_{M+1}\|\mu'_{M+1})$ and the claimed right-hand side are infinite) and
using the induction hypothesis gives the identity for $M + 1$.

\emph{Second equality.} By the integrated chain rule of Lemma~\ref{lem:pinfo_kl_step} applied
to the finite measure $\Prob|_{\{t<\tau\}}\circ X_t^{-1}$ on the finite set $\As$, the $t$-th
term is $\sum_a \Prob(t < \tau, X_t = a)\,\KL(\nu_a\|\nu'_a)$, and
$\sum_{t<M}\Prob(t<\tau, X_t = a) = \E_\nu[\sum_{t<M}\indic\{t<\tau, X_t = a\}] = \E_\nu[N_a(\tau\wedge M)]$.
\end{proof}

\textbf{Lean remark.} This is the sister project's
\texttt{Learning.IsAlgEnvSeq.klDiv\_map\_hittingProcess\_sigmaHistory\_compProd} (first
equality, composition-product form, for \texttt{hittingProcess (sigmaHistory O X Y) S 0 M})
and \texttt{klDiv\_map\_hittingProcess\_sigmaHistory} (integral form), both derived from the
general chain rule for a process with conditional laws given its prefixes,
\texttt{InformationTheory.klDiv\_map\_stoppedProcess\_sigmaPrefix\_compProd}
(\texttt{Maiti2026Power/Mathlib/InformationTheory/KLStoppedPrefix.lean}), transported to the
two spaces through the laws of the paths. The splitting of the law of $Z_{M+1}$ is
\texttt{map\_stoppedProcess\_sigmaPrefix\_succ\_eq\_add}, the set $B$ comes from
\texttt{IsStoppingTime.exists\_measurableSet\_preimage\_finPrefix}, and the restriction of the
conditional law of the step to $\{M < \tau\}$ is \texttt{HasCondDistrib.restrict\_lt\_of\_isStoppingTime}.
These files should be copied to \texttt{Essakine2026Tight/LeanMachineLearning/SequentialLearning/}
(and the \texttt{SigmaPrefix} layer to \texttt{Essakine2026Tight/Mathlib/}) until they reach LML.
The reindexing by actions is as in Lemma~\ref{lem:kl_history_le}, with
$N_a(\tau\wedge M) = \sum_{t < M}\indic\{t < \tau\}\indic\{X_t = a\}$ (the sum
\texttt{∑ t ∈ range M, if (t : ℕ∞) < τ ω then ... else 0}, sister project
\texttt{sum\_ite\_lt\_eq\_sum\_range\_toNat\_min}).

\begin{lemma}[Divergence decomposition at a stopping time]\label{lem:kl_stopped_hist_le}
\uses{lem:pinfo_stopped_chain_rule, lem:pinfo_kl_restrict}
In the standing setting, let $S$ be a stopping rule whose stopping times $\tau$ and $\tau'$
are almost surely finite under $\Prob$ and under $\Prob'$ respectively. Then
\[
  \KL\big(\mathrm{law}_\nu(\hist_\tau)\,\big\|\,\mathrm{law}_{\nu'}(\hist'_{\tau'})\big)
  = \sum_{t=0}^\infty \KL\big(\Prob|_{\{t < \tau\}}\circ X_t^{-1}\otimes\nu\,\big\|\,
      \Prob|_{\{t<\tau\}}\circ X_t^{-1}\otimes\nu'\big)
  = \sum_{a\in\As}\E_\nu[N_a(\tau)]\,\KL(\nu_a\|\nu'_a) \qquad\text{in } [0,\infty];
\]
in particular $\KL(\mathrm{law}_\nu(\hist_\tau)\,\|\,\mathrm{law}_{\nu'}(\hist'_{\tau'}))
\le \sum_a \E_\nu[N_a(\tau)]\KL(\nu_a\|\nu'_a)$, and the right-hand side is finite as soon as
$\E_\nu[\tau] < \infty$ and every $\KL(\nu_a\|\nu'_a)$ is finite.
\end{lemma}

\begin{proof}
\uses{lem:pinfo_stopped_chain_rule, lem:pinfo_kl_restrict}
Let $\mu := \mathrm{law}_\nu(\hist_\tau)$ and $\mu' := \mathrm{law}_{\nu'}(\hist'_{\tau'})$,
probability measures on $\mathcal H$. The series is the supremum over $M$ of its partial sums
(\texttt{ENNReal.tsum\_eq\_iSup\_nat}), which by Lemma~\ref{lem:pinfo_stopped_chain_rule} are
$\KL(\mu_M\|\mu'_M)$ with $\mu_M := \mathrm{law}_\nu(\hist_{\tau\wedge M})$,
$\mu'_M := \mathrm{law}_{\nu'}(\hist'_{\tau'\wedge M})$. Since $\tau < \infty$ almost surely,
$\hist_{\tau\wedge M} = \mathrm{trunc}_M(\hist_\tau)$ almost surely, so
$\mu_M = (\mathrm{trunc}_M)_*\mu$; likewise $\mu'_M = (\mathrm{trunc}_M)_*\mu'$ (this is where
almost sure finiteness under $\Prob'$ is used).

\emph{Upper bound $\KL(\mu\|\mu') \ge \sup_M\KL(\mu_M\|\mu'_M)$.} Data processing
(\texttt{klDiv\_map\_le}) for the measurable map $\mathrm{trunc}_M$.

\emph{Lower bound.} The sets $\mathcal H_{<M} := \{\text{length} < M\}$ increase to $\mathcal H$
(every history has finite length), so by Lemma~\ref{lem:pinfo_kl_restrict}(iii)
$\KL(\mu\|\mu') = \sup_M \KL(\mu|_{\mathcal H_{<M}}\,\|\,\mu'|_{\mathcal H_{<M}})$. Now
$\mathrm{trunc}_M$ is the identity on $\mathcal H_{<M}$ and
$\mathrm{trunc}_M^{-1}(\mathcal H_{<M}) = \mathcal H_{<M}$ (a truncated history has length
$< M$ iff the original one has), so $\mu|_{\mathcal H_{<M}} = ((\mathrm{trunc}_M)_*\mu)|_{\mathcal H_{<M}}
= \mu_M|_{\mathcal H_{<M}}$ and similarly for $\mu'$. By Lemma~\ref{lem:pinfo_kl_restrict}(i),
$\KL(\mu_M|_{\mathcal H_{<M}}\|\mu'_M|_{\mathcal H_{<M}}) \le \KL(\mu_M\|\mu'_M)$. Taking suprema
gives $\KL(\mu\|\mu') \le \sup_M\KL(\mu_M\|\mu'_M)$.

\emph{Third form.} By Lemma~\ref{lem:pinfo_stopped_chain_rule} the partial sums are
$\sum_a\E_\nu[N_a(\tau\wedge M)]\KL(\nu_a\|\nu'_a)$, and $N_a(\tau\wedge M) \uparrow N_a(\tau)$
as $M \to \infty$ (a nondecreasing sequence of nonnegative variables), so monotone convergence
(\texttt{lintegral\_iSup}) and the finiteness of the sum over $a$ give
$\sup_M \sum_a\E_\nu[N_a(\tau\wedge M)]\KL(\nu_a\|\nu'_a) = \sum_a\E_\nu[N_a(\tau)]\KL(\nu_a\|\nu'_a)$
(the suprema of finitely many nondecreasing sequences may be exchanged with the finite sum,
\texttt{ENNReal.finset\_sum\_iSup\_of\_monotone}). Finally $N_a(\tau) \le \tau$, so
$\E_\nu[N_a(\tau)] \le \E_\nu[\tau]$.
\end{proof}

\textbf{Lean remark.} Sister project
\texttt{Learning.IsAlgEnvSeq.klDiv\_map\_hittingValue\_sigmaHistory\_compProd} and
\texttt{klDiv\_map\_hittingValue\_sigmaHistory}, from
\texttt{InformationTheory.klDiv\_map\_stoppedValue\_sigmaPrefix\_compProd}, whose proof is
exactly the argument above (\texttt{HasLaw.stoppedProcess\_sigmaPrefix} for
$\mu_M = (\mathrm{trunc}_M)_*\mu$, \texttt{restrict\_map\_truncFin} for the restriction
identity, \texttt{klDiv\_eq\_iSup\_restrict} and \texttt{klDiv\_restrict\_le}). The hypotheses
are \texttt{∀ᵐ ω ∂P, τ ω ≠ ⊤} and \texttt{∀ᵐ ω ∂P', τ' ω ≠ ⊤}. For an identification algorithm
$A$ the stopping rule is \texttt{A.stopSet}, the stopping times are \texttt{A.stoppingTime O X Y}
and the stopped histories \texttt{A.stoppedHist O X Y}, with laws
\texttt{stoppedHistMeasure A.alg env A.stopSet} (LML
\texttt{IsAlgEnvSeq.hasLaw\_stoppedValue\_sigmaHistory}).

\begin{lemma}[Appending the output does not change the divergence]\label{lem:kl_output_pair}
Let $\mu, \mu'$ be finite measures on a measurable space $\mathcal H$ and $\rho$ a Markov kernel
from $\mathcal H$ to a measurable space $\mathcal D$. Then
$\KL(\mu\otimes\rho\,\|\,\mu'\otimes\rho) = \KL(\mu\|\mu')$. In the standing setting, for an
identification algorithm $A$ with sampling rule $\mathrm{alg}$, stopping rule $S$ and output
kernel $\rho$, and two runs (LML \texttt{IdentAlg.IsRun}) of $A$ in $\mathrm{env}$ on
$(\Omega, \Prob)$ with output $\mathrm{out}$ and in $\mathrm{env}'$ on $(\Omega', \Prob')$ with
output $\mathrm{out}'$,
\[
  \KL\big(\mathrm{law}_\nu(\hist_\tau, \mathrm{out})\,\big\|\,\mathrm{law}_{\nu'}(\hist'_{\tau'}, \mathrm{out}')\big)
  = \KL\big(\mathrm{law}_\nu(\hist_\tau)\,\big\|\,\mathrm{law}_{\nu'}(\hist'_{\tau'})\big),
\]
and $\KL(\mathrm{law}_\nu(\mathrm{out})\,\|\,\mathrm{law}_{\nu'}(\mathrm{out}')) \le$ the same
quantity.
\end{lemma}

\begin{proof}
The equality is Mathlib \texttt{klDiv\_compProd\_left} (same kernel on both sides). In a run,
the conditional law of the output given the stopped history is $\rho$
(\texttt{IsRun.hasCondDistrib\_output}), so the law of $(\hist_\tau, \mathrm{out})$ is
$\mathrm{law}_\nu(\hist_\tau)\otimes\rho$ (\texttt{HasCondDistrib.map\_eq}), and similarly under
$\Prob'$. The last claim is data processing (\texttt{klDiv\_map\_le}) along the second
projection.
\end{proof}

\textbf{Lean remark.} Sister project \texttt{Learning.IdentAlg.IsRun.klDiv\_map\_stoppedHist\_out}
and \texttt{IsRun.klDiv\_map\_out\_le\_stoppedHist} (\texttt{RunDivergence.lean}), to be copied
to \texttt{Essakine2026Tight/LeanMachineLearning/SequentialLearning/ChangeOfMeasure.lean}. No
finiteness of the stopping time is needed: when $\tau = \infty$ the stopped history is an
arbitrary element of $\mathcal H$ and the output is still drawn from $\rho$ applied to it.

\section{The change of measure}
\label{sec:pinfo_change_of_measure}

\begin{theorem}[Change of measure at a stopping time]\label{thm:change_of_measure}
\uses{def:klbin, lem:kl_stopped_hist_le, lem:kl_output_pair}
In the standing setting, let $A$ be an identification algorithm with sampling rule
$\mathrm{alg}$, stopping rule $S$ and output kernel $\rho$, and consider a run of $A$ in
$\mathrm{env}$ on $(\Omega,\Prob)$ and a run in $\mathrm{env}'$ on $(\Omega',\Prob')$, with
stopping times $\tau$, $\tau'$ and outputs $\mathrm{out}$, $\mathrm{out}'$. Assume that
$\tau < \infty$ $\Prob$-almost surely. Then for every measurable set $C \subseteq \mathcal H\times\mathcal D$,
with $E := \{(\hist_\tau, \mathrm{out}) \in C\}$ and $E' := \{(\hist'_{\tau'}, \mathrm{out}') \in C\}$,
\[
  \klbin\big(\Prob(E), \Prob'(E')\big) \le \sum_{a\in\As}\E_\nu[N_a(\tau)]\,\KL(\nu_a\|\nu'_a)
  \qquad\text{in } [0, \infty].
\]
Equivalently, $\sup_C \klbin(\Prob(E), \Prob'(E')) \le \sum_a \E_\nu[N_a(\tau)]\KL(\nu_a\|\nu'_a)$,
the form of Lemma~17 of \cite{essakine2026tight}. When moreover $\tau' < \infty$
$\Prob'$-almost surely, the same bound holds with $\klbin(\Prob(E), \Prob'(E'))$ replaced by
$\KL(\mathrm{law}_\nu(\hist_\tau,\mathrm{out})\,\|\,\mathrm{law}_{\nu'}(\hist'_{\tau'},\mathrm{out}'))$.
\end{theorem}

\begin{proof}
\uses{lem:klbin_le_kl, lem:kl_output_pair, lem:kl_stopped_hist_le, lem:pinfo_stopped_chain_rule, lem:pinfo_kl_restrict, lem:pinfo_klbin_convex_right, lem:pinfo_klbin_lsc, def:klbin}
Write $D := \sum_a\E_\nu[N_a(\tau)]\KL(\nu_a\|\nu'_a)$; if $D = \infty$ there is nothing to
prove, so assume $D < \infty$.

\emph{Two-sided case.} If $\tau' < \infty$ $\Prob'$-a.s., then by Lemma~\ref{lem:kl_output_pair}
and Lemma~\ref{lem:kl_stopped_hist_le},
$\KL(\mathrm{law}_\nu(\hist_\tau,\mathrm{out})\,\|\,\mathrm{law}_{\nu'}(\hist'_{\tau'},\mathrm{out}'))
= \KL(\mathrm{law}_\nu(\hist_\tau)\|\mathrm{law}_{\nu'}(\hist'_{\tau'})) \le D$, and
Lemma~\ref{lem:klbin_le_kl} applied to the two laws of the pair and the set $C$ gives the claim.

\emph{General case.} We reduce to bounded stopping times. Fix $M \in \N$. Let
$F_S := \{x \in \mathcal H : x \in S \text{ and no proper prefix of } x \text{ is in } S\}$ be the
(measurable) set of histories at which the rule fires; by definition of $\tau$,
$\hist_{\tau\wedge M} \in F_S$ iff $\tau \le M$, and then $\hist_{\tau\wedge M} = \hist_\tau$.
Consider the identification algorithm $A_M$ with the same sampling rule and output kernel and
the stopping rule $S_M := S \cup \{\text{length} = M\}$, whose stopping time is $\tau \wedge M$,
bounded, hence almost surely finite under both $\Prob$ and $\Prob'$. Runs of $A_M$ in
$\mathrm{env}$ and $\mathrm{env}'$ exist (LML's canonical runs on the Ionescu-Tulcea spaces),
with outputs $\mathrm{out}_M$, $\mathrm{out}'_M$, and the laws of their (stopped history,
output) pairs are $\mu_M\otimes\rho$ and $\mu'_M\otimes\rho$ with
$\mu_M := \mathrm{law}_\nu(\hist_{\tau\wedge M})$, $\mu'_M := \mathrm{law}_{\nu'}(\hist'_{\tau'\wedge M})$
(Lemma~\ref{lem:kl_output_pair}; these laws are determined by $A_M$ and the environments).
The two-sided case for $A_M$ and the sets $C_M := C \cap (F_S\times\mathcal D)$ and
$C_M^c := (\mathcal H\times\mathcal D \setminus C)\cap(F_S\times\mathcal D)$, together with
Lemma~\ref{lem:pinfo_stopped_chain_rule} and $N_a(\tau\wedge M) \le N_a(\tau)$, gives
\[
  \klbin(x_M, y_M) \le D \quad\text{and}\quad \klbin(\bar x_M, \bar y_M) \le D,
\]
where $x_M := (\mu_M\otimes\rho)(C_M) = \Prob(E, \tau \le M)$,
$y_M := (\mu'_M\otimes\rho)(C_M) = \Prob'(E', \tau'\le M)$,
$\bar x_M := \Prob(E^c, \tau\le M)$ and $\bar y_M := \Prob'(E'^c, \tau'\le M)$: indeed, on
$\{\tau \le M\}$ the pair $(\hist_{\tau\wedge M}, \cdot)$ has its first component in $F_S$ and
equals $(\hist_\tau, \cdot)$, so $(\mu_M\otimes\rho)(C_M) = \int_{F_S}\rho(x, C_x)\,\mu_M(dx)
= \int_{F_S}\rho(x, C_x)\,\mathrm{law}_\nu(\hist_\tau)(dx) = \Prob((\hist_\tau,\mathrm{out})\in C, \tau\le M)$
(with $C_x$ the section of $C$ at $x$; the law of $(\hist_\tau, \mathrm{out})$ is
$\mathrm{law}_\nu(\hist_\tau)\otimes\rho$ by Lemma~\ref{lem:kl_output_pair}, and
$\mu_M$ and $\mathrm{law}_\nu(\hist_\tau)$ agree on $F_S \cap \{\text{length}\le M\}$), and
likewise under $\Prob'$.

\emph{Passage to the limit.} As $M \to \infty$: $x_M \uparrow \Prob(E, \tau<\infty) = \Prob(E) =: x$
and $\bar x_M \uparrow \Prob(E^c) = 1 - x$ (almost sure finiteness under $\Prob$);
$y_M \uparrow \Prob'(E', \tau'<\infty) =: y_1$ and $\bar y_M \uparrow \Prob'(E'^c, \tau'<\infty)
= 1 - c - y_1$ with $c := \Prob'(\tau' = \infty)$. Lemma~\ref{lem:pinfo_klbin_lsc} gives
$\klbin(x, y_1) \le D$ and $\klbin(1 - x, 1 - c - y_1) \le D$, i.e.\ (by the symmetry
$\klbin(1-u, 1-v) = \klbin(u, v)$ of Definition~\ref{def:klbin}) $\klbin(x, y_1 + c) \le D$.
Finally $y := \Prob'(E') = y_1 + \Prob'(E', \tau' = \infty) \in [y_1, y_1 + c]$, so
Lemma~\ref{lem:pinfo_klbin_convex_right} gives $\klbin(x, y) \le \max\{\klbin(x,y_1), \klbin(x, y_1+c)\} \le D$.
\end{proof}

\textbf{Lean remark.} The statement is for two runs
\texttt{h : A.IsRun (stationaryEnv ν) O X Y out P}, \texttt{h' : A.IsRun (stationaryEnv ν')
O' X' Y' out' P'} with \texttt{[Fintype 𝓐]}, \texttt{[MeasurableSpace.CountablyGenerated 𝓨]},
\texttt{hτ : ∀ᵐ ω ∂P, A.stoppingTime O X Y ω ≠ ⊤} and a measurable
\texttt{C : Set ((Σ n, Hist Unit 𝓐 𝓨 n) × 𝓓)}; the conclusion is
\texttt{klDiv (bernoulliMeasure true false (P.real E)) (bernoulliMeasure true false (P'.real E'))
≤ ∑ a, (∫⁻ ω, (pullCount X a (A.stoppingTime O X Y ω).toNat ω : ℝ≥0∞) ∂P) * klDiv (ν a) (ν' a)},
with the real form \texttt{klBin (P.real E) (P'.real E') ≤ (…).toReal} when
$\Prob'(E') \in (0,1)$ and the right-hand side is finite. The two-sided case is
\texttt{IsRun.klDiv\_map\_out\_le\_lintegral} of the sister project (there for the output alone;
the pair version is the same proof through \texttt{IsRun.klDiv\_map\_stoppedHist\_out}), and
the reduction to $\tau\wedge M$ uses the identification algorithm
\texttt{\{A with stopSet := A.stopSet ∪ \{x | x.1 = M\}\}}, whose stopping time is
\texttt{min (A.stoppingTime O X Y) M} (LML \texttt{hittingAfter} of a union), and the canonical
runs \texttt{IdentAlg.isRun\_canonical} on \texttt{trajMeasure}. File:
\texttt{Essakine2026Tight/LeanMachineLearning/SequentialLearning/ChangeOfMeasure.lean}.

\begin{remark}[The paper's Lemma~17]\label{rem:pinfo_paper_lemma}
\uses{thm:change_of_measure, def:episode_env}
Lemma~17 of \cite{essakine2026tight} is Lemma~1 of \cite{kaufmann2016complexity}: for two
bandit models $\nu, \nu'$ with $K$ arms whose arm distributions are mutually absolutely
continuous, an algorithm $((A_t)_t, \tau, \hat S)$ and an almost surely finite stopping time
$\sigma$, $\sum_a \E_\nu[N_a(\sigma)]\KL(\nu_a\|\nu'_a) \ge \sup_{\mathcal E\in\mathcal F_\sigma}
\klbin(\Prob_\nu(\mathcal E), \Prob_{\nu'}(\mathcal E))$. This is
Theorem~\ref{thm:change_of_measure} with $\As = [K]$, $\mathcal Y = \R$, the stationary
environments of the two bandit models, $S$ the stopping rule of $\sigma$, and $\mathcal E$ an
event of $(\hist_\sigma, \hat S)$: the events of $\mathcal F_\sigma$ (the stopped
$\sigma$-algebra) are exactly the events of the stopped history on $\{\sigma < \infty\}$, and the
recommendation is a function of the stopped history (or of it and independent randomness, which
is the output kernel $\rho$). Mutual absolute continuity is not needed: if some
$\KL(\nu_a\|\nu'_a)$ is infinite the bound is trivial unless $\E_\nu[N_a(\sigma)] = 0$, in which
case the convention $0\cdot\infty = 0$ is the right one (an arm that is never pulled does not
help to distinguish the models). The paper applies the lemma to the episode environment of the
hard instances (Definition~\ref{def:episode_env}): the ``arms'' are the policies
$\pi \in \As^{[H]\times\Ss}$, the ``reward'' of an episode is its state sequence, and
$\KL(\nu_\pi\|\nu'_\pi)$ is the divergence of the two laws of the state sequence under $\pi$
(lem:hard\_kl\_episode in Chapter~\ref{chap:lower_bound}); the stopping time is the number of
episodes $\tau$ of the BPI algorithm and the event is ``the output realizes $u$''.
\end{remark}

\begin{remark}[One-sided finiteness]\label{rem:pinfo_one_sided}
\uses{thm:change_of_measure}
Theorem~\ref{thm:change_of_measure} only assumes that $\tau$ is almost surely finite under
the \emph{first} environment. This matters for the lower bound: the hypothesis
$\E_{\mathcal M_0}[\tau] < \infty$ (Chapter~\ref{chap:lower_bound}) says nothing about the
stopping time under the alternative instances $\mathcal M_u$, and almost sure finiteness under
$\nu$ does not transfer to $\nu'$ even when all one-step divergences are finite (the laws of
the whole trajectories are typically mutually singular). The proof of
\cite{kaufmann2016complexity}, through the likelihood ratio $L_\sigma$ and the identity
$\Prob_{\nu'}(\mathcal E) = \E_\nu[\indic_{\mathcal E}/L_\sigma]$, silently uses
$\mathcal E \subseteq \{\sigma < \infty\}$ under $\nu'$; the truncation argument above (bounded
stopping times $\tau\wedge M$, then the convexity of $\klbin(x,\cdot)$, which handles the
mass $\Prob'(\tau' = \infty)$ that the events $E'$ may or may not contain) removes this
restriction. The divergence identity of Lemma~\ref{lem:kl_stopped_hist_le} itself does need
finiteness under both environments, because the stopped history is not defined on
$\{\tau = \infty\}$.
\end{remark}

\chapter{Elementary inequalities}
\label{chap:pre_analysis}

This chapter collects the real-analysis facts used by the analysis of Entropic-BPI
(Chapters~\ref{chap:analysis_pos}, \ref{chap:analysis_neg} and~\ref{chap:sample_complexity}):
two square-root inequalities used to split the concentration terms, the bound
$(1 + a/H)^h \le e^a$ that removes the factor $\kappa^{h-1}$ of the unrolled certificate
recursion, the Bhatia--Davis inequality and the normalized variance bound for exponential
transforms (the paper's Lemma~24, \cite{essakine2026tight}), the logarithmic sum bound (the
paper's Lemma~28) used by the counting argument, and the self-bounding inequality (the paper's
Lemma~29, Lemma~13 of \cite{menard2021fast}) that turns the recursive bound on the number of
episodes into an explicit one. Everything is stated over $\R$ with no reference to the MDP.

\textbf{Constants.} Every constant is verified below. Two of the paper's statements are
adjusted: the logarithmic sum bound (Lemma~\ref{lem:log_sum_bound}) holds with the constant
$3$ in place of $4$ (both forms are stated, since the form with $4$ is the one used by the
counting argument); and the self-bounding inequality
(Lemma~\ref{lem:self_bounding}) is \emph{false} with the constant $C_1$ printed in the paper
(a counterexample is given in Remark~\ref{rem:pana_self_bounding_paper}), so it is stated and
proved with a corrected $C_1$.

\textbf{What Mathlib provides.} The real square root \texttt{Real.sqrt} with
\texttt{Real.sqrt\_le\_sqrt}, \texttt{Real.sqrt\_le\_left}, \texttt{Real.le\_sqrt},
\texttt{Real.sq\_sqrt}, \texttt{Real.sqrt\_mul}; the exponential and logarithm with
\texttt{Real.add\_one\_le\_exp}, \texttt{Real.log\_le\_sub\_one\_of\_pos},
\texttt{Real.one\_sub\_inv\_le\_log\_of\_pos}, \texttt{Real.log\_le\_log\_iff},
\texttt{Real.log\_rpow}, \texttt{Real.rpow\_natCast}; the variance
\texttt{ProbabilityTheory.variance} with \texttt{variance\_eq\_sub} (variance as
$\E X^2 - (\E X)^2$ for $X \in L^2$), \texttt{memLp\_of\_bounded} and the Bhatia--Davis
inequality itself (\texttt{ProbabilityTheory.variance\_le\_sub\_mul\_sub},
Lemma~\ref{lem:bhatia_davis}); the AM--GM inequality in the form \texttt{two\_mul\_le\_add\_sq};
the weak logarithm-versus-power bound \texttt{Real.log\_le\_rpow\_div}.

\textbf{What is missing.} The normalized variance bound, the logarithmic sum bound and the
self-bounding inequality, and the small lemmas around them. They are in
\texttt{Essakine2026Tight/Mathlib/} (\texttt{Algebra/Order/Field/Basic.lean},
\texttt{Analysis/Real/Sqrt.lean}, \texttt{Analysis/SpecialFunctions/Exp.lean},
\texttt{Analysis/SpecialFunctions/Log/Basic.lean}, \texttt{Analysis/SpecialFunctions/Pow/Real.lean},
\texttt{Analysis/SpecialFunctions/Log/SelfBounding.lean},
\texttt{Probability/Moments/Variance.lean}).

\section{Square roots, exponentials and logarithms}
\label{sec:pre_analysis_basic}

\begin{lemma}[Subadditivity of the square root]\label{lem:sqrt_add_le}
\lean{Real.sqrt_add_le}\leanok
For $a, b \ge 0$, $\sqrt{a + b} \le \sqrt a + \sqrt b$.
\end{lemma}

\begin{proof}
\leanok
Both sides are nonnegative, so it suffices to compare squares (Mathlib
\texttt{Real.sqrt\_le\_left} with $y = \sqrt a + \sqrt b \ge 0$):
$(\sqrt a + \sqrt b)^2 = a + b + 2\sqrt a\sqrt b \ge a + b$ because $\sqrt a\sqrt b \ge 0$
(\texttt{Real.sq\_sqrt}, \texttt{Real.sqrt\_nonneg}).
\end{proof}

\begin{lemma}[Weighted AM--GM for a square root]\label{lem:sqrt_mul_le}
\lean{Real.two_mul_sqrt_mul_le, Real.sqrt_mul_le_half_add_half, Real.sqrt_mul_le_add}\leanok
For $a, b \ge 0$ and $\eta > 0$,
\[
  \sqrt{ab} \le \frac{\eta a}{2} + \frac{b}{2\eta} \le \eta a + \frac{b}{\eta} .
\]
\end{lemma}

\begin{proof}
\leanok
Let $x := \sqrt{\eta a}$ and $y := \sqrt{b/\eta}$, so that $x, y \ge 0$, $x^2 = \eta a$,
$y^2 = b/\eta$ and $xy = \sqrt{ab}$ (\texttt{Real.sqrt\_mul}, \texttt{Real.sq\_sqrt}; $\eta a \ge 0$
and $b/\eta \ge 0$). From $(x - y)^2 \ge 0$ we get $2xy \le x^2 + y^2$
(\texttt{two\_mul\_le\_add\_sq}), i.e.\ $2\sqrt{ab} \le \eta a + b/\eta$
(\texttt{Real.two\_mul\_sqrt\_mul\_le}), which is the first
inequality. The second holds because $\eta a \ge 0$ and $b/\eta \ge 0$.
\end{proof}

\begin{lemma}[A power bound for $(1 + a/H)^h$]\label{lem:one_add_div_pow_le_exp}
\lean{Real.one_add_div_pow_le_exp}\leanok
For $a \ge 0$, integers $H \ge 1$ and $0 \le h \le H$: $(1 + a/H)^h \le e^{a}$.
\end{lemma}

\begin{proof}
\leanok
By \texttt{Real.add\_one\_le\_exp}, $1 + a/H \le e^{a/H}$, and $1 + a/H \ge 0$. Since
$x \mapsto x^h$ is monotone on $[0, \infty)$ (\texttt{pow\_le\_pow\_left₀}),
$(1 + a/H)^h \le (e^{a/H})^h = e^{ah/H}$ (\texttt{Real.exp\_nat\_mul}), and $ah/H \le a$
because $0 \le h \le H$ and $a \ge 0$, so $e^{ah/H} \le e^{a}$ (\texttt{Real.exp\_le\_exp}).
\end{proof}

\begin{lemma}[{$2 - x \le 1/x$ on $(0, \infty)$}]\label{lem:two_sub_le_inv}
\lean{two_sub_le_inv}\leanok
For $x > 0$ (in any linearly ordered field), $2 - x \le 1/x$. (It is used for
$x = e^{\beta\eps} \in (0, 1]$.)
\end{lemma}

\begin{proof}
\leanok
$1/x - (2 - x) = (1 - x)^2/x \ge 0$ (\texttt{field\_simp}, \texttt{positivity}).
\end{proof}

\section{The Bhatia--Davis inequality and the normalized variance bound}
\label{sec:pre_analysis_variance}

\begin{lemma}[Bhatia--Davis inequality]\label{lem:bhatia_davis}
\lean{ProbabilityTheory.variance_le_sub_mul_sub}\leanok
Let $(\Omega, \Prob)$ be a probability space, $Y : \Omega \to \R$ an a.e.-measurable random
variable and $m, M$ real numbers with $Y \in [m, M]$ almost surely. Then, with $\mu := \E Y$,
\[
  \Var(Y) \le (M - \mu)(\mu - m) .
\]
\end{lemma}

\begin{proof}
\leanok
This is Mathlib's \texttt{ProbabilityTheory.variance\_le\_sub\_mul\_sub}. The argument:
$Y$ is bounded a.s.\ on a probability space, hence in $L^2$
(\texttt{memLp\_of\_bounded}), and $\Var(Y) = \E Y^2 - \mu^2$ (\texttt{variance\_eq\_sub}).
Almost surely $(M - Y)(Y - m) \ge 0$, so
\[
  0 \le \E[(M - Y)(Y - m)] = \E[-Y^2 + (M + m)Y - Mm] = -\E Y^2 + (M + m)\mu - Mm ,
\]
i.e.\ $\E Y^2 \le (M + m)\mu - Mm$. Therefore
$\Var(Y) = \E Y^2 - \mu^2 \le (M + m)\mu - Mm - \mu^2 = (M - \mu)(\mu - m)$.
This is the inequality of \cite{bhatia2000better}.
\end{proof}

\begin{lemma}[The mean of a bounded random variable]\label{lem:pana_integral_mem_Icc}
\lean{ProbabilityTheory.integral_mem_Icc_of_ae_mem_Icc}\leanok
Let $(\Omega, \Prob)$ be a probability space and $Y : \Omega \to \R$ a.e.-strongly measurable
with $Y \in [m, M]$ almost surely. Then $\E Y \in [m, M]$.
\end{lemma}

\begin{proof}
\leanok
$Y$ is integrable (\texttt{memLp\_of\_bounded} with $p = 1$), and the integral is monotone
(\texttt{integral\_mono\_ae}) with $\E[m] = m$, $\E[M] = M$ on a probability space.
\end{proof}

\textbf{Lean remark.} The finite-vector version $\Var_p(f) \le (M - pf)(pf - m)$ for a
probability vector $p$ on a finite set and $f \in [m, M]$ is the same statement for the measure
$\sum_s p(s)\delta_s$ (\texttt{weightedMeasure}), or is proved directly by the same computation
with finite sums.

\begin{lemma}[Maximization of the normalized Bhatia--Davis bound]\label{lem:pana_var_ratio_max}
\lean{sub_mul_sub_div_sq_le}\leanok
For $m, M, \mu > 0$ (in any linearly ordered field),
\[
  \frac{(M - \mu)(\mu - m)}{\mu^2} \le \frac{(M - m)^2}{4Mm}
  = \frac{(M/m - 1)^2}{4\,M/m},
\]
with equality if and only if $\mu = 2Mm/(M + m)$.
\end{lemma}

\begin{proof}
\leanok
Since $\mu > 0$ and $4Mm > 0$, the inequality is equivalent to
$4Mm(M - \mu)(\mu - m) \le (M - m)^2\mu^2$ (\texttt{div\_le\_div\_iff₀}). Expanding both sides,
\[
  (M - m)^2\mu^2 - 4Mm(M - \mu)(\mu - m)
  = (M + m)^2\mu^2 - 4Mm(M + m)\mu + 4M^2m^2
  = \big((M + m)\mu - 2Mm\big)^2 \ge 0 ,
\]
where we used $(M - m)^2 + 4Mm = (M + m)^2$ (\texttt{nlinarith} with this square as hint).
Equality holds iff $(M + m)\mu = 2Mm$. The last
expression of the statement is obtained by dividing numerator and denominator by $m^2$.
(This is the calculus maximization of the paper's proof of Lemma~24 made algebraic: the
function $\mu \mapsto (M - \mu)(\mu - m)/\mu^2$ increases on $[m, 2Mm/(M + m)]$ and decreases
on $[2Mm/(M + m), M]$, with maximum $(M - m)^2/(4Mm)$. The restriction $\mu \in [m, M]$ is not
needed.)
\end{proof}

\begin{lemma}[Normalized variance bound for an exponential transform]\label{lem:normalized_variance_bound}
\uses{lem:bhatia_davis, lem:pana_var_ratio_max, lem:pana_integral_mem_Icc}
\lean{ProbabilityTheory.variance_exp_div_sq_integral_le}\leanok
Let $(\Omega, \Prob)$ be a probability space, $R \ge 0$, $\beta \in \R$ and $f : \Omega \to \R$
an a.e.-measurable random variable with $f \in [0, R]$ almost surely. Let $Y := e^{\beta f}$
and $\mu := \E Y$. Then
\[
  \frac{\Var(Y)}{\mu^2} \le \frac{(e^{\abs\beta R} - 1)^2}{4\,e^{\abs\beta R}} .
\]
\end{lemma}

\begin{proof}
\uses{lem:bhatia_davis, lem:pana_var_ratio_max, lem:pana_integral_mem_Icc}
\leanok
Let $m := \min\{1, e^{\beta R}\}$ and $M := \max\{1, e^{\beta R}\}$.
Almost surely $\beta f$ lies between $0$ and $\beta R$ (if $\beta \ge 0$ then
$0 \le \beta f \le \beta R$; if $\beta < 0$ then $\beta R \le \beta f \le 0$), so by
monotonicity of $\exp$, $m \le Y \le M$ a.s.\ (\texttt{ProbabilityTheory.exp\_mul\_mem\_Icc\_min\_max}).
Lemma~\ref{lem:pana_integral_mem_Icc} gives
$\mu \in [m, M]$, and $m > 0$, hence $\mu > 0$; Lemma~\ref{lem:bhatia_davis} gives
$\Var(Y) \le (M - \mu)(\mu - m)$.
Dividing by $\mu^2$ and applying Lemma~\ref{lem:pana_var_ratio_max},
\[
  \frac{\Var(Y)}{\mu^2} \le \frac{(M - \mu)(\mu - m)}{\mu^2} \le \frac{(M - m)^2}{4Mm} .
\]
Finally $(M - m)^2/(4Mm) = (e^{\abs\beta R} - 1)^2/(4e^{\abs\beta R})$
(\texttt{ProbabilityTheory.max\_sub\_min\_sq\_div\_eq}): if $\beta \ge 0$ then $m = 1$,
$M = e^{\beta R}$; if $\beta < 0$ then $m = e^{\beta R}$, $M = 1$ and
$(1 - e^{\beta R})^2/(4e^{\beta R}) = (e^{-\beta R} - 1)^2/(4e^{-\beta R})$.
\end{proof}

\textbf{Lean remark.} This is the paper's Lemma~24 (Appendix~F). The bound is used in
Lemma~\ref{lem:variance_ratio_le} with $f$ the return of an episode and $R = \Gmax$; the
paper's constant $V_G = (e^{\abs\beta G} - 1)^2/e^{\abs\beta G}$ (Definition~\ref{def:upper_bound_const})
is four times the bound above, so the factor $4$ is slack that the formal proof may drop.
The case split on the sign of $\beta$ only appears in the two auxiliary lemmas (the a.s.\ bounds
on $Y$ and the identification of $(M - m)^2/(4Mm)$).

\section{The logarithmic sum bound}
\label{sec:pre_analysis_log_sum}

\begin{lemma}[One step of the logarithmic sum bound]\label{lem:pana_log_sum_step}
\lean{Real.div_max_one_le_three_mul_log}\leanok
For $U \ge 0$ and $u \in [0, 1]$,
\[
  \frac{u}{\max\{U, 1\}} \le 3\log\frac{U + u + 1}{U + 1} .
\]
\end{lemma}

\begin{proof}
\leanok
Let $x := u/(U + 1) \ge 0$. Since $1 + x > 0$, Mathlib \texttt{Real.one\_sub\_inv\_le\_log\_of\_pos}
gives $\log(1 + x) \ge 1 - \frac{1}{1 + x} = \frac{x}{1 + x} = \frac{u}{U + 1 + u}$, so the
right-hand side is at least $3u/(U + 1 + u)$. It
suffices to show $u(U + 1 + u) \le 3u\max\{U, 1\}$, which follows from $u \ge 0$ and
$U + 1 + u \le 3\max\{U, 1\}$: if $U \ge 1$ then
$U + 1 + u \le U + 2 \le 3U$; if $U < 1$ then $U + 1 + u < 3 = 3\max\{U, 1\}$. Hence
$\frac{u}{\max\{U,1\}} \le \frac{3u}{U + 1 + u} \le 3\log\frac{U + u + 1}{U + 1}$.
\end{proof}

\begin{lemma}[Logarithmic sum bound]\label{lem:log_sum_bound}
\uses{lem:pana_log_sum_step}
\lean{Real.sum_div_max_one_le_three_mul_log, Real.sum_div_max_one_le_four_mul_log}\leanok
Let $(u_t)_{t \ge 1}$ be a sequence with $u_t \in [0, 1]$ and $U_t := \sum_{i = 1}^t u_i$
($U_0 = 0$). For every integer $T \ge 0$,
\[
  \sum_{t = 0}^{T - 1}\frac{u_{t + 1}}{\max\{U_t, 1\}} \le 3\log(U_{T} + 1) \le 4\log(U_{T} + 1) .
\]
\end{lemma}

\begin{proof}
\uses{lem:pana_log_sum_step}
\leanok
For each $t$, $U_t \ge 0$, $u_{t+1} \in [0,1]$ and $U_{t+1} = U_t + u_{t+1}$, so
Lemma~\ref{lem:pana_log_sum_step} gives
\[
  \frac{u_{t + 1}}{\max\{U_t, 1\}} \le 3\log\frac{U_{t+1} + 1}{U_t + 1}
  = 3\big(\log(U_{t+1} + 1) - \log(U_t + 1)\big)
\]
(\texttt{Real.log\_div}, both arguments being positive). Summing over $t = 0, \dots, T - 1$, the
right-hand side telescopes (\texttt{Finset.sum\_range\_sub}) to
$3(\log(U_{T} + 1) - \log(U_0 + 1)) = 3\log(U_{T} + 1)$. Finally $U_{T} + 1 \ge 1$, so
$\log(U_{T} + 1) \ge 0$ (\texttt{Real.log\_nonneg}) and $3\log(\cdot) \le 4\log(\cdot)$.
\end{proof}

\begin{remark}[On the constant]\label{rem:pana_log_sum_const}
This is the paper's Lemma~28 (Appendix~F), stated there with the constant $4$ and without
proof (it is Lemma~8 of \cite{menard2021fast}). The one-step bound holds with the constant
$3$, hence so does the sum bound; the constant $2$ fails at $U = u = 1$ ($1 > 2\log(3/2)$).
The form with the constant $4$ is the one consumed by
Lemma~\ref{lem:sum_counts_bound}. \textbf{Lean remark.} The sequence is indexed by $\N$ with
$u_{t+1} = $ \texttt{u t}, so that $U_t = $ \texttt{∑ i ∈ Finset.range t, u i} and the sum is
\texttt{∑ t ∈ Finset.range T, u t / max (∑ i ∈ Finset.range t, u i) 1}.
\end{remark}

\section{The self-bounding inequality}
\label{sec:pre_analysis_self_bounding}

The counting argument of Chapter~\ref{chap:sample_complexity} produces an inequality of the
form $\tau \le C\sqrt{\tau\,\ell(\tau)} + D\,\ell'(\tau)$ with $\ell, \ell'$ polylogarithmic in
$\tau$, and one must solve it for $\tau$. We first solve the quadratic part, then bound the
logarithm of a quantity that is bounded by a multiple of the square of its own logarithm.

\begin{lemma}[Root of a quadratic inequality]\label{lem:pana_quadratic_root}
\lean{Real.le_add_sqrt_of_sq_le}\leanok
Let $b, c \ge 0$ and $x$ real with $x^2 \le bx + c$. Then $x \le b + \sqrt c$.
\end{lemma}

\begin{proof}
\leanok
Suppose $x > b + \sqrt c$. Then $x > 0$ (as $b + \sqrt c \ge 0$), $x > b$ and $x > \sqrt c \ge 0$,
so
\[
  x^2 = x\cdot x > x(b + \sqrt c) = bx + x\sqrt c \ge bx + \sqrt c\sqrt c = bx + c ,
\]
using $x\sqrt c \ge \sqrt c\sqrt c$ (from $x \ge \sqrt c$ and $\sqrt c \ge 0$) and
$\sqrt c\sqrt c = c$ (\texttt{Real.mul\_self\_sqrt}). This contradicts $x^2 \le bx + c$.
(No sign assumption on $x$ is needed.)
\end{proof}

\begin{lemma}[Logarithm versus powers]\label{lem:pana_log_le_rpow}
\lean{Real.log_le_rpow_div_exp_one_mul}\leanok
For $y > 0$ and $\gamma > 0$, $\log y \le \dfrac{y^{\gamma}}{e\gamma}$.
\end{lemma}

\begin{proof}
\leanok
Let $z := y^\gamma > 0$. Mathlib \texttt{Real.log\_le\_sub\_one\_of\_pos} applied to $z/e > 0$
gives $\log(z/e) \le z/e - 1$, and $\log(z/e) = \log z - 1$ (\texttt{Real.log\_div},
\texttt{Real.log\_exp}), so $\log z \le z/e$. Since $\log z = \gamma\log y$
(\texttt{Real.log\_rpow}), $\gamma\log y \le y^\gamma/e$; divide by $\gamma > 0$.
(Mathlib's \texttt{Real.log\_le\_rpow\_div} is the same bound without the factor $1/e$.)
\end{proof}

\begin{lemma}[A quantity bounded by the square of its logarithm]\label{lem:pana_log_self_bound}
\uses{lem:pana_log_le_rpow}
\lean{Real.log_le_of_le_mul_sq_log}\leanok
Let $y \ge 1$ and $M > 0$ with $y \le M(\log y)^2$. Then
\[
  \log y \le \tfrac85\log(4M) .
\]
\end{lemma}

\begin{proof}
\uses{lem:pana_log_le_rpow}
\leanok
Since $y \ge 1$, $\log y \ge 0$, and Lemma~\ref{lem:pana_log_le_rpow} with $\gamma = 3/16$ gives
$0 \le \log y \le \frac{16}{3e}\,y^{3/16}$. Squaring (both sides are nonnegative,
\texttt{pow\_le\_pow\_left₀}) and using $y^{3/16}\cdot y^{3/16} = y^{3/8}$ (\texttt{Real.rpow\_add}),
\[
  (\log y)^2 \le \frac{256}{9e^2}\,y^{3/8} \le 4\,y^{3/8},
\]
because $9e^2 \ge 64$ (indeed $e \ge 2.7$, \texttt{Real.exp\_one\_gt\_d9}, so $e^2 \ge 7.29 > 64/9$).
Hence $y \le 4M\,y^{3/8}$, and since $y = y^{5/8}\,y^{3/8}$ with $y^{3/8} > 0$
(\texttt{Real.rpow\_add}, \texttt{le\_of\_mul\_le\_mul\_right}), $y^{5/8} \le 4M$. Both sides are
positive, so taking logarithms (\texttt{Real.log\_le\_log}, \texttt{Real.log\_rpow}) yields
$\frac58\log y \le \log(4M)$, i.e.\ $\log y \le \frac85\log(4M)$.
\end{proof}

\begin{lemma}[Reduction of the self-bounding inequality to a quadratic one]\label{lem:pana_self_bounding_quadratic}
\uses{lem:pana_quadratic_root}
\lean{Real.le_mul_sq_log_of_le_sqrt_mul_add}\leanok
Let $A, B, C, D \ge 0$, $E \ge B$, $\alpha \in \R$ and $\tau \ge 0$ with $L := \log(\alpha\tau) \ge 1$
satisfy
\begin{equation}\label{eq:pana_self_bounding_hyp}
  \tau \le C\sqrt{\tau\big(A\log(\alpha\tau) + B\log(\alpha\tau)^2\big)}
  + D\big(A\log(\alpha\tau) + E\log(\alpha\tau)^2\big) .
\end{equation}
Let $K := C^2(A + B) + (D + 2\sqrt D\,C)(A + E)$. Then $\tau \le K L^2$.
\end{lemma}

\begin{proof}
\uses{lem:pana_quadratic_root}
\leanok
\begin{enumerate}
  \item[(i)] \emph{Reduction to a quadratic inequality.} Since $L \ge 1$ and $A \ge 0$,
        $AL \le AL^2$, hence $AL + BL^2 \le (A + B)L^2$ and $AL + EL^2 \le (A + E)L^2$. Both
        terms of~\eqref{eq:pana_self_bounding_hyp} are monotone in these quantities ($C, D \ge 0$,
        $\sqrt{\cdot}$ monotone), and $\sqrt{\tau(A + B)L^2} = L\sqrt{A+B}\sqrt\tau$ as $L \ge 0$
        (\texttt{Real.sqrt\_mul}, \texttt{Real.sqrt\_sq}); so
        $\tau \le CL\sqrt{A + B}\,\sqrt\tau + D(A + E)L^2$.
  \item[(ii)] \emph{Solving it.} Put $x := \sqrt\tau$, $b := CL\sqrt{A+B} \ge 0$ and
        $c := D(A+E)L^2 \ge 0$; then $x^2 = \tau \le bx + c$ and Lemma~\ref{lem:pana_quadratic_root}
        gives $x \le b + \sqrt c = CL\sqrt{A+B} + \sqrt D\sqrt{A+E}\,L$. Squaring (both sides
        nonnegative),
        \[
          \tau \le L^2\Big(C^2(A+B) + D(A+E)\Big) + 2C\sqrt D\,L^2\,\sqrt{A+B}\sqrt{A+E} .
        \]
        As $B \le E$, $\sqrt{A+B} \le \sqrt{A+E}$, so $\sqrt{A+B}\sqrt{A+E} \le A + E$ and
        the last term is at most $2C\sqrt D\,(A+E)L^2$. Hence $\tau \le KL^2$.
\end{enumerate}
\end{proof}

\begin{lemma}[The logarithm of the self-bounding inequality]\label{lem:pana_self_bounding_log}
\uses{lem:pana_self_bounding_quadratic, lem:pana_log_self_bound}
\lean{Real.one_le_log_mul_of_exp_one_le, Real.log_le_of_le_sqrt_mul_add}\leanok
Let $A, B, C, D \ge 0$, $E \ge B$, $\alpha \ge e$ and $\tau > 1$
satisfy~\eqref{eq:pana_self_bounding_hyp}, and let $K$ be as in
Lemma~\ref{lem:pana_self_bounding_quadratic}. Then $L := \log(\alpha\tau) \ge 1$, $K > 0$ and
$L \le \frac85\log(4\alpha K)$.
\end{lemma}

\begin{proof}
\uses{lem:pana_self_bounding_quadratic, lem:pana_log_self_bound}
\leanok
$\alpha\tau \ge \alpha \ge e$, so $L \ge 1$ (\texttt{Real.le\_log\_iff\_exp\_le};
\texttt{Real.one\_le\_log\_mul\_of\_exp\_one\_le}). Lemma~\ref{lem:pana_self_bounding_quadratic}
gives $\tau \le KL^2$; since $\tau > 0$, this forces $K > 0$. Let $y := \alpha\tau \ge 1$ and
$M := \alpha K > 0$. Multiplying $\tau \le KL^2$ by $\alpha$ gives $y \le M(\log y)^2$, and
Lemma~\ref{lem:pana_log_self_bound} yields $L = \log y \le \frac85\log(4M) = \frac85\log(4\alpha K)$.
\end{proof}

\begin{lemma}[Self-bounding inequality, sharp form]\label{lem:pana_self_bounding_core}
\uses{lem:pana_self_bounding_quadratic, lem:pana_self_bounding_log}
\lean{Real.le_mul_sq_log_add_one_of_le_sqrt_mul_add}\leanok
Let $A, B, C, D \ge 0$, $E \ge B$, $\alpha \ge e$ and $\tau \ge 0$
satisfy~\eqref{eq:pana_self_bounding_hyp}, let $K := C^2(A + B) + (D + 2\sqrt D\,C)(A + E)$ and
let $M \ge 4\alpha K$. Then
\[
  \tau \le K\Big(\tfrac85\log M\Big)^2 + 1 .
\]
\end{lemma}

\begin{proof}
\uses{lem:pana_self_bounding_quadratic, lem:pana_self_bounding_log}
\leanok
If $\tau \le 1$ the claim holds because $K \ge 0$. Assume $\tau > 1$. By
Lemma~\ref{lem:pana_self_bounding_log}, $K > 0$, $L := \log(\alpha\tau) \ge 1$ and
$L \le \frac85\log(4\alpha K) \le \frac85\log M$ ($0 < 4\alpha K \le M$ and $\log$ is monotone);
by Lemma~\ref{lem:pana_self_bounding_quadratic}, $\tau \le KL^2$.
Then $0 \le L \le \frac85\log M$ gives $L^2 \le (\frac85\log M)^2$, and
$\tau \le KL^2 \le K(\frac85\log M)^2 \le K(\frac85\log M)^2 + 1$.
\end{proof}

\begin{lemma}[The constant of the self-bounding inequality]\label{lem:pana_four_alpha_K_le}
\lean{Real.four_mul_mul_le}\leanok
Let $A, B, C, D \ge 0$, $E \ge B$ and $\alpha \ge e$, and $K := C^2(A + B) + (D + 2\sqrt D\,C)(A + E)$.
Then $4\alpha K \le 11\alpha^2(A + E)(C + \sqrt D)^2$.
\end{lemma}

\begin{proof}
\leanok
$K \le C^2(A+E) + (D + 2\sqrt D C)(A+E) = (A+E)\big(C^2 + 2C\sqrt D + D\big) = (A+E)(C + \sqrt D)^2$
(using $B \le E$ and $\sqrt D^2 = D$), and $4\alpha \le 11\alpha^2$ because $\alpha \ge e \ge 4/11$.
Both factors are nonnegative, so the products compare.
\end{proof}

\begin{lemma}[Self-bounding inequality]\label{lem:self_bounding}
\uses{lem:pana_self_bounding_core, lem:pana_four_alpha_K_le}
\lean{Real.le_of_le_sqrt_mul_add, Real.le_of_le_sqrt_mul_add'}\leanok
Let $A, B, C, D \ge 0$, $E \ge B$, $\alpha \ge e$ and $\tau \ge 0$ satisfy
\[
  \tau \le C\sqrt{\tau\big(A\log(\alpha\tau) + B\log(\alpha\tau)^2\big)}
  + D\big(A\log(\alpha\tau) + E\log(\alpha\tau)^2\big) .
\]
Then
\[
  \tau \le C^2(A + B)\,C_1^2 + \big(D + 2\sqrt D\,C\big)(A + E)\,C_1^2 + 1,
  \qquad
  C_1 := \frac85\log\Big(11\,\alpha^2\,(A + E)\,\big(C + \sqrt D\big)^2\Big) .
\]
If moreover $D \ge 1$, the same holds with $C_1 := \frac85\log\big(11\alpha^2(A+E)(C + D)^2\big)$.
\end{lemma}

\begin{proof}
\uses{lem:pana_self_bounding_core, lem:pana_four_alpha_K_le}
\leanok
Let $K := C^2(A+B) + (D + 2\sqrt D C)(A+E)$, so that the claimed bound reads
$\tau \le KC_1^2 + 1$, which is Lemma~\ref{lem:pana_self_bounding_core} with
$M := 11\alpha^2(A+E)(C+\sqrt D)^2 \ge 4\alpha K$ (Lemma~\ref{lem:pana_four_alpha_K_le}).
If $D \ge 1$ then $\sqrt D \le D$ (\texttt{Real.sqrt\_le\_left}), so
$(C + \sqrt D)^2 \le (C + D)^2$ and $M' := 11\alpha^2(A+E)(C+D)^2 \ge M \ge 4\alpha K$ also
qualifies.
\end{proof}

\begin{remark}[The paper's constant is wrong]\label{rem:pana_self_bounding_paper}
The paper's Lemma~29 (Appendix~F, quoted from Lemma~13 of \cite{menard2021fast}) states the
conclusion with $C_1 = \frac85\log(11\alpha^2(A+E)(C+D))$, i.e.\ with $(C + D)$ instead of
$(C + \sqrt D)^2$, for positive $A, B, C, D, E$ with $1 \le B \le E$. That statement is false.
Take $\alpha = e$, $B = E = 1$, $C = 1000$, $A = D = 10^{-6}$ and $\tau = 4\cdot 10^8$. Then
$L = \log(\alpha\tau) \approx 20.81$ and the right-hand side of the hypothesis is
$C\sqrt{\tau(AL + L^2)} + D(AL + L^2) \approx 1000\cdot\sqrt{4\cdot10^8}\cdot 20.81
\approx 4.16\cdot 10^8 \ge \tau$, so the hypothesis holds; but the paper's
$C_1 = \frac85\log(11e^2(1 + 10^{-6})(1000 + 10^{-6})) \approx 18.09$ gives the bound
$C^2(A+B)C_1^2 + (D + 2\sqrt D C)(A+E)C_1^2 + 1 \approx 10^6\cdot 327.2 + 655 \approx 3.27\cdot 10^8 < \tau$.
(The largest $\tau$ satisfying the hypothesis is about $4.37\cdot 10^8$; the corrected $C_1$
of Lemma~\ref{lem:self_bounding} is $\approx 29.1$ and gives the bound $\approx 8.5\cdot10^8$.)
The defect is in the exponent of $C + D$: for $D \to 0$ and $A \to 0$ the hypothesis
$\tau \le C\sqrt\tau\,L$ only gives $\tau \le C^2L^2$ with $L \approx 2\log C + \log\alpha +
2\log L$, which is not below $\frac85\log C + O(1)$ for large $C$. The correct exponent is
$2$, and this is what Lemma~\ref{lem:self_bounding} proves, with the weaker hypotheses
$A, B, C, D \ge 0$, $B \le E$ (positivity and $B \ge 1$ are not used). The sharper bound
$C_1 = \frac85\log(4\alpha K)$ of Lemma~\ref{lem:pana_self_bounding_core} is what the proof
actually gives. In the application (Definition~\ref{def:upper_bound_const}), $D \ge 1$, so
the form with $(C + D)^2$ applies; the asymptotic order of the bound is unchanged (only the
logarithmic factor changes).

\textbf{Lean remark.} \texttt{Real.log 0 = 0}, so the hypothesis and the conclusion make sense
for $\tau = 0$ and $K = 0$ without side conditions; the case $\tau \le 1$ is treated first and
$K > 0$ is derived in the other case. The powers $y^{3/16}$, $y^{3/8}$, $y^{5/8}$ are
\texttt{Real.rpow}; $\sqrt\tau$ is \texttt{Real.sqrt} with \texttt{Real.sq\_sqrt},
\texttt{Real.sqrt\_le\_sqrt}. The numerical facts are $e \ge 2.7$
(\texttt{Real.exp\_one\_gt\_d9}) and $4/11 \le e$.
\end{remark}

